\documentclass[12pt,openany]{book}

% ── Packages ──────────────────────────────────────────────────
\usepackage[margin=1in]{geometry}
\usepackage{amsmath,amssymb,amsthm,mathtools}
\usepackage{stmaryrd}
\usepackage{tikz-cd}
\usepackage{enumitem}
\usepackage{hyperref}
\usepackage{cleveref}
\usepackage{tocloft}
\usepackage{fancyhdr}
\usepackage{titlesec}
\usepackage[expansion=false]{microtype}

% ── Theorem environments ─────────────────────────────────────
\theoremstyle{plain}
\newtheorem{theorem}{Theorem}[chapter]
\newtheorem{lemma}[theorem]{Lemma}
\newtheorem{proposition}[theorem]{Proposition}
\newtheorem{corollary}[theorem]{Corollary}
\theoremstyle{definition}
\newtheorem{definition}[theorem]{Definition}
\newtheorem{example}[theorem]{Example}
\newtheorem{construction}[theorem]{Construction}
\theoremstyle{remark}
\newtheorem{remark}[theorem]{Remark}
\newtheorem{notation}[theorem]{Notation}

% ── JuGeo macros ──────────────────────────────────────────────
\newcommand{\Grade}{\mathsf{Grade}}
\newcommand{\opG}{\oplus}
\newcommand{\otG}{\otimes}
\newcommand{\gzero}{\mathbf{0}}
\newcommand{\gone}{\mathbf{1}}

% ── The 4 objects (Fraktur) ───────────────────────────────────
\newcommand{\STop}{\mathfrak{S}}     % Semantic Topos
\newcommand{\PStr}{\mathfrak{P}}     % Proof Structure
\newcommand{\CSpec}{\mathfrak{H}}    % Cohomological Spectrum
\newcommand{\GEval}{\mathfrak{Q}}    % Graded Evaluation

% ── Judgment shorthand ────────────────────────────────────────
\newcommand{\JJ}{\mathcal{J}}
\newcommand{\Sh}{\mathsf{Sh}}
\newcommand{\PSh}{\mathsf{PSh}}
\newcommand{\Hom}{\mathrm{Hom}}
\newcommand{\id}{\mathrm{id}}
\newcommand{\ob}{\mathrm{ob}}
\newcommand{\op}{\mathrm{op}}
\newcommand{\Set}{\mathsf{Set}}
\newcommand{\Type}{\mathsf{Type}}
\newcommand{\Prop}{\mathsf{Prop}}
\newcommand{\Cech}{\check{C}}
\newcommand{\coker}{\mathrm{coker}}

% ── CTT / NL macros ──────────────────────────────────────────
\newcommand{\phon}{\mathsf{phon}}
\newcommand{\morph}{\mathsf{morph}}
\newcommand{\synt}{\mathsf{synt}}
\newcommand{\sem}{\mathsf{sem}}
\newcommand{\prag}{\mathsf{prag}}
\newcommand{\disc}{\mathsf{disc}}
\newcommand{\info}{\mathsf{info}}
\newcommand{\pros}{\mathsf{pros}}
\newcommand{\poet}{\mathsf{poet}}
\newcommand{\aesth}{\mathsf{aesth}}
\newcommand{\GId}{\mathsf{GId}}
\newcommand{\FMLM}{\mathsf{FMLM}}
\newcommand{\HPLF}{\mathsf{HPLF}}

% ── Page style ────────────────────────────────────────────────
\pagestyle{fancy}
\fancyhf{}
\fancyhead[LE,RO]{\thepage}
\fancyhead[LO]{\nouppercase{\rightmark}}
\fancyhead[RE]{\nouppercase{\leftmark}}

% ══════════════════════════════════════════════════════════════
\title{\Huge\textbf{Judgment Geometry}\\[0.5em]
  \Large A Four-Object Framework with\\
  Compositional Theory of Theories}
\author{Halley Young}
\date{2025}

\begin{document}
\frontmatter
\maketitle

\chapter*{Preface}
\addcontentsline{toc}{chapter}{Preface}

This monograph develops Judgment Geometry (JuGeo), a framework in
which every judgment---about a program, a sentence, a poem, a
proof---is captured by four deeply structured mathematical objects:
the Semantic Topos~$\STop$, the Proof Structure~$\PStr$, the
Cohomological Spectrum~$\CSpec$, and the Graded Evaluation~$\GEval$.

These four objects arise from a single question that every domain
asks:
\begin{quote}
  \emph{Over this space~($\STop$), how well~($\GEval$) can we
  construct~($\PStr$) a global solution, and what prevents
  it~($\CSpec$)?}
\end{quote}
\noindent Each object corresponds to one irreducible aspect of
judgment: \textbf{Where} (the context), \textbf{What} (the
construction), \textbf{Why not} (the obstructions), and
\textbf{How good} (the quality).  Together they form the judgment
4-tuple $\JJ = (\STop, \PStr, \CSpec, \GEval)$, and together they
are exactly the data of a \emph{graded section of a fibration}---one
operation, viewed from four angles.

Parts~I--II develop the core four-object theory and establish the
mathematical foundations.  Part~III extends JuGeo to natural
language and poetry via the Compositional Theory of Theories (CTT).
Part~IV contains detailed proofs.  Part~V presents worked examples
across all domains JuGeo targets.  Part~VI explores applications,
both existing and speculative.

\tableofcontents

\mainmatter

% ── Part I: Foundations ───────────────────────────────────────
\part{Foundations}
% Part I: Foundations
% This file is \include'd from main.tex

\chapter{Introduction}
\label{ch:intro}

\section{The fundamental problem}
\label{sec:intro:problem}

Consider the following four artifacts, drawn from four different domains.

\begin{enumerate}[label=(\roman*)]
\item A Python function implementing binary search uses the loop guard
  \texttt{while low < high} instead of \texttt{while low <= high}.
  Every branch of the loop body is locally correct; the error is
  visible only at the \emph{interface} between the guard and the body.

\item The English sentence ``The horse raced past the barn fell'' is
  locally grammatical in every sub-phrase---``the horse raced past the
  barn'' parses as a relative clause, and ``fell'' is a well-formed
  verb---but the global parse fails because the two fragments cannot
  be composed into a single syntactic tree without additional
  structure.

\item A sonnet whose three quatrains are metrically and imagistically
  perfect but whose volta (the turn between the octave and the sestet)
  introduces a non-sequitur.  Each stanza, read in isolation, is
  accomplished verse; the poem as a whole is incoherent.

\item A mathematical proof in which each lemma is impeccably
  established but the overall argument is circular: Lemma~3 depends on
  Lemma~1, Lemma~1 depends on Lemma~5, and Lemma~5 depends on
  Lemma~3.  Every \emph{local} step is valid; the \emph{global}
  derivation is unsound.
\end{enumerate}

In every case the pattern is the same: \emph{local correctness does
not imply global correctness}.  Locally verified parts fail to compose
into a globally verified whole.  The bug, the ambiguity, the
incoherence, and the circularity are all invisible when one inspects
the parts in isolation; they emerge only at the \emph{overlaps}
between parts.

This observation is the starting point of \textbf{Judgment Geometry}
(JuGeo).  We develop a mathematical framework in which:

\begin{itemize}
\item The artifact under analysis is modeled as a \emph{site}---a
  category of regions equipped with a Grothendieck topology specifying
  which families of sub-regions jointly observe everything about a
  larger region.

\item Properties of the artifact are modeled as \emph{presheaves} (or
  \emph{sheaves}) on the site: functors that assign to each region a
  set of local data and to each inclusion a restriction map.

\item The question ``do locally verified properties hold globally?''
  is answered by the \emph{descent condition}: local sections glue
  into a global section if and only if the first \v{C}ech cohomology
  group vanishes, $\check{H}^1(\mathcal{U}, \mathcal{F}) = 0$.

\item When the cohomology does not vanish, the non-trivial classes are
  \emph{obstructions}---first-class mathematical objects that identify
  exactly which overlaps between regions fail to reconcile.

\item Every claim carries a \emph{grade} drawn from a semiring
  $\Grade$, measuring the quality or confidence of the evidence
  supporting it.
\end{itemize}

The result is a \emph{judgment}: a structured datum that records not
only what is claimed and where, but how well the claim is supported,
what evidence underlies it, what obstacles prevent global validity, and
how those obstacles might be resolved.


\section{Why existing frameworks fail}
\label{sec:intro:existing}

Three broad families of tools address the problem of artifact
correctness, and each captures only a fragment of the full picture.

\paragraph{Type theory and proof assistants.}
Systems such as Lean, Coq, and F$^\star$ are founded on dependent type
theory.  They provide absolute certainty about the claims they
verify---if a term type-checks, the proposition it inhabits is
true---but they operate in a binary regime: either the proof
type-checks or it does not.  There is no room for partial evidence,
for claims supported at different confidence levels by different
sources, or for graded quality.  A flawless proof and a dubious
conjecture are equally rejected when they fail to type-check; a
beautiful proof and a tedious case-bash are equally accepted when they
do.  Moreover, the local-to-global composition problem is solved by
the programmer manually: the system checks individual lemmas but does
not automatically determine whether the lemma graph is globally
coherent.

\paragraph{Testing and runtime analysis.}
Test suites, fuzzing, and runtime monitoring sample the behavior of an
artifact on finitely many inputs.  They can detect bugs but cannot
certify their absence.  They have no internal notion of geometric
structure: a test suite does not know that two test cases exercise
overlapping code regions, or that a passing test on the loop body and
a passing test on the loop guard jointly cover the loop only if the
guard--body interface is itself tested.  There is no notion of
\emph{covering} or of \emph{descent}.

\paragraph{Machine learning and LLM-based analysis.}
Statistical models can classify artifacts (``this code is likely
buggy,'' ``this sentence is likely grammatical'') and even propose
repairs, but they provide confidence scores rather than proofs.  The
scores are not compositional: the confidence that function $f$ is
correct and the confidence that function $g$ is correct do not combine
in any principled way into a confidence that the module $\{f, g\}$ is
correct.  There is no algebra of trust, no sheaf-theoretic gluing, and
no cohomological obstruction theory.

\medskip

Judgment Geometry fills the gap by integrating geometric context
(sites and covers), logical evidence (typed proofs and witnesses),
cohomological invariants (\v{C}ech cohomology classes measuring
failure to glue), and quality measures (grades from a semiring) into a
single, unified mathematical framework.


\section{The core insight}
\label{sec:intro:insight}

The central idea can be stated in a single sentence:

\begin{quote}
\itshape
A judgment is a section of a presheaf over a semantic site, carrying
proof-relevant evidence graded by quality, whose global validity is
governed by the vanishing of \v{C}ech cohomology.
\end{quote}

More precisely, we associate to any structured artifact---a program, a
sentence, a poem, a proof---a small category $\mathcal{C}$ whose
objects are the meaningful sub-regions of the artifact and whose
morphisms are the structural relationships (containment, restriction,
transport) between regions.  We equip $\mathcal{C}$ with a
Grothendieck topology $J$ specifying which families of sub-regions
constitute \emph{covers}.  The pair $(\mathcal{C}, J)$ is a
\emph{site}.

Over this site, properties of the artifact become \emph{presheaves}:
contravariant functors $\mathcal{F} : \mathcal{C}^{\op} \to \Set$
assigning to each region $U$ a set $\mathcal{F}(U)$ of local claims,
with restriction maps that specialize claims to sub-regions.  A
\emph{local section} is an element $s \in \mathcal{F}(U)$---a claim
about a particular region.  A \emph{global section} is a compatible
family of local sections: one for each member of a covering family,
agreeing on all overlaps.

The \v{C}ech cohomology group $\check{H}^1(\mathcal{U}, \mathcal{F})$
measures the obstruction to gluing: it vanishes if and only if every
compatible family of local sections extends to a global section.  When
it does not vanish, its non-trivial elements identify precisely the
overlaps at which gluing fails.

Finally, each claim carries a \emph{grade} $T \in \Grade$, where
$\Grade$ is a semiring measuring quality along multiple dimensions.
The grade algebra provides the quantitative layer: not merely ``is
this artifact correct?'' but ``how good is it, along which axes, and
how does quality compose?''


\section{The 4-object approach}
\label{sec:intro:4objects}

The full data of a judgment can be organized into four maximally rich
mathematical objects, each drawn from a well-established branch of
mathematics:

\[
\JJ = (\STop,\; \PStr,\; \CSpec,\; \GEval).
\]

\begin{description}[leftmargin=2.5em]
\item[$\STop$ --- the Semantic Topos.]
  The site $(\mathcal{C}, J)$ together with the category of sheaves
  $\Sh(\mathcal{C}, J)$.  This is where the artifact lives: its
  regions, its covers, and the gluing conditions that govern
  local-to-global transitions.

\item[$\PStr$ --- the Proof Structure.]
  The dependent-type-theoretic datum: the context, the proposition,
  the evidence term, and the typing judgment $\Gamma \vdash e : A$.
  This is the logical content of the claim: what is asserted, and what
  evidence supports it.

\item[$\CSpec$ --- the Cohomological Spectrum.]
  The \v{C}ech cohomology groups $\check{H}^n(\mathcal{U},
  \mathcal{F})$ and the obstruction classes they contain.  This is the
  diagnostic layer: it measures failure to glue, localizes
  inconsistencies, and guides repair.

\item[$\GEval$ --- the Graded Evaluation.]
  The grade semiring $\Grade$ together with the evaluation functor
  $\mu$ that maps derivations to quality measures.  This is the
  quantitative layer: it assigns a multi-dimensional quality score to
  every judgment and every composition of judgments.
\end{description}

The central thesis of this monograph is that these four objects are
not merely a convenient packaging of the data involved in judgment
but are \emph{exactly the right level of abstraction}: each
corresponds to a deep mathematical structure (topos theory,
type theory, cohomology, quantale theory), and together they
capture everything needed for judgment across domains.

Part~I of this monograph (the present Part) develops the mathematical
foundations: sites, presheaves, sheaves, \v{C}ech cohomology, and the
Grade semiring.  It introduces the judgment and its individual
components.  Part~II then develops each of the four objects
$(\STop, \PStr, \CSpec, \GEval)$ in depth, showing how they unify
these foundations into a single coherent framework.


\section{Motivating examples}
\label{sec:intro:examples}

We return to the four artifacts from \S\ref{sec:intro:problem} and
sketch how Judgment Geometry analyzes each.

\begin{example}[A buggy program]
\label{ex:intro:program}
The binary search function with guard \texttt{while low < high}
determines a site with three coordinates: $c_{\mathrm{guard}}$ (the
loop guard), $c_{\mathrm{body}}$ (the loop body), and
$c_{\mathrm{fn}}$ (the function).  The covering family is
$\{c_{\mathrm{guard}} \to c_{\mathrm{fn}},\; c_{\mathrm{body}} \to
c_{\mathrm{fn}}\}$.  The guard's local section claims ``when
$\mathit{low} = \mathit{high}$, the loop exits without checking
$\mathit{arr}[\mathit{low}]$.''  The body's local section assumes
``$\mathit{arr}[\mathit{mid}]$ is always checked.''  These disagree on
the overlap $c_{\mathrm{guard}} \cap c_{\mathrm{body}}$: the
coboundary $\delta^0 \neq 0$, so $\check{H}^1 \neq 0$.  The
obstruction class identifies the guard--body interface, and the repair
frontier suggests changing \texttt{<} to \texttt{<=}.
\end{example}

\begin{example}[An ambiguous sentence]
\label{ex:intro:sentence}
The garden-path sentence ``The horse raced past the barn fell'' determines
a site whose coordinates are the morphemes, phrases, and the full
sentence.  At the phrase level, ``the horse raced past the barn''
parses as a noun phrase with a reduced relative clause, and ``fell'' is
a well-formed verb phrase.  But on the overlap---the interface where
the parsed noun phrase must serve as the subject of ``fell''---the two
local parses are incompatible: the noun phrase has already consumed the
verb ``raced,'' leaving no room for ``fell'' as the main verb.  The
\v{C}ech cohomology does not vanish; the obstruction localizes the
ambiguity to the phrase--sentence boundary.
\end{example}

\begin{example}[A poem with forced rhyme]
\label{ex:intro:poem}
Consider a sonnet in which the first quatrain is metrically perfect
and imagistically vivid, but the second quatrain distorts its syntax
to force a rhyme: ``And through the misty vale the river \emph{gleams}
/ Reflecting back her silver-threaded \emph{beams}.''  Each quatrain,
analyzed as a local section over the stanza coordinate, receives a
high grade for meter and rhyme but a lower grade for naturalness.  The
grade semiring captures this: $T = T_{\text{meter}} \otG
T_{\text{rhyme}} \otG T_{\text{imagery}} \otG T_{\text{naturalness}}$,
and the forced rhyme depresses $T_{\text{naturalness}}$, pulling the
overall grade down through the bottleneck principle.
\end{example}

\begin{example}[A flawed proof]
\label{ex:intro:proof}
A mathematical argument consists of five lemmas, each individually
correct, but with a circular dependency: Lemma~3 uses Lemma~1,
Lemma~1 uses Lemma~5, and Lemma~5 uses Lemma~3.  The site has
coordinates for each lemma and for the overall theorem.  Each local
section (lemma) is valid; the \v{C}ech computation reveals a
non-trivial 1-cocycle supported on the triangle of overlaps
$\{1 \cap 3,\; 3 \cap 5,\; 5 \cap 1\}$.  More precisely, the
dependency cycle creates a failure of the cocycle condition on triple
overlaps, yielding a non-trivial class in $\check{H}^1$.  The
obstruction pinpoints the circular dependency and the repair frontier
suggests breaking the cycle.
\end{example}


\section{Overview of the monograph}
\label{sec:intro:overview}

This monograph is organized into six Parts.

\paragraph{Part~I: Foundations
  (Chapters~\ref{ch:intro}--\ref{ch:grade}).}
The present Part develops the mathematical foundations from scratch.
Chapter~\ref{ch:sites} introduces category theory, Grothendieck
topologies, presheaves, sheaves, and the Yoneda lemma.
Chapter~\ref{ch:judgment} defines the judgment and its components.
Chapter~\ref{ch:cech} develops \v{C}ech cohomology and descent.
Chapter~\ref{ch:grade} introduces the Grade semiring and the
evaluation functor.

\paragraph{Part~II: The Judgment Framework.}
Develops $\STop$, $\PStr$, $\CSpec$, and $\GEval$ as standalone
mathematical structures, explaining the design principles behind
the 4-tuple and establishing the judgment fibration.  Also shows
that the component-wise flattened view
$(c, \varphi, A, E, O, B, T, \Pi)$ loses no information.

\paragraph{Part~III: Compositional Theory of Theories.}
Extends Judgment Geometry to natural language and poetry via a
multi-stratum architecture in which each linguistic level
(phonological, morphological, syntactic, semantic, pragmatic,
discourse, poetic, aesthetic) is a coordinate in a master site.

\paragraph{Part~IV: Formal Proofs.}
Collects the detailed proofs of all major theorems stated in
Parts~I--III.

\paragraph{Part~V: Worked Examples.}
Presents extended worked examples across all four domains: program
verification, natural language analysis, poetry evaluation, and
mathematical reasoning.

\paragraph{Part~VI: Applications and Future Directions.}
Explores applications of Judgment Geometry to automated repair,
synthesis, equivalence checking, incremental analysis, documentation
alignment, and multi-agent verification.

\chapter{Semantic Sites and Presheaves}
\label{ch:sites}

This chapter develops the categorical foundations on which Judgment
Geometry rests.  We begin with the basic language of category theory,
introduce Grothendieck topologies and sites, and then build presheaves
and sheaves over sites.  The Yoneda lemma, the central organizing
principle of the theory, receives a full proof.  We close with
concrete examples of sites drawn from all four target domains.

The treatment is self-contained: no prior knowledge of category theory
is assumed beyond elementary set theory.

%% ─── Section 1: Category Theory Preliminaries ─────────────────────

\section{Categories, functors, and natural transformations}
\label{sec:sites:categories}

\begin{definition}[Category]
\label{def:category}
A \emph{category} $\mathcal{C}$ consists of:
\begin{enumerate}[label=(\roman*)]
\item A collection $\ob(\mathcal{C})$ of \emph{objects}.
\item For each pair of objects $A, B \in \ob(\mathcal{C})$, a set
  $\Hom_{\mathcal{C}}(A, B)$ of \emph{morphisms} (or \emph{arrows})
  from $A$ to $B$.
\item For each triple $A, B, C$, a \emph{composition} map
  \[
    \circ : \Hom_{\mathcal{C}}(B, C) \times \Hom_{\mathcal{C}}(A, B)
    \to \Hom_{\mathcal{C}}(A, C).
  \]
\item For each object $A$, an \emph{identity} morphism
  $\id_A \in \Hom_{\mathcal{C}}(A, A)$.
\end{enumerate}
These data satisfy:
\begin{itemize}
\item \emph{Associativity}: $h \circ (g \circ f) = (h \circ g) \circ f$
  whenever the compositions are defined.
\item \emph{Identity}: $f \circ \id_A = f = \id_B \circ f$ for every
  $f : A \to B$.
\end{itemize}
A category is \emph{small} if $\ob(\mathcal{C})$ is a set (not a
proper class).
\end{definition}

\begin{example}[Basic categories]
\label{ex:basic-categories}
\begin{enumerate}[label=(\alph*)]
\item $\Set$: the category whose objects are sets and whose morphisms are functions.
\item For any partially ordered set $(P, \leq)$: the category with one morphism $a \to b$ whenever $a \leq b$.
\item For any group $G$: the category with one object $\ast$ and $\Hom(\ast, \ast) = G$, composition given by group multiplication.
\end{enumerate}
\end{example}

\begin{definition}[Functor]
\label{def:functor}
Let $\mathcal{C}$ and $\mathcal{D}$ be categories.  A \emph{functor}
$F : \mathcal{C} \to \mathcal{D}$ consists of:
\begin{enumerate}[label=(\roman*)]
\item A map $F : \ob(\mathcal{C}) \to \ob(\mathcal{D})$.
\item For each pair $A, B \in \ob(\mathcal{C})$, a map
  $F_{A,B} : \Hom_{\mathcal{C}}(A, B) \to \Hom_{\mathcal{D}}(F(A), F(B))$.
\end{enumerate}
These maps preserve identities ($F(\id_A) = \id_{F(A)}$) and
composition ($F(g \circ f) = F(g) \circ F(f)$).

A functor $F : \mathcal{C}^{\op} \to \mathcal{D}$, where
$\mathcal{C}^{\op}$ denotes the \emph{opposite category} (same
objects, reversed morphisms), is called a \emph{contravariant functor}
from $\mathcal{C}$ to $\mathcal{D}$.
\end{definition}

\begin{definition}[Natural transformation]
\label{def:nat-trans}
Let $F, G : \mathcal{C} \to \mathcal{D}$ be functors.  A \emph{natural
transformation} $\eta : F \Rightarrow G$ is a family of morphisms
$\{\eta_A : F(A) \to G(A)\}_{A \in \ob(\mathcal{C})}$ such that for
every morphism $f : A \to B$ in $\mathcal{C}$, the diagram
\[
\begin{tikzcd}
F(A) \ar[r, "\eta_A"] \ar[d, "F(f)"'] & G(A) \ar[d, "G(f)"] \\
F(B) \ar[r, "\eta_B"']                 & G(B)
\end{tikzcd}
\]
commutes: $G(f) \circ \eta_A = \eta_B \circ F(f)$.  If each $\eta_A$
is an isomorphism, $\eta$ is a \emph{natural isomorphism}.
\end{definition}


%% ─── Section 2: Sites ─────────────────────────────────────────────

\section{Grothendieck topologies and sites}
\label{sec:sites:grothendieck}

The notion of a Grothendieck topology generalizes the idea of an open
cover in point-set topology to an arbitrary category.  It specifies,
for each object, which families of morphisms should be regarded as
``covers.''

\begin{definition}[Sieve]
\label{def:sieve}
Let $\mathcal{C}$ be a category and $X \in \ob(\mathcal{C})$.  A
\emph{sieve} on $X$ is a collection $S$ of morphisms with codomain
$X$ that is closed under precomposition: if $(f : U \to X) \in S$ and
$g : V \to U$ is any morphism, then $(f \circ g : V \to X) \in S$.

Equivalently, $S$ is a subfunctor of the representable presheaf
$\Hom_{\mathcal{C}}(-, X) : \mathcal{C}^{\op} \to \Set$.
\end{definition}

\begin{definition}[Grothendieck topology]
\label{def:grothendieck-topology}
A \emph{Grothendieck topology} $J$ on a category $\mathcal{C}$ assigns
to each object $X \in \ob(\mathcal{C})$ a collection $J(X)$ of sieves
on $X$, called \emph{covering sieves}, satisfying three axioms:
\begin{enumerate}
\item \textbf{Identity (maximality).}
  The maximal sieve
  $t_X = \{f : U \to X \mid f \text{ is any morphism to } X\}$
  is in $J(X)$.

\item \textbf{Stability (base change).}
  If $S \in J(X)$ and $g : Y \to X$ is any morphism, then the
  \emph{pullback sieve}
  \[
    g^* S = \{h : Z \to Y \mid g \circ h \in S\}
  \]
  is in $J(Y)$.

\item \textbf{Transitivity (local character).}
  If $S \in J(X)$ and $R$ is a sieve on $X$ such that for every
  $(f : U \to X) \in S$ the pullback sieve $f^* R$ is in $J(U)$, then
  $R \in J(X)$.
\end{enumerate}
A pair $(\mathcal{C}, J)$ is called a \emph{site}.
\end{definition}

\begin{remark}
\label{rem:covering-families}
In practice, one often specifies a topology via \emph{covering
families} rather than sieves.  A covering family of $X$ is a set
$\{f_i : U_i \to X\}_{i \in I}$ of morphisms with common codomain
$X$.  The covering family \emph{generates} the sieve
$S = \{g : V \to X \mid g \text{ factors through some } f_i\}$.
When we say a family ``is a cover,'' we mean the sieve it generates is
a covering sieve.
\end{remark}

\begin{example}[The classical site]
\label{ex:classical-site}
Let $X$ be a topological space and $\mathcal{O}(X)$ the poset category
of open sets (morphisms are inclusions $U \hookrightarrow V$ whenever
$U \subseteq V$).  The Grothendieck topology is defined by declaring
$\{U_i \hookrightarrow U\}_{i \in I}$ to be a covering family whenever
$\bigcup_{i} U_i = U$.  The three axioms are immediate from the
properties of open sets.
\end{example}


%% ─── Section 3: Presheaves ────────────────────────────────────────

\section{Presheaves and the free cocompletion}
\label{sec:sites:presheaves}

\begin{definition}[Presheaf]
\label{def:presheaf}
Let $\mathcal{C}$ be a small category.  A \emph{presheaf} on
$\mathcal{C}$ is a functor $F : \mathcal{C}^{\op} \to \Set$.  The
category of presheaves on $\mathcal{C}$, denoted
\[
  \PSh(\mathcal{C}) = [\mathcal{C}^{\op}, \Set],
\]
has presheaves as objects and natural transformations as morphisms.
\end{definition}

For a morphism $f : U \to V$ in $\mathcal{C}$, a presheaf $F$ gives a
\emph{restriction map} $F(f) : F(V) \to F(U)$.  If $s \in F(V)$ is a
``section over $V$,'' then $F(f)(s) \in F(U)$ is its ``restriction to
$U$.''  We often write $s|_U$ for $F(f)(s)$ when the morphism $f$ is
clear from context.

\begin{definition}[Representable presheaf (Yoneda embedding)]
\label{def:yoneda-embedding}
For each object $c \in \ob(\mathcal{C})$, the \emph{representable
presheaf} $\mathrm{y}(c) = \Hom_{\mathcal{C}}(-, c)$ is defined by:
\begin{align*}
  \mathrm{y}(c)(U) &= \Hom_{\mathcal{C}}(U, c), \\
  \mathrm{y}(c)(f : U \to V) &= (- \circ f) :
    \Hom_{\mathcal{C}}(V, c) \to \Hom_{\mathcal{C}}(U, c).
\end{align*}
The assignment $c \mapsto \mathrm{y}(c)$ defines the \emph{Yoneda
embedding} $\mathrm{y} : \mathcal{C} \to \PSh(\mathcal{C})$.
\end{definition}

\begin{theorem}[The Yoneda lemma]
\label{thm:yoneda}
Let $\mathcal{C}$ be a small category, $c \in \ob(\mathcal{C})$, and
$F \in \PSh(\mathcal{C})$.  There is a bijection, natural in both $c$
and $F$:
\[
  \Hom_{\PSh(\mathcal{C})}(\mathrm{y}(c), F) \;\cong\; F(c).
\]
\end{theorem}

\begin{proof}
We construct explicit maps in both directions and verify they are
mutually inverse.

\medskip\noindent\textbf{From left to right.}  Given a natural
transformation $\eta : \mathrm{y}(c) \Rightarrow F$, define
$\Phi(\eta) = \eta_c(\id_c) \in F(c)$.  This is an element of $F(c)$
because $\eta_c : \mathrm{y}(c)(c) = \Hom(c, c) \to F(c)$ and
$\id_c \in \Hom(c, c)$.

\medskip\noindent\textbf{From right to left.}  Given $x \in F(c)$,
define a natural transformation $\Psi(x) : \mathrm{y}(c) \Rightarrow F$
by setting, for each object $U$ and each morphism $f \in \mathrm{y}(c)(U) = \Hom(U, c)$:
\[
  \Psi(x)_U(f) = F(f)(x) \in F(U).
\]
We must check that $\Psi(x)$ is natural.  Let $g : U \to V$ be a
morphism in $\mathcal{C}$.  For any $f \in \Hom(V, c)$:
\begin{align*}
  (\Psi(x)_U \circ \mathrm{y}(c)(g))(f)
  &= \Psi(x)_U(f \circ g) = F(f \circ g)(x), \\
  (F(g) \circ \Psi(x)_V)(f)
  &= F(g)(F(f)(x)) = F(f \circ g)(x),
\end{align*}
where the last equality uses the functoriality of $F$ (since $F$ is
contravariant, $F(g) \circ F(f) = F(f \circ g)$).  Hence the
naturality square commutes.

\medskip\noindent\textbf{Mutual inverses.}
$\Phi(\Psi(x)) = \Psi(x)_c(\id_c) = F(\id_c)(x) = x$ since $F$
preserves identities.

Conversely, for any $\eta : \mathrm{y}(c) \Rightarrow F$ and any
$f \in \Hom(U, c)$:
\[
  \Psi(\Phi(\eta))_U(f) = F(f)(\eta_c(\id_c))
  = \eta_U(\mathrm{y}(c)(f)(\id_c))
  = \eta_U(f),
\]
where the second equality uses naturality of $\eta$ applied to $f : U \to c$:
\[
  F(f) \circ \eta_c = \eta_U \circ \mathrm{y}(c)(f).
\]
Evaluating both sides at $\id_c$ gives $F(f)(\eta_c(\id_c)) = \eta_U(f)$.
Hence $\Psi(\Phi(\eta)) = \eta$.
\end{proof}

\begin{corollary}[Full faithfulness of Yoneda]
\label{cor:yoneda-ff}
The Yoneda embedding $\mathrm{y} : \mathcal{C} \to \PSh(\mathcal{C})$
is full and faithful: for all $c, d \in \ob(\mathcal{C})$,
\[
  \Hom_{\mathcal{C}}(c, d) \;\cong\;
  \Hom_{\PSh(\mathcal{C})}(\mathrm{y}(c), \mathrm{y}(d)).
\]
\end{corollary}

\begin{proof}
Apply the Yoneda lemma with $F = \mathrm{y}(d)$:
$\Hom(\mathrm{y}(c), \mathrm{y}(d)) \cong \mathrm{y}(d)(c) = \Hom(c, d)$.
\end{proof}

\begin{remark}[Presheaves as the free cocompletion]
\label{rem:free-cocompletion}
The category $\PSh(\mathcal{C})$ is the free cocompletion of
$\mathcal{C}$: it is the universal category with all small colimits,
equipped with a functor $\mathrm{y} : \mathcal{C} \to \PSh(\mathcal{C})$
such that every functor $F : \mathcal{C} \to \mathcal{D}$ into a
cocomplete category $\mathcal{D}$ extends essentially uniquely to a
colimit-preserving functor $\hat{F} : \PSh(\mathcal{C}) \to \mathcal{D}$.
This universal property justifies calling presheaves ``generalized
objects'' of $\mathcal{C}$.
\end{remark}


%% ─── Section 4: Sheaves ──────────────────────────────────────────

\section{Sheaves and sheafification}
\label{sec:sites:sheaves}

A sheaf is a presheaf satisfying a gluing condition with respect to a
Grothendieck topology.  The gluing condition is the categorical
abstraction of the familiar fact that a function defined on a
topological space can be constructed by specifying it on each open set
of a cover, provided the definitions agree on overlaps.

\begin{definition}[Matching family and amalgamation]
\label{def:matching-family}
Let $(\mathcal{C}, J)$ be a site, $F \in \PSh(\mathcal{C})$, and
$\{f_i : U_i \to X\}_{i \in I}$ a covering family in $J(X)$.  A
\emph{matching family} (or \emph{compatible family}) for the cover is
a collection $\{s_i \in F(U_i)\}_{i \in I}$ such that for every pair
$i, j \in I$ and every pair of morphisms
$g : W \to U_i$, $h : W \to U_j$ with $f_i \circ g = f_j \circ h$,
we have $F(g)(s_i) = F(h)(s_j)$.

An \emph{amalgamation} of such a family is a section $s \in F(X)$ with
$F(f_i)(s) = s_i$ for all $i \in I$.
\end{definition}

\begin{definition}[Sheaf]
\label{def:sheaf}
A presheaf $F \in \PSh(\mathcal{C})$ is a \emph{sheaf} for the
topology $J$ if for every covering family $\{f_i : U_i \to X\}$ in $J(X)$,
every matching family has a unique amalgamation.  Equivalently, $F$
is a sheaf if the diagram
\[
  F(X) \;\xrightarrow{\;\sim\;}\;
  \mathrm{eq}\!\left(
    \prod_{i \in I} F(U_i) \;\rightrightarrows\;
    \prod_{(i,j) \in I \times I} F(U_i \times_X U_j)
  \right)
\]
is an isomorphism for every cover, where the two parallel arrows are
induced by the two projections from the fiber product
$U_i \times_X U_j$ to $U_i$ and to $U_j$.

The full subcategory of $\PSh(\mathcal{C})$ consisting of sheaves is
denoted $\Sh(\mathcal{C}, J)$.
\end{definition}

\begin{proposition}[Sheafification]
\label{prop:sheafification}
The inclusion functor $\iota : \Sh(\mathcal{C}, J) \hookrightarrow
\PSh(\mathcal{C})$ has a left adjoint
\[
  a : \PSh(\mathcal{C}) \to \Sh(\mathcal{C}, J),
\]
called \emph{sheafification} (or the \emph{associated sheaf functor}).
For every presheaf $F$, there is a canonical morphism
$\eta_F : F \to \iota(a(F))$ that is universal among morphisms from
$F$ to sheaves.  The adjunction
$a \dashv \iota : \PSh(\mathcal{C}) \rightleftarrows \Sh(\mathcal{C}, J)$
is exact: $a$ preserves all finite limits.
\end{proposition}

\begin{proof}[Proof sketch]
The sheafification $a(F)$ is constructed in two steps.  First, one
forms the ``separated presheaf'' $F^+$ by quotienting each $F(X)$ by
the equivalence relation that identifies two sections $s, t \in F(X)$
whenever there exists a cover $\{U_i \to X\}$ such that $s|_{U_i} = t|_{U_i}$
for all $i$.  This yields a separated presheaf (injectivity of
amalgamation).  Second, one repeats the construction: $(F^+)^+$ is a
sheaf (existence of amalgamation).  Thus $a(F) = (F^+)^+$.  The
adjunction and exactness are standard; see \cite{johnstone2002} for
the full proof.
\end{proof}

\begin{remark}
\label{rem:topos}
The category $\Sh(\mathcal{C}, J)$ is a \emph{Grothendieck topos}: it
is complete, cocomplete, has a subobject classifier, and satisfies
the Giraud axioms.  In Judgment Geometry, the topos $\Sh(\mathcal{C}, J)$
is the first of the four objects, $\STop = \Sh(\mathcal{C}, J)$, and
provides the geometric universe in which judgments live.
\end{remark}


%% ─── Section 5: Examples of sites ─────────────────────────────────

\section{Examples of semantic sites}
\label{sec:sites:examples}

We now illustrate how the abstract machinery of sites applies to the
four domains motivating Judgment Geometry.

\begin{example}[Program verification]
\label{ex:site:programs}
Given a program $P$, the \emph{semantic site}
$\mathcal{S}(P) = (\mathcal{C}_P, J_P)$ is defined as follows:

\medskip\noindent\textbf{Objects (coordinates).}  The objects of
$\mathcal{C}_P$ are the structural units of $P$:
\begin{center}
\begin{tabular}{ll}
\emph{Coordinate kind} & \emph{Program construct} \\
\hline
Module   & top-level file \\
Class    & \texttt{class} definition \\
Function & \texttt{def} / \texttt{async def} \\
Scope    & lexical scope (loop body, \texttt{with}-block) \\
Branch   & \texttt{if}/\texttt{else}/\texttt{elif}/\texttt{match} case
\end{tabular}
\end{center}

\medskip\noindent\textbf{Morphisms.}  Morphisms encode four
structural relationships:
\begin{description}[nosep,leftmargin=2em]
\item[Restriction:] function $\to$ loop body, class $\to$ method.
\item[Inclusion:] branch $\hookrightarrow$ conditional, method $\hookrightarrow$ class.
\item[Transport:] call site $\to$ callee (equivalence-preserving).
\item[Refinement:] coordinate $\to$ finer decomposition.
\end{description}

\medskip\noindent\textbf{Topology.}  Covering families arise from:
\begin{itemize}[nosep]
\item \emph{Control-flow covers}: the branches of an \texttt{if/else}
  cover the conditional.
\item \emph{Scope covers}: a function's local scopes cover the function.
\item \emph{Module covers}: the top-level definitions cover the module.
\end{itemize}

The topology axioms are verified:
\emph{Identity} holds because every coordinate covers itself.
\emph{Stability} holds because restricting a control-flow decomposition
to a sub-scope yields a valid decomposition of the sub-scope.
\emph{Transitivity} holds because nested decompositions compose: if a
module is covered by its functions, and each function is covered by its
branches, the composition covers the module.
\end{example}

\begin{example}[Natural language semantics]
\label{ex:site:nl}
A sentence $S$ determines a site $\mathcal{S}(S)$ whose objects are
linguistic units at multiple granularities:
\begin{center}
\begin{tabular}{ll}
\emph{Coordinate kind} & \emph{Linguistic unit} \\
\hline
Morpheme  & smallest meaningful unit (\emph{un-}, \emph{-ing}) \\
Word      & lexical item \\
Phrase    & syntactic constituent (NP, VP, PP) \\
Clause    & finite or non-finite clause \\
Sentence  & complete sentence \\
Discourse & multi-sentence passage
\end{tabular}
\end{center}
Morphisms are the compositional structure maps: morpheme $\to$ word
(affixation), word $\to$ phrase (phrase-structure rules), phrase $\to$
clause (predication), clause $\to$ sentence (embedding/coordination),
sentence $\to$ discourse (discourse relations).  A word is covered by
its morphemes; a phrase is covered by its immediate constituents; a
sentence is covered by its clauses.

A presheaf over this site assigns to each level a set of semantic
representations (e.g., $\lambda$-terms in the Montagovian tradition),
and the restriction maps implement compositional semantics: the
meaning of a phrase is determined by---and restricts to---the meanings
of its parts.
\end{example}

\begin{example}[Poetry]
\label{ex:site:poetry}
A poem $P$ determines a site $\mathcal{S}(P)$ with a hierarchy of
prosodic and structural units:
\begin{center}
\begin{tabular}{ll}
\emph{Coordinate kind} & \emph{Poetic unit} \\
\hline
Syllable & individual syllable (stressed/unstressed) \\
Foot     & metrical unit (iamb, trochee, dactyl, \ldots) \\
Line     & a line of verse \\
Couplet/Tercet & grouping of 2 or 3 lines \\
Stanza   & stanza (quatrain, sestet, \ldots) \\
Poem     & the entire poem
\end{tabular}
\end{center}
Morphisms encode prosodic containment: syllable $\to$ foot $\to$ line
$\to$ stanza $\to$ poem.  A line is covered by its feet; a stanza is
covered by its lines.  The presheaf assigns to each unit the relevant
quality data: metrical regularity, rhyme adherence, imagery, emotional
tone.  Restriction maps propagate constraints: the meter of a line is
determined by the meters of its feet, and the rhyme scheme of a stanza
is determined by the end-words of its lines.
\end{example}

\begin{example}[Mathematics]
\label{ex:site:math}
A mathematical theory $\mathcal{T}$ determines a site whose objects
are the components of the theory:
\begin{center}
\begin{tabular}{ll}
\emph{Coordinate kind} & \emph{Mathematical unit} \\
\hline
Definition  & definition of a concept \\
Lemma       & auxiliary result \\
Theorem     & main result \\
Proof step  & individual step in a proof \\
Theory      & the entire theory
\end{tabular}
\end{center}
Morphisms encode logical dependency: a theorem depends on (receives
morphisms from) the lemmas it cites; a lemma depends on the
definitions it uses.  A theorem is covered by the lemmas and
definitions that constitute its proof.  The presheaf assigns to each
coordinate the set of logical conditions (axioms, hypotheses) active
at that coordinate, and restriction maps propagate hypotheses down the
dependency tree.
\end{example}


%% ─── Section 6: Restriction and Extension ──────────────────────────

\section{Restriction and extension along morphisms}
\label{sec:sites:restriction}

A fundamental operation in sheaf theory is the \emph{restriction} of
data along morphisms.  Given a presheaf $F$ on $\mathcal{C}$ and a
morphism $f : U \to V$, the restriction map $F(f) : F(V) \to F(U)$
``pulls back'' data from $V$ to $U$.  In all four domains, restriction
has a concrete interpretation:

\begin{itemize}
\item \textbf{Programs:} restricting a module-level specification to a
  function yields the function's local verification conditions.

\item \textbf{Natural language:} restricting a sentence-level semantic
  representation to a phrase yields the phrase's meaning in the
  sentential context.

\item \textbf{Poetry:} restricting a poem-level quality assessment to a
  stanza yields the stanza's contribution to the poem's quality.

\item \textbf{Mathematics:} restricting a theory-level consistency
  requirement to a lemma yields the lemma's proof obligation.
\end{itemize}

\begin{proposition}[Restriction preserves compatibility]
\label{prop:restriction-compatibility}
Let $\{f_i : U_i \to X\}$ be a covering family and
$\{s_i \in F(U_i)\}$ a matching family.  If $g : Y \to X$ is any
morphism, then the restricted family $\{F(\mathrm{pr}_i)(s_i)\}$, where
$\mathrm{pr}_i : U_i \times_X Y \to U_i$ is the projection, is a matching
family for the pullback cover $\{U_i \times_X Y \to Y\}$.
\end{proposition}

\begin{proof}
The pullback cover $\{U_i \times_X Y \to Y\}_{i \in I}$ is a cover by
the stability axiom.  The compatibility of the restricted family follows
from the functoriality of $F$: if $s_i|_W = s_j|_W$ on the overlap
$W = U_i \times_X U_j$, then restricting further along $W \times_X Y$
preserves this agreement because $F$ respects composition of morphisms.
\end{proof}

\begin{remark}[Extension]
\label{rem:extension}
The dual operation, \emph{extension} (or \emph{pushforward}) along a
morphism, is more subtle: not every local section extends to a larger
region.  The obstruction to extension is precisely the content of
\v{C}ech cohomology, which we develop in Chapter~\ref{ch:cech}.
The distinction between restriction (always possible) and extension
(possible only when cohomological obstructions vanish) is the
mathematical heart of the local-to-global problem.
\end{remark}

\chapter{The Judgment and Its Components}
\label{ch:judgment}

With the language of sites and presheaves in hand, we now define the
central algebraic object of Judgment Geometry: the \emph{judgment}.
A judgment is not merely a logical claim (``$A$ is true'') nor merely
a typed term (``$e : A$'').  It is a structured datum that records
\emph{where} the claim is made, \emph{what} data is being judged,
\emph{what} is claimed about it, \emph{what evidence} supports the
claim, \emph{what obstacles} prevent global validity, \emph{what
cohomological invariants} measure failure to glue, \emph{how good}
the evidence is, and \emph{how the quality assessment was computed}.

We define a judgment as a 4-tuple $\JJ = (\STop, \PStr, \CSpec,
\GEval)$ of structured objects.  Each of these objects has internal
components, and when we flatten the structured objects to their most
visible constituents, we obtain an 8-component view
$(c, \varphi, A, E, O, B, T, \Pi)$ that is useful for concrete
computation.  This chapter introduces the components; Part~II develops
each structured object in full depth.

\section{Definition of the judgment}
\label{sec:judgment:def}

\begin{definition}[Judgment]
\label{def:judgment}
Let $(\mathcal{C}, J)$ be a site and $\mathcal{F}$ a presheaf on
$\mathcal{C}$.  A \emph{judgment} is a 4-tuple
\[
  \JJ = (\STop,\; \PStr,\; \CSpec,\; \GEval)
\]
whose four objects encode the \textbf{Where}, \textbf{What},
\textbf{Why not}, and \textbf{How good} of the claim (developed
fully in Part~II).

Each object has internal components.  When we flatten these to their
most visible constituents, we obtain the \emph{component-wise view}:
\[
  (c,\; \varphi,\; A,\; E,\; O,\; B,\; T,\; \Pi)
\]
where:
\begin{enumerate}[label=(\roman*),itemsep=4pt]
\item $c \in \ob(\mathcal{C})$ is the \textbf{site component}: the
  coordinate (region, locus) in the site to which the judgment
  applies.

\item $\varphi \in \mathcal{F}(c)$ is the \textbf{presheaf component}:
  the data being judged, as a section of the presheaf $\mathcal{F}$
  at coordinate $c$.

\item $A$ is the \textbf{type component}: the specification,
  proposition, or type against which $\varphi$ is judged.  In a
  program verification context, $A$ might be a postcondition; in
  natural language, a grammaticality constraint; in poetry, a metrical
  template; in mathematics, a theorem statement.

\item $E$ is the \textbf{evidence component}: the proof, witness, or
  test result supporting the claim $\varphi : A$.  Evidence is
  \emph{typed}: we write $E : A$ in context $c$ to indicate that $E$
  is a witness for $A$ at coordinate $c$.  Evidence may be a formal
  proof, an SMT solver certificate, a runtime test trace, or an
  expert assertion.

\item $O$ is the \textbf{obstruction component}: an open cover
  $\mathcal{U} = \{U_i \to c\}$ witnessing non-trivial topology.
  The cover $O$ records the decomposition of $c$ into sub-regions
  over which local verification was performed.  If $c$ is atomic
  (admits only the trivial cover), we set $O = \{c \to c\}$.

\item $B \in \check{H}^1(\mathcal{U}, \mathcal{F})$ is the
  \textbf{cohomology class}: the \v{C}ech 1-class measuring failure
  of the local sections to glue over the cover $O$.  If
  $B = 0$, the local verifications compose into a global
  verification.  If $B \neq 0$, the class identifies which overlaps
  fail and how.

\item $T \in \Grade$ is the \textbf{grade}: a quality measure drawn
  from the Grade semiring (\Cref{ch:grade}), recording how well the
  evidence supports the claim.  The grade is \emph{not} a probability;
  it is an element of an algebraic structure that supports
  composition, comparison, and residuation.

\item $\Pi$ is the \textbf{proof trail}: a record of how $T$ was
  computed---which evidence channels were consulted, which rules were
  applied, which compositions were performed.  The proof trail
  provides auditability: any grade can be traced back to its
  evidential sources.
\end{enumerate}

\noindent The components cluster naturally into four pairs---$(c,
\varphi)$, $(A, E)$, $(O, B)$, $(T, \Pi)$---each pair being the
visible face of one of the four structured objects $\STop$, $\PStr$,
$\CSpec$, $\GEval$.
\end{definition}

\begin{remark}[Relation to Martin-L\"of judgments]
\label{rem:mlof}
Martin-L\"of's type theory uses judgments of the form
$\Gamma \vdash a : A$ (three components: context, term, type).
F$^\star$ extends this to $\Gamma \vdash e : M\;a\;\mathit{wp}$ (four
components: context, term, monad, weakest precondition).  The JuGeo
judgment subsumes both: setting $O$ to the trivial cover, $B = 0$,
$T = \gone$, and $\Pi = \emptyset$ recovers a Martin-L\"of judgment
from the remaining triple $(c, \varphi, A)$ together with the
evidence $E$.  The additional components---obstruction, cohomology
class, grade, and proof trail---carry the geometric, diagnostic, and
qualitative information that type theory alone does not track.
\end{remark}


\section{The components in detail}
\label{sec:judgment:components}

\subsection{The site component \texorpdfstring{$c$}{c}}
\label{sec:judgment:site-component}

The coordinate $c$ answers the question: \emph{where does this
judgment apply?}  A judgment without a coordinate is meaningless---it
is a claim about nothing in particular.  By anchoring every judgment
to a coordinate in the site, we gain:

\begin{itemize}
\item \textbf{Locality:} we can restrict a judgment to sub-regions
  and ask whether it holds there.
\item \textbf{Compositionality:} we can compose judgments along shared
  boundaries (overlaps between coordinates).
\item \textbf{Geometric structure:} the covering families of the site
  organize the local-to-global analysis.
\end{itemize}

\subsection{The presheaf component \texorpdfstring{$\varphi$}{φ}}
\label{sec:judgment:presheaf-component}

The section $\varphi \in \mathcal{F}(c)$ is the data being judged.
In different domains:

\begin{itemize}
\item \textbf{Programs:} $\varphi$ is the verification condition at
  coordinate $c$---a logical formula expressing what must hold for the
  code at $c$ to satisfy its specification.
\item \textbf{Natural language:} $\varphi$ is the semantic
  representation at $c$---a $\lambda$-term, discourse referent, or
  feature bundle.
\item \textbf{Poetry:} $\varphi$ is the quality profile at $c$---a
  record of metrical stress, rhyme, imagery, and emotional tone.
\item \textbf{Mathematics:} $\varphi$ is the logical content at
  $c$---the statement of a lemma, the hypothesis set of a proof step.
\end{itemize}

The key property of $\varphi$ is that it transforms covariantly under
restriction: if $f : U \to c$ is a morphism (so $U$ is a sub-region of
$c$), then $\varphi|_U = \mathcal{F}(f)(\varphi) \in \mathcal{F}(U)$
is the restriction of $\varphi$ to $U$.

\subsection{The type component \texorpdfstring{$A$}{A}}
\label{sec:judgment:type-component}

The type $A$ specifies \emph{what it means for $\varphi$ to be
correct} at coordinate $c$.  Following the Curry-Howard correspondence,
$A$ can be viewed simultaneously as:
\begin{itemize}
\item A \emph{type} that evidence must inhabit: $E : A$.
\item A \emph{proposition} that evidence must prove.
\item A \emph{specification} that the artifact must satisfy.
\end{itemize}

Types in Judgment Geometry are classified into five \emph{proposition
kinds}:

\begin{description}[leftmargin=2.5em]
\item[Structural:] properties of shape---type signatures, arity,
  return types, syntactic well-formedness.
\item[Behavioral:] properties of execution---postconditions,
  preconditions, loop invariants, termination.
\item[Relational:] properties relating multiple coordinates---equivalence,
  commutativity, idempotence, refinement.
\item[Resource:] properties about resource usage---purity, effect
  tracking, memory safety.
\item[Semantic:] higher-level properties---liveness, confluence,
  coherence, aesthetic quality.
\end{description}

The proposition kind determines which evidence channels are applicable,
which descent strategies are appropriate, and how the grade is computed.

\subsection{The evidence component \texorpdfstring{$E$}{E}}
\label{sec:judgment:evidence}

Evidence is a \emph{multi-channel} datum: a judgment may be supported
by evidence from several sources simultaneously.  Formally, $E$ is a
multiset of evidence items, each tagged with a \emph{channel}:

\begin{description}[leftmargin=2.5em]
\item[Proof:] a formal derivation in a proof calculus.
\item[Solver:] a certificate from an SMT solver (Z3, CVC5, etc.)
  operating within a decidable fragment.
\item[Runtime:] an execution trace on a declared set of inputs.
\item[Oracle:] an assertion from an external agent (human expert, LLM).
\end{description}

The typing judgment $E : A$ in context $c$ asserts that $E$ is
\emph{admissible} evidence for $A$ at $c$.  Different channels produce
evidence at different confidence levels, and these levels are
tracked by the grade $T$.

\subsection{The obstruction component \texorpdfstring{$O$}{O}}
\label{sec:judgment:obstruction}

The cover $O = \{U_i \to c\}_{i \in I}$ records the decomposition of
the coordinate $c$ into sub-regions over which local analysis was
performed.  The cover serves two purposes:

\begin{enumerate}
\item It specifies the \emph{granularity} of the analysis: a fine cover
  (many small patches) gives more detailed information but requires
  checking more overlaps; a coarse cover (few large patches) is cheaper
  but may miss local inconsistencies.

\item It provides the input to the \v{C}ech complex
  (\Cref{ch:cech}): the cohomology class $B$ is computed relative to
  the cover $O$.  Different covers may yield different cohomology
  groups; the \emph{true} cohomology is the colimit over all covers
  (the direct limit over refinements).
\end{enumerate}

\subsection{The cohomology class \texorpdfstring{$B$}{B}}
\label{sec:judgment:cohomology}

The class $B \in \check{H}^1(\mathcal{U}, \mathcal{F})$ is the central
diagnostic datum of the judgment.  It answers the question:
\emph{do the local verifications compose into a global verification?}

\begin{itemize}
\item If $B = 0$: the local sections glue.  The artifact is globally
  correct (relative to the specification $A$ and the evidence $E$).

\item If $B \neq 0$: the local sections fail to glue.  The class $B$
  identifies exactly which overlaps are inconsistent and how.  Its
  \emph{support} $\mathrm{supp}(B) = \{(i, j) \mid B_{ij} \neq 0\}$
  localizes the failure; its structure suggests repairs.
\end{itemize}

The full development of \v{C}ech cohomology and its role in Judgment
Geometry is the subject of \Cref{ch:cech}.

\subsection{The grade \texorpdfstring{$T$}{T}}
\label{sec:judgment:grade}

The grade $T \in \Grade$ is a quality measure.  Unlike a probability
(which satisfies the Kolmogorov axioms and lives in $[0, 1]$ with
additive structure), the grade lives in a \emph{semiring}
$(\Grade, \opG, \otG, \gzero, \gone)$ with two operations:

\begin{itemize}
\item $\opG$ (competition/alternative): takes the better of two grades.
\item $\otG$ (composition/conjunction): combines grades along a chain,
  typically by taking the minimum (bottleneck).
\end{itemize}

The full development of the Grade semiring is the subject of
\Cref{ch:grade}.

\subsection{The proof trail \texorpdfstring{$\Pi$}{Π}}
\label{sec:judgment:trail}

The proof trail $\Pi$ is an append-only record documenting:
\begin{enumerate}[nosep]
\item Which evidence channels were consulted and in what order.
\item Which composition rules were applied (restriction, transport,
  gluing).
\item Which grade computations were performed (and the intermediate
  values).
\item Whether any promotions or demotions of grade occurred, and their
  justifications.
\end{enumerate}

The proof trail provides \emph{auditability}: any stakeholder can
trace a judgment's grade back to its evidential roots.  In security-
and safety-critical applications, the proof trail is the certification
artifact.


\section{Composition of judgments}
\label{sec:judgment:composition}

Judgments compose along morphisms in the site.  This composition is the
algebraic content of the gluing operation.

\begin{definition}[Restriction of judgments]
\label{def:judgment-restriction}
Let $\JJ = (c, \varphi, A, E, O, B, T, \Pi)$ be a judgment and
$f : U \to c$ a morphism in the site.  The \emph{restriction} of
$\JJ$ along $f$ is
\[
  \JJ|_U = (U,\; \varphi|_U,\; A|_U,\; E|_U,\; f^*O,\; f^*B,\;
  T',\; \Pi \cdot [\text{restrict along } f])
\]
where:
\begin{itemize}[nosep]
\item $\varphi|_U = \mathcal{F}(f)(\varphi)$,
\item $A|_U$ is the specification specialized to $U$,
\item $E|_U$ is the evidence restricted to $U$,
\item $f^*O$ and $f^*B$ are the pulled-back cover and cohomology class,
\item $T' \leq T$ (restriction may \emph{attenuate} the grade but never
  increase it),
\item $\Pi \cdot [\text{restrict along } f]$ appends the restriction
  step to the proof trail.
\end{itemize}
\end{definition}

\begin{definition}[Composition of judgments]
\label{def:judgment-composition}
Let $\JJ_1 = (c_1, \varphi_1, A_1, E_1, O_1, B_1, T_1, \Pi_1)$ and
$\JJ_2 = (c_2, \varphi_2, A_2, E_2, O_2, B_2, T_2, \Pi_2)$ be
judgments whose coordinates share an overlap: there exist morphisms
$g_1 : W \to c_1$ and $g_2 : W \to c_2$ with
$\varphi_1|_W = \varphi_2|_W$.  The \emph{composition}
$\JJ_1 \circ \JJ_2$ is the judgment
\[
  \JJ_1 \circ \JJ_2 = (c_1 \cup c_2,\; \varphi,\; A_1 \wedge A_2,\;
  E_1 \sqcup E_2,\; O,\; B,\; T_1 \otG T_2,\;
  \Pi_1 \cdot \Pi_2 \cdot [\text{compose}])
\]
where $\varphi$ is the glued section (existing precisely when $B = 0$
for the two-element cover $\{c_1, c_2\}$), $O$ is the union cover,
and $B$ is the \v{C}ech class of the composed cover.
\end{definition}

\begin{proposition}[Associativity of composition]
\label{prop:composition-assoc}
Judgment composition is associative: $(\JJ_1 \circ \JJ_2) \circ \JJ_3
= \JJ_1 \circ (\JJ_2 \circ \JJ_3)$ whenever both sides are defined.
\end{proposition}

\begin{proof}
Associativity follows from the associativity of section gluing (which
is a consequence of the sheaf condition), the associativity of $\otG$
in the Grade semiring, and the associativity of evidence union.
\end{proof}


\section{The judgment category}
\label{sec:judgment:category}

Judgments over a fixed site $(\mathcal{C}, J)$ and presheaf
$\mathcal{F}$ organize into a category.

\begin{definition}[Category of judgments]
\label{def:judgment-category}
The \emph{judgment category} $\mathbf{Jud}(\mathcal{C}, J, \mathcal{F})$
has:
\begin{itemize}
\item \textbf{Objects:} judgments
  $\JJ = (c, \varphi, A, E, O, B, T, \Pi)$.

\item \textbf{Morphisms:} a morphism $\JJ_1 \to \JJ_2$ (called a
  \emph{refinement}) consists of:
  \begin{enumerate}[nosep,label=(\alph*)]
  \item A morphism $f : c_1 \to c_2$ in $\mathcal{C}$ (the judgment
    moves to a finer or coarser coordinate).
  \item A refinement of evidence: $E_2$ extends or strengthens $E_1$
    relative to $f$.
  \item A grade inequality: $T_2 \geq T_1$ (refinement improves or
    preserves quality).
  \item Compatibility: $\varphi_2|_{c_1} = \varphi_1$ and
    $A_2|_{c_1}$ implies $A_1$.
  \end{enumerate}
\end{itemize}
\end{definition}

\begin{proposition}[Well-definedness]
\label{prop:judgment-cat-welldef}
$\mathbf{Jud}(\mathcal{C}, J, \mathcal{F})$ is a category: composition
of refinements is associative and has identities.
\end{proposition}

\begin{proof}
The identity refinement at $\JJ$ is $(\id_c, \id_E, T = T, \id_\varphi)$.
Composition of refinements $(f, E_2, T_2, \ldots) \circ (g, E_3, T_3, \ldots)
= (f \circ g, E_3, T_3, \ldots)$ is associative because composition in
$\mathcal{C}$ is associative and $\geq$ on $\Grade$ is transitive.
\end{proof}

\begin{remark}[Enrichment]
\label{rem:enrichment}
The judgment category is naturally \emph{enriched} over $\Grade$: the
``hom-set'' between two judgments carries a grade measuring the
quality of the refinement.  This enrichment reflects the fact that some
refinements are better than others---a refinement supported by a formal
proof is better than one supported only by an oracle assertion.
The enriched perspective becomes central in Part~II when we define
the Graded Evaluation object $\GEval$.
\end{remark}


\section{The four structured objects}
\label{sec:judgment:structured}

\begin{remark}[From components to structure]
\label{rem:components-to-structure}
The component-wise view $(c, \varphi, A, E, O, B, T, \Pi)$
is useful for concrete calculations, but its eight entries are
\emph{not independent}.  They cluster into four natural groups:

\begin{enumerate}
\item \textbf{Geometric:} $c$ and $O$ (the site component and the
  cover) belong to the site and its topology.  Together with the
  presheaf $\mathcal{F}$ (which supplies $\varphi$), they constitute the
  \emph{Semantic Topos} $\STop$.

\item \textbf{Logical:} $A$ and $E$ (the type and the evidence)
  constitute a dependent typing judgment $E : A$ in context $c$.
  Together they form the \emph{Proof Structure} $\PStr$.

\item \textbf{Cohomological:} $O$ and $B$ (the cover and the
  cohomology class) together with the full \v{C}ech complex
  constitute the \emph{Cohomological Spectrum} $\CSpec$.

\item \textbf{Evaluative:} $T$ and $\Pi$ (the grade and the proof
  trail) together with the Grade semiring constitute the
  \emph{Graded Evaluation} $\GEval$.
\end{enumerate}

These four objects are developed in full depth in Part~II.  Each is
\emph{maximally rich}: $\STop$ is a Grothendieck topos, $\PStr$
is a comprehension category, $\CSpec$ is a spectral sequence, and
$\GEval$ is a quantale with evaluation functor.  The component-wise
view is the shadow of these four objects; the 4-tuple
$\JJ = (\STop, \PStr, \CSpec, \GEval)$ is the natural, canonical
formulation.
\end{remark}

\chapter{\v{C}ech Cohomology and Descent}
\label{ch:cech}

This chapter develops the cohomological heart of Judgment Geometry.
The \v{C}ech complex provides a computable invariant that detects
whether locally verified properties compose into globally verified
ones.  When they do not, the cohomology classes identify exactly where
and how the composition fails.

We develop the theory from first principles, without assuming prior
knowledge of algebraic topology.

%% ─── Section 1: Open covers ──────────────────────────────────────

\section{Open covers and refinements}
\label{sec:cech:covers}

\begin{definition}[Open cover]
\label{def:open-cover}
Let $(\mathcal{C}, J)$ be a site and $X \in \ob(\mathcal{C})$.  An
\emph{open cover} of $X$ is a covering family
$\mathcal{U} = \{f_i : U_i \to X\}_{i \in I}$ belonging to $J(X)$.
\end{definition}

\begin{definition}[Overlaps]
\label{def:overlaps}
Given a cover $\mathcal{U} = \{U_i \to X\}$, the \emph{double overlap}
(or \emph{pairwise intersection}) of $U_i$ and $U_j$ is the fiber
product
\[
  U_{ij} = U_i \times_X U_j,
\]
equipped with projection morphisms $\mathrm{pr}_1 : U_{ij} \to U_i$
and $\mathrm{pr}_2 : U_{ij} \to U_j$.  The \emph{triple overlap} is
\[
  U_{ijk} = U_i \times_X U_j \times_X U_k,
\]
and higher overlaps are defined analogously.
\end{definition}

\begin{definition}[Refinement of covers]
\label{def:cover-refinement}
A cover $\mathcal{V} = \{V_\alpha \to X\}$ is a \emph{refinement} of
$\mathcal{U} = \{U_i \to X\}$ if there exists a map
$\lambda : \mathcal{V} \to \mathcal{U}$ (on index sets) and morphisms
$\rho_\alpha : V_\alpha \to U_{\lambda(\alpha)}$ compatible with the
structure maps to $X$.  The collection of all covers of $X$, ordered
by refinement, forms a directed system.
\end{definition}


%% ─── Section 2: The Čech complex ──────────────────────────────────

\section{The \v{C}ech complex}
\label{sec:cech:complex}

\begin{definition}[\v{C}ech cochain groups]
\label{def:cech-cochain}
Let $\mathcal{U} = \{U_i \to X\}_{i \in I}$ be a cover and
$\mathcal{F}$ a presheaf on $(\mathcal{C}, J)$.  The \emph{\v{C}ech
cochain groups} are:
\begin{align}
  \Cech^0(\mathcal{U}, \mathcal{F})
    &= \prod_{i \in I} \mathcal{F}(U_i)
    && \text{(local sections on each patch)}, \\
  \Cech^1(\mathcal{U}, \mathcal{F})
    &= \prod_{i < j} \mathcal{F}(U_{ij})
    && \text{(data on double overlaps)}, \\
  \Cech^2(\mathcal{U}, \mathcal{F})
    &= \prod_{i < j < k} \mathcal{F}(U_{ijk})
    && \text{(data on triple overlaps)}.
\end{align}
More generally, $\Cech^n(\mathcal{U}, \mathcal{F}) =
\prod_{i_0 < \cdots < i_n} \mathcal{F}(U_{i_0 \cdots i_n})$.
\end{definition}

The key to the entire theory is the coboundary map, which measures
\emph{discrepancy} between overlapping local sections.

\begin{definition}[Coboundary maps]
\label{def:coboundary}
The \emph{zeroth coboundary map}
$\delta^0 : \Cech^0(\mathcal{U}, \mathcal{F})
\to \Cech^1(\mathcal{U}, \mathcal{F})$
sends a family of local sections $(s_i)_{i \in I}$ to the family of
discrepancies:
\[
  (\delta^0 s)_{ij}
  = s_i|_{U_{ij}} - s_j|_{U_{ij}}
  \in \mathcal{F}(U_{ij}).
\]
Here the ``subtraction'' is the comparison operation in $\mathcal{F}$:
$(\delta^0 s)_{ij}$ measures how much the restrictions of $s_i$ and $s_j$
to their overlap disagree.  If $\mathcal{F}$ takes values in abelian
groups, this is literal subtraction; in general, it is the natural
comparison in the fiber.

The \emph{first coboundary map}
$\delta^1 : \Cech^1(\mathcal{U}, \mathcal{F})
\to \Cech^2(\mathcal{U}, \mathcal{F})$
sends a family $(\alpha_{ij})$ to:
\[
  (\delta^1 \alpha)_{ijk}
  = \alpha_{jk}|_{U_{ijk}} - \alpha_{ik}|_{U_{ijk}} + \alpha_{ij}|_{U_{ijk}}.
\]
\end{definition}

\begin{proposition}[$\delta^1 \circ \delta^0 = 0$]
\label{prop:delta-squared-zero}
The composition $\delta^1 \circ \delta^0 = 0$: for any family
$(s_i) \in \Cech^0$,
\[
  (\delta^1 \delta^0 s)_{ijk} = 0.
\]
\end{proposition}

\begin{proof}
Compute directly:
\begin{align*}
  (\delta^1 \delta^0 s)_{ijk}
  &= (\delta^0 s)_{jk}|_{U_{ijk}}
   - (\delta^0 s)_{ik}|_{U_{ijk}}
   + (\delta^0 s)_{ij}|_{U_{ijk}} \\
  &= (s_j|_{U_{jk}} - s_k|_{U_{jk}})|_{U_{ijk}}
   - (s_i|_{U_{ik}} - s_k|_{U_{ik}})|_{U_{ijk}}
   + (s_i|_{U_{ij}} - s_j|_{U_{ij}})|_{U_{ijk}} \\
  &= s_j|_{U_{ijk}} - s_k|_{U_{ijk}}
   - s_i|_{U_{ijk}} + s_k|_{U_{ijk}}
   + s_i|_{U_{ijk}} - s_j|_{U_{ijk}} \\
  &= 0. \qedhere
\end{align*}
\end{proof}

The sequence $\Cech^0 \xrightarrow{\delta^0} \Cech^1 \xrightarrow{\delta^1} \Cech^2 \to \cdots$
is therefore a \emph{cochain complex}: the image of each coboundary
map is contained in the kernel of the next.

\begin{remark}[Interpretation of the degrees]
\label{rem:cech-degrees}
\begin{itemize}
\item An element of $\Cech^0$ is a \emph{choice of local data}:
  one section on each patch.
\item An element of $\Cech^1$ is \emph{overlap data}: one datum on
  each pairwise intersection.
\item An element of $\Cech^2$ is \emph{higher compatibility data}:
  one datum on each triple intersection.
\end{itemize}
The coboundary $\delta^0$ measures \emph{pairwise discrepancy}; the
coboundary $\delta^1$ measures \emph{failure of the cocycle condition}
on triple overlaps.
\end{remark}


%% ─── Section 3: Čech cohomology ──────────────────────────────────

\section{\v{C}ech cohomology groups}
\label{sec:cech:cohomology}

\begin{definition}[\v{C}ech cohomology]
\label{def:cech-cohomology}
The \emph{\v{C}ech cohomology groups} of $\mathcal{F}$ with respect to
the cover $\mathcal{U}$ are:
\[
  \check{H}^n(\mathcal{U}, \mathcal{F})
  = \frac{\ker(\delta^n : \Cech^n \to \Cech^{n+1})}
         {\operatorname{im}(\delta^{n-1} : \Cech^{n-1} \to \Cech^n)},
\]
with the convention that $\Cech^{-1} = 0$.  The elements of
$\ker \delta^n$ are called \emph{$n$-cocycles}; the elements of
$\operatorname{im} \delta^{n-1}$ are called \emph{$n$-coboundaries}.
\end{definition}

The first three cohomology groups have particularly clear
interpretations.

\subsection{\texorpdfstring{$\check{H}^0$}{H0}: Global sections}
\label{sec:cech:h0}

\begin{proposition}[$H^0$ is the group of global sections]
\label{prop:h0-global-sections}
$\check{H}^0(\mathcal{U}, \mathcal{F}) = \ker \delta^0$
consists of those families of local sections $(s_i)_{i \in I}$
that agree on all pairwise overlaps:
\[
  s_i|_{U_{ij}} = s_j|_{U_{ij}} \quad \text{for all } i, j.
\]
When $\mathcal{F}$ is a sheaf, such a compatible family glues to a
unique global section $s \in \mathcal{F}(X)$, so
$\check{H}^0(\mathcal{U}, \mathcal{F}) \cong \mathcal{F}(X)$.
\end{proposition}

\begin{proof}
The condition $\delta^0(s) = 0$ means $s_i|_{U_{ij}} - s_j|_{U_{ij}} = 0$
for all $i, j$, which is exactly the matching family condition of
\Cref{def:matching-family}.  By the sheaf condition
(\Cref{def:sheaf}), every matching family has a unique amalgamation,
giving $\check{H}^0 \cong \mathcal{F}(X)$.
\end{proof}

\subsection{\texorpdfstring{$\check{H}^1$}{H1}: Obstructions to gluing}
\label{sec:cech:h1}

This is the most important cohomology group for Judgment Geometry.

\begin{proposition}[$H^1$ measures failure to glue]
\label{prop:h1-obstructions}
A cocycle $\alpha \in \ker \delta^1 \subset \Cech^1$ is a family
$(\alpha_{ij})$ satisfying the \emph{cocycle condition} on triple overlaps:
\[
  \alpha_{ij}|_{U_{ijk}} + \alpha_{jk}|_{U_{ijk}} = \alpha_{ik}|_{U_{ijk}}
  \quad \text{for all } i, j, k.
\]
The cocycle $\alpha$ is a \emph{coboundary} if there exist sections
$t_i \in \mathcal{F}(U_i)$ with $\alpha_{ij} = t_i|_{U_{ij}} - t_j|_{U_{ij}}$,
i.e., $\alpha = \delta^0(t)$.

The quotient $\check{H}^1(\mathcal{U}, \mathcal{F})
= \ker \delta^1 / \operatorname{im} \delta^0$ measures the
\emph{essential} obstructions to gluing: two cocycles that differ by a
coboundary represent the same obstruction, because the coboundary
component can be absorbed by modifying the local sections.

If $\check{H}^1 = 0$, every cocycle is a coboundary, and every
compatible family of local sections (after appropriate modification)
glues to a global section.  If $\check{H}^1 \neq 0$, there exist
genuine obstructions to gluing that cannot be removed by modifying
individual local sections.
\end{proposition}

\subsection{\texorpdfstring{$\check{H}^2$}{H2}: Obstructions to modifying gluings}
\label{sec:cech:h2}

\begin{remark}[$H^2$ and structural obstructions]
\label{rem:h2}
The group $\check{H}^2(\mathcal{U}, \mathcal{F})$ measures obstructions
at a higher level: the failure of attempted gluings to be consistent
across triple overlaps.  In Judgment Geometry, $\check{H}^2 \neq 0$
signals a \emph{structural} obstruction---the cover itself is
inadequate, and the decomposition of the artifact must be redesigned.
This corresponds to architectural refactoring (in programs), paragraph
restructuring (in text), formal restructuring (in poetry), or
reorganization of the proof structure (in mathematics).

In practice, $\check{H}^2$ obstructions are rarer than $\check{H}^1$
obstructions but harder to resolve, because they require changing the
cover rather than merely modifying local sections.
\end{remark}


%% ─── Section 4: Descent ──────────────────────────────────────────

\section{Descent: the local-to-global principle}
\label{sec:cech:descent}

The central theorem of Judgment Geometry connects the vanishing of
\v{C}ech cohomology to the existence of global sections.

\begin{definition}[Descent datum]
\label{def:descent-datum}
A \emph{descent datum} for a presheaf $\mathcal{F}$ relative to a
cover $\mathcal{U} = \{U_i \to X\}$ is a 1-cocycle
$\alpha \in \ker \delta^1 \subset \Cech^1(\mathcal{U}, \mathcal{F})$.
The descent datum is \emph{effective} if $\alpha$ is a coboundary:
$\alpha = \delta^0(s)$ for some $(s_i) \in \Cech^0$.
\end{definition}

\begin{theorem}[Fundamental Theorem of Judgment Geometry]
\label{thm:fundamental}
Let $(\mathcal{C}, J)$ be a site, $\mathcal{F}$ a presheaf on
$\mathcal{C}$, and $X \in \ob(\mathcal{C})$.  Let $A$ be a
specification for the artifact at $X$.  Then:

\begin{enumerate}
\item The artifact at $X$ satisfies $A$ if and only if there exists a
  global section $s \in \mathcal{F}(X)$ witnessing $A$.

\item A global section exists if and only if, for some (equivalently,
  every sufficiently fine) cover $\mathcal{U}$ of $X$, the first
  \v{C}ech cohomology vanishes:
  \[
    \check{H}^1(\mathcal{U}, \mathcal{F}) = 0.
  \]

\item When $\check{H}^1(\mathcal{U}, \mathcal{F}) \neq 0$, each
  non-trivial class $[{\alpha}] \in \check{H}^1$ is an
  \emph{obstruction}: it identifies the overlaps at which local
  verifications fail to compose, and its support
  $\mathrm{supp}([\alpha]) = \{(i,j) \mid \alpha_{ij} \neq 0\}$
  localizes the failure.
\end{enumerate}
\end{theorem}

\begin{proof}
(1)~follows from the definition: a global section is a witness.

(2)~$(\Rightarrow)$: If $s \in \mathcal{F}(X)$ is a global section,
restricting to each $U_i$ gives a compatible family:
$s_i = s|_{U_i}$.  Then $\delta^0(s)_{ij} = s_i|_{U_{ij}} - s_j|_{U_{ij}}
= s|_{U_{ij}} - s|_{U_{ij}} = 0$, so every 1-cocycle is the zero
cocycle (since any cocycle $\alpha$ can be written as $\alpha = \alpha - 0
= \alpha - \delta^0(s)$, and if the presheaf is separated, the coboundary
structure forces $[\alpha] = 0$ in $\check{H}^1$).

$(\Leftarrow)$: If $\check{H}^1 = 0$, every 1-cocycle is a coboundary.
Given local sections $s_i \in \mathcal{F}(U_i)$, the discrepancy
$\delta^0(s)$ is a 1-cocycle (by \Cref{prop:delta-squared-zero}),
hence a coboundary.  This means there exist modified local sections
$s'_i$ with $\delta^0(s') = 0$---a compatible family.  For a sheaf
$\mathcal{F}$, this compatible family has a unique amalgamation:
a global section $s \in \mathcal{F}(X)$.

(3)~follows from the definition of support: $\alpha_{ij} \neq 0$
means the restrictions of $s_i$ and $s_j$ to $U_{ij}$ disagree, and
this disagreement cannot be eliminated by modifying individual sections
(since $[\alpha] \neq 0$).
\end{proof}


%% ─── Section 5: The exact sequence ───────────────────────────────

\section{The \v{C}ech exact sequence}
\label{sec:cech:exact}

The \v{C}ech cohomology groups fit into a fundamental exact sequence
that organizes the local-to-global analysis.

\begin{proposition}[The \v{C}ech exact sequence]
\label{prop:cech-exact}
For any cover $\mathcal{U} = \{U_i \to X\}$ and presheaf $\mathcal{F}$,
there is an exact sequence:
\[
  0 \to \check{H}^0(\mathcal{U}, \mathcal{F})
  \xrightarrow{\;\iota\;}
  \prod_{i} \mathcal{F}(U_i)
  \xrightarrow{\;\delta^0\;}
  \prod_{i<j} \mathcal{F}(U_{ij})
  \xrightarrow{\;\delta^1\;}
  \prod_{i<j<k} \mathcal{F}(U_{ijk})
  \to \cdots
\]
The map $\iota$ sends a global section $s$ to the family of
restrictions $(s|_{U_i})$.  Exactness at $\prod \mathcal{F}(U_i)$
states that a family $(s_i)$ is in the image of $\iota$ if and only if
$\delta^0(s) = 0$---precisely the sheaf condition.
\end{proposition}


\section{Long exact sequences from short exact sequences}
\label{sec:cech:les}

\begin{theorem}[Long exact sequence in \v{C}ech cohomology]
\label{thm:cech-les}
Let $0 \to \mathcal{F}' \to \mathcal{F} \to \mathcal{F}'' \to 0$ be a
short exact sequence of presheaves on $(\mathcal{C}, J)$.  Then for
any cover $\mathcal{U}$ there is a long exact sequence:
\[
\begin{tikzcd}[column sep=small]
  0 \ar[r] &
  \check{H}^0(\mathcal{U}, \mathcal{F}') \ar[r] &
  \check{H}^0(\mathcal{U}, \mathcal{F}) \ar[r] &
  \check{H}^0(\mathcal{U}, \mathcal{F}'') \ar[dll, "\delta"',
    out=-10, in=170, looseness=0.7] \\
  & \check{H}^1(\mathcal{U}, \mathcal{F}') \ar[r] &
  \check{H}^1(\mathcal{U}, \mathcal{F}) \ar[r] &
  \check{H}^1(\mathcal{U}, \mathcal{F}'') \ar[dll, "\delta"',
    out=-10, in=170, looseness=0.7] \\
  & \check{H}^2(\mathcal{U}, \mathcal{F}') \ar[r] &
  \cdots &
\end{tikzcd}
\]
The connecting homomorphism $\delta : \check{H}^n(\mathcal{U}, \mathcal{F}'')
\to \check{H}^{n+1}(\mathcal{U}, \mathcal{F}')$ is the obstruction to
lifting: a cohomology class in $\mathcal{F}''$ lifts to $\mathcal{F}$
if and only if its image under $\delta$ vanishes.
\end{theorem}

\begin{proof}[Proof sketch]
The short exact sequence of presheaves induces a short exact sequence
of cochain complexes:
$0 \to \Cech^\bullet(\mathcal{U}, \mathcal{F}') \to
\Cech^\bullet(\mathcal{U}, \mathcal{F}) \to
\Cech^\bullet(\mathcal{U}, \mathcal{F}'') \to 0$.
The long exact sequence in cohomology then follows from the standard
snake lemma argument in homological algebra.  The connecting
homomorphism $\delta$ is constructed by diagram chasing: given a
cocycle $\alpha'' \in \Cech^n(\mathcal{U}, \mathcal{F}'')$, lift it
to $\alpha \in \Cech^n(\mathcal{U}, \mathcal{F})$, apply $\delta^n$
to get an element of $\Cech^{n+1}(\mathcal{U}, \mathcal{F})$, and
verify that this element lies in the image of
$\Cech^{n+1}(\mathcal{U}, \mathcal{F}')$.  The class of this element
is $\delta[\alpha'']$, and it is independent of the choice of lift.
\end{proof}

\begin{remark}[Operational significance]
\label{rem:les-operational}
The long exact sequence has immediate practical consequences.  If
$\mathcal{F}'$ is the presheaf of structural properties (type
signatures, arity) and $\mathcal{F}''$ is the presheaf of behavioral
properties (postconditions, invariants), then $\mathcal{F} = \mathcal{F}'
\oplus \mathcal{F}''$ is the full judgment presheaf.  The exact
sequence tells us:

\begin{enumerate}[nosep]
\item If structural verification succeeds ($\check{H}^1(\mathcal{F}') = 0$)
  but behavioral verification fails ($\check{H}^1(\mathcal{F}'') \neq 0$),
  then full verification fails:
  $\check{H}^1(\mathcal{F}) \neq 0$.

\item The converse need not hold: behavioral obstructions may cancel
  with structural data via the connecting homomorphism.

\item Structural verification can be performed first (it is cheaper),
  and behavioral verification deferred until structural verification
  succeeds.  The exact sequence guarantees this staging is sound.
\end{enumerate}
\end{remark}


%% ─── Section 6: Examples ─────────────────────────────────────────

\section{Examples of descent and obstruction}
\label{sec:cech:examples}

We now present four worked examples, one from each domain, illustrating
how \v{C}ech cohomology detects and localizes failures of global
coherence.

\begin{example}[Programs: circular imports]
\label{ex:cech:programs}
Consider a Python project with three modules $M_1$, $M_2$, $M_3$:
\begin{itemize}[nosep]
\item $M_1$ imports a function from $M_2$.
\item $M_2$ imports a class from $M_3$.
\item $M_3$ imports a constant from $M_1$.
\end{itemize}
Each module type-checks in isolation (the local sections are valid).
The site has coordinates $\{M_1, M_2, M_3, P\}$ where $P$ is the
project, with covering family $\{M_i \to P\}_{i=1}^{3}$.  The
overlaps $M_{ij} = M_i \cap M_j$ are the shared interfaces (imports).

The \v{C}ech 1-cocycle $\alpha$ is defined by:
\begin{align*}
  \alpha_{12} &= \text{``$M_1$ expects function $f$ from $M_2$''}, \\
  \alpha_{23} &= \text{``$M_2$ expects class $C$ from $M_3$''}, \\
  \alpha_{13} &= \text{``$M_3$ expects constant $k$ from $M_1$''}.
\end{align*}
The cocycle condition on the triple overlap $M_{123}$ requires
$\alpha_{12} + \alpha_{23} = \alpha_{13}$, which fails because the
dependency chain $M_1 \to M_2 \to M_3 \to M_1$ is circular.  Hence
$\check{H}^1 \neq 0$, and the obstruction class is supported on the
triangle $\{(1,2), (2,3), (1,3)\}$, identifying the circular import.
\end{example}

\begin{example}[Natural language: garden-path sentences]
\label{ex:cech:nl}
The sentence ``The horse raced past the barn fell'' has a site with
coordinates for each phrase and the full sentence $S$.  The cover
$\{P_1, P_2\} \to S$ decomposes the sentence into ``the horse raced
past the barn'' ($P_1$, parsed as a relative clause modification of
``horse'') and ``fell'' ($P_2$, the main predicate).

At $P_1$: the local section assigns ``horse'' the semantic role of
patient (the horse that was raced).  At $P_2$: the local section
requires a subject.  On the overlap $P_1 \cap P_2$: $P_1$ has
consumed the subject position (``horse'' is the patient of ``raced''),
but $P_2$ needs it as the agent of ``fell.''  The discrepancy
$\delta^0 \neq 0$: the two local parses make incompatible demands on
the subject.  $\check{H}^1 \neq 0$, localizing the ambiguity to the
phrase--sentence boundary.
\end{example}

\begin{example}[Poetry: failed volta]
\label{ex:cech:poetry}
A Petrarchan sonnet has an octave (lines 1--8) and a sestet
(lines 9--14), connected by the \emph{volta} (the ``turn'' at line 9).
The site has coordinates $\{Q_1, Q_2, V, Q_3, C, P\}$: two quatrains,
the volta, a sestet ``quatrain,'' a couplet, and the poem.  The cover
$\{Q_1, Q_2, V, Q_3, C\} \to P$ decomposes the poem.

Suppose each quatrain develops the theme of loss, but the volta
introduces, without transition, the topic of gardening.  The local
section at $Q_2$ (thematic content: loss) and the local section at
$Q_3$ (thematic content: gardening) disagree on the overlap
$Q_2 \cap Q_3 = V$: the volta should mediate the thematic transition,
but instead introduces a non-sequitur.  The coboundary
$\delta^0 \neq 0$, and $\check{H}^1 \neq 0$, with the obstruction
supported at the volta.
\end{example}

\begin{example}[Mathematics: circular proof]
\label{ex:cech:math}
A proof of Theorem $T$ uses Lemmas $L_1, L_3, L_5$, with
dependencies $L_1 \to L_3 \to L_5 \to L_1$.  The site has coordinates
$\{L_1, L_3, L_5, T\}$ with cover $\{L_i \to T\}$.  Each lemma is
individually valid (local sections exist), but the overlaps encode
the dependency structure.

The 1-cocycle $\alpha$ records the logical dependency on each overlap:
$\alpha_{13}$ records that $L_1$ uses $L_3$; $\alpha_{35}$ records
that $L_3$ uses $L_5$; $\alpha_{15}$ records that $L_5$ uses $L_1$.
The cocycle condition demands $\alpha_{13} + \alpha_{35} = \alpha_{15}$
(transitivity of dependency), but this fails because the dependency
is circular rather than transitive.  $\check{H}^1 \neq 0$; the
obstruction identifies the circular dependency and the repair frontier
suggests breaking the cycle at any one edge.
\end{example}


\section{Descent as verification}
\label{sec:cech:descent-verification}

We summarize the connection between descent and verification.

\begin{theorem}[Descent criterion for verification]
\label{thm:descent-verification}
Let $(\mathcal{C}, J)$ be the semantic site of an artifact, and let
$\mathcal{F}$ be the presheaf of verification conditions.  Checking
$\check{H}^1(\mathcal{U}, \mathcal{F}) = 0$ for a cover $\mathcal{U}$
is equivalent to checking that the artifact is globally correct
relative to the specification.

More precisely, the descent procedure is:
\begin{enumerate}[nosep]
\item \textbf{Local verification:} compute local sections
  $s_i \in \mathcal{F}(U_i)$ for each patch $U_i$.
\item \textbf{Discrepancy computation:} compute the coboundary
  $\delta^0(s)$ to measure pairwise disagreement on overlaps.
\item \textbf{Cohomological classification:}
  \begin{itemize}[nosep]
  \item If $\delta^0(s) = 0$: the local sections agree; glue them into
    a global section.  The artifact is verified.
  \item If $[\delta^0(s)] = 0$ in $\check{H}^1$ but $\delta^0(s) \neq 0$:
    the discrepancy is a coboundary; modify local sections to achieve
    agreement, then glue.
  \item If $[\delta^0(s)] \neq 0$ in $\check{H}^1$: the artifact has
    genuine obstructions.  Report the obstruction classes, their
    supports, and repair frontiers.
  \end{itemize}
\end{enumerate}
\end{theorem}

\begin{remark}
The descent procedure is \emph{algorithmic}: given a finite site (as
in all practical applications), the \v{C}ech complex is finite, and
the cohomology groups are computable.  The computational cost is
$O(|I|^2 \cdot T_{\mathrm{restrict}})$, where $|I|$ is the number of
patches and $T_{\mathrm{restrict}}$ is the cost of a single restriction
operation.  For typical artifacts, $|I|$ is small (2--20 patches), so
the quadratic factor is negligible; the dominant cost is the
verification of individual patches.
\end{remark}

\chapter{The Grade Semiring}
\label{ch:grade}

The preceding chapters developed the geometric (sites, covers) and
cohomological (\v{C}ech complex, descent) layers of Judgment Geometry.
These layers answer the qualitative question: \emph{does the artifact
pass or fail?}  This chapter introduces the quantitative layer: the
\emph{Grade semiring}, which answers the question \emph{how good is
the artifact, along which axes, and how does quality compose?}

The Grade semiring is not an ad hoc scoring system bolted onto the
framework.  It is a principled algebraic structure---a quantale---with
two operations (competition and composition), a partial order, and a
residuation operation that internalizes implication.  The semiring
interacts tightly with the site and the cohomology: grades restrict
along morphisms, compose along chains, and attenuate through descent.


%% ─── Section 1: Definition ────────────────────────────────────────

\section{Definition of the Grade semiring}
\label{sec:grade:def}

\begin{definition}[Grade semiring]
\label{def:grade-semiring}
The \emph{Grade semiring} is a tuple
\[
  \Grade = ([0,1],\; \opG,\; \otG,\; \gzero,\; \gone)
\]
where:
\begin{enumerate}[label=(\roman*),itemsep=3pt]
\item The \emph{carrier} is the closed unit interval $[0, 1]$.

\item $\opG : [0,1] \times [0,1] \to [0,1]$ is the
  \textbf{competition} (or \textbf{alternative}) operation:
  $a \opG b = \max(a, b)$.  Intuitively, $a \opG b$ is the grade of
  ``the better of two alternatives.''

\item $\otG : [0,1] \times [0,1] \to [0,1]$ is the
  \textbf{composition} (or \textbf{conjunction}) operation:
  $a \otG b = \min(a, b)$.  Intuitively, $a \otG b$ is the grade of
  ``both $a$ and $b$ together''---the bottleneck principle, where the
  weakest link determines the overall quality.

\item $\gzero = 0$ is the \textbf{zero} element: the grade of an
  artifact that is completely unsatisfactory.

\item $\gone = 1$ is the \textbf{unit} element: the grade of a
  perfect artifact.
\end{enumerate}
\end{definition}

\begin{proposition}[Semiring axioms]
\label{prop:semiring-axioms}
$(\Grade, \opG, \otG, \gzero, \gone)$ satisfies the semiring axioms:
\begin{enumerate}
\item $(\Grade, \opG, \gzero)$ is a commutative monoid:
  \begin{itemize}[nosep]
  \item Associativity: $(a \opG b) \opG c = a \opG (b \opG c)$.
  \item Commutativity: $a \opG b = b \opG a$.
  \item Identity: $a \opG \gzero = a$.
  \end{itemize}

\item $(\Grade, \otG, \gone)$ is a commutative monoid:
  \begin{itemize}[nosep]
  \item Associativity: $(a \otG b) \otG c = a \otG (b \otG c)$.
  \item Commutativity: $a \otG b = b \otG a$.
  \item Identity: $a \otG \gone = a$.
  \end{itemize}

\item Distributivity: $a \otG (b \opG c) = (a \otG b) \opG (a \otG c)$.

\item Absorption: $a \otG \gzero = \gzero$.
\end{enumerate}
\end{proposition}

\begin{proof}
Each axiom follows from the properties of $\max$ and $\min$ on $[0,1]$.
\begin{enumerate}
\item $\max$ is associative, commutative, and $\max(a, 0) = a$ for
  $a \in [0,1]$.

\item $\min$ is associative, commutative, and $\min(a, 1) = a$ for
  $a \in [0,1]$.

\item $\min(a, \max(b, c)) = \max(\min(a, b), \min(a, c))$.
  This is the distributive law for lattices, which holds in any
  distributive lattice.

\item $\min(a, 0) = 0$.  \qedhere
\end{enumerate}
\end{proof}


%% ─── Section 2: Partial order ─────────────────────────────────────

\section{The partial order and lattice structure}
\label{sec:grade:order}

\begin{proposition}[Induced partial order]
\label{prop:grade-order}
The semiring $\Grade$ carries a natural partial order defined by:
\[
  a \leq b \quad \Longleftrightarrow \quad a \opG b = b.
\]
This coincides with the usual order on $[0, 1]$: $a \leq b$ iff
$\max(a, b) = b$ iff $a \leq b$ in $\mathbb{R}$.
\end{proposition}

\begin{proof}
$a \opG b = \max(a, b) = b$ iff $a \leq b$ in the usual order on
$\mathbb{R}$.
\end{proof}

\begin{proposition}[Bounded distributive lattice]
\label{prop:grade-lattice}
Under the order $\leq$, $\Grade$ is a bounded distributive lattice
with:
\begin{itemize}[nosep]
\item Join: $a \vee b = a \opG b = \max(a, b)$.
\item Meet: $a \wedge b = a \otG b = \min(a, b)$.
\item Bottom: $\bot = \gzero = 0$.
\item Top: $\top = \gone = 1$.
\end{itemize}
Distributivity ($a \wedge (b \vee c) = (a \wedge b) \vee (a \wedge c)$)
holds, as verified in \Cref{prop:semiring-axioms}.
\end{proposition}


%% ─── Section 3: Quantale structure ────────────────────────────────

\section{Quantale structure}
\label{sec:grade:quantale}

The Grade semiring is not merely a lattice; it has the richer structure
of a \emph{quantale}: a complete lattice equipped with an associative
multiplication that distributes over arbitrary joins.

\begin{definition}[Quantale]
\label{def:quantale}
A \emph{quantale} is a triple $(Q, \bigvee, \otimes)$ where:
\begin{enumerate}[label=(\roman*),nosep]
\item $(Q, \leq)$ is a complete lattice (every subset has a supremum).
\item $\otimes : Q \times Q \to Q$ is an associative binary operation.
\item $\otimes$ distributes over arbitrary suprema:
  $a \otimes \bigvee S = \bigvee \{a \otimes s \mid s \in S\}$
  and
  $(\bigvee S) \otimes a = \bigvee \{s \otimes a \mid s \in S\}$.
\end{enumerate}
A quantale is \emph{commutative} if $\otimes$ is commutative, and
\emph{unital} if there exists $e \in Q$ with $e \otimes a = a = a \otimes e$
for all $a$.
\end{definition}

\begin{theorem}[$\Grade$ is a commutative unital quantale]
\label{thm:grade-quantale}
$([0,1], \bigvee, \otG)$ is a commutative unital quantale with
unit $\gone = 1$.
\end{theorem}

\begin{proof}
The interval $[0, 1]$ is a complete lattice under the usual order,
with $\bigvee S = \sup S$ for any $S \subseteq [0, 1]$.  The
operation $\otG = \min$ is associative and commutative with unit $1$.
Distributivity of $\min$ over arbitrary suprema: for any $a \in [0,1]$
and $S \subseteq [0,1]$,
\[
  \min(a, \sup S)
  = \sup\{\min(a, s) \mid s \in S\}.
\]
This holds because $\min(a, -)$ preserves directed suprema (it is
Scott-continuous) and, more strongly, preserves all suprema since $[0,1]$
is a complete Heyting algebra.
\end{proof}


%% ─── Section 4: Residuation ──────────────────────────────────────

\section{Residuation: grade implication}
\label{sec:grade:residuation}

The quantale structure on $\Grade$ provides a natural notion of
\emph{implication} via residuation.

\begin{definition}[Residuation]
\label{def:residuation}
The \emph{residual} (or \emph{grade implication}) is defined by:
\[
  a \multimap b
  = \bigvee \{c \in [0,1] \mid a \otG c \leq b\}
  = \bigvee \{c \mid \min(a, c) \leq b\}.
\]
Intuitively, $a \multimap b$ is the \emph{best grade that, combined
with $a$, still does not exceed $b$}.
\end{definition}

\begin{proposition}[Computation of residuation]
\label{prop:residuation-compute}
For $a, b \in [0, 1]$:
\[
  a \multimap b =
  \begin{cases}
    1 & \text{if } a \leq b, \\
    b & \text{if } a > b.
  \end{cases}
\]
\end{proposition}

\begin{proof}
If $a \leq b$: then $\min(a, c) \leq a \leq b$ for all $c$, so
$\{c \mid \min(a, c) \leq b\} = [0, 1]$, and $\bigvee [0,1] = 1$.

If $a > b$: then $\min(a, c) \leq b$ iff $c \leq b$ (since
$a > b$, the minimum is determined by $c$ when $c \leq b$, and equals
$a > b$ when $c > a$, and equals $c$ when $b < c \leq a$).  More
precisely: if $c \leq b$, then $\min(a, c) = c \leq b$; if $c > b$,
then $\min(a, c) \geq \min(a, c) > b$ when $c > b$ and $c \leq a$ gives
$\min(a,c) = c > b$.  So $\{c \mid \min(a, c) \leq b\} = [0, b]$, and
$\bigvee [0, b] = b$.
\end{proof}

\begin{proposition}[Adjunction property]
\label{prop:residuation-adjunction}
Residuation satisfies the adjunction:
\[
  a \otG c \leq b \quad \Longleftrightarrow \quad c \leq (a \multimap b).
\]
This makes $\Grade$ a \emph{closed monoidal category} (more precisely,
a closed symmetric monoidal poset).
\end{proposition}

\begin{proof}
$(\Rightarrow)$: If $a \otG c \leq b$, then
$c \in \{c' \mid a \otG c' \leq b\}$, so $c \leq \bigvee\{c' \mid a \otG c' \leq b\} = a \multimap b$.

$(\Leftarrow)$: If $c \leq a \multimap b$, then
$a \otG c \leq a \otG (a \multimap b)$.  We claim
$a \otG (a \multimap b) \leq b$.  If $a \leq b$: $a \otG 1 = a \leq b$.
If $a > b$: $a \otG b = \min(a, b) = b$.  In both cases,
$a \otG (a \multimap b) \leq b$, so $a \otG c \leq b$.
\end{proof}


%% ─── Section 5: The evaluation functor ───────────────────────────

\section{The evaluation functor}
\label{sec:grade:evaluation}

The Grade semiring acquires its semantic content through the
\emph{evaluation functor}, which maps derivations (proofs, evidence
chains) to grades.

\begin{definition}[Evaluation functor]
\label{def:evaluation-functor}
Let $\mathbf{Prf}$ be the category whose objects are evidence items
(proofs, solver certificates, runtime traces, oracle assertions) and
whose morphisms are evidence transformations (composition, weakening,
restriction).  The \emph{evaluation functor} is a functor
\[
  \mu : \mathbf{Prf} \to \Grade
\]
(viewing $\Grade$ as a category via its partial order) that assigns to
each evidence item a grade and preserves the ordering: stronger
evidence receives a higher grade.
\end{definition}

\begin{remark}[Channel-dependent evaluation]
\label{rem:channel-evaluation}
In practice, $\mu$ is channel-dependent:
\begin{itemize}[nosep]
\item Formal proofs: $\mu(\text{proof}) = \gone = 1$ (maximal confidence).
\item Solver certificates: $\mu(\text{solver}) \in [0.9, 1]$
  (high confidence, modulated by fragment coverage).
\item Runtime witnesses: $\mu(\text{runtime}) \in [0.5, 0.9]$
  (confidence depends on test coverage).
\item Oracle assertions: $\mu(\text{oracle}) \in [0, 0.5]$
  (low baseline confidence, subject to promotion with justification).
\end{itemize}
These ranges are illustrative; the specific evaluation function is a
parameter of the framework, not a fixed choice.
\end{remark}


%% ─── Section 6: Harmony ──────────────────────────────────────────

\section{Harmony: the bottleneck principle}
\label{sec:grade:harmony}

A central application of the Grade semiring is the \emph{harmony}
functional, which computes the overall quality of an artifact by
composing grades across strata.

\begin{definition}[Harmony]
\label{def:harmony}
Let $\varphi$ be an artifact analyzed across $n$ strata
$s_1, \ldots, s_n$ (e.g., in a program: type safety, behavioral
correctness, performance, documentation; in a poem: meter, rhyme,
imagery, emotion, unity).  The \emph{harmony} of $\varphi$ is:
\[
  h(\varphi) = \bigotimes_{s \in \text{strata}} T_s(\varphi)
  = T_{s_1}(\varphi) \otG T_{s_2}(\varphi) \otG \cdots \otG T_{s_n}(\varphi)
  = \min(T_{s_1}(\varphi), \ldots, T_{s_n}(\varphi)).
\]
This is the \textbf{bottleneck principle}: the overall quality is
determined by the weakest stratum.
\end{definition}

\begin{proposition}[Properties of harmony]
\label{prop:harmony-properties}
\begin{enumerate}
\item $h(\varphi) = \gone$ if and only if $T_s(\varphi) = \gone$ for
  every stratum $s$.  (An artifact is perfect only if every aspect is
  perfect.)

\item $h(\varphi) = \gzero$ if and only if $T_s(\varphi) = \gzero$
  for some stratum $s$.  (A single catastrophic failure in any
  stratum renders the artifact worthless.)

\item $h$ is monotone in each argument: improving any stratum's grade
  cannot decrease the harmony.

\item $h$ is idempotent on uniform grades: if $T_s(\varphi) = t$ for
  all $s$, then $h(\varphi) = t$.
\end{enumerate}
\end{proposition}

\begin{proof}
All four properties follow from the corresponding properties of $\min$
on $[0, 1]$.
\end{proof}


%% ─── Section 7: OT as degenerate case ────────────────────────────

\section{Optimality Theory as a degenerate case}
\label{sec:grade:ot}

Optimality Theory (OT), the dominant framework in generative phonology,
can be recovered as a degenerate case of the Grade semiring.

\begin{proposition}[OT as a special case of $\Grade$]
\label{prop:ot-special-case}
In classical OT, constraints are ranked $C_1 \gg C_2 \gg \cdots \gg C_n$,
and each constraint assigns a binary pass/fail to a candidate.  This is
the special case of $\Grade$ where:
\begin{enumerate}[nosep]
\item The grade set is $\{0, 1\}$ (binary).
\item The strata are the OT constraints $C_1, \ldots, C_n$.
\item The harmony functional $h$ evaluates candidates by the
  lexicographic order induced by the constraint ranking.
\end{enumerate}
\end{proposition}

\begin{proof}[Proof sketch]
In binary $\Grade$, $\otG = \min$ on $\{0, 1\}$ is conjunction
(both must pass) and $\opG = \max$ on $\{0, 1\}$ is disjunction
(at least one must pass).  The OT evaluation of a candidate is
the lexicographically ordered tuple $(C_1(\varphi), \ldots, C_n(\varphi))$,
which corresponds to the graded harmony $h(\varphi) =
C_1(\varphi) \otG \cdots \otG C_n(\varphi)$ in the binary Grade
with the lexicographic tiebreaking convention.

The key difference is that classical OT is restricted to binary grades
and strict ranking, while $\Grade$ allows continuous grades and the
quantale structure provides residuation (``how much violation of
$C_i$ can be tolerated given the grade on $C_j$?''), which OT
cannot express.
\end{proof}

\begin{remark}
\label{rem:ot-generalization}
This embedding of OT into $\Grade$ is significant for linguistic
applications of Judgment Geometry.  It means that the JuGeo framework
can express everything OT can express (ranked constraints with binary
evaluations) while also expressing graded constraints, continuous
quality measures, and the residuation calculus.  The multi-stratum
architecture of the Compositional Theory of Theories (Part~III)
exploits this generalization.
\end{remark}


%% ─── Section 8: The trust algebra ────────────────────────────────

\section{The trust algebra: grades along chains}
\label{sec:grade:trust}

When judgments are composed along chains of evidence---one source
produces evidence that is consumed by another---the grades compose
multiplicatively.

\begin{definition}[Trust composition]
\label{def:trust-composition}
If $A$ and $B$ are judgments with grades $T(A)$ and $T(B)$, and $B$
depends on $A$ (i.e., $B$ uses $A$ as evidence), then the grade of
the composite judgment is:
\[
  T(A \circ B) = T(A) \otG T(B) = \min(T(A), T(B)).
\]
This is the \textbf{conservative principle}: a chain of evidence is
only as strong as its weakest link.
\end{definition}

\begin{proposition}[Properties of trust composition]
\label{prop:trust-properties}
\begin{enumerate}
\item \textbf{Monotonicity:} strengthening any link in the chain
  cannot weaken the composite:
  $T(A) \leq T(A')$ implies $T(A \circ B) \leq T(A' \circ B)$.

\item \textbf{Conservative join:} composing a high-grade judgment
  with a low-grade judgment yields the low grade:
  $\gone \otG t = t$ for any $t$.

\item \textbf{Associativity:} $T((A \circ B) \circ C) = T(A \circ (B \circ C))$.

\item \textbf{No silent promotion:} there is no operation that
  increases the grade of a composite beyond the minimum of its
  components, unless an explicit promotion step is recorded in the
  proof trail $\Pi$.  Any system claiming to promote trust must
  justify the promotion and record it.
\end{enumerate}
\end{proposition}

\begin{proof}
Properties (1)--(3) follow from the corresponding properties of $\min$
and the associativity of $\otG$.  Property (4) is a design invariant
of the framework: promotion is a distinguished operation, not a
consequence of the algebra, and its occurrence is always logged in the
proof trail.
\end{proof}

\begin{remark}[The trust algebra as enrichment]
\label{rem:trust-enrichment}
The trust composition law makes the judgment category
$\mathbf{Jud}(\mathcal{C}, J, \mathcal{F})$ into a category
\emph{enriched} over $\Grade$: the hom-objects are not merely sets but
carry grades, and composition of morphisms respects the Grade
multiplication.  This enrichment is the bridge to the Graded
Evaluation object $\GEval$ of Part~II.
\end{remark}


%% ─── Section 9: Examples ──────────────────────────────────────────

\section{Examples of graded evaluation}
\label{sec:grade:examples}

We close with concrete examples of graded evaluation across the four
domains.

\begin{example}[Programs]
\label{ex:grade:programs}
A program is evaluated across three strata:
\begin{itemize}[nosep]
\item $T_{\text{type}} = 1.0$: the program type-checks (formal proof).
\item $T_{\text{behav}} = 0.85$: the program passes 85\% of behavioral
  test cases (runtime evidence).
\item $T_{\text{perf}} = 0.72$: the program meets 72\% of performance
  benchmarks (runtime evidence).
\end{itemize}
The harmony is $h = \min(1.0, 0.85, 0.72) = 0.72$.  The bottleneck
is performance; improving performance would have the greatest impact
on overall quality.  The residual $T_{\text{perf}} \multimap T_{\text{type}} = 1.0$
tells us that performance quality does not constrain type safety---they are
independent strata.
\end{example}

\begin{example}[Natural language]
\label{ex:grade:nl}
A sentence is evaluated across four strata:
\begin{itemize}[nosep]
\item $T_{\phon} = 0.95$: phonological well-formedness.
\item $T_{\synt} = 1.0$: syntactic grammaticality.
\item $T_{\sem} = 0.80$: semantic coherence.
\item $T_{\prag} = 0.60$: pragmatic felicity.
\end{itemize}
The harmony is $h = \min(0.95, 1.0, 0.80, 0.60) = 0.60$.  The
bottleneck is pragmatic felicity: the sentence is grammatical and
meaningful but contextually odd.  The residual
$T_{\prag} \multimap T_{\synt} = 1.0$ confirms that pragmatic
infelicity does not impugn syntactic grammaticality.
\end{example}

\begin{example}[Poetry]
\label{ex:grade:poetry}
A sonnet is evaluated across five strata:
\begin{itemize}[nosep]
\item $T_{\text{meter}} = 0.90$: metrical regularity.
\item $T_{\text{rhyme}} = 0.95$: rhyme scheme adherence.
\item $T_{\text{imagery}} = 0.85$: vividness and originality of imagery.
\item $T_{\text{emotion}} = 0.70$: emotional resonance.
\item $T_{\text{unity}} = 0.55$: thematic unity across stanzas.
\end{itemize}
The harmony is $h = \min(0.90, 0.95, 0.85, 0.70, 0.55) = 0.55$.
The bottleneck is thematic unity: the individual stanzas are
accomplished but the poem lacks a coherent arc.  This is precisely
the kind of defect detected by the \v{C}ech cohomology of
\Cref{ch:cech}---the local sections (stanza-level quality) are
high, but they fail to glue (the volta fails, or the thematic
thread is broken).
\end{example}

\begin{example}[Mathematics]
\label{ex:grade:math}
A mathematical proof is evaluated across four strata:
\begin{itemize}[nosep]
\item $T_{\text{rigor}} = 1.0$: every step is formally justified.
\item $T_{\text{elegance}} = 0.65$: the proof is correct but
  ungainly (case-bashing rather than conceptual argument).
\item $T_{\text{novelty}} = 0.80$: the result extends known work in
  a non-trivial way.
\item $T_{\text{significance}} = 0.90$: the result has important
  consequences.
\end{itemize}
The harmony is $h = \min(1.0, 0.65, 0.80, 0.90) = 0.65$.  The
bottleneck is elegance.  The proof is unimpeachably correct but fails
to achieve the aesthetic standard expected of a publishable result.
The residual $T_{\text{elegance}} \multimap T_{\text{rigor}} = 1.0$
confirms that lack of elegance does not compromise rigor.
\end{example}

\begin{remark}[The interaction of grade and cohomology]
\label{rem:grade-cohomology-interaction}
The Grade semiring and the \v{C}ech cohomology interact in a
fundamental way.  When $\check{H}^1 \neq 0$, the obstruction class
$B$ depresses the grade: the overall harmony cannot exceed the grade
of the worst overlap.  Conversely, when $\check{H}^1 = 0$ and all
local grades are high, the global grade is high.  This interaction
is formalized in Part~II via the Graded Evaluation object $\GEval$,
which pairs the quantale structure of $\Grade$ with the spectral
sequence of $\CSpec$ to produce a single coherent quality assessment.
\end{remark}


% ── Part II: The Four Objects ─────────────────────────────────
\part{The Judgment Framework}
% Part II: The Judgment Framework
% This file is \include'd from main.tex

\chapter{Design Principles --- Why Four Objects}
\label{ch:motivation}

Every judgment answers a single question:
\begin{quote}
  \emph{Over this space, how well can we construct a global solution,
  and what prevents it?}
\end{quote}
This question decomposes into exactly four irreducible aspects:
\begin{itemize}[nosep]
\item \textbf{Where} --- the context, the space over which the
  judgment is made;
\item \textbf{What} --- the construction, the content being asserted;
\item \textbf{Why not} --- the obstructions that prevent local pieces
  from assembling into a global whole;
\item \textbf{How good} --- the quality of the construction, measured
  in a graded algebraic structure.
\end{itemize}
Each aspect corresponds to a deep, well-studied branch of mathematics:
topos theory, type theory, cohomology, and quantale theory.  This
chapter explains why a judgment naturally has exactly four aspects
and why each demands its own mathematical universe.

\section{Why structured objects, not flat components}

When judgment components are listed as independent entries in a flat
tuple, their intricate dependencies must be imposed as external
axioms.  But the components of a judgment are not independent: the
coordinate~$c$ determines which
propositions~$\varphi$ make sense; the type~$A$ constrains which
terms~$E$ can serve as evidence; the cover encoded in~$O$ determines
which cohomology classes~$B$ can arise; the grade~$T$ summarises the
trail~$\Pi$.

Encoding these dependencies externally is
possible but unnatural.  It is analogous to defining a group as a
set~$G$ with a binary operation~$\mu$, an identity element~$e$, and an
inversion map~$\iota$, then imposing associativity, identity, and
inverse axioms externally.  The definition \emph{works}, but it hides
the fact that the axioms are not arbitrary constraints---they are the
defining properties of a single algebraic concept.

The 4-tuple formulation avoids this: each of the four objects
\emph{internalises} the relevant constraints, just as a group
internalises its axioms.

\section{The first aspect: $(c,\varphi)$ is a presheaf evaluation}
\label{sec:obs-topos}

Let $C$ be the site category whose objects are the coordinates of the
judgment, equipped with a Grothendieck topology~$J$ encoding the
coverage relation.  A proposition~$\varphi$ defined at coordinate~$c$
is a section of a presheaf $\varphi \in \ob(\PSh(C))$ evaluated at the
object~$c$:
\[
  \varphi(c) \;\in\; \Set.
\]
The pair $(c,\varphi)$ is therefore a \emph{pointed presheaf}---an
object of the slice category $\PSh(C)/\mathbf{y}(c)$, where
$\mathbf{y}$ is the Yoneda embedding.

But a pointed presheaf lives inside a much richer world: the topos
$\Sh(C,J)$ of sheaves on the site.  This topos has a subobject
classifier~$\Omega$, an internal Heyting algebra, exponential objects,
a sheafification functor, geometric morphisms to and from other topoi,
and---crucially---an \emph{internal logic} in which the proposition
$\varphi$ can be interpreted.

None of this structure is visible in the pair $(c,\varphi)$ taken
alone.  Yet all of it is \emph{determined} by the site~$(C,J)$ and
participates in the judgment: the internal logic governs which
propositions are well-formed; the sheaf condition controls how local
data glue; the subobject classifier measures truth values.

\begin{remark}
The passage from $(c,\varphi)$ to a full topos is not a gratuitous
abstraction.  Every topos-theoretic concept we introduce in
Chapter~\ref{ch:stopos} will be exercised in the examples of
Part~\ref{part:examples}.
\end{remark}

\section{The second aspect: $(A,E)$ is a typing judgment}
\label{sec:obs-type}

A carrier type~$A$ together with an evidence term~$E$ satisfying
$\Gamma \vdash E : A$ is a \emph{typing judgment} in the sense of
Martin-L\"of type theory.  It is not merely a pair of a set and an
element.  The derivation witnessing~$\Gamma \vdash E : A$ is a tree
whose internal structure records every logical step: variable lookups,
$\lambda$-abstractions, applications, eliminations.

This derivation tree lives inside a rich proof-theoretic universe:
\begin{itemize}[nosep]
\item Dependent function types $\Pi(x:A).\,B(x)$ and dependent pair
  types $\Sigma(x:A).\,B(x)$.
\item Identity types $\mathrm{Id}_A(a,b)$ with path induction and, in
  the homotopical setting, higher identity types.
\item A universe hierarchy
  $\Type_0 : \Type_1 : \Type_2 : \cdots$.
\item $\beta$-reduction (computation) and $\eta$-expansion
  (extensionality), satisfying subject reduction, confluence, and
  strong normalisation.
\item W-types for general inductive definitions.
\end{itemize}
The pair $(A,E)$, extracted from this universe, discards the context,
the derivation, the reduction theory, and the type formers.  Restoring
them gives the \emph{Proof Structure}~$\PStr$.

\section{The third aspect: $(O,B)$ is a cohomological datum}
\label{sec:obs-cohom}

The obligation cover~$O$ and the obstruction class~$B$ encode the
failure of local sections to glue.  In the language of \v{C}ech
cohomology:
\begin{itemize}[nosep]
\item $O = \{U_i \to X\}$ is an open cover of the site object~$X$.
\item $B \in \check{H}^1(\mathcal{U},\mathcal{F})$ is a first
  cohomology class measuring the obstruction to global gluing.
\end{itemize}
The fundamental theorem of Judgment Geometry states:
\[
  P \models \varphi
  \quad\Longleftrightarrow\quad
  \check{H}^1(\mathcal{U},\mathcal{F}_\varphi) = 0.
\]
But $\check{H}^1$ is only the first layer of a much deeper structure.
The full \v{C}ech complex $C^\bullet(\mathcal{U},\mathcal{F})$ yields
cohomology groups in every degree.  The short exact sequence of sheaves
produces a long exact sequence in cohomology.  The Mayer--Vietoris
sequence relates the cohomology of a union to the cohomology of its
parts.  Spectral sequences connect \v{C}ech cohomology to derived
functor cohomology.  Characteristic classes detect higher-order
twisting.

All of this machinery is implicit in the pair $(O,B)$ but invisible in
the component-wise view.  Making it explicit gives the \emph{Cohomological
Spectrum}~$\CSpec$.

\section{The fourth aspect: $(T,\Pi)$ is a graded derivation}
\label{sec:obs-grade}

The trust grade~$T$ and the provenance trail~$\Pi$ record the quality
of a judgment and the history of how that quality was established.  At
the component level, $T$ ranges over a linearly ordered set
\[
  T_{\mathrm{contra}} \prec T_{\mathrm{unver}} \prec T_{\mathrm{copilot}}
  \prec T_{\mathrm{oracle}} \prec T_{\mathrm{runtime}} \prec
  T_{\mathrm{solver}} \prec T_{\mathrm{proof}},
\]
with operations of conservative join, attenuation, promotion, and
demotion.

This is the shadow of a full algebraic structure.  The grade set~$G$,
equipped with two binary operations $\opG$ (join) and $\otG$ (tensor),
a zero~$\gzero$, and a unit~$\gone$, forms an \emph{idempotent
commutative semiring}---in fact, a \emph{commutative unital quantale}.
The quantale carries:
\begin{itemize}[nosep]
\item A complete lattice structure with $\opG = \bigvee$.
\item A residuation operation $\multimap$ right adjoint to $\otG$,
  making the quantale \emph{closed monoidal}.
\item An evaluation functor from the proof category to the quantale,
  making it a \emph{graded monad}.
\item Stability and convergence conditions governing iterated
  re-evaluation.
\end{itemize}
The pair $(T,\Pi)$ records a single grade and a single trail.
Embedding them in the full quantale-enriched structure gives the
\emph{Graded Evaluation}~$\GEval$.

\section{The thesis: four maximally rich objects}

We have identified four pairs and four ambient structures:
\begin{center}
\renewcommand{\arraystretch}{1.3}
\begin{tabular}{@{}cccc@{}}
\textbf{Pair} & \textbf{Ambient structure} & \textbf{Rich object} &
\textbf{Chapter} \\
\hline
$(c,\varphi)$ & elementary topos with site &
$\STop$ & \ref{ch:stopos} \\
$(A,E)$ & proof-relevant type theory &
$\PStr$ & \ref{ch:pstr} \\
$(O,B)$ & derived \v{C}ech cohomology &
$\CSpec$ & \ref{ch:cspec} \\
$(T,\Pi)$ & quantale-enriched evaluation &
$\GEval$ & \ref{ch:geval}
\end{tabular}
\end{center}

A judgment is a compatible quadruple $(\STop,\PStr,\CSpec,\GEval)$.
The compatibility conditions, formulated as a fibration structure on the
total space of judgments, are the subject of
Chapter~\ref{ch:fibration}.

\section{Why ``rich''}

Each of the four objects is \emph{maximally rich} in a precise sense: it
carries all the internal structure that its ambient mathematical
theory provides.  Concretely:

\begin{enumerate}[label=(\roman*),nosep]
\item \textbf{Internal coherence.}  Each object satisfies the axioms of
  its ambient category (topos axioms, type-theoretic rules, cochain
  complex identities, semiring axioms) without external enforcement.
\item \textbf{Own morphisms.}  Each object lives in a category with a
  well-defined notion of morphism: geometric morphisms for~$\STop$,
  context morphisms for~$\PStr$, chain maps for~$\CSpec$,
  grade-preserving maps for~$\GEval$.
\item \textbf{Own logic.}  Each object carries an internal logic:
  the Mitchell--B\'enabou language of the topos, the type-theoretic
  derivation calculus, the long exact sequence in cohomology, the
  residuated implication $\multimap$ of the quantale.
\end{enumerate}

The component-wise view retains only the \emph{data} of each pair; the
4-tuple retains both the data and the \emph{ambient structure}.  Every
theorem can be stated in the 4-tuple language, and new
theorems---invisible from the flattened component
perspective---become accessible.

\begin{remark}[Enrichment is canonical]
\label{rem:canonical}
Given a set of components $(c,\varphi)$, there is a smallest topos
containing them; given $(A,E)$, a minimal type theory; and so on.
This canonicity is made precise in
Chapter~\ref{ch:equivalence}, where we exhibit an adjoint equivalence
between the structured and flattened views.
\end{remark}

\chapter{The Semantic Topos \texorpdfstring{$\STop$}{S}}
\label{ch:stopos}

The first of our four objects encodes \emph{where} a judgment lives and
\emph{what} it asserts.  At the component level, this information
appears as the pair $(c,\varphi)$---a coordinate and a presheaf
section.  We now develop it into a full elementary topos equipped with
a distinguished presheaf and the complete internal logical apparatus
that the topos provides.

\section{Sites and Grothendieck topologies}
\label{sec:sites}

\begin{definition}[Site]
\label{def:site}
A \emph{site} is a pair $(C,J)$ where $C$ is a small category and
$J$ assigns to each object $U \in \ob(C)$ a collection $J(U)$ of
\emph{covering sieves} on~$U$---families of morphisms with codomain
$U$ closed under precomposition---satisfying:
\begin{enumerate}[label=(\alph*),nosep]
\item \textbf{Maximality.}  The maximal sieve
  $\{f : V \to U \mid V \in \ob(C)\}$ is in $J(U)$.
\item \textbf{Stability.}  If $S \in J(U)$ and $g : V \to U$ is any
  morphism, then the pullback sieve
  $g^*(S) = \{h : W \to V \mid g \circ h \in S\}$ is in $J(V)$.
\item \textbf{Transitivity.}  If $S \in J(U)$ and $R$ is a sieve on
  $U$ such that for every $f : V \to U$ in~$S$ the pullback
  $f^*(R) \in J(V)$, then $R \in J(U)$.
\end{enumerate}
\end{definition}

\begin{example}[Program site]
\label{ex:program-site}
Let $P$ be a program.  The \emph{semantic site} $\mathcal{S}_P$ has:
\begin{itemize}[nosep]
\item Objects: semantic coordinates---modules, classes, functions,
  branches, expressions.
\item Morphisms: scope inclusions $f \hookrightarrow g$ when $f$ is
  defined in the scope of~$g$.
\item Topology: a family $\{U_i \to U\}$ covers~$U$ if the $U_i$
  jointly exhaust every execution path through~$U$.
\end{itemize}
\end{example}

\begin{example}[Discourse site]
\label{ex:discourse-site}
A natural-language discourse~$D$ determines a site whose objects are
discourse segments (sentences, clauses, referring expressions) and
whose coverings reflect anaphoric accessibility: $\{s_i\}$ covers a
discourse referent~$x$ if the $s_i$ jointly determine the reference
of~$x$.
\end{example}

\section{Presheaves and the Yoneda embedding}

\begin{definition}[Presheaf category]
For a small category~$C$, the \emph{presheaf category}
$\PSh(C) = [C^{\op}, \Set]$ has:
\begin{itemize}[nosep]
\item Objects: functors $\mathcal{F} : C^{\op} \to \Set$.
\item Morphisms: natural transformations.
\end{itemize}
\end{definition}

The Yoneda embedding $\mathbf{y} : C \hookrightarrow \PSh(C)$ sends
$U \mapsto \Hom_C(-,U)$ and is fully faithful.  The Yoneda lemma gives
$\Hom_{\PSh(C)}(\mathbf{y}(U), \mathcal{F}) \cong \mathcal{F}(U)$.

\begin{definition}[The datum $\varphi$]
\label{def:phi}
The \emph{judgment presheaf} is a presheaf
$\varphi \in \ob(\PSh(C))$.  For each coordinate
$U \in \ob(C)$, the set $\varphi(U)$ records the local data of
the proposition being judged at stage~$U$.  For each inclusion
$f : V \to U$, the restriction $\varphi(f) : \varphi(U) \to
\varphi(V)$ specialises data from a larger to a smaller scope.
\end{definition}

\section{The sheaf condition and sheafification}
\label{sec:sheafification}

\begin{definition}[Sheaf]
A presheaf $\mathcal{F} \in \PSh(C)$ is a \emph{sheaf} for the
topology~$J$ if, for every covering sieve $S \in J(U)$, the natural
map
\[
  \mathcal{F}(U) \;\longrightarrow\;
  \lim_{(V \to U) \in S}\, \mathcal{F}(V)
\]
is an isomorphism.  Equivalently, every compatible family of local
sections has a unique amalgamation.
\end{definition}

\begin{construction}[Plus construction]
\label{con:plus}
For a presheaf $\mathcal{F}$, the \emph{plus construction} yields
$\mathcal{F}^+$ with
\[
  \mathcal{F}^+(U) \;=\;
  \varinjlim_{S \in J(U)}\;
  \mathrm{Match}(S, \mathcal{F}),
\]
where $\mathrm{Match}(S,\mathcal{F})$ is the set of matching families
for the sieve~$S$.  Iterating twice---$\mathcal{F}^{++}$---yields the
\emph{associated sheaf} (sheafification), and the canonical map
$\mathcal{F} \to \mathcal{F}^{++}$ is universal.
\end{construction}

\begin{definition}[Sheafification adjunction]
\label{def:sheafification}
Let $i : \Sh(C,J) \hookrightarrow \PSh(C)$ be the inclusion.  The
\emph{sheafification functor} $a : \PSh(C) \to \Sh(C,J)$ is the left
adjoint of~$i$:
\[
  a \;\dashv\; i,
  \qquad
  \Hom_{\Sh}(a\mathcal{F},\, \mathcal{G})
  \;\cong\;
  \Hom_{\PSh}(\mathcal{F},\, i\mathcal{G}).
\]
The inclusion~$i$ is fully faithful, so the counit
$a \circ i \xrightarrow{\;\sim\;} \id_{\Sh}$ is an isomorphism.
\end{definition}

\section{The subobject classifier and internal logic}
\label{sec:omega}

\begin{definition}[Subobject classifier]
\label{def:subobj}
A \emph{subobject classifier} in a topos~$\mathcal{E}$ is a monomorphism
$\mathrm{true} : 1 \to \Omega$ such that for every monomorphism
$m : S \hookrightarrow X$ there is a unique \emph{characteristic map}
$\chi_m : X \to \Omega$ making the square
\[
\begin{tikzcd}
  S \arrow[r] \arrow[d, hook, "m"'] & 1 \arrow[d, "\mathrm{true}"] \\
  X \arrow[r, "\chi_m"'] & \Omega
\end{tikzcd}
\]
a pullback.
\end{definition}

\begin{theorem}
\label{thm:omega-exists}
In the sheaf topos $\Sh(C,J)$, the subobject classifier exists and is
given by
\[
  \Omega(U) \;=\; \{\text{closed sieves on } U\},
\]
where a sieve~$S$ on~$U$ is \emph{$J$-closed} if, whenever
$f^*(S) \in J(V)$, then $f \in S$.  The map
$\mathrm{true}_U : \{*\} \to \Omega(U)$ picks the maximal sieve.
\end{theorem}

\begin{proposition}[Heyting algebra structure]
\label{prop:heyting}
For each $U \in \ob(C)$, the set $\Omega(U)$ carries a Heyting algebra
structure:
\begin{align*}
  (S_1 \wedge S_2)(U) &= S_1 \cap S_2, \\
  (S_1 \vee S_2)(U) &= j(S_1 \cup S_2), \\
  (S_1 \Rightarrow S_2)(U) &= \{f : V \to U \mid
    \text{for all } g : W \to V,\;
    f \circ g \in S_1 \;\Longrightarrow\;
    f \circ g \in S_2 \},
\end{align*}
where $j$ denotes the closure operation associated to~$J$.  In
particular, $\Omega$ is in general \emph{not} Boolean:
$S \vee \neg S$ need not equal the maximal sieve.
\end{proposition}

\begin{definition}[Internal quantifiers]
\label{def:quantifiers}
For each morphism $f : Y \to X$ in the topos, the pullback functor
$f^* : \mathrm{Sub}(X) \to \mathrm{Sub}(Y)$ on subobject lattices has:
\begin{itemize}[nosep]
\item a left adjoint $\exists_f \dashv f^*$ ---
  the \emph{internal existential} along~$f$;
\item a right adjoint $f^* \dashv \forall_f$ ---
  the \emph{internal universal} along~$f$.
\end{itemize}
These satisfy the Beck--Chevalley condition: for every pullback square,
$\exists$ and $\forall$ commute with pullback.
\end{definition}

\begin{remark}[The Mitchell--B\'enabou language]
\label{rem:mitchell-benabou}
The internal logic of a topos~$\mathcal{E}$ is formalised by the
\emph{Mitchell--B\'enabou language}: terms are generalised elements
(morphisms into objects of~$\mathcal{E}$), formulas are morphisms into
$\Omega$, and the logical connectives $\wedge, \vee, \Rightarrow,
\neg, \forall, \exists$ are the natural transformations described
above.  The semantics is Kripke--Joyal: a formula
$\varphi$ \emph{holds at stage~$U$} if the corresponding sieve is
$J$-dense at~$U$.  This means the topos provides a
\emph{constructive higher-order logic} in which the judgment
presheaf~$\varphi$ can be interpreted without any external logical
framework.
\end{remark}

\section{Internal equality and Lawvere--Tierney topologies}

\begin{definition}[Internal equality]
For each object $X$ in the topos, the \emph{equality predicate}
$=_X : X \times X \to \Omega$ is the characteristic map of the diagonal
$\Delta_X : X \hookrightarrow X \times X$.
\end{definition}

\begin{definition}[Lawvere--Tierney topology]
\label{def:lt-topology}
A \emph{Lawvere--Tierney topology} (or \emph{local operator}) on a
topos~$\mathcal{E}$ is an endomorphism $j : \Omega \to \Omega$
satisfying:
\begin{enumerate}[label=(\alph*),nosep]
\item $j \circ \mathrm{true} = \mathrm{true}$,
\item $j \circ j = j$,
\item $j \circ \wedge = \wedge \circ (j \times j)$.
\end{enumerate}
\end{definition}

\begin{proposition}
Every Lawvere--Tierney topology~$j$ determines a sub-topos
$\Sh_j(\mathcal{E}) \hookrightarrow \mathcal{E}$ whose objects are the
$j$-sheaves, and conversely every sub-topos arises from a unique~$j$.
The operator~$j$ provides a \emph{modal operator} on truth values:
$j(\varphi)$ is the $j$-closure of~$\varphi$, interpretable as
``$\varphi$ holds $j$-locally.''
\end{proposition}

In Judgment Geometry, different Lawvere--Tierney topologies correspond to
different notions of \emph{local truth}---e.g., ``true on all
branches,'' ``true up to syntactic $\alpha$-equivalence,'' or ``true
modulo discourse context.''

\section{Points and stalks}

\begin{definition}[Point of a topos]
\label{def:point}
A \emph{point} of $\Sh(C,J)$ is a geometric morphism
$p : \Set \to \Sh(C,J)$, i.e., an adjoint pair
$p^* \dashv p_*$ with $p^*$ left exact (preserving finite limits).
The \emph{stalk} of a sheaf~$\mathcal{F}$ at~$p$ is the set
$\mathcal{F}_p = p^*(\mathcal{F})$.
\end{definition}

\begin{theorem}[Deligne]
\label{thm:deligne}
If the site $(C,J)$ is \emph{coherent}---every covering sieve
contains a finite sub-cover---then the sheaf topos $\Sh(C,J)$ has
\emph{enough points}: a morphism of sheaves is an isomorphism if and
only if it induces isomorphisms on all stalks.
\end{theorem}

\begin{remark}
Program sites are typically coherent (functions have finitely many
branches), so Deligne's theorem applies: local-to-global arguments via
stalks are valid.  Discourse sites require more care; infinite
anaphoric chains may violate coherence.
\end{remark}

\section{Slice topoi and dependent types}

\begin{theorem}
\label{thm:slice-topos}
For any object $X$ of a topos~$\mathcal{E}$, the slice category
$\mathcal{E}/X$ is again a topos.  Its subobject classifier is
$\Omega_X = \Omega \times X \to X$ (the projection), and the
forgetful functor $\Sigma_X : \mathcal{E}/X \to \mathcal{E}$ has
both a left adjoint $X^* = (-) \times X$ and a right adjoint
$\Pi_X = (-)^X$ (internal hom over~$X$).
\end{theorem}

In Judgment Geometry, the slice topos $\STop/\varphi$ encodes
\emph{dependent judgments}: judgments that depend on the data
of~$\varphi$.  An object of $\STop/\varphi$ is a sheaf $\mathcal{G}$
together with a map $\mathcal{G} \to \varphi$, representing a family of
types indexed by the local sections of~$\varphi$.

\section{Geometric morphisms}
\label{sec:geom-morph}

\begin{definition}[Geometric morphism]
\label{def:geom-morph}
A \emph{geometric morphism} $f : \mathcal{E} \to \mathcal{F}$ between
topoi is an adjoint pair $f^* \dashv f_*$ where the \emph{inverse
image} $f^*$ preserves finite limits.  A geometric morphism is:
\begin{itemize}[nosep]
\item \emph{essential} if $f^*$ has a further left adjoint $f_!$,
  giving $f_! \dashv f^* \dashv f_*$;
\item \emph{surjective} (or a \emph{surjection}) if $f^*$ is faithful;
\item \emph{an inclusion} (or \emph{embedding}) if $f_*$ is fully
  faithful.
\end{itemize}
\end{definition}

\begin{theorem}[Hyperconnected--localic factorisation]
\label{thm:factorisation}
Every geometric morphism $f : \mathcal{E} \to \mathcal{F}$ factors
uniquely as
\[
  \mathcal{E} \xrightarrow{\;f_h\;} \mathcal{G}
  \xrightarrow{\;f_l\;} \mathcal{F},
\]
where $f_h$ is hyperconnected (connected and locally connected) and
$f_l$ is localic (the inverse image is a Heyting algebra homomorphism
on subobject lattices).
\end{theorem}

In practice, geometric morphisms encode \emph{change of site}: if the
program is refactored (renaming modules, splitting functions), the
old and new sites are related by a geometric morphism, and the
factorisation theorem decomposes the refactoring into a ``topological''
part (the localic factor, changing the underlying space) and a
``logical'' part (the hyperconnected factor, changing the fibres).

\section{The complete definition}

We now assemble the full definition of the Semantic Topos.

\begin{definition}[Semantic Topos $\STop$]
\label{def:stop}
A \emph{Semantic Topos} is a tuple
\[
  \STop \;=\;
  \bigl(C,\; J,\; \varphi,\; \Omega,\; \mathrm{Sub},\;
  \wedge,\, \vee,\, \Rightarrow,\,
  \forall,\, \exists,\; {=},\;
  a,\; \mathrm{Pts},\; \tau \bigr)
\]
where:
\begin{enumerate}[label=(\roman*),nosep]
\item $(C,J)$ is a site (Definition~\ref{def:site});
\item $\varphi \in \ob(\PSh(C))$ is the judgment presheaf
  (Definition~\ref{def:phi});
\item $\Omega$ is the subobject classifier of $\Sh(C,J)$
  (Theorem~\ref{thm:omega-exists});
\item $\mathrm{Sub}$ is the subobject functor,
  $\wedge, \vee, \Rightarrow$ are the Heyting operations on~$\Omega$
  (Proposition~\ref{prop:heyting});
\item $\forall, \exists$ are the internal quantifiers
  (Definition~\ref{def:quantifiers});
\item ${=}$ is internal equality via the diagonal;
\item $a : \PSh(C) \to \Sh(C,J)$ is sheafification
  (Definition~\ref{def:sheafification});
\item $\mathrm{Pts}$ is the class of points of $\Sh(C,J)$
  (Definition~\ref{def:point});
\item $\tau$ is a choice of Lawvere--Tierney topology
  $j : \Omega \to \Omega$ (Definition~\ref{def:lt-topology}),
  determining the operative notion of local truth.
\end{enumerate}
\end{definition}

\begin{remark}
The data (iii)--(viii) are \emph{determined} by the site $(C,J)$; they
are not additional choices.  We list them to emphasise that $\STop$
\emph{carries} all this structure, making it available for the internal
logic of judgments.  Only~$\varphi$ (the datum being judged) and
$\tau$ (the chosen notion of local truth) involve genuine choices.
\end{remark}

\begin{definition}[Morphism of Semantic Topoi]
\label{def:stop-morph}
A \emph{morphism} $f : \STop_1 \to \STop_2$ is a geometric morphism
$f = (f^*,f_*) : \Sh(C_1,J_1) \to \Sh(C_2,J_2)$
(Definition~\ref{def:geom-morph}) together with a natural isomorphism
$\alpha : f^*(\varphi_2) \xrightarrow{\;\sim\;} \varphi_1$ that is
compatible with the Lawvere--Tierney topologies:
$f^* \circ j_2 = j_1 \circ f^*$ on~$\Omega_1$.
\end{definition}

We write $\mathbf{STop}$ for the category of Semantic Topoi and their
morphisms.

\section{Examples}
\label{sec:stop-examples}

\begin{example}[Program AST topos]
\label{ex:ast-topos}
For a typed functional program, the site objects are AST nodes
(module, function, let-binding, pattern-match branch, expression),
morphisms are containment, and covering families are exhaustive case
splits.  The presheaf $\varphi$ assigns to each node the type
specification it must satisfy.  The subobject classifier $\Omega$
assigns to each node the lattice of ``locally valid partial
specifications.''  Sheafification forces local specifications to be
globally consistent.
\end{example}

\begin{example}[Natural-language discourse topos]
\label{ex:discourse-topos}
For a discourse, the site objects are Discourse Representation
Structures (DRSs) in the sense of Kamp--Reyle, morphisms are
accessibility relations, and covers are exhaustive sets of accessible
antecedents for a pronoun.  The presheaf $\varphi$ assigns to each DRS
the set of possible referents.  The internal logic is intuitionistic:
$\varphi \vee \neg\varphi$ fails when a pronoun is genuinely ambiguous.
\end{example}

\begin{example}[Poetic form topos]
\label{ex:poetry-topos}
For a poem, site objects are formal positions (line, foot, syllable);
covers require that metrical constraints be checked on all feet within
a line.  The presheaf $\varphi$ assigns to each position the set of
admissible stress patterns.  Sheafification forces local metrical
constraints to be globally consistent, detecting enjambment and
intentional metrical violations.
\end{example}

\begin{example}[Mathematical theory topos]
For a mathematical theory~$T$, the site is the syntactic category
$\mathcal{C}_T$ whose objects are formulas-in-context and whose
morphisms are provable entailments.  The topology is generated by
finite disjunctions.  The presheaf $\varphi$ is the definable set
corresponding to the formula under scrutiny.  Sheafification yields the
classifying topos $\Set[T]$, and geometric morphisms
$\Set \to \Set[T]$ are precisely the models of~$T$.
\end{example}

\chapter{The Proof Structure \texorpdfstring{$\PStr$}{P}}
\label{ch:pstr}

The second object encodes the \emph{evidence} supporting a judgment.
At the component level, the pair $(A,E)$---a type and a term inhabiting
it---records the bare fact that evidence exists.  We now develop the
full proof-relevant type-theoretic universe in which this pair lives:
dependent types, derivation trees, reduction theory,
identity types, and inductive definitions.

\section{Contexts and well-formedness}
\label{sec:contexts}

\begin{definition}[Context]
\label{def:context}
A \emph{context} is a telescope of typed variable declarations
\[
  \Gamma \;=\; (x_1 : A_1,\; x_2 : A_2(x_1),\;
  \ldots,\; x_n : A_n(x_1,\ldots,x_{n-1})),
\]
where each $A_k$ may depend on the preceding variables.  The empty
context is written $\cdot$.
\end{definition}

\begin{definition}[Context well-formedness]
Well-formedness $\vdash \Gamma \;\mathsf{ctx}$ is defined inductively:
\begin{enumerate}[label=(\alph*),nosep]
\item $\vdash \cdot \;\mathsf{ctx}$.
\item If $\vdash \Gamma \;\mathsf{ctx}$ and
  $\Gamma \vdash A : \Type_i$ for some~$i$, then
  $\vdash (\Gamma, x : A) \;\mathsf{ctx}$, provided
  $x \notin \mathrm{dom}(\Gamma)$.
\end{enumerate}
\end{definition}

Standard structural rules govern the manipulation of contexts:

\begin{definition}[Structural rules]
\label{def:structural}
Let $\Gamma$ be a well-formed context.
\begin{itemize}[nosep]
\item \textbf{Weakening.} If $\Gamma \vdash \mathcal{J}$ and
  $\vdash (\Gamma, x:A) \;\mathsf{ctx}$, then
  $\Gamma, x:A \vdash \mathcal{J}$.
\item \textbf{Contraction.} If $\Gamma, x:A, y:A \vdash \mathcal{J}$,
  then $\Gamma, x:A \vdash \mathcal{J}[y := x]$.
\item \textbf{Exchange.} If $x$ does not occur in~$B$ and $y$ does not
  occur in~$A$, then $\Gamma, x:A, y:B \vdash \mathcal{J}$ implies
  $\Gamma, y:B, x:A \vdash \mathcal{J}$.
\end{itemize}
\end{definition}

\section{Types and universe hierarchy}
\label{sec:universes}

\begin{definition}[Universe hierarchy]
\label{def:universe}
The \emph{universe hierarchy} is a countable tower
\[
  \Type_0 \;:\; \Type_1 \;:\; \Type_2 \;:\; \cdots
\]
with a cumulativity rule: if $\Gamma \vdash A : \Type_i$ then
$\Gamma \vdash A : \Type_{i+1}$.  We adopt the Russell convention:
a type $A : \Type_i$ is simultaneously a term of $\Type_i$ and a
type whose terms $a : A$ can be formed.
\end{definition}

\begin{remark}[Universe polymorphism]
In practice, we work \emph{universe-polymorphically}: definitions and
theorems are quantified over universe levels, avoiding the need to fix
a specific $\Type_i$.  This is essential when the type~$A$ in a
judgment may itself be a universe (e.g., $A = \Type_0$).
\end{remark}

\section{Evidence terms and the $\lambda$-calculus}

\begin{definition}[Evidence terms]
\label{def:terms}
The \emph{evidence terms} are generated by the grammar
\[
  E \;::=\; x \;\mid\; \lambda x.\, E \;\mid\; E_1\; E_2
  \;\mid\; (E_1, E_2) \;\mid\; \pi_1\, E \;\mid\; \pi_2\, E
  \;\mid\; \mathrm{let}\; x = E_1 \;\mathrm{in}\; E_2
  \;\mid\; \cdots
\]
where $x$ ranges over variables and ``$\cdots$'' stands for the
constructors and eliminators of the specific type formers (identity
types, W-types, etc.\ introduced below).  Each term is assigned a type
by a derivation of the form $\Gamma \vdash E : A$.
\end{definition}

\begin{definition}[Explicit substitution]
The \emph{substitution} $E[x := N]$ replaces every free occurrence
of~$x$ in~$E$ by~$N$, with the standard capture-avoiding convention.
Substitution satisfies:
\begin{enumerate}[label=(\alph*),nosep]
\item $x[x := N] = N$.
\item $y[x := N] = y$ if $y \neq x$.
\item $(\lambda y.\, E)[x := N] = \lambda y.\, E[x := N]$ if
  $y \neq x$ and $y \notin \mathrm{FV}(N)$ (rename otherwise).
\item $(E_1 \; E_2)[x := N] = E_1[x := N] \; E_2[x := N]$.
\end{enumerate}
\end{definition}

\section{Typing derivations}
\label{sec:derivations}

\begin{definition}[Typing derivation]
\label{def:derivation}
A \emph{typing derivation} $\mathcal{D}$ is a finite tree whose nodes
are typing judgments and whose edges connect premises to conclusions
according to the inference rules of the type theory.  Each node is
labelled by the name of the rule applied.

Formally, $\mathcal{D}$ is an element of the inductively defined set
of derivation trees:
\[
  \mathcal{D} \;=\;
  \frac{\mathcal{D}_1 \quad \cdots \quad \mathcal{D}_k}
  {\Gamma \vdash E : A}\;\; (\text{rule-name}),
\]
where $\mathcal{D}_1, \ldots, \mathcal{D}_k$ are the sub-derivations
for the premises.
\end{definition}

\begin{remark}[Proof relevance]
\label{rem:proof-relevance}
Two terms $E$ and $E'$ may both inhabit type~$A$ in context~$\Gamma$
but have \emph{different} derivation trees.  The derivation is
part of the judgment's evidence: it records not just \emph{that} $A$ is
inhabited, but \emph{how}.  This is the essence of
\emph{proof relevance}: the proof matters, not just its existence.
In Judgment Geometry, the derivation tree~$\mathcal{D}$ is the
structural skeleton of the evidence.
\end{remark}

\section{Reduction theory}
\label{sec:reduction}

\begin{definition}[$\beta$-reduction]
\label{def:beta}
The \emph{$\beta$-reduction} relation $\to_\beta$ is generated by:
\[
  (\lambda x.\, M)\; N \;\to_\beta\; M[x := N].
\]
This is the fundamental \emph{computation} rule: applying a function
to an argument computes by substitution.
\end{definition}

\begin{definition}[$\eta$-expansion]
\label{def:eta}
For function types, the \emph{$\eta$-expansion} rule is:
\[
  f \;\to_\eta\; \lambda x.\, f\; x
  \qquad (x \notin \mathrm{FV}(f)).
\]
This enforces \emph{extensionality}: two functions that agree on all
inputs are equal.  Analogous $\eta$-rules exist for pair types
($p \to_\eta (\pi_1\, p,\, \pi_2\, p)$) and other type formers.
\end{definition}

\begin{theorem}[Subject reduction]
\label{thm:subject-reduction}
If $\Gamma \vdash E : A$ and $E \to_\beta E'$, then
$\Gamma \vdash E' : A$.  That is, $\beta$-reduction preserves typing.
\end{theorem}

\begin{theorem}[Confluence]
\label{thm:confluence}
The $\beta$-reduction relation is confluent (Church--Rosser): if
$E \to_\beta^* E_1$ and $E \to_\beta^* E_2$, then there exists~$E_3$
with $E_1 \to_\beta^* E_3$ and $E_2 \to_\beta^* E_3$.
\end{theorem}

\begin{theorem}[Strong normalisation]
\label{thm:sn}
Every well-typed term in the type-theoretic fragment (without general
recursion) is strongly normalising: every reduction sequence from~$E$
terminates.
\end{theorem}

\section{Dependent type formers}
\label{sec:type-formers}

\begin{definition}[Dependent product]
\label{def:pi}
Given $\Gamma \vdash A : \Type_i$ and
$\Gamma, x : A \vdash B(x) : \Type_j$, the \emph{dependent product}
$\Pi(x : A).\, B(x) : \Type_{\max(i,j)}$ has:
\begin{itemize}[nosep]
\item \textbf{Introduction:}
  $\Gamma \vdash \lambda x.\, e : \Pi(x:A).\, B(x)$
  when $\Gamma, x : A \vdash e : B(x)$.
\item \textbf{Elimination:}
  $\Gamma \vdash f\; a : B(a)$
  when $\Gamma \vdash f : \Pi(x:A).\, B(x)$ and $\Gamma \vdash a : A$.
\item \textbf{Computation:}
  $(\lambda x.\, e)\; a \;=_\beta\; e[x := a]$.
\item \textbf{Uniqueness:}
  $f \;=_\eta\; \lambda x.\, f\; x$.
\end{itemize}
When $B$ does not depend on~$x$, we recover the simple function type
$A \to B$.
\end{definition}

\begin{definition}[Dependent sum]
\label{def:sigma}
Given $\Gamma \vdash A : \Type_i$ and
$\Gamma, x : A \vdash B(x) : \Type_j$, the \emph{dependent sum}
$\Sigma(x : A).\, B(x) : \Type_{\max(i,j)}$ has:
\begin{itemize}[nosep]
\item \textbf{Introduction:}
  $\Gamma \vdash (a, b) : \Sigma(x:A).\, B(x)$
  when $\Gamma \vdash a : A$ and $\Gamma \vdash b : B(a)$.
\item \textbf{Elimination:}
  $\pi_1 : \Sigma(x:A).\, B(x) \to A$ and
  $\pi_2 : (p : \Sigma(x:A).\, B(x)) \to B(\pi_1\, p)$.
\item \textbf{Computation:}
  $\pi_1(a,b) =_\beta a$ and $\pi_2(a,b) =_\beta b$.
\item \textbf{Uniqueness:}
  $p =_\eta (\pi_1\, p,\, \pi_2\, p)$.
\end{itemize}
When $B$ does not depend on~$x$, we recover the product type
$A \times B$.
\end{definition}

\section{Identity types}
\label{sec:identity}

\begin{definition}[Identity type]
\label{def:identity}
For $\Gamma \vdash a, b : A$, the \emph{identity type}
$\mathrm{Id}_A(a,b) : \Type_i$ (where $A : \Type_i$) has:
\begin{itemize}[nosep]
\item \textbf{Introduction:}
  $\mathrm{refl}_a : \mathrm{Id}_A(a,a)$ --- the reflexivity witness.
\item \textbf{Elimination (J-rule):}  Given a type family
  $C : \Pi(x,y:A).\, \mathrm{Id}_A(x,y) \to \Type_j$ and a term
  $d : \Pi(x:A).\, C(x,x,\mathrm{refl}_x)$, there is a term
  \[
    \mathrm{J}(C,d,a,b,p) : C(a,b,p)
  \]
  for every $p : \mathrm{Id}_A(a,b)$, with computation rule
  $\mathrm{J}(C,d,a,a,\mathrm{refl}_a) =_\beta d(a)$.
\end{itemize}
\end{definition}

\begin{remark}[Higher identity types]
The identity type can be iterated: given $p, q : \mathrm{Id}_A(a,b)$,
we can form $\mathrm{Id}_{\mathrm{Id}_A(a,b)}(p,q)$---a type whose
inhabitants are \emph{paths between paths}.  This gives each type~$A$
the structure of an $\infty$-groupoid, connecting type theory to
homotopy theory.
\end{remark}

\begin{remark}[Univalence]
The \emph{univalence axiom} states that the canonical map
$(A =_{\Type} B) \to (A \simeq B)$, sending a path to the induced
equivalence, is itself an equivalence.  Under univalence, equivalent
types are identical as elements of the universe.  We do not require
univalence in the base theory but note that it is compatible with all
our constructions.
\end{remark}

\section{W-types and general induction}
\label{sec:wtypes}

\begin{definition}[W-type]
\label{def:wtype}
Given $\Gamma \vdash A : \Type_i$ and
$\Gamma, x : A \vdash B(x) : \Type_i$, the \emph{W-type}
$W(x : A).\, B(x) : \Type_i$ is the type of well-founded trees with:
\begin{itemize}[nosep]
\item \textbf{Constructor:}
  $\mathrm{sup} : \Pi(a : A).\, (B(a) \to W(x:A).\, B(x))
  \to W(x:A).\, B(x)$.
  Intuitively: a tree consists of a label $a : A$ at the root and a
  $B(a)$-indexed family of subtrees.
\item \textbf{Eliminator:}
  $\mathrm{wrec}$ --- the dependent recursor, which maps out of the
  W-type by providing, for each $a : A$ and each family of subtrees
  $f : B(a) \to W$, a value in the target type, given recursive access
  to the values on subtrees.
\item \textbf{Computation:}
  $\mathrm{wrec}(h, \mathrm{sup}(a,f))
  =_\beta h(a, f, \lambda b.\, \mathrm{wrec}(h, f(b)))$.
\end{itemize}
\end{definition}

\begin{proposition}
W-types suffice to encode all finitary inductive types: natural
numbers ($A = \mathbf{2}$, $B(\mathtt{zero}) = \mathbf{0}$,
$B(\mathtt{succ}) = \mathbf{1}$), lists, ordinal notations, and the
syntax trees of typed $\lambda$-calculi.
\end{proposition}

\section{The proof category}
\label{sec:proof-category}

\begin{construction}[Proof category]
\label{con:proof-category}
Given a context~$\Gamma$, we construct the \emph{proof category}
$\mathbf{Proof}_\Gamma$:
\begin{itemize}[nosep]
\item \textbf{Objects:} types $A$ with $\Gamma \vdash A : \Type_i$
  for some~$i$.
\item \textbf{Morphisms:} $\Hom(A,B) = \{E \mid
  \Gamma \vdash E : A \to B\} / {=_{\beta\eta}}$---terms of function
  type modulo $\beta\eta$-equivalence.
\item \textbf{Composition:} $g \circ f = \lambda x.\, g\;(f\;x)$.
\item \textbf{Identity:} $\id_A = \lambda x.\, x$.
\end{itemize}
\end{construction}

\begin{theorem}
\label{thm:lccc}
The proof category $\mathbf{Proof}_\Gamma$ is a locally cartesian
closed category (LCCC):
\begin{enumerate}[label=(\alph*),nosep]
\item It has a terminal object~$\mathbf{1}$ (the unit type).
\item It has pullbacks (constructed via dependent sums and identity
  types).
\item For every morphism $f : A \to B$, the pullback functor
  $f^* : \mathbf{Proof}_\Gamma / B \to
  \mathbf{Proof}_\Gamma / A$ has a right adjoint $\Pi_f$
  (dependent product along~$f$) and a left adjoint $\Sigma_f$
  (dependent sum along~$f$).
\end{enumerate}
Moreover, there is an equivalence of categories between LCCCs with
additional structure (universes, W-types) and dependent type theories
with the corresponding type formers.
\end{theorem}

\section{The complete definition}

\begin{definition}[Proof Structure $\PStr$]
\label{def:pstr}
A \emph{Proof Structure} is a tuple
\[
  \PStr \;=\;
  (\Gamma,\; A,\; E,\; \mathcal{D},\; {\to_\beta},\;
  {\to_\eta},\; \mathcal{U},\;
  \mathrm{Id},\; \Sigma,\; \Pi,\; W,\; \mathrm{Elim})
\]
where:
\begin{enumerate}[label=(\roman*),nosep]
\item $\Gamma$ is a well-formed context (Definition~\ref{def:context});
\item $A$ is a type with $\Gamma \vdash A : \Type_i$ for some~$i$;
\item $E$ is an evidence term with $\Gamma \vdash E : A$
  (Definition~\ref{def:terms});
\item $\mathcal{D}$ is a derivation tree witnessing
  $\Gamma \vdash E : A$ (Definition~\ref{def:derivation});
\item $\to_\beta$ is the $\beta$-reduction relation
  (Definition~\ref{def:beta});
\item $\to_\eta$ is the $\eta$-expansion relation
  (Definition~\ref{def:eta});
\item $\mathcal{U} = (\Type_i)_{i \in \mathbb{N}}$ is the universe
  hierarchy (Definition~\ref{def:universe});
\item $\mathrm{Id}$ denotes the identity type former
  (Definition~\ref{def:identity});
\item $\Sigma$ denotes the dependent sum type former
  (Definition~\ref{def:sigma});
\item $\Pi$ denotes the dependent product type former
  (Definition~\ref{def:pi});
\item $W$ denotes the W-type former
  (Definition~\ref{def:wtype});
\item $\mathrm{Elim}$ is the collection of eliminators (induction
  principles) for all type formers.
\end{enumerate}
\end{definition}

\begin{definition}[Morphism of Proof Structures]
\label{def:pstr-morph}
A \emph{morphism} $\sigma : \PStr_1 \to \PStr_2$ is a
\emph{context morphism} (substitution)
$\sigma : \Gamma_2 \to \Gamma_1$---a family of terms
$(\sigma(x_k))_{x_k \in \mathrm{dom}(\Gamma_2)}$ with
$\Gamma_1 \vdash \sigma(x_k) : A_k[\sigma]$---such that
$A_2[\sigma] = A_1$ and $E_2[\sigma] =_{\beta\eta} E_1$, and the
derivation~$\mathcal{D}_2$ transports along~$\sigma$ to a derivation
of $\Gamma_1 \vdash E_1 : A_1$.
\end{definition}

We write $\mathbf{PStr}$ for the category of Proof Structures and their
morphisms.

\section{Examples}

\begin{example}[Type-checking derivation]
\label{ex:type-checking}
A type-checker for a dependently typed language, given a term~$E$ and a
type~$A$ in context~$\Gamma$, constructs a Proof Structure: the
derivation tree~$\mathcal{D}$ is the trace of the type-checking
algorithm.  The $\beta$-reduction relation governs normalisation during
type-checking (e.g., reducing $(\lambda x.\,x)\;3$ to~$3$ to
check that $3 : \mathbb{N}$).
\end{example}

\begin{example}[Semantic parse]
\label{ex:semantic-parse}
A Montague-style semantic parse of a sentence produces a $\PStr$:
$\Gamma$ contains typed discourse referents;
$A = \mathrm{Prop}$ (the type of propositions);
$E$ is the $\lambda$-term
representing the sentence's meaning (e.g., $\lambda x.\,
\mathsf{dog}(x) \wedge \mathsf{runs}(x)$);
$\mathcal{D}$ is the derivation tree of the combinatory categorial
grammar.
\end{example}

\begin{example}[Proof of a theorem]
\label{ex:formal-proof}
A formal proof in a proof assistant (Coq, Lean, Agda) is a Proof
Structure.  The context~$\Gamma$ contains the hypotheses; $A$ is the
theorem statement; $E$ is the proof term; $\mathcal{D}$ is the
elaborated derivation; $\to_\beta$ governs reduction during proof
checking.
\end{example}

\chapter{The Cohomological Spectrum \texorpdfstring{$\CSpec$}{H}}
\label{ch:cspec}

The third object encodes the \emph{obstructions} that prevent local
judgments from assembling into a global one.  At the component level,
the pair $(O,B)$ records a cover and a first cohomology class.  We now
develop the full \v{C}ech cohomological apparatus: the cochain
complex in every degree, the long exact sequence, spectral sequences,
and the derived category.

\section{Covers and their combinatorics}
\label{sec:covers}

\begin{definition}[Cover]
\label{def:cover}
Let $X \in \ob(C)$ be an object of the site.  A \emph{cover} of~$X$
is a family $\mathcal{U} = \{U_i \to X\}_{i \in I}$ that generates a
covering sieve in the topology~$J$.  A \emph{refinement}
$\mathcal{V} \to \mathcal{U}$ is a cover $\mathcal{V} = \{V_j \to X\}$
together with a function $\lambda : J' \to I$ and morphisms
$V_j \to U_{\lambda(j)}$ over~$X$.
\end{definition}

\begin{definition}[Nerve]
\label{def:nerve}
The \emph{nerve} $N(\mathcal{U})$ of a cover $\mathcal{U}$ is the
simplicial set with
\[
  N(\mathcal{U})_n \;=\;
  \coprod_{i_0, \ldots, i_n \in I}
  (U_{i_0} \times_X U_{i_1} \times_X \cdots \times_X U_{i_n}),
\]
with face maps given by omitting a factor and degeneracy maps given by
repeating a factor.  When the $U_i$ are open subsets of a space, the
fibre products are intersections:
$U_{i_0} \times_X \cdots \times_X U_{i_n} = U_{i_0} \cap \cdots \cap U_{i_n}$.
\end{definition}

\section{The \v{C}ech complex}
\label{sec:cech-complex}

\begin{definition}[\v{C}ech complex]
\label{def:cech}
Let $\mathcal{F}$ be a presheaf on~$C$ and $\mathcal{U} = \{U_i \to X\}$
a cover.  The \emph{\v{C}ech complex}
$\Cech^\bullet(\mathcal{U}, \mathcal{F})$ is the cochain complex:
\begin{align*}
  \Cech^0(\mathcal{U}, \mathcal{F})
  &= \prod_{i \in I} \mathcal{F}(U_i),
  \\[3pt]
  \Cech^1(\mathcal{U}, \mathcal{F})
  &= \prod_{i < j} \mathcal{F}(U_i \cap U_j),
  \\[3pt]
  \Cech^2(\mathcal{U}, \mathcal{F})
  &= \prod_{i < j < k} \mathcal{F}(U_i \cap U_j \cap U_k),
  \\[3pt]
  \Cech^n(\mathcal{U}, \mathcal{F})
  &= \prod_{i_0 < \cdots < i_n}
  \mathcal{F}(U_{i_0} \cap \cdots \cap U_{i_n}).
\end{align*}
\end{definition}

\begin{definition}[Coboundary maps]
\label{def:coboundary-p2}
The \emph{coboundary} $\delta^n : \Cech^n \to \Cech^{n+1}$ is defined
by alternating restrictions:
\[
  (\delta^n \alpha)_{i_0, \ldots, i_{n+1}}
  \;=\;
  \sum_{k=0}^{n+1} (-1)^k \;
  \alpha_{i_0, \ldots, \widehat{i_k}, \ldots, i_{n+1}}
  \Big|_{U_{i_0} \cap \cdots \cap U_{i_{n+1}}},
\]
where $\widehat{i_k}$ indicates omission. In additive notation, when
$\mathcal{F}$ takes values in abelian groups, the coboundary for $n=0$
is the classical formula
\[
  (\delta^0 s)_{ij} \;=\;
  s_j\big|_{U_i \cap U_j} - s_i\big|_{U_i \cap U_j}.
\]
\end{definition}

\begin{proposition}
\label{prop:d-squared}
$\delta^{n+1} \circ \delta^n = 0$ for all~$n \geq 0$.
\end{proposition}

\begin{proof}
Each term in the double sum
$(\delta^{n+1} \circ \delta^n)(\alpha)$ appears exactly twice with
opposite signs, via the identity
\[
  \sum_{k < l} (-1)^{k+l} + \sum_{l \leq k} (-1)^{k+l+1} \;=\; 0.
  \qedhere
\]
\end{proof}

\section{\v{C}ech cohomology groups}
\label{sec:cech-cohomology}

\begin{definition}[\v{C}ech cohomology]
\label{def:cech-cohom}
The \emph{$n$-th \v{C}ech cohomology group} of $\mathcal{F}$ with
respect to the cover $\mathcal{U}$ is
\[
  \check{H}^n(\mathcal{U}, \mathcal{F})
  \;=\;
  \frac{\ker(\delta^n : \Cech^n \to \Cech^{n+1})}
       {\mathrm{im}(\delta^{n-1} : \Cech^{n-1} \to \Cech^n)},
\]
with the convention $\Cech^{-1} = 0$.
\end{definition}

We now interpret each degree in the context of Judgment Geometry.

\begin{proposition}[Degree-by-degree interpretation]
\label{prop:degree-interp}
\hfill
\begin{enumerate}[label=(\alph*),nosep]
\item $\check{H}^0(\mathcal{U}, \mathcal{F}) = \ker\, \delta^0
  = \Gamma(X, \mathcal{F})$: the space of \emph{global sections},
  i.e., locally defined data that are already globally consistent.
\item $\check{H}^1(\mathcal{U}, \mathcal{F})
  = \ker\, \delta^1 / \mathrm{im}\, \delta^0$: classes of
  \emph{$1$-cocycles modulo coboundaries}.  A non-trivial class is an
  obstruction to gluing---exactly the JuGeo obstruction~$B$.
\item $\check{H}^2(\mathcal{U}, \mathcal{F})$: obstructions to
  \emph{modifying} gluings.  When $H^1 \neq 0$, one may ask whether the
  cocycle can be trivialized by changing the cover; $H^2$ measures
  the secondary obstruction.  This degree illustrates the advantage
  of the full spectrum over isolated components.
\item Higher $\check{H}^n$: higher-order obstructions, each
  controlling the ``deformability'' of lower-degree data.
\end{enumerate}
\end{proposition}

\begin{theorem}[Independence of cover]
\label{thm:cover-independence}
Let $\mathrm{Cov}(X)$ denote the filtered category of covers of~$X$
with refinements as morphisms.  Then the \emph{absolute \v{C}ech
cohomology}
\[
  \check{H}^n(X, \mathcal{F})
  \;=\;
  \varinjlim_{\mathcal{U} \in \mathrm{Cov}(X)}
  \check{H}^n(\mathcal{U}, \mathcal{F})
\]
is independent of the choice of cover: any two covers have a common
refinement, and the refinement maps induce isomorphisms on the
colimit.
\end{theorem}

\section{Long exact sequences}
\label{sec:les}

\begin{theorem}[Long exact sequence in cohomology]
\label{thm:les}
Let $0 \to \mathcal{F}' \to \mathcal{F} \to \mathcal{F}'' \to 0$ be a
short exact sequence of sheaves on~$(C,J)$.  There is a long exact
sequence of abelian groups:
\[
  0 \to H^0(\mathcal{F}') \to H^0(\mathcal{F}) \to H^0(\mathcal{F}'')
  \xrightarrow{\;\partial^0\;}
  H^1(\mathcal{F}') \to H^1(\mathcal{F}) \to H^1(\mathcal{F}'')
  \xrightarrow{\;\partial^1\;} H^2(\mathcal{F}') \to \cdots
\]
The \emph{connecting homomorphisms}~$\partial^n$ are constructed by
diagram chasing: given an $n$-cocycle in $\Cech^n(\mathcal{F}'')$,
lift it to $\Cech^n(\mathcal{F})$, apply the coboundary~$\delta^n$
(which lands in the image of $\mathcal{F}'$), and project to
$\Cech^{n+1}(\mathcal{F}')$.  One verifies that the result is a
cocycle and that its class is independent of the lift.
\end{theorem}

\begin{remark}
In Judgment Geometry, the long exact sequence governs how obstructions
propagate through sub-judgments.  If a judgment decomposes as
$\mathcal{F}' \hookrightarrow \mathcal{F} \twoheadrightarrow
\mathcal{F}''$ (e.g., splitting a program specification into interface
and implementation), then the connecting homomorphism~$\partial$ maps
obstructions in the quotient ($\mathcal{F}''$) to obstructions in the
sub-object ($\mathcal{F}'$).
\end{remark}

\section{Mayer--Vietoris}
\label{sec:mayer-vietoris}

\begin{theorem}[Mayer--Vietoris sequence]
\label{thm:mv}
Let $X = U \cup V$ be a decomposition of the site object into two
sub-objects.  There is a long exact sequence:
\[
  \cdots \to H^n(X,\mathcal{F})
  \xrightarrow{\;\rho\;}
  H^n(U,\mathcal{F}) \oplus H^n(V,\mathcal{F})
  \xrightarrow{\;\delta_{\mathrm{MV}}\;}
  H^n(U \cap V,\mathcal{F})
  \xrightarrow{\;\partial\;}
  H^{n+1}(X,\mathcal{F}) \to \cdots
\]
where $\rho$ is restriction and $\delta_{\mathrm{MV}}$ is the
difference of restrictions.
\end{theorem}

\begin{example}[Modular decomposition]
\label{ex:modular-mv}
A program with two interacting modules $M_1$ and $M_2$ has
$X = M_1 \cup M_2$ with $M_1 \cap M_2$ the shared interface.  The
Mayer--Vietoris sequence gives:
\[
  H^0(M_1) \oplus H^0(M_2) \to H^0(M_1 \cap M_2)
  \xrightarrow{\;\partial\;} H^1(X).
\]
Thus $H^1(X) \neq 0$---a global bug---arises precisely when the
interface specifications are individually satisfied but mutually
inconsistent.
\end{example}

\section{Spectral sequences}
\label{sec:spectral}

\begin{definition}[Spectral sequence]
\label{def:spectral}
A \emph{cohomological spectral sequence} in an abelian category is a
collection $\{E_r^{p,q},\; d_r^{p,q}\}_{r \geq 0}$ where:
\begin{itemize}[nosep]
\item Each $E_r^{p,q}$ is an object (the $(p,q)$-entry on
  page~$r$).
\item Each $d_r^{p,q} : E_r^{p,q} \to E_r^{p+r,\, q-r+1}$ is a
  differential with $d_r \circ d_r = 0$.
\item $E_{r+1}^{p,q} = H^{p,q}(E_r, d_r) =
  \ker\, d_r^{p,q} / \mathrm{im}\, d_r^{p-r,\, q+r-1}$.
\end{itemize}
The spectral sequence \emph{converges} to a graded object
$H^n = \bigoplus_{p+q=n} E_\infty^{p,q}$ if it degenerates at a
finite page: $E_r = E_\infty$ for some~$r$.
\end{definition}

\begin{theorem}[\v{C}ech-to-derived functor spectral sequence]
\label{thm:cech-derived}
For a cover $\mathcal{U}$ of~$X$ and a sheaf $\mathcal{F}$, there is a
spectral sequence with
\[
  E_2^{p,q}
  \;=\;
  \check{H}^p\!\left(\mathcal{U},\;
  \underline{H}^q(\mathcal{F})\right)
  \;\Longrightarrow\;
  H^{p+q}(X, \mathcal{F}),
\]
where $\underline{H}^q(\mathcal{F})$ is the presheaf
$U \mapsto H^q(U, \mathcal{F})$ (derived functor cohomology on open
sub-objects).  If $\mathcal{U}$ is a \emph{Leray cover}---meaning
$H^q(U_{i_0} \cap \cdots \cap U_{i_p}, \mathcal{F}) = 0$ for
$q > 0$---the spectral sequence degenerates at $E_2$ and
\[
  \check{H}^n(\mathcal{U}, \mathcal{F})
  \;\cong\; H^n(X, \mathcal{F}).
\]
\end{theorem}

\begin{remark}
In Judgment Geometry, convergence of the spectral sequence means that
\v{C}ech obstructions (computable from a cover) faithfully capture all
derived obstructions.  The Leray condition identifies ``good''
covers---those for which the combinatorial \v{C}ech computation suffices.
\end{remark}

\section{Obstruction theory}
\label{sec:obstruction}

\begin{construction}[Obstruction tower]
\label{con:obstruction-tower}
Let $X$ be built by attaching cells: $X^{(0)} \subseteq X^{(1)}
\subseteq \cdots \subseteq X$.  Given a sheaf~$\mathcal{F}$ and a
section $s_n \in \Gamma(X^{(n)}, \mathcal{F})$ over the $n$-skeleton,
the obstruction to extending $s_n$ to a section over $X^{(n+1)}$ is a
class
\[
  o_{n+1}(s_n) \;\in\; H^{n+1}(X^{(n+1)},\; X^{(n)};\; \mathcal{F}).
\]
If $o_{n+1}(s_n) = 0$, the section extends; otherwise, $s_n$ must be
modified on the $n$-skeleton (if possible) to kill the obstruction.
\end{construction}

\begin{proposition}[Systematic obstruction theory]
\label{prop:systematic-obstruction}
A global section $s \in \Gamma(X, \mathcal{F})$ exists if and only if
all obstructions $o_1, o_2, o_3, \ldots$ can be successively killed.
More precisely, an inductive construction produces:
\begin{enumerate}[label=(\alph*),nosep]
\item $s_0 \in \Gamma(X^{(0)}, \mathcal{F})$ exists freely
  (sections on a discrete set).
\item $s_n$ extends to $s_{n+1}$ iff $o_{n+1}(s_n) = 0$.
\item If $o_{n+1} \neq 0$ but represents the zero class in the
  colimit $\check{H}^{n+1}(X, \mathcal{F})$, then $s_n$ can be
  refined to kill~$o_{n+1}$.
\end{enumerate}
\end{proposition}

\section{The derived category}
\label{sec:derived}

\begin{definition}[Derived category]
\label{def:derived}
The \emph{bounded derived category} $D^b(\Sh(C,J))$ is obtained from
the category of bounded cochain complexes of sheaves by formally
inverting quasi-isomorphisms (maps inducing isomorphisms on all
cohomology sheaves).  It is a \emph{triangulated category} with
shift functor $[1]$ and distinguished triangles.
\end{definition}

\begin{remark}
The \v{C}ech complex $\Cech^\bullet(\mathcal{U}, \mathcal{F})$ is an
object of~$D^b$.  Its cohomology is $R\Gamma(X, \mathcal{F})$---the
derived global sections.  The advantage of working in $D^b$ is that all
spectral sequence arguments become statements about distinguished
triangles, and the derived Hom
$\mathrm{RHom}(\mathcal{F}, \mathcal{G})$ encodes extension classes
$\mathrm{Ext}^n(\mathcal{F}, \mathcal{G})$ that measure higher
obstructions to morphisms between sheaves.
\end{remark}

\section{Characteristic classes}
\label{sec:char-classes}

In classical topology, characteristic classes attach cohomological
invariants to fibre bundles, measuring how ``twisted'' the bundle is.
In Judgment Geometry, the role of the fibre bundle is played by
the sheaf~$\mathcal{F}$ over the site, and the characteristic classes
measure the complexity of the judgment's non-trivial topology.

\begin{definition}[Characteristic class in JuGeo]
\label{def:char-class}
Let $\mathcal{F}$ be a locally free sheaf of rank~$r$ on the site
$(C,J)$.  The \emph{$k$-th characteristic class}
$c_k(\mathcal{F}) \in H^{2k}(X, \mathbb{Z})$ (in the classical case)
or $c_k(\mathcal{F}) \in \check{H}^k(X, \mathcal{F}^\times)$ (in
the general sheaf-theoretic case) measures the degree-$k$ obstruction
to trivialising~$\mathcal{F}$.
\begin{itemize}[nosep]
\item $c_0 = 1$: the rank is locally constant.
\item $c_1$: the first obstruction to choosing a global frame---in
  JuGeo, whether local evidence can be coherently oriented.
\item Higher $c_k$: successive obstructions to higher-order
  trivialisations.
\end{itemize}
\end{definition}

\section{The complete definition}

\begin{definition}[Cohomological Spectrum $\CSpec$]
\label{def:cspec}
A \emph{Cohomological Spectrum} is a tuple
\[
  \CSpec \;=\;
  \bigl(\mathcal{U},\; \Cech^\bullet,\;
  (\check{H}^n)_{n \geq 0},\; (\delta^n),\;
  \partial,\; \mathrm{LES},\; \mathrm{MV},\;
  E_r^{p,q},\; \mathrm{char},\; D^b\bigr)
\]
where:
\begin{enumerate}[label=(\roman*),nosep]
\item $\mathcal{U}$ is a cover of the site object
  (Definition~\ref{def:cover});
\item $\Cech^\bullet = \Cech^\bullet(\mathcal{U}, \mathcal{F})$ is the
  \v{C}ech cochain complex (Definition~\ref{def:cech});
\item $(\check{H}^n)_{n \geq 0}$ are the \v{C}ech cohomology groups
  (Definition~\ref{def:cech-cohom});
\item $(\delta^n)$ are the coboundary maps
  (Definition~\ref{def:coboundary-p2});
\item $\partial$ denotes the connecting homomorphisms of the long exact
  sequence (Theorem~\ref{thm:les});
\item $\mathrm{LES}$ is the long exact sequence datum
  (Theorem~\ref{thm:les});
\item $\mathrm{MV}$ is the Mayer--Vietoris sequence datum
  (Theorem~\ref{thm:mv});
\item $E_r^{p,q}$ is the \v{C}ech-to-derived spectral sequence
  (Theorem~\ref{thm:cech-derived});
\item $\mathrm{char}$ denotes the characteristic classes
  (Definition~\ref{def:char-class});
\item $D^b = D^b(\Sh(C,J))$ is the bounded derived category
  (Definition~\ref{def:derived}).
\end{enumerate}
\end{definition}

\begin{definition}[Morphism of Cohomological Spectra]
\label{def:cspec-morph}
A \emph{morphism} $f : \CSpec_1 \to \CSpec_2$ consists of:
\begin{enumerate}[label=(\alph*),nosep]
\item a refinement of covers $\lambda : \mathcal{U}_2 \to \mathcal{U}_1$;
\item induced cochain maps
  $\lambda^* : \Cech^\bullet(\mathcal{U}_1, \mathcal{F}_1) \to
  \Cech^\bullet(\mathcal{U}_2, \mathcal{F}_2)$;
\item compatibility with all connecting homomorphisms:
  $\lambda^* \circ \partial_1 = \partial_2 \circ \lambda^*$.
\end{enumerate}
\end{definition}

We write $\mathbf{CSpec}$ for the category of Cohomological Spectra.

\section{Examples across domains}
\label{sec:cspec-examples}

\begin{example}[Program bug as $H^1$ class]
\label{ex:program-bug}
A function \texttt{merge\_sort} with two branches---one for lists of
length $\leq 1$, one for longer lists---has cover
$\mathcal{U} = \{U_{\mathrm{base}},\, U_{\mathrm{rec}}\}$.
A $0$-cochain assigns a sorting invariant to each branch; a $1$-cocycle
on the overlap $U_{\mathrm{base}} \cap U_{\mathrm{rec}}$ records
whether the invariants match at the boundary.
If $H^1 \neq 0$, the local invariants are inconsistent: there is a bug
at the case split.
\end{example}

\begin{example}[NL ambiguity as $H^1$ class]
\label{ex:nl-ambiguity}
In the sentence ``Every farmer who owns a donkey beats it,'' the
discourse site has coordinates for ``every farmer,'' ``a donkey,'' and
``it.''  The presheaf assigns possible referents at each coordinate.
The $1$-cocycle measures whether the pronoun ``it'' can be coherently
resolved across all accessible antecedents.  When
$H^1 \neq 0$, the pronoun is genuinely ambiguous---no single
assignment of referents is globally consistent.
\end{example}

\begin{example}[Poetic tension as $H^1$ class]
\label{ex:poetic-tension}
In a sonnet with an ABAB rhyme scheme, the cover consists of
quatrains.  A $0$-cochain assigns a consistent rhyme pattern to each
quatrain; a $1$-cocycle on the overlap between quatrains~1 and~2
measures whether the shared rhyme (the B-rhyme of quatrain~1 and the
A-rhyme of quatrain~2) is phonologically consistent.  When $H^1 \neq
0$, there is metrical or phonological tension between quatrains---a
feature exploited deliberately in English sonnets.
\end{example}

\begin{example}[Proof gap as $H^1$ class]
\label{ex:proof-gap}
A mathematical proof split into lemmas $L_1, \ldots, L_k$ has a cover
$\mathcal{U} = \{L_i\}$.  Each lemma's proof is a $0$-cochain; the
$1$-cocycle on the overlap $L_i \cap L_{i+1}$ records whether the
conclusion of~$L_i$ matches the hypothesis of~$L_{i+1}$.  When
$H^1 \neq 0$, there is a \emph{proof gap}---the lemmas do not chain
together into a complete proof.
\end{example}

\chapter{The Graded Evaluation \texorpdfstring{$\GEval$}{Q}}
\label{ch:geval}

The fourth object encodes the \emph{quality} of a judgment and the
\emph{history} by which that quality was established.  At the
component level, the pair $(T,\Pi)$ records a single grade and a
provenance trail.  We now develop the full algebraic structure that
governs how grades compose, how they respond to new evidence, and how
they converge under iterated re-evaluation.

\section{The grade semiring}
\label{sec:grade-semiring}

\begin{definition}[Grade semiring]
\label{def:grade-semiring-p2}
The \emph{grade semiring} is the algebraic structure
\[
  G \;=\; ([0,1],\; \opG,\; \otG,\; \gzero,\; \gone)
\]
where $[0,1]$ is the real unit interval and the operations are:
\begin{itemize}[nosep]
\item $\opG$ (join): $a \opG b = \max(a,b)$,
\item $\otG$ (tensor): $a \otG b = a \cdot b$ (multiplication),
\item $\gzero = 0$ (additive identity / worst grade),
\item $\gone = 1$ (multiplicative identity / perfect grade).
\end{itemize}
\end{definition}

\begin{proposition}[Semiring axioms]
\label{prop:semiring}
$(G, \opG, \otG, \gzero, \gone)$ is a commutative semiring:
\begin{enumerate}[label=(\alph*),nosep]
\item $\opG$ is associative, commutative, and $\gzero$ is its identity:
  $a \opG \gzero = a$.
\item $\otG$ is associative, commutative, and $\gone$ is its identity:
  $a \otG \gone = a$.
\item $\otG$ distributes over $\opG$:
  $a \otG (b \opG c) = (a \otG b) \opG (a \otG c)$.
\item $\gzero$ annihilates: $a \otG \gzero = \gzero$.
\end{enumerate}
Moreover, $\opG$ is \emph{idempotent}: $a \opG a = a$.
\end{proposition}

\begin{proof}
All properties follow from the standard properties of $\max$ and
$\cdot$ on $[0,1]$.  Idempotency of~$\opG$ is immediate:
$\max(a,a) = a$.  Distributivity: $a \cdot \max(b,c) = \max(a \cdot b,
a \cdot c)$ since multiplication by $a \geq 0$ preserves order.
\end{proof}

\section{Lattice and quantale structure}
\label{sec:quantale}

\begin{definition}[Partial order]
\label{def:partial-order}
Define $a \leq b$ iff $a \opG b = b$, i.e., $\max(a,b) = b$, i.e.,
$a \leq b$ in the usual order on $[0,1]$.
\end{definition}

\begin{theorem}[Complete lattice]
\label{thm:lattice}
$(G, \leq)$ is a complete lattice.  For any subset
$S \subseteq [0,1]$:
\[
  \bigvee S = \sup S, \qquad
  \bigwedge S = \inf S.
\]
In particular, $\opG = \bigvee$ (binary join), $\gzero = \bot$, and
$\gone = \top$.
\end{theorem}

\begin{theorem}[Quantale]
\label{thm:quantale}
$(G, \otG, \gone, \leq)$ is a \emph{commutative unital quantale}:
$\otG$ distributes over arbitrary joins:
\[
  a \otG \bigvee S \;=\; \bigvee \{a \otG s : s \in S\}.
\]
This follows from the continuity of multiplication on $[0,1]$.
\end{theorem}

\section{Residuation}
\label{sec:residuation}

\begin{definition}[Residuation]
\label{def:residuation-p2}
The \emph{residual} (or \emph{internal hom}) $\multimap$ is the right
adjoint to $\otG$ in the lattice order:
\[
  a \otG c \leq b
  \quad\Longleftrightarrow\quad
  c \leq a \multimap b.
\]
Explicitly:
\[
  a \multimap b \;=\;
  \bigvee\, \{c \in [0,1] : a \otG c \leq b\}
  \;=\;
  \begin{cases}
    1 & \text{if } a \leq b, \\
    b / a & \text{if } a > b > 0, \\
    0 & \text{if } a > 0,\; b = 0.
  \end{cases}
\]
\end{definition}

\begin{theorem}[Closed monoidal structure]
\label{thm:closed-monoidal}
$(G, \otG, \multimap, \gone)$ is a symmetric closed monoidal structure
on the lattice $(G, \leq)$.  The residuation satisfies:
\begin{enumerate}[label=(\alph*),nosep]
\item $\gone \multimap a = a$ for all~$a$.
\item $a \multimap (b \multimap c) = (a \otG b) \multimap c$
  (currying).
\item $a \leq b$ implies $b \multimap c \leq a \multimap c$
  (contravariance in the first argument).
\item $a \leq b$ implies $c \multimap a \leq c \multimap b$
  (covariance in the second argument).
\end{enumerate}
\end{theorem}

\begin{remark}[Interpretation]
The residual $a \multimap b$ answers the question: ``Given that our
current evidence has quality~$a$, what is the best quality~$b$ that
can be achieved?''  If $a \leq b$ (the evidence already meets the
target), the answer is $\gone$ (perfect).  If $a > b$, the answer
degrades proportionally.
\end{remark}

\section{Grade and proof trail}
\label{sec:grade-trail}

\begin{definition}[Grade]
\label{def:grade}
The \emph{grade} of a judgment is an element $T \in G$.  It represents
the current assessed quality of the evidence supporting the judgment.
\end{definition}

\begin{definition}[Proof trail]
\label{def:trail}
The \emph{proof trail} is a finite sequence
\[
  \Pi \;=\; [(T_1, r_1),\; (T_2, r_2),\; \ldots,\; (T_k, r_k)]
\]
where each $T_i \in G$ is a grade and each $r_i$ is a \emph{reason}---
a symbolic tag identifying the source of the assessment (e.g.,
``type-checker,'' ``runtime test,'' ``SMT solver,'' ``human review'').
\end{definition}

\begin{definition}[Trail aggregation]
\label{def:aggregation}
The grade $T$ is determined from the trail~$\Pi$ by one of two
canonical aggregation modes:
\begin{itemize}[nosep]
\item \textbf{Optimistic (join):}
  $T = \bigoplus_{i=1}^{k} T_i = \max(T_1, \ldots, T_k)$ ---
  the best assessment wins.
\item \textbf{Conservative (tensor):}
  $T = \bigotimes_{i=1}^{k} T_i = T_1 \cdot T_2 \cdots T_k$ ---
  all assessments must concur; any low grade drags the aggregate down.
\end{itemize}
The choice of mode is part of the judgment's policy.
\end{definition}

\section{The evaluation functor}
\label{sec:eval-functor}

\begin{definition}[Evaluation functor]
\label{def:eval-functor}
Let $\mathbf{Proof}_\Gamma$ be the proof category
(Construction~\ref{con:proof-category}).
The \emph{evaluation functor}
$\mu : \mathbf{Proof}_\Gamma \to G$ is a function assigning a grade to
each morphism (proof term) such that:
\begin{enumerate}[label=(\alph*),nosep]
\item $\mu(\id_A) = \gone$ for every type~$A$.
\item $\mu(g \circ f) = \mu(g) \otG \mu(f)$ for composable
  morphisms~$f$ and~$g$.
\end{enumerate}
That is, $\mu$ is a \emph{lax monoidal functor} from
$(\mathbf{Proof}_\Gamma, \circ, \id)$ to $(G, \otG, \gone)$.
\end{definition}

\begin{remark}
Condition~(b) says that the grade of a composite proof is the tensor
of the grades of its parts.  Since $a \otG b \leq \min(a,b)$ on
$[0,1]$, composing proofs can only maintain or reduce quality---the
weakest link constrains the chain.
\end{remark}

\section{Stability and convergence}
\label{sec:stability-convergence}

\begin{definition}[Stability]
\label{def:stability}
The evaluation functor~$\mu$ is \emph{$L$-stable} with respect to a
metric $d$ on the presheaf space if
\[
  |T(\varphi) - T(\varphi')| \;\leq\; L \cdot d(\varphi, \varphi')
\]
for all presheaves $\varphi, \varphi'$.  Small perturbations of the
input data produce at most proportionally small changes in the grade.
\end{definition}

\begin{definition}[Convergence]
\label{def:convergence}
Define the \emph{re-evaluation sequence}
\[
  T^{(0)} = T, \qquad
  T^{(n+1)} = \mu\!\left(\mathrm{refine}(E, T^{(n)})\right),
\]
where $\mathrm{refine}(E, T^{(n)})$ re-evaluates the evidence~$E$ in
light of the grade $T^{(n)}$.  The evaluation is \emph{convergent} if
this sequence has a limit $T^* = \lim_{n \to \infty} T^{(n)}$.
\end{definition}

\begin{theorem}[Fixpoint convergence]
\label{thm:fixpoint}
If the re-evaluation map $T \mapsto \mu(\mathrm{refine}(E, T))$ is
monotone and continuous on the complete lattice $(G, \leq)$, then by
the Kleene fixed-point theorem the re-evaluation sequence converges
to the least fixpoint:
\[
  T^* \;=\; \bigvee_{n \geq 0}\, T^{(n)}.
\]
This is the \FMLM{} fixpoint property: iterated re-evaluation
stabilises.
\end{theorem}

\begin{remark}
Convergence is essential in practice: a judgment's grade should not
oscillate indefinitely under re-evaluation.  The monotonicity
condition says that higher grades cannot lead to lower re-evaluated
grades---a natural requirement when refinement improves evidence.
\end{remark}

\section{Entropy and information content}
\label{sec:entropy}

\begin{definition}[Grade entropy]
\label{def:entropy}
The \emph{entropy} of a grade $T \in (0,1)$ is
\[
  \mathcal{E}(T) \;=\;
  -T \ln T \;-\; (1-T) \ln(1-T),
\]
with $\mathcal{E}(0) = \mathcal{E}(1) = 0$ by continuity.
\end{definition}

\begin{proposition}[Entropy properties]
\label{prop:entropy}
\hfill
\begin{enumerate}[label=(\alph*),nosep]
\item $\mathcal{E}(T) \geq 0$ for all $T \in [0,1]$.
\item $\mathcal{E}$ is maximised at $T = 1/2$ (maximal uncertainty).
\item $\mathcal{E}(T) \to 0$ as $T \to 0$ or $T \to 1$ (confident
  grades carry low entropy).
\item $\mathcal{E}$ is concave.
\end{enumerate}
The entropy measures the \emph{information content} of the grade: a
grade near $0$ or $1$ is informative (high confidence), while a grade
near $1/2$ is uninformative (maximal uncertainty).
\end{proposition}

\section{Bayesian update}
\label{sec:bayesian}

\begin{definition}[Grade update]
\label{def:bayes}
When new evidence $e$ arrives with likelihood $\ell(e \mid \varphi)$
and prior grade~$T$, the \emph{updated grade} is
\[
  T' \;=\;
  \frac{T \otG \ell(e \mid \varphi)}
       {T \otG \ell(e \mid \varphi) \;+\;
        (1 - T) \otG \ell(e \mid \neg\varphi)},
\]
which is the standard Bayesian update when $\otG$ is multiplication.
The proof trail extends by one entry: $\Pi' = \Pi \cdot [(T', e)]$.
\end{definition}

\begin{proposition}
\label{prop:bayes-coherence}
Bayesian update is coherent with the quantale structure:
\begin{enumerate}[label=(\alph*),nosep]
\item If $\ell(e \mid \varphi) = \gone$, then $T' \geq T$: perfectly
  supporting evidence can only increase the grade.
\item If $\ell(e \mid \varphi) = \gzero$, then $T' = \gzero$:
  contradicting evidence kills the grade.
\item The update is monotone in both $T$ and $\ell$.
\end{enumerate}
\end{proposition}

\section{The trust algebra}
\label{sec:trust-algebra}

\begin{definition}[Trust composition]
\label{def:trust-composition-p2}
For a chain of judgments $J_1 \circ J_2$ (judgment $J_2$ depends on the
conclusion of~$J_1$), the composite grade is
\[
  T(J_1 \circ J_2) \;=\; T(J_1) \otG T(J_2).
\]
Since $\otG$ is multiplication on $[0,1]$, trust degrades along chains.
\end{definition}

\begin{proposition}[Sub-transitivity]
\label{prop:sub-transitive}
Trust is \emph{sub-transitive} in general:
\[
  T(A \to B) \otG T(B \to C) \;\leq\; T(A \to C).
\]
Equality holds when the intermediate judgment~$B$ introduces no
additional uncertainty (i.e., $T(B \to C)$ already accounts for the
quality of~$B$).  In general, the direct path $A \to C$ may have
higher trust than the indirect path through~$B$, because it avoids
compounding errors.
\end{proposition}

\begin{definition}[Grade polynomials]
\label{def:grade-poly}
A \emph{grade polynomial} is an algebraic expression built from grades
using $\opG$ and $\otG$:
\[
  T \;=\; (T_1 \otG T_2) \opG (T_3 \otG T_4).
\]
Such expressions arise when a judgment has multiple independent proof
paths (joined with~$\opG$), each of which is a chain of sub-judgments
(tensored with~$\otG$).  Evaluation of grade polynomials in the
quantale gives the aggregate grade.
\end{definition}

\section{The Eilenberg--Moore category}
\label{sec:eilenberg-moore}

\begin{definition}[$G$-modules]
\label{def:g-module}
A \emph{$G$-module} is a set $M$ equipped with a $G$-action
$\bullet : G \times M \to M$ satisfying:
\begin{enumerate}[label=(\alph*),nosep]
\item $\gone \bullet m = m$ for all $m \in M$.
\item $(a \otG b) \bullet m = a \bullet (b \bullet m)$.
\item $a \bullet (m \opG_M n) = (a \bullet m) \opG_M (a \bullet n)$,
  where $\opG_M$ is a join operation on~$M$.
\end{enumerate}
The category $G\text{-}\mathbf{Mod}$ of $G$-modules and equivariant
maps is the natural home for graded data: each element of~$M$
represents a datum, and the $G$-action scales its quality.
\end{definition}

\begin{remark}
Graded judgments form a $G$-module: the underlying set is the set of
judgments, and the action $T \bullet \JJ$ re-grades $\JJ$ by
tensoring with~$T$.  Morphisms of $G$-modules are grade-compatible
judgment transformations.
\end{remark}

\section{The complete definition}

\begin{definition}[Graded Evaluation $\GEval$]
\label{def:geval}
A \emph{Graded Evaluation} is a tuple
\[
  \GEval \;=\;
  \bigl(G,\; \opG,\; \otG,\; \gzero,\; \gone,\;
  {\leq},\; T,\; \Pi,\; {\multimap},\;
  \mu,\; \sigma,\; \kappa,\; \mathcal{E},\;
  \mathrm{Bayes}\bigr)
\]
where:
\begin{enumerate}[label=(\roman*),nosep]
\item $(G, \opG, \otG, \gzero, \gone)$ is the grade semiring
  (Definition~\ref{def:grade-semiring-p2});
\item $\leq$ is the induced partial order
  (Definition~\ref{def:partial-order});
\item $T \in G$ is the current grade (Definition~\ref{def:grade});
\item $\Pi$ is the proof trail (Definition~\ref{def:trail});
\item $\multimap$ is the residuation
  (Definition~\ref{def:residuation-p2});
\item $\mu$ is the evaluation functor
  (Definition~\ref{def:eval-functor});
\item $\sigma$ is the stability constant
  (Definition~\ref{def:stability});
\item $\kappa$ is the convergence datum
  (Definition~\ref{def:convergence});
\item $\mathcal{E}$ is the entropy function
  (Definition~\ref{def:entropy});
\item $\mathrm{Bayes}$ is the Bayesian update rule
  (Definition~\ref{def:bayes}).
\end{enumerate}
\end{definition}

\begin{definition}[Morphism of Graded Evaluations]
\label{def:geval-morph}
A \emph{morphism} $f : \GEval_1 \to \GEval_2$ is a semiring
homomorphism $f_G : G_1 \to G_2$ together with a map of trails
$f_\Pi : \Pi_1 \to \Pi_2$ such that:
\begin{enumerate}[label=(\alph*),nosep]
\item $f_G(T_1) = T_2$;
\item $f_G$ preserves the residuation:
  $f_G(a \multimap b) = f_G(a) \multimap f_G(b)$;
\item $f_\Pi$ is order-preserving and respects the evaluation functor:
  $\mu_2 \circ f_\Pi = f_G \circ \mu_1$.
\end{enumerate}
\end{definition}

We write $\mathbf{GEval}$ for the category of Graded Evaluations.

\section{Examples}

\begin{example}[Test coverage]
\label{ex:test-coverage}
A program's test suite produces a Graded Evaluation.  Each test~$t_i$
yields a grade $T_i \in \{0, 1\}$ (pass/fail); the aggregate is
$T = \bigotimes_i T_i$ under conservative aggregation (all tests must
pass) or $T = |\{i : T_i = 1\}| / k$ (coverage fraction) under a
softer policy.  The trail
$\Pi = [(T_1, t_1), \ldots, (T_k, t_k)]$ records which tests
contributed.  The entropy $\mathcal{E}(T)$ is high for $T \approx
0.5$ (half the tests fail), indicating maximal uncertainty about
correctness.
\end{example}

\begin{example}[Semantic adequacy]
\label{ex:semantic-adequacy}
A semantic parse of a sentence receives a grade $T$ from a model that
estimates how well the $\lambda$-term captures the intended meaning.
The trail records: the parser's confidence ($T_1$), human evaluation
($T_2$), and entailment test ($T_3$).  Bayesian update incorporates
each source sequentially, and the entropy decreases as evidence
accumulates.
\end{example}

\begin{example}[Poetic quality]
\label{ex:poetic-quality}
A poem's quality is graded by multiple criteria: meter regularity
($T_1$), rhyme consistency ($T_2$), semantic coherence ($T_3$), and
aesthetic judgement ($T_4$).  The grade polynomial
$T = (T_1 \otG T_2) \opG (T_3 \otG T_4)$ says: the overall quality
is the best of two independent assessments---formal (meter $\otG$
rhyme) and semantic (coherence $\otG$ aesthetics).
\end{example}

\begin{example}[Proof confidence]
\label{ex:proof-confidence}
A formally verified theorem has $T = \gone$ by a single trail entry
$(1, \text{proof-checker})$.  A theorem verified by an SMT solver with
a timeout has $T = 0.95$ with reason ``solver-with-timeout.''  If later
a formal proof is produced, the trail grows by $(\gone,
\text{proof-checker})$ and the optimistic aggregate becomes $T = \gone$.
\end{example}

\chapter{The Judgment Fibration}
\label{ch:fibration}

The four objects $\STop$, $\PStr$, $\CSpec$, $\GEval$ do not stand
independently: they are bound together by compatibility conditions that
express how the semantic, proof-theoretic, cohomological, and
evaluative aspects of a judgment constrain one another.  This chapter
formulates these constraints as a \emph{fibration} over the category
of Semantic Topoi, making precise the idea that the other three objects
are ``fibred'' over the site.

\section{Compatibility conditions}
\label{sec:compatibility}

\begin{definition}[Compatible quadruple]
\label{def:compatible}
A \emph{compatible quadruple} $(\STop, \PStr, \CSpec, \GEval)$
consists of one object from each of the four categories, subject to the
following conditions:

\medskip
\noindent\textbf{(C1) Topos determines context.}
The context $\Gamma$ of the Proof Structure is drawn from the internal
language of the Semantic Topos:
\begin{itemize}[nosep]
\item Each variable declaration $x_i : A_i$ in~$\Gamma$ corresponds to
  a generalised element $x_i : 1 \to \varphi_i$ in the
  Mitchell--B\'enabou language of~$\STop$.
\item The type $A$ is an object of the internal type theory of the
  topos.
\item Well-formedness of~$\Gamma$ is equivalent to well-definedness
  of the corresponding string of internal terms.
\end{itemize}

\medskip
\noindent\textbf{(C2) Proof determines cohomology.}
The cover $\mathcal{U}$ of the Cohomological Spectrum arises from
the proof search for~$E : A$:
\begin{itemize}[nosep]
\item The cover $\mathcal{U} = \{U_i\}$ consists of the sub-goals into
  which the proof obligation $\Gamma \vdash ? : A$ is decomposed.
\item A $1$-cocycle in $\check{H}^1(\mathcal{U}, \mathcal{F})$ records
  a failure to glue sub-proofs: the boundary
  $(\delta^0 E)_{ij} = E_i|_{U_{ij}} - E_j|_{U_{ij}}$ does not
  vanish.
\item In particular: $\check{H}^1 = 0$ if and only if the sub-proofs
  glue into a global proof $E$, recovering the fundamental theorem of
  Judgment Geometry.
\end{itemize}

\medskip
\noindent\textbf{(C3) Cohomology constrains grade.}
The grade~$T$ of the Graded Evaluation is bounded by the cohomological
data:
\begin{itemize}[nosep]
\item If $\check{H}^1(\mathcal{U}, \mathcal{F}) = 0$, then the grade
  can be high: $T$ depends only on the quality of~$E$.
\item If $\check{H}^1 \neq 0$, then $T$ is bounded:
  \[
    T \;\leq\; \gone \;-\; \|[\check{H}^1]\|
  \]
  for some norm $\|\cdot\|$ on the cohomology group.  Non-trivial
  obstructions cap the achievable quality.
\item Higher cohomology provides finer bounds:
  $\check{H}^2 \neq 0$ implies $T$ is further reduced, reflecting
  the impossibility of even modifying the attempted gluing.
\end{itemize}

\medskip
\noindent\textbf{(C4) Grade feeds back to topos.}
The Graded Evaluation informs the construction of future Semantic
Topoi:
\begin{itemize}[nosep]
\item Low-grade regions of the site may be refined (adding new
  coordinates) or coarsened (merging coordinates) to improve
  future judgments.
\item The Lawvere--Tierney topology $\tau$ may be adjusted:
  replacing~$\tau$ by a finer topology (more covering sieves)
  demands stronger evidence, while a coarser topology is more
  permissive.
\item This feedback loop gives Judgment Geometry its
  \emph{adaptive} character.
\end{itemize}
\end{definition}

\section{The fibration}
\label{sec:fibration-def}

\begin{construction}[Total category of judgments]
\label{con:total-category}
Define the category $\JJ$ of \emph{judgments} as follows:
\begin{itemize}[nosep]
\item \textbf{Objects:} compatible quadruples
  $(\STop, \PStr, \CSpec, \GEval)$ as in
  Definition~\ref{def:compatible}.
\item \textbf{Morphisms:} a morphism
  $(f, \sigma, \lambda, h) :
  (\STop_1, \PStr_1, \CSpec_1, \GEval_1) \to
  (\STop_2, \PStr_2, \CSpec_2, \GEval_2)$
  consists of:
  \begin{enumerate}[label=(\alph*),nosep]
  \item a morphism $f : \STop_1 \to \STop_2$ of Semantic Topoi
    (Definition~\ref{def:stop-morph});
  \item a morphism $\sigma : \PStr_1 \to \PStr_2$ of Proof Structures
    (Definition~\ref{def:pstr-morph}), compatible with~$f$ via~(C1);
  \item a morphism $\lambda : \CSpec_1 \to \CSpec_2$ of Cohomological
    Spectra (Definition~\ref{def:cspec-morph}), compatible with
    $\sigma$ via~(C2);
  \item a morphism $h : \GEval_1 \to \GEval_2$ of Graded Evaluations
    (Definition~\ref{def:geval-morph}), compatible with $\lambda$
    via~(C3).
  \end{enumerate}
\item \textbf{Composition:} componentwise.
\item \textbf{Identity:} $(\id_\STop, \id_\PStr, \id_\CSpec,
  \id_\GEval)$.
\end{itemize}
\end{construction}

\begin{definition}[Judgment fibration]
\label{def:judgment-fibration}
The \emph{judgment fibration} is the functor
\[
  \pi : \JJ \;\longrightarrow\; \mathbf{STop}
\]
defined by $\pi(\STop, \PStr, \CSpec, \GEval) = \STop$ on objects and
$\pi(f, \sigma, \lambda, h) = f$ on morphisms.  This functor
``projects'' a judgment onto its underlying Semantic Topos, forgetting
the proof, cohomological, and evaluative data.
\end{definition}

\begin{theorem}[Fibration structure]
\label{thm:fibration}
The functor $\pi : \JJ \to \mathbf{STop}$ is a Grothendieck fibration:
for every morphism $f : \STop_1 \to \STop_2$ in~$\mathbf{STop}$ and
every object $(\STop_2, \PStr_2, \CSpec_2, \GEval_2)$ in the fibre
over~$\STop_2$, there exists a \emph{Cartesian lift}---a morphism
in~$\JJ$ over~$f$ that is universal among all lifts.
\end{theorem}

\begin{proof}[Proof sketch]
The Cartesian lift is constructed by pulling back the proof, cohomological,
and evaluative data along the geometric morphism~$f$.
Specifically:
\begin{enumerate}[label=(\roman*),nosep]
\item The context $\Gamma_1$ is obtained by applying $f^*$ to the
  types in~$\Gamma_2$.
\item The cover $\mathcal{U}_1$ is the pullback cover
  $f^*(\mathcal{U}_2)$.
\item The grade $T_1 = T_2$ is unchanged (grades are scalar).
\end{enumerate}
The universality follows from the universality of pullback in each
component category.
\end{proof}

\section{Fibres and sections}
\label{sec:fibres-sections}

\begin{definition}[Fibre]
The \emph{fibre} over a Semantic Topos $\STop$ is the subcategory
\[
  \pi^{-1}(\STop)
  \;=\;
  \{(\PStr, \CSpec, \GEval) :
  (\STop, \PStr, \CSpec, \GEval) \in \ob(\JJ)\}
\]
with morphisms $(\id_\STop, \sigma, \lambda, h)$.
\end{definition}

The fibre $\pi^{-1}(\STop)$ is the space of all possible judgments
\emph{over a fixed site and presheaf}.  It encodes the freedom in
choosing evidence, decomposing into sub-goals, and assigning grades.

\begin{definition}[Section]
\label{def:section}
A \emph{section} of the fibration~$\pi$ is a functor
$s : \mathbf{STop} \to \JJ$ with $\pi \circ s = \id_{\mathbf{STop}}$.
That is, $s$ assigns to each Semantic Topos a compatible judgment
over it, coherently across all morphisms of topoi.
\end{definition}

\begin{remark}
A section is a \emph{global judgment policy}: for every site and
presheaf, it specifies how to construct evidence, decompose into
sub-goals, and assign grades, in a way that is compatible with
all changes of site.
\end{remark}

\section{Harmony as section}
\label{sec:harmony-section}

\begin{definition}[Harmonious judgment]
\label{def:harmonious}
A judgment $(\STop, \PStr, \CSpec, \GEval)$ is \emph{harmonious} if:
\begin{enumerate}[label=(\alph*),nosep]
\item It is compatible (Definition~\ref{def:compatible}).
\item $\check{H}^1(\mathcal{U}, \mathcal{F}) = 0$: no gluing
  obstructions.
\item $T \geq T_{\mathrm{threshold}}$ for a chosen threshold
  $T_{\mathrm{threshold}} \in G$.
\end{enumerate}
A section $s : \mathbf{STop} \to \JJ$ is \emph{harmonious} if $s(\STop)$
is harmonious for every~$\STop$.
\end{definition}

\begin{proposition}
\label{prop:harmony-section}
A program~$P$ satisfies its specification~$\varphi$ if and only if there
exists a harmonious section $s$ of the judgment fibration with
$s(\STop_P) = (\STop_P, \PStr, \CSpec, \GEval)$ where $\PStr$ contains
a valid proof $E : A$ and $T = \gone$.
\end{proposition}

This recovers the fundamental theorem of Judgment Geometry: correctness
is the existence of a harmonious section with maximal grade.

\section{The bottleneck principle}
\label{sec:bottleneck-principle}

\begin{theorem}[Bottleneck principle]
\label{thm:bottleneck-p2}
For a section $s$ over a cover $\{U_i\}$ of the site, the global grade
is bounded by:
\[
  T_{\mathrm{global}}
  \;=\;
  \bigotimes_{i} T_{s(U_i)}
  \;\leq\;
  \min_i \, T_{s(U_i)}.
\]
That is, the global grade is dominated by the \emph{worst local grade}.
\end{theorem}

\begin{proof}
Since $\otG$ is multiplication on $[0,1]$ and all factors are in
$[0,1]$, we have $\prod_i T_i \leq \min_i T_i$.  The product
$\bigotimes_i T_{s(U_i)}$ computes the conservative aggregate
(Definition~\ref{def:aggregation}), which is the relevant measure for
global validity.
\end{proof}

\begin{corollary}
\label{cor:bottleneck}
A single stratum with $T_{s(U_j)} = \gzero$ forces
$T_{\mathrm{global}} = \gzero$, regardless of the grades on all other
strata.  In practice: a single failing module blocks the entire
program's verification.
\end{corollary}

\section{Descent as optimisation}
\label{sec:descent-optimisation}

The judgment fibration frames the central problem of Judgment Geometry
as an \emph{optimisation problem over sections}:

\begin{definition}[Judgment optimisation]
\label{def:judgment-opt}
Given a Semantic Topos $\STop$, the \emph{judgment optimisation
problem} is to find a section $s^*$ of $\pi$ that maximises the
global grade subject to the cohomological constraint:
\[
  s^* \;=\;
  \arg\max_{s \in \Gamma(\pi)}\;
  T_{\mathrm{global}}(s)
  \qquad\text{subject to}\qquad
  \check{H}^1(\mathcal{U}_{s(\STop)},\; \mathcal{F}) = 0.
\]
\end{definition}

\begin{remark}
The constraint $\check{H}^1 = 0$ is a necessary condition for
correctness; the optimisation over~$T$ selects the highest-quality
proof among all valid ones.  When no section satisfies the constraint,
the optimum is $T^* = \gzero$, signalling that the judgment is
\emph{not} harmonious.
\end{remark}

\begin{remark}[Descent = gluing sections]
The name ``descent'' is apt: in algebraic geometry, descent theory asks
when a local datum (defined on a cover) can be glued into a global
datum.  In Judgment Geometry, a ``local judgment'' is a section defined
on each $U_i$; \emph{descent} is the process of gluing these local
judgments into a global section.  The obstruction to descent is precisely
$\check{H}^1$, and successful descent is equivalent to finding an
optimal global section.
\end{remark}

\section{Summary}

The judgment fibration $\pi : \JJ \to \mathbf{STop}$ provides a
unified geometric framework for the four objects:

\begin{center}
\renewcommand{\arraystretch}{1.3}
\begin{tabular}{@{}ll@{}}
\textbf{Concept} & \textbf{Fibration interpretation} \\
\hline
Judgment & Object of the total category $\JJ$ \\
Semantic context & Object of the base $\mathbf{STop}$ \\
Evidence + obstructions + grade & Object in the fibre $\pi^{-1}(\STop)$ \\
Correctness & Existence of a harmonious section \\
Verification & Search for a section with $T \geq T_{\mathrm{threshold}}$ \\
Bug & Obstruction: $\check{H}^1 \neq 0$ or $T < T_{\mathrm{threshold}}$ \\
Refactoring & Geometric morphism in $\mathbf{STop}$ \\
Quality & Global grade $T_{\mathrm{global}} = \bigotimes_i T_i$
\end{tabular}
\end{center}

\chapter{The Component-Wise View}
\label{ch:equivalence}

This chapter makes precise the relationship between the structured
4-tuple formulation $(\STop, \PStr, \CSpec, \GEval)$ and its
component-wise flattening $(c, \varphi, A, E, O, B, T, \Pi)$.
We construct two functors---a
flattening functor that extracts the eight components from a compatible
quadruple, and a structuring functor that assembles components into their
canonical structured objects---and prove that they form an adjoint
equivalence of categories.  The upshot: the component-wise view is a
faithful, lossless projection of the structured view.

\section{The categories}

\begin{definition}[The category $\mathbf{Judg}_8$]
\label{def:judg8}
The category $\mathbf{Judg}_8$ of \emph{component-wise judgments} has:
\begin{itemize}[nosep]
\item \textbf{Objects:} component tuples
  $\mathcal{J} = (c, \varphi, A, E, O, B, T, \Pi)$ satisfying the
  coherence axioms of Judgment Geometry.
\item \textbf{Morphisms:} a morphism
  $\mathcal{J}_1 \to \mathcal{J}_2$ is a tuple of component maps
  $(f_c, f_\varphi, f_A, f_E, f_O, f_B, f_T, f_\Pi)$ preserving
  all compatibility conditions: $f_c$ maps coordinates to coordinates,
  $f_\varphi$ is a natural transformation of presheaves compatible
  with $f_c$, $f_A$ is a type map compatible with $f_E$, and so on.
\end{itemize}
\end{definition}

\begin{definition}[The category $\mathbf{Judg}_4$]
\label{def:judg4}
The category $\mathbf{Judg}_4$ is the category $\JJ$ of
Construction~\ref{con:total-category}: compatible quadruples
$(\STop, \PStr, \CSpec, \GEval)$ with the morphisms defined there.
\end{definition}

\section{The flattening functor}
\label{sec:forgetful-functor}

\begin{construction}[Flattening functor $U$]
\label{con:forgetful-functor}
The \emph{flattening functor}
$U : \mathbf{Judg}_4 \to \mathbf{Judg}_8$ extracts the eight
components from their rich ambient structures:
\[
  U(\STop, \PStr, \CSpec, \GEval)
  \;=\;
  (c,\; \varphi,\; A,\; E,\; O,\; B,\; T,\; \Pi)
\]
where:
\begin{enumerate}[label=(\roman*),nosep]
\item $c$ is any chosen object of the underlying site category~$C$
  of~$\STop$ (or the terminal object of~$C$ if it exists).
\item $\varphi$ is the judgment presheaf from~$\STop$
  (Definition~\ref{def:phi}).
\item $A$ is the type from~$\PStr$ (component~(ii) of
  Definition~\ref{def:pstr}).
\item $E$ is the evidence term from~$\PStr$ (component~(iii)).
\item $O$ is the cover $\mathcal{U}$ from~$\CSpec$
  (component~(i) of Definition~\ref{def:cspec}).
\item $B$ is the first cohomology class
  $[\check{H}^1(\mathcal{U}, \mathcal{F})]$ from~$\CSpec$
  (component~(iii), degree $n=1$).
\item $T$ is the grade from~$\GEval$
  (component~(iii) of Definition~\ref{def:geval}).
\item $\Pi$ is the proof trail from~$\GEval$ (component~(iv)).
\end{enumerate}
On morphisms, $U(f, \sigma, \lambda, h)$ extracts the component maps
in the obvious way.
\end{construction}

\begin{proposition}
\label{prop:U-well-defined}
The functor~$U$ is well-defined: the component tuple
$U(\STop, \PStr, \CSpec, \GEval)$ satisfies all coherence axioms of
Judgment Geometry.
\end{proposition}

\begin{proof}
The compatibility conditions (C1)--(C4) of
Definition~\ref{def:compatible} imply:
\begin{itemize}[nosep]
\item $\varphi$ is a presheaf on a site containing~$c$: the pair
  $(c, \varphi)$ is well-formed.
\item $\Gamma \vdash E : A$: the pair $(A, E)$ is a valid typing
  judgment.
\item The cover~$O$ and class~$B$ arise from the same site and
  sheaf: $(O, B)$ is a valid cohomological datum.
\item The grade~$T$ is in the grade semiring with trail~$\Pi$:
  $(T, \Pi)$ is a valid graded evaluation.
\item The cross-component axioms (e.g., $B = 0 \Leftrightarrow$
  $E$ is globally defined) are guaranteed by~(C2).
  \qedhere
\end{itemize}
\end{proof}

\section{The structuring functor}
\label{sec:enrichment-functor}

\begin{construction}[Structuring functor $R$]
\label{con:enrichment-functor}
The \emph{structuring functor}
$R : \mathbf{Judg}_8 \to \mathbf{Judg}_4$ assembles each component
tuple into its canonical rich structure:
\[
  R(c, \varphi, A, E, O, B, T, \Pi)
  \;=\;
  (\STop_{\min},\; \PStr_{\min},\; \CSpec_{\min},\; \GEval_{\min})
\]
where each component is the \emph{minimal} (or \emph{free}) rich object
containing the given data:

\medskip
\noindent\textbf{(i) Minimal topos $\STop_{\min}$.}
Let $C_{c,\varphi}$ be the full subcategory of the site generated by
the coordinate~$c$ and all coordinates appearing in the support
of~$\varphi$.  Equip it with the induced Grothendieck topology.  Then
$\STop_{\min}$ is the sheaf topos $\Sh(C_{c,\varphi}, J)$ with
presheaf~$\varphi$, together with all the data of
Definition~\ref{def:stop} (which are determined by the site).
The Lawvere--Tierney topology is the trivial one ($j = \id_\Omega$).

\medskip
\noindent\textbf{(ii) Minimal proof structure $\PStr_{\min}$.}
The context~$\Gamma$ is the free context generated by the variables
occurring in~$E$, with types assigned by the minimal type theory in
which $E : A$ is derivable.  The derivation~$\mathcal{D}$ is the
canonical derivation of $\Gamma \vdash E : A$.  The reduction theory
is the standard $\beta\eta$-reduction.  The universe level is the
smallest $i$ such that $A : \Type_i$.  Type formers ($\Pi$, $\Sigma$,
$\mathrm{Id}$, $W$) are included precisely when required by~$E$.

\medskip
\noindent\textbf{(iii) Minimal cohomological spectrum $\CSpec_{\min}$.}
The cover is~$O$.  The \v{C}ech complex $\Cech^\bullet(O, \mathcal{F})$
is computed directly.  The cohomology groups $\check{H}^n$ are the
cohomology of this complex.  The long exact sequence, Mayer--Vietoris,
spectral sequence, characteristic classes, and derived category are all
determined by the site and sheaf.

\medskip
\noindent\textbf{(iv) Minimal graded evaluation $\GEval_{\min}$.}
The grade semiring is the standard $([0,1], \max, \cdot, 0, 1)$.
The grade is~$T$, the trail is~$\Pi$.  The residuation, evaluation
functor, stability, convergence, entropy, and Bayesian update are
determined by the semiring and the data.

\medskip
\noindent On morphisms, $R$ maps a component-wise morphism of tuples
to the induced morphism of rich structures.
\end{construction}

\begin{proposition}
\label{prop:R-well-defined}
The enrichment functor~$R$ is well-defined: the quadruple
$R(c, \varphi, A, E, O, B, T, \Pi)$ satisfies the compatibility
conditions (C1)--(C4) of Definition~\ref{def:compatible}.
\end{proposition}

\begin{proof}
The coherence axioms of the component tuple ensure that the components
are mutually
consistent.  The minimal constructions preserve this consistency:
(C1)~holds because the context is read off from the same site and
presheaf; (C2)~holds because the cover $O$ arose from the same proof
decomposition; (C3)~holds because the coherence axiom
$B = 0 \Leftrightarrow$ global proof exists translates directly;
(C4)~holds trivially since no feedback is active in the minimal
construction.
\end{proof}

\section{The equivalence theorem}
\label{sec:equivalence-theorem}

\begin{theorem}[Adjoint equivalence]
\label{thm:equivalence}
The functors $R : \mathbf{Judg}_8 \to \mathbf{Judg}_4$ and
$U : \mathbf{Judg}_4 \to \mathbf{Judg}_8$ form an adjoint equivalence:
\[
  R \;\dashv\; U,
  \qquad
  U \circ R \;\cong\; \id_{\mathbf{Judg}_8},
  \qquad
  R \circ U \;\cong\; \id_{\mathbf{Judg}_4}.
\]
\end{theorem}

\begin{proof}
We construct the unit and counit natural isomorphisms.

\medskip
\noindent\textbf{Unit: $\eta : \id_{\mathbf{Judg}_8}
\xrightarrow{\;\sim\;} U \circ R$.}
For a component tuple
$\mathcal{J} = (c, \varphi, A, E, O, B, T, \Pi)$:
\[
  U(R(\mathcal{J}))
  = U(\STop_{\min}, \PStr_{\min}, \CSpec_{\min}, \GEval_{\min})
  = (c', \varphi, A, E, O, B', T, \Pi).
\]
The coordinate $c' = c$ (since $c \in \ob(C_{c,\varphi})$) and the
obstruction class $B' = B$ (since the \v{C}ech cohomology of the
minimal complex with cover~$O$ yields the same $\check{H}^1$).
Thus $U \circ R(\mathcal{J}) = \mathcal{J}$, and $\eta$ is the
identity natural transformation.

\medskip
\noindent\textbf{Counit: $\varepsilon : R \circ U
\xrightarrow{\;\sim\;} \id_{\mathbf{Judg}_4}$.}
For a quadruple $\mathcal{Q} = (\STop, \PStr, \CSpec, \GEval)$:
\[
  R(U(\mathcal{Q}))
  = R(c, \varphi, A, E, O, B, T, \Pi)
  = (\STop_{\min}, \PStr_{\min}, \CSpec_{\min}, \GEval_{\min}).
\]
The minimal topos $\STop_{\min}$ is a sub-topos of~$\STop$; the
inclusion $\iota : \STop_{\min} \hookrightarrow \STop$ is a geometric
morphism.  Similarly, $\PStr_{\min}$ is a sub-structure of~$\PStr$
(the minimal type theory embeds into the ambient one), $\CSpec_{\min}$
embeds into~$\CSpec$, and $\GEval_{\min}$ embeds into~$\GEval$.

The counit $\varepsilon_\mathcal{Q} : R(U(\mathcal{Q})) \to \mathcal{Q}$
is the canonical inclusion of the minimal enrichment into the full
enrichment.  It is an isomorphism up to the choice of ambient structure:
the additional data in~$\STop$ beyond~$\STop_{\min}$ (extra objects in
the site, a non-trivial Lawvere--Tierney topology) are determined by
the site and play no role in the component-level data.  Formally,
$\varepsilon$ is a natural isomorphism because any two enrichments of
the same component data are canonically isomorphic as objects of
$\mathbf{Judg}_4$ (the enrichment is determined up to unique
isomorphism by the data it enriches).

\medskip
\noindent\textbf{Triangle identities.}
We verify $U \varepsilon \circ \eta U = \id_U$ and
$\varepsilon R \circ R \eta = \id_R$.  Both follow immediately from
$\eta = \id$ and $\varepsilon$ being the canonical inclusion.
\end{proof}

\section{Transfer of theorems}

\begin{corollary}[Theorem transfer]
\label{cor:theorem-transfer}
Every theorem about component-wise judgments transfers to structured
judgments and vice versa.  Specifically, for any property~$P$ definable
in the language of either category:
\begin{enumerate}[label=(\alph*),nosep]
\item $\mathcal{J} \in \mathbf{Judg}_8$ satisfies~$P$ if and only if
  $R(\mathcal{J}) \in \mathbf{Judg}_4$ satisfies the transported
  property $R^*(P)$.
\item $\mathcal{Q} \in \mathbf{Judg}_4$ satisfies~$Q$ if and only if
  $U(\mathcal{Q}) \in \mathbf{Judg}_8$ satisfies $U^*(Q)$.
\end{enumerate}
\end{corollary}

\begin{proof}
An adjoint equivalence is, in particular, an equivalence of categories.
Equivalent categories satisfy the same categorical properties: any
diagram that commutes in one commutes in the other, any limit or
colimit that exists in one has a counterpart in the other, and any
functor defined on one extends to the other.
\end{proof}

\begin{corollary}[Fundamental theorem, 4-tuple version]
\label{cor:fundamental}
A program~$P$ satisfies its specification~$\varphi$ if and only if
the Cohomological Spectrum of the associated judgment satisfies
$\check{H}^1(\mathcal{U}, \mathcal{F}_\varphi) = 0$.
\end{corollary}

\begin{proof}
Apply the theorem transfer to the fundamental theorem of
Judgment Geometry.
\end{proof}

\section{Strict gain in expressiveness}
\label{sec:strict-gain}

\begin{remark}[The structured view sees more]
\label{rem:strict-gain}
While the two formulations are categorically equivalent (every theorem
transfers), the structured 4-tuple is \emph{strictly more expressive in
practice} because it makes visible:

\begin{enumerate}[label=(\roman*),nosep]
\item \textbf{Morphisms.}  The component-wise view has no intrinsic
  notion of ``morphism of judgments''; the 4-tuple inherits geometric
  morphisms, context morphisms, chain maps, and grade homomorphisms
  from its ambient categories.  These morphisms enable the study of
  \emph{judgment transformations}: refactoring, generalisation,
  specialisation, composition.

\item \textbf{Internal logic.}  The Semantic Topos carries the
  Mitchell--B\'enabou language; the Proof Structure carries the
  type-theoretic derivation calculus.  The flattened components
  $(c, \varphi)$ and $(A, E)$ carry no intrinsic logic.

\item \textbf{Higher cohomology.}  The component-wise view records only
  $H^1$ (the obstruction class~$B$); the Cohomological Spectrum
  carries $H^n$ for all~$n$, plus the long exact sequence, spectral
  sequences, and the derived category.  These detect obstructions
  invisible to~$H^1$ alone.

\item \textbf{Residuation.}  The Graded Evaluation carries the
  residual $\multimap$, which answers ``what quality is achievable
  given the current evidence?''  The flattened grade~$T$ is a bare
  number with no such algebraic structure.

\item \textbf{Fibration structure.}  The judgment fibration
  $\pi : \JJ \to \mathbf{STop}$ (Chapter~\ref{ch:fibration}) provides
  a unified geometric framework---fibres, sections, descent,
  optimisation---that has no counterpart in the flat component-wise
  view.
\end{enumerate}

These are not additional axioms but \emph{consequences} of the
structured formulation: they are present in the ambient mathematical
structures and become available in the 4-tuple.  The equivalence
theorem guarantees backward compatibility---no theorem is lost---while
the structured view provides forward capability.
\end{remark}


% ── Part III: CTT Extension ───────────────────────────────────
\part{Compositional Theory of Theories}
% Part III: Compositional Theory of Theories
% This file is \include'd from main.tex

\chapter{The Theory of Theories}
\label{ch:theory-of-theories}

Parts~I and~II developed Judgment Geometry as a four-object framework
$(\STop, \PStr, \CSpec, \GEval)$ for universal program verification.
In this Part we extend the framework to its most demanding application
domain: natural language, poetry, and aesthetics.  The central
technical device is the \emph{Compositional Theory of Theories}
(CTT)\index{CTT}, which treats each linguistic theory---phonology,
syntax, semantics, pragmatics, and so on---as a typed, graded
component that plugs into the judgment fibration at a designated
stratum.

\section{The Problem: 450+ Theories}
\label{sec:ctt-problem}

Natural language is not governed by a single formal system.
A complete account of even a single English sentence may require
contributions from phonology (segmental features, syllable weight,
stress assignment), morphology (inflectional paradigms, derivational
productivity), syntax (constituency, dependency, binding), semantics
(lambda terms, event structure, possible worlds), pragmatics (speech
acts, implicature, presupposition accommodation), discourse structure
(coherence relations, QUD management), information structure
(focus--topic partitioning, givenness), prosody (intonational
phrasing, pitch accent placement), and---for literary
language---poetics (meter, rhyme scheme, figuration) and aesthetics
(emotional resonance, rasa, sublimity).

Each of these ten \emph{strata} hosts dozens of competing and
complementary theories.  Within syntax alone one finds
Transformational Grammar, HPSG, LFG, CCG, TAG, Minimalism,
Dependency Grammar, and Construction Grammar, among others.  Within
semantics: Montague Grammar, DRT, SDRT, situation semantics, event
semantics, dynamic semantics, distributional semantics, frame
semantics.  Within pragmatics: Gricean theory, Relevance Theory,
speech act theory, game-theoretic pragmatics.  A conservative count
across all ten strata yields over 450 distinct formalisms that have
been proposed, debated, refined, and---in many cases---implemented.

The traditional response has been to pick one theory per stratum and
ignore the rest, yielding monolithic architectures (Montague Grammar,
HPSG, LFG) that are deep in one dimension and shallow in all others.
An alternative is the ``plug-and-play'' approach of constraint-based
unification architectures, but these lack a principled account of
quality and inter-stratal interaction.

\begin{remark}[The modularity challenge]
\label{rem:modularity}
Any framework that aspires to cover all of natural language must solve
two problems simultaneously:
\begin{enumerate}[label=(\roman*)]
\item \textbf{Horizontal compositionality}: theories at the same
  stratum must compete or cooperate gracefully (OT-style ranking,
  Bayesian model comparison, ensemble methods).
\item \textbf{Vertical compositionality}: theories at different
  strata must constrain one another (phonology--syntax interface,
  syntax--semantics interface, semantics--pragmatics interface).
\end{enumerate}
CTT solves both problems by embedding all theories into a single
graded algebraic structure.
\end{remark}

\section{Theories as Typed Objects}
\label{sec:theory-typed}

We now formalize the notion of a \emph{theory} within the
four-object framework.

\begin{definition}[Theory]
\label{def:theory}
A \textbf{theory} is a quadruple
\[
  \mathcal{T} \;=\; \bigl(s,\; d,\; \alpha,\; \kappa\bigr)
\]
where:
\begin{enumerate}[label=(\alph*)]
\item $s \in \mathrm{Strata} = \{\phon, \morph, \synt, \sem, \prag,
  \disc, \info, \pros, \poet, \aesth\}$ is the \textbf{stratum};
\item $d \in \mathbb{N}$ is the \textbf{depth} (sub-level within the
  stratum; e.g.\ segmental phonology has depth~$0$, prosodic
  phonology depth~$1$, intonational phonology depth~$2$);
\item $\alpha : \mathrm{Expr} \times \mathrm{GluingData}
  \to \GEval$ is the \textbf{analysis function}, which takes an
  expression and the current gluing context and returns a graded
  evaluation;
\item $\kappa : \mathrm{List}(\mathrm{Constraint})$ is the list of
  \textbf{constraints} the theory imposes.
\end{enumerate}
\end{definition}

The key insight is that every linguistic theory, however different in
origin and formalism, can be cast as an instance of
Definition~\ref{def:theory}:\ it lives at a stratum, it analyses
expressions, and it imposes constraints.  The analysis function
$\alpha$ returns a value in $\GEval$---not a binary accept/reject,
but a graded evaluation that records \emph{how well} the expression
satisfies the theory's expectations.

\begin{definition}[Theory registry]
\label{def:registry}
The \textbf{theory registry} is a dependent type
\[
  \mathrm{Reg} \;:\; \prod_{s \in \mathrm{Strata}}
  \mathrm{List}\bigl(\mathrm{Theory}(s)\bigr)
\]
assigning to each stratum a list of registered theories.
Registration is the act of adding a theory to this list.
\end{definition}

Each registered theory ``plugs in'' to the judgment fibration
(Chapter~\ref{ch:fibration}) at its stratum and depth.  The
fibration's fiber over stratum~$s$ is precisely the category of all
local sections produced by the theories registered at~$s$.

\begin{example}[Montague Grammar as a theory]
\label{ex:montague-theory}
Montague Grammar~\cite{montague1973} is registered as
$\mathcal{T}_{\mathrm{MG}} = (\sem, 0, \alpha_{\mathrm{MG}},
\kappa_{\mathrm{MG}})$ where $\alpha_{\mathrm{MG}}$ maps a
sentence to its denotation in a simply-typed lambda calculus with
generalized quantifiers, and $\kappa_{\mathrm{MG}}$ enforces
compositionality (the meaning of a complex expression is determined
by the meanings of its parts and their mode of combination).
\end{example}

\begin{example}[Optimality Theory as a theory]
\label{ex:ot-theory}
OT~\cite{prince1993} is registered as $\mathcal{T}_{\mathrm{OT}} =
(\phon, 0, \alpha_{\mathrm{OT}}, \kappa_{\mathrm{OT}})$ where
$\alpha_{\mathrm{OT}}$ maps an input--output pair to its harmony
value computed by evaluating ranked, violable constraints, and
$\kappa_{\mathrm{OT}}$ is the ranked constraint set itself.
Crucially, the OT harmony value is naturally a grade: each violation
reduces the grade multiplicatively via $\otG$.
\end{example}

\section{The Free Graded Semiring of Theories}
\label{sec:free-semiring}

Theories do not exist in isolation.  They compose in two ways:
\emph{competition} (alternative analyses at the same stratum) and
\emph{cooperation} (joint constraints across strata or within a
stratum).  These two modes correspond precisely to the two operations
of the grade semiring $(\Grade, \opG, \otG, \gzero, \gone)$
introduced in Chapter~\ref{ch:grade}.

\begin{definition}[Theory composition]
\label{def:theory-composition}
Let $\mathcal{T}_1, \mathcal{T}_2$ be theories.
\begin{enumerate}[label=(\alph*)]
\item \textbf{Cooperative composition} ($\otG$):  both theories must
  be satisfied.  The combined analysis is
  \[
    \alpha_{\mathcal{T}_1 \otG \mathcal{T}_2}(e, G)
    \;=\; \alpha_{\mathcal{T}_1}(e, G)
    \;\otG\; \alpha_{\mathcal{T}_2}(e, G).
  \]
  Since $\otG$ is addition in log-space, this multiplies
  probabilities: the joint grade is the product of individual grades.

\item \textbf{Competitive composition} ($\opG$):  the best analysis
  wins.  The combined analysis is
  \[
    \alpha_{\mathcal{T}_1 \opG \mathcal{T}_2}(e, G)
    \;=\; \alpha_{\mathcal{T}_1}(e, G)
    \;\opG\; \alpha_{\mathcal{T}_2}(e, G).
  \]
  Since $\opG$ is log-sum-exp, this is a smooth maximum: the better
  analysis dominates.
\end{enumerate}
\end{definition}

\begin{theorem}[CTT as free graded semiring]
\label{thm:free-semiring}
Let $\mathfrak{T} = \{\mathcal{T}_i\}_{i \in I}$ be a set of
theories.  The \textbf{Compositional Theory of Theories} generated
by $\mathfrak{T}$ is the free commutative semiring
\[
  \mathrm{CTT}(\mathfrak{T})
  \;=\; \Grade\langle \mathcal{T}_i : i \in I \rangle
  \;\cong\; \bigoplus_{S \subseteq I}
  \bigotimes_{i \in S} \mathcal{T}_i
\]
over the grade semiring $(\Grade, \opG, \otG, \gzero, \gone)$.
The Harmony of an expression $e$ under the full CTT is
\[
  \mathrm{Harm}(e) \;=\; \bigotimes_{s \in \mathrm{Strata}}
  \bigoplus_{\mathcal{T} \in \mathrm{Reg}(s)}
  \alpha_{\mathcal{T}}(e, G).
\]
\end{theorem}

\begin{proof}
The semiring axioms for $(\Grade, \opG, \otG, \gzero, \gone)$ were
established in Chapter~\ref{ch:grade}.  The free construction is
standard: we adjoin formal generators $\mathcal{T}_i$ and quotient
by the semiring axioms.  The Harmony formula arises from the
structure of the stratal fibration: at each stratum, competing
theories combine via $\opG$; across strata, the chosen analyses
combine via $\otG$.
\end{proof}

\begin{remark}[Interpretation of the Harmony formula]
The outer $\bigotimes$ enforces \emph{all strata must be satisfied}:
a fatal semantic violation cannot be rescued by phonological beauty.
The inner $\bigoplus$ implements \emph{best-theory-wins}: at each
stratum, the analysis with the highest grade dominates.
This is the formal content of the bottleneck principle.
\end{remark}

\section{The Open World Property}
\label{sec:open-world}

A critical design requirement for any framework covering natural
language is \emph{extensibility}: new phenomena are discovered, new
theories are proposed, and the framework must accommodate them
without requiring modification of existing components.

\begin{theorem}[Open world property]
\label{thm:open-world}
Let $\mathrm{CTT}(\mathfrak{T})$ be the theory of theories
generated by $\mathfrak{T}$, and let $\mathcal{T}_{\mathrm{new}}$ be
a new theory at stratum~$s$ with depth~$d$.  Then:
\begin{enumerate}[label=(\roman*)]
\item Registration of $\mathcal{T}_{\mathrm{new}}$ extends the
  Harmony by exactly one $\opG$-summand in stratum~$s$:
  \[
    \mathrm{Harm}_{\mathrm{new}}(e)
    \;=\; \mathrm{Harm}_{\mathrm{old}}(e)
    \;\otG\; \frac{
      \bigoplus_{\mathcal{T} \in \mathrm{Reg}(s)}
      \alpha_{\mathcal{T}}(e,G) \;\opG\;
      \alpha_{\mathcal{T}_{\mathrm{new}}}(e,G)
    }{
      \bigoplus_{\mathcal{T} \in \mathrm{Reg}(s)}
      \alpha_{\mathcal{T}}(e,G)
    }.
  \]
\item No existing theory's analysis function or constraint set is
  modified.
\item The time to register $\mathcal{T}_{\mathrm{new}}$ is
  $O(1)$---independent of $|\mathfrak{T}|$.
\end{enumerate}
\end{theorem}

\begin{proof}
Property~(i) follows from the distributivity of $\otG$ over $\opG$
in the grade semiring: the new theory adds one term to the $\opG$
sum at its stratum, and the $\otG$ product over strata factors
accordingly.  Properties~(ii) and~(iii) follow from the free
construction: generators are independent, and registration is simply
appending to the list $\mathrm{Reg}(s)$.
\end{proof}

\begin{corollary}[Scalability]
\label{cor:scalability}
The system scales linearly in the total number of theories.  Adding
the $n$-th theory costs $O(1)$ registration and $O(1)$ additional
per-expression evaluation time (one call to $\alpha_{\mathcal{T}_n}$).
\end{corollary}

\section{The Ten Strata}
\label{sec:ten-strata}

The stratal architecture is fixed at ten strata, ordered by the
classical linguistic hierarchy from form to meaning to use to art:

\begin{definition}[Strata]
\label{def:strata}
The set of strata is
\[
  \mathrm{Strata} = \{\phon, \morph, \synt, \sem, \prag, \disc,
  \info, \pros, \poet, \aesth\}
\]
with the partial order
\[
  \phon \preceq \morph \preceq \synt \preceq \sem \preceq \prag
  \preceq \disc, \qquad
  \info,\; \pros,\; \poet,\; \aesth
\]
where the first six strata form a total order (the classical
``pipeline''), while the last four are partially ordered, reflecting
the fact that information structure, prosody, poetics, and aesthetics
interact with all levels of the pipeline.
\end{definition}

Each stratum is equipped with its own internal structure:

\begin{enumerate}[label=(\arabic*)]
\item \textbf{Sub-site} $(\mathcal{C}_s, J_s)$:\ the category of
  coordinates relevant to stratum~$s$ (e.g.\ for $\phon$, objects
  are phonemes, syllables, feet, and prosodic words; for $\sem$,
  objects are lambda terms, event variables, and possible worlds).

\item \textbf{Sub-types} $\Type_s$:\ the universe of types at
  stratum~$s$ (e.g.\ $\Type_{\phon}$ contains feature bundles,
  syllable types, stress patterns; $\Type_{\sem}$ contains entity
  types, truth-value types, event types).

\item \textbf{Sub-cohomology} $H^*_s$:\ the cohomology of the
  sub-site, whose classes detect local gluing failures within the
  stratum (e.g.\ $H^1_{\phon}$ detects stress clashes;
  $H^1_{\sem}$ detects type mismatches in semantic composition).

\item \textbf{Inter-stratal constraints} $\kappa_{s,s'}$:\ the
  constraints governing the interface between strata $s$ and $s'$
  (e.g.\ $\kappa_{\synt,\sem}$ is the syntax--semantics interface;
  $\kappa_{\sem,\prag}$ governs presupposition projection;
  $\kappa_{\phon,\synt}$ governs prosodic phrasing relative to
  syntactic constituency).
\end{enumerate}

\begin{remark}[Why exactly ten?]
The choice of ten strata is empirical, not mathematical.  The
framework is parameterized by the stratal set; one could use fewer
(collapsing information structure into pragmatics, say) or more
(splitting semantics into lexical and compositional sub-strata).
Ten strata represent a consensus granularity that aligns with the
major subdivisions of modern linguistics and literary theory.
\end{remark}

\section{Theory Registration and Plugging In}
\label{sec:registration}

\begin{construction}[Registration protocol]
\label{con:registration}
To register a theory $\mathcal{T} = (s, d, \alpha, \kappa)$:
\begin{enumerate}[label=\textbf{Step \arabic*:}]
\item \textbf{Stratum assignment.}  Verify $s \in \mathrm{Strata}$
  and $d \in \mathbb{N}$.  The theory is filed at position $(s, d)$
  in the stratal fibration.
\item \textbf{Type declaration.}  Declare the input and output types
  of $\alpha$ within the sub-type universe $\Type_s$.
  Specifically, $\alpha$ must have signature
  $\alpha : \mathrm{Expr}_s \times \mathrm{GluingData} \to \GEval_s$
  where $\mathrm{Expr}_s$ and $\GEval_s$ are the expression type
  and evaluation type at stratum~$s$.
\item \textbf{Constraint declaration.}  Each constraint
  $c \in \kappa$ specifies:
  \begin{itemize}
  \item a \emph{scope}: the set of strata it references
    (intra-stratal if scope $= \{s\}$, inter-stratal otherwise);
  \item a \emph{weight} $w_c \in \Grade$;
  \item a \emph{predicate} $p_c : \mathrm{GluingData} \to \Grade$.
  \end{itemize}
\item \textbf{Append.}  Add $\mathcal{T}$ to $\mathrm{Reg}(s)$.
  This is an $O(1)$ operation.
\end{enumerate}
\end{construction}

After registration, the theory participates automatically in Harmony
computation: the stratal fibration includes its local sections, the
$\opG$ competition at its stratum accounts for its analyses, and the
$\otG$ composition across strata respects its inter-stratal
constraints.

\section{Classification of Representative Theories}
\label{sec:classification}

We now survey representative theories across all ten strata,
indicating for each its CTT formalization as
$(s, d, \alpha, \kappa)$.  This is not exhaustive---each stratum
hosts dozens of theories---but illustrates the range of the framework.

\medskip

\noindent\textbf{Phonology} ($\phon$).
\begin{itemize}[nosep]
\item \emph{SPE} (Chomsky--Halle):\ $({\phon}, 0, \alpha_{\mathrm{SPE}}, \kappa_{\mathrm{SPE}})$.
  Rules apply sequentially; $\alpha$ traces rule application,
  $\kappa$ is the ordered rule set.
\item \emph{OT Phonology}~\cite{prince1993}:\ $({\phon}, 0,
  \alpha_{\mathrm{OT}}, \kappa_{\mathrm{OT}})$.
  Parallel evaluation of candidates; $\alpha$ returns the harmony
  of the optimal candidate.
\item \emph{Autosegmental Phonology}:\ $({\phon}, 1,
  \alpha_{\mathrm{auto}}, \kappa_{\mathrm{auto}})$.
  Tier-based representations; $\alpha$ checks association-line
  well-formedness.
\end{itemize}

\noindent\textbf{Morphology} ($\morph$).
\begin{itemize}[nosep]
\item \emph{Distributed Morphology}:\ $({\morph}, 0,
  \alpha_{\mathrm{DM}}, \kappa_{\mathrm{DM}})$.
  Late insertion; $\alpha$ evaluates vocabulary item competition.
\item \emph{Word-and-Paradigm Morphology}:\ $({\morph}, 0,
  \alpha_{\mathrm{WP}}, \kappa_{\mathrm{WP}})$.
  Paradigm-based; $\alpha$ computes analogical similarity.
\item \emph{Finite-State Morphology}:\ $({\morph}, 1,
  \alpha_{\mathrm{FSM}}, \kappa_{\mathrm{FSM}})$.
  Transducer-based; $\alpha$ is transducer output weight.
\end{itemize}

\noindent\textbf{Syntax} ($\synt$).
\begin{itemize}[nosep]
\item \emph{Minimalism}:\ $({\synt}, 0, \alpha_{\mathrm{Min}},
  \kappa_{\mathrm{Min}})$.  Merge-based derivation; $\alpha$
  counts economy violations.
\item \emph{HPSG}:\ $({\synt}, 0, \alpha_{\mathrm{HPSG}},
  \kappa_{\mathrm{HPSG}})$.  Typed feature structures; $\alpha$
  checks unification success with graded constraints.
\item \emph{CCG}:\ $({\synt}, 0, \alpha_{\mathrm{CCG}},
  \kappa_{\mathrm{CCG}})$.  Combinatory rules; $\alpha$ counts
  type-raising and composition steps.
\item \emph{TAG}:\ $({\synt}, 1, \alpha_{\mathrm{TAG}},
  \kappa_{\mathrm{TAG}})$.  Tree-adjoining; $\alpha$ evaluates
  derivation tree well-formedness.
\item \emph{LFG}:\ $({\synt}, 0, \alpha_{\mathrm{LFG}},
  \kappa_{\mathrm{LFG}})$.  Parallel c-structure and f-structure;
  $\alpha$ checks completeness, coherence, consistency.
\end{itemize}

\noindent\textbf{Semantics} ($\sem$).
\begin{itemize}[nosep]
\item \emph{Montague Grammar}~\cite{montague1973}:\ $({\sem}, 0,
  \alpha_{\mathrm{MG}}, \kappa_{\mathrm{MG}})$.  IL denotations;
  $\alpha$ computes compositional truth conditions.
\item \emph{DRT}~\cite{kamp1993}:\ $({\sem}, 0,
  \alpha_{\mathrm{DRT}}, \kappa_{\mathrm{DRT}})$.  Discourse
  Representation Structures; $\alpha$ builds and verifies DRSs.
\item \emph{Event Semantics}:\ $({\sem}, 1,
  \alpha_{\mathrm{ev}}, \kappa_{\mathrm{ev}})$.  Neo-Davidsonian
  events; $\alpha$ checks thematic role assignment.
\item \emph{Frame Semantics}:\ $({\sem}, 1,
  \alpha_{\mathrm{frame}}, \kappa_{\mathrm{frame}})$.
  FrameNet-style frames; $\alpha$ evaluates frame element coverage.
\end{itemize}

\noindent\textbf{Pragmatics} ($\prag$).
\begin{itemize}[nosep]
\item \emph{Gricean Pragmatics}:\ $({\prag}, 0,
  \alpha_{\mathrm{Grice}}, \kappa_{\mathrm{Grice}})$.
  Maxim compliance; $\alpha$ grades adherence to Quality, Quantity,
  Relation, Manner.
\item \emph{Relevance Theory}:\ $({\prag}, 0,
  \alpha_{\mathrm{RT}}, \kappa_{\mathrm{RT}})$.
  Cost--benefit; $\alpha$ computes cognitive effect per processing
  effort.
\item \emph{Speech Act Theory}:\ $({\prag}, 1,
  \alpha_{\mathrm{SA}}, \kappa_{\mathrm{SA}})$.
  Illocutionary force; $\alpha$ checks felicity conditions.
\end{itemize}

\noindent\textbf{Discourse} ($\disc$).
\begin{itemize}[nosep]
\item \emph{RST (Rhetorical Structure Theory)}:\ $({\disc}, 0,
  \alpha_{\mathrm{RST}}, \kappa_{\mathrm{RST}})$.
  Rhetorical relations; $\alpha$ evaluates tree well-formedness.
\item \emph{SDRT}:\ $({\disc}, 0, \alpha_{\mathrm{SDRT}},
  \kappa_{\mathrm{SDRT}})$.  Segmented DRSs; $\alpha$ checks
  discourse attachment and update.
\item \emph{QUD Theory}:\ $({\disc}, 1, \alpha_{\mathrm{QUD}},
  \kappa_{\mathrm{QUD}})$.  Questions Under Discussion; $\alpha$
  grades relevance to the current QUD stack.
\end{itemize}

\noindent\textbf{Information Structure} ($\info$).
\begin{itemize}[nosep]
\item \emph{Focus Theory} (Rooth, Alternative Semantics):\
  $({\info}, 0, \alpha_{\mathrm{foc}}, \kappa_{\mathrm{foc}})$.
  Focus alternatives; $\alpha$ grades congruence with the QUD.
\item \emph{Givenness Hierarchy} (Gundel et al.):\
  $({\info}, 1, \alpha_{\mathrm{given}}, \kappa_{\mathrm{given}})$.
  Cognitive status; $\alpha$ evaluates referential form against
  givenness level.
\end{itemize}

\noindent\textbf{Prosody} ($\pros$).
\begin{itemize}[nosep]
\item \emph{Metrical Phonology}:\
  $({\pros}, 0, \alpha_{\mathrm{met}}, \kappa_{\mathrm{met}})$.
  Stress grids; $\alpha$ grades rhythmic regularity.
\item \emph{ToBI Intonation}:\
  $({\pros}, 1, \alpha_{\mathrm{ToBI}}, \kappa_{\mathrm{ToBI}})$.
  Pitch accent and boundary tone annotation; $\alpha$ evaluates
  prosodic phrasing.
\end{itemize}

\noindent\textbf{Poetics} ($\poet$).
\begin{itemize}[nosep]
\item \emph{Generative Metrics}:\
  $({\poet}, 0, \alpha_{\mathrm{GM}}, \kappa_{\mathrm{GM}})$.
  Metrical rules; $\alpha$ grades scansion quality.
\item \emph{Jakobsonian Poetics}:\
  $({\poet}, 1, \alpha_{\mathrm{Jak}}, \kappa_{\mathrm{Jak}})$.
  Parallelism and the poetic function; $\alpha$ measures structural
  recurrence.
\item \emph{Cognitive Poetics}:\
  $({\poet}, 2, \alpha_{\mathrm{CP}}, \kappa_{\mathrm{CP}})$.
  Conceptual metaphor, blending; $\alpha$ evaluates figurative
  coherence.
\end{itemize}

\noindent\textbf{Aesthetics} ($\aesth$).
\begin{itemize}[nosep]
\item \emph{Rasa Theory} (Bharata, Abhinavagupta):\
  $({\aesth}, 0, \alpha_{\mathrm{rasa}}, \kappa_{\mathrm{rasa}})$.
  Nine rasas; $\alpha$ grades emotional evocation.
\item \emph{Kantian Aesthetics}:\
  $({\aesth}, 1, \alpha_{\mathrm{Kant}}, \kappa_{\mathrm{Kant}})$.
  Purposiveness without purpose; $\alpha$ evaluates free play of
  imagination and understanding.
\item \emph{Sublimity Theory}:\
  $({\aesth}, 2, \alpha_{\mathrm{sub}}, \kappa_{\mathrm{sub}})$.
  The Burkean/Kantian sublime; $\alpha$ grades the tension between
  comprehension and apprehension.
\end{itemize}

\medskip

This survey instantiates 30 of the 450+ theories.  The full
registry---maintained computationally in the \texttt{jugeo.ctt}
package---grows as new theories are formalized and registered.
The open world property (Theorem~\ref{thm:open-world}) guarantees
that each new theory extends the system without modifying existing
components.

\begin{remark}[Theories as generators]
\label{rem:generators}
In the free semiring picture, each theory $\mathcal{T}_i$ is a
\emph{generator}.  An analysis of an expression is a \emph{monomial}
$\bigotimes_{i \in S} \mathcal{T}_i$ specifying which theories are
active and how they interact.  The full CTT is the \emph{polynomial
semiring} over these generators---the space of all possible
multi-theory analyses, with $\opG$ and $\otG$ providing the algebra.
\end{remark}

\chapter{Graded Dependent Types for Natural Language}
\label{ch:gdt}

The theories registered in Chapter~\ref{ch:theory-of-theories} analyse
expressions and return grades.  But what \emph{are} these grades,
formally?  In this chapter we develop \emph{Graded Dependent Type
Theory} (GDT), the type-theoretic backbone of CTT.  Every typing
judgment carries an intrinsic quality measure; every identity
assertion admits degrees; and every type constructor is
parameterized not only by its arguments but by the grade to which
those arguments are satisfied.

\section{Graded Judgments}
\label{sec:graded-judgments}

\begin{definition}[Graded typing judgment]
\label{def:graded-judgment}
A \textbf{graded typing judgment} has the form
\[
  \Gamma \vdash_g e : A
  \qquad\text{where } g \in \Grade
\]
and is read:\ ``in context $\Gamma$, the term $e$ inhabits type $A$
at grade $g$.''
\end{definition}

The grade $g$ measures \emph{how well} the term $e$ satisfies the
specification $A$.  Three distinguished values organize the grading
scale:

\begin{itemize}[nosep]
\item $g = \gone$ (perfect satisfaction):\ $e$ is a canonical
  inhabitant of $A$ with no violations.
\item $g = \gzero$ (complete failure):\ $e$ does not inhabit $A$
  in any meaningful sense.
\item $\gzero < g < \gone$ (partial match):\ $e$ approximately
  inhabits $A$, with the grade quantifying the degree of
  approximation.
\end{itemize}

\begin{remark}[Relation to binary typing]
When $\Grade = \{\gzero, \gone\}$, the graded judgment
$\Gamma \vdash_g e : A$ reduces to the classical binary judgment
$\Gamma \vdash e : A$ (when $g = \gone$) or its negation (when
$g = \gzero$).  Thus GDT is a \emph{conservative extension} of
Martin-L\"of Type Theory~\cite{mlm1984}: every derivation in MLTT
is a derivation in GDT with all grades equal to~$\gone$.
\end{remark}

The structural rules of GDT are:

\begin{definition}[Graded structural rules]
\label{def:structural-rules}
\leavevmode
\begin{enumerate}[label=(\alph*)]
\item \textbf{Weakening.} If $\Gamma \vdash_g e : A$ and
  $x : B$ is fresh, then $\Gamma, x : B \vdash_g e : A$.
  Grade is preserved.
\item \textbf{Substitution.} If $\Gamma, x : A \vdash_g e : B$
  and $\Gamma \vdash_{g'} a : A$, then
  $\Gamma \vdash_{g \otG g'} e[a/x] : B[a/x]$.
  Grades compose via $\otG$.
\item \textbf{Grade monotonicity.} If $\Gamma \vdash_g e : A$
  and $g' \leq g$, then $\Gamma \vdash_{g'} e : A$.
  Lower grades are always available.
\end{enumerate}
\end{definition}

The substitution rule is crucial: when we substitute a term of
grade~$g'$ into a judgment of grade~$g$, the result has grade
$g \otG g'$.  Since $\otG$ is addition in log-space (multiplication
of probabilities), the combined grade is at most as good as the
worst component---the algebraic expression of the bottleneck
principle.

\section{Graded Identity Types}
\label{sec:gid}

The most important new type constructor in GDT is the
\emph{graded identity type}, which subsumes a remarkable range of
linguistic phenomena under a single notion of ``sameness at a
grade.''

\begin{definition}[Graded identity type]
\label{def:gid}
Given a type $A$ and terms $a, b : A$, the \textbf{graded identity
type} is the family
\[
  \GId_g(A, a, b) \qquad (g \in \Grade)
\]
whose inhabitants are evidence that $a$ and $b$ are ``the same at
grade~$g$.''
\end{definition}

The graded identity type has the following rules:

\begin{definition}[Rules for $\GId$]
\label{def:gid-rules}
\leavevmode
\begin{enumerate}[label=(\alph*)]
\item \textbf{Formation.}
  \[
    \frac{\Gamma \vdash A : \Type \qquad
      \Gamma \vdash a : A \qquad
      \Gamma \vdash b : A \qquad
      g \in \Grade}{
      \Gamma \vdash \GId_g(A, a, b) : \Type}
  \]

\item \textbf{Introduction (reflexivity).}
  \[
    \frac{\Gamma \vdash a : A}{
      \Gamma \vdash_{\gone} \mathsf{refl}_a : \GId_{\gone}(A, a, a)}
  \]

\item \textbf{Elimination (transport).}
  \[
    \frac{\Gamma \vdash_g p : \GId_g(A, a, b) \qquad
      \Gamma, x : A \vdash C(x) : \Type \qquad
      \Gamma \vdash_{g'} c : C(a)}{
      \Gamma \vdash_{g \otG g'} \mathsf{tr}(p, c) : C(b)}
  \]

\item \textbf{Computation ($\beta$-rule).}
  \[
    \mathsf{tr}(\mathsf{refl}_a, c) \equiv c
    \qquad\text{at grade } \gone \otG g' = g'.
  \]
\end{enumerate}
\end{definition}

\begin{proposition}[Graded groupoid structure]
\label{prop:graded-groupoid}
The graded identity types endow each type with a \emph{graded
groupoid} structure:
\begin{enumerate}[label=(\roman*)]
\item \textbf{Reflexivity:}\ $\mathsf{refl}_a : \GId_{\gone}(A,a,a)$.
\item \textbf{Symmetry:}\ if $p : \GId_g(A,a,b)$, then
  $p^{-1} : \GId_g(A,b,a)$.
\item \textbf{Transitivity with grade degradation:}\ if
  $p : \GId_{g_1}(A,a,b)$ and $q : \GId_{g_2}(A,b,c)$, then
  $q \cdot p : \GId_{g_1 \otG g_2}(A,a,c)$.
\end{enumerate}
\end{proposition}

\begin{proof}
Reflexivity is the introduction rule.  Symmetry follows from
transport along the family $C(x) = \GId_g(A, x, a)$.  Transitivity
follows from transport of $p$ along $q$; the grades compose via
$\otG$ by the substitution rule.
\end{proof}

The key linguistic insight is that different grades of identity
correspond to different linguistic relations:

\begin{example}[Linguistic identity spectrum]
\label{ex:identity-spectrum}
Let $A = \mathrm{Expression}$ be the type of English expressions.
\begin{center}
\renewcommand{\arraystretch}{1.2}
\begin{tabular}{lll}
\textbf{Grade range} & \textbf{Relation} & \textbf{Example} \\
\hline
$g = \gone$ & Definitional equality & ``cat'' $=$ ``cat'' \\
$g > 0.9$ & Synonymy & ``big'' $\sim$ ``large'' \\
$g > 0.8$ & Paraphrase & ``X killed Y'' $\sim$ ``Y died because of X'' \\
$g > 0.7$ & Metaphor & ``Juliet is the sun'' \\
$g > 0.6$ & Translation & ``cat'' $\sim$ ``chat'' (Fr.) \\
$g > 0.5$ & Allusion & ``To be or not to be'' $\leadsto$ Hamlet \\
$g > 0.4$ & Rhyme & ``cat'' $\sim$ ``bat'' (phonological $\GId$) \\
$g > 0.3$ & Metrical equivalence & iamb $\sim$ pyrrhic (prosodic $\GId$) \\
\end{tabular}
\end{center}
The continuous nature of $\Grade$ allows smooth interpolation
between these discrete categories.
\end{example}

\begin{remark}[Relation to HoTT identity]
When $g = \gone$, the graded identity type $\GId_{\gone}(A,a,b)$
is precisely the identity type $\mathrm{Id}_A(a,b)$ of Homotopy
Type Theory~\cite{hott2013}.  The graded groupoid structure
at $g = \gone$ recovers the $\infty$-groupoid structure of HoTT.
At lower grades, we obtain a novel structure: a graded $\infty$-groupoid
in which paths have quality.
\end{remark}

\begin{definition}[Graded univalence]
\label{def:graded-univalence}
For types $A, B$, define the type of \textbf{graded equivalences}
\[
  A \simeq_g B \;:\equiv\;
  \Sigma(f : A \to B).\,\Sigma(h : B \to A).\,
  \GId_g\bigl(A \to A,\; h \circ f,\; \id_A\bigr)
  \times
  \GId_g\bigl(B \to B,\; f \circ h,\; \id_B\bigr).
\]
\textbf{Graded univalence} asserts
$(A \simeq_g B) \simeq_{\gone} \GId_g(\Type, A, B)$.
\end{definition}

Graded univalence states that two types are ``the same at grade~$g$''
precisely when there is a graded equivalence between them.
At $g = \gone$ this is ordinary univalence; at lower grades it
captures \emph{approximate type equivalence}---a type-theoretic
formalization of ``close enough.''

\section{Graded Dependent Products and Sums}
\label{sec:graded-pi-sigma}

\begin{definition}[Graded dependent product]
\label{def:graded-pi}
Given $\Gamma \vdash A : \Type$ and $\Gamma, x : A \vdash B(x) : \Type$,
the \textbf{graded dependent product} is
\[
  \Pi_g(x : A).\, B(x)
\]
whose inhabitants are functions $f$ such that for all $a : A$,
we have $\Gamma \vdash_{g'} f(a) : B(a)$ with $g' \geq g$.  The
overall grade of $f$ is
\[
  \mathrm{grade}(f) \;=\; \bigotimes_{a : A} g'(a),
\]
the $\otG$-product of grades over all arguments.
\end{definition}

\begin{definition}[Graded dependent sum]
\label{def:graded-sigma}
The \textbf{graded dependent sum}
\[
  \Sigma_g(x : A).\, B(x)
\]
is inhabited by pairs $(a, b)$ with $a : A$ at some grade $g_1 \geq g$
and $b : B(a)$ at some grade $g_2 \geq g$.  The grade of the pair is
$g_1 \otG g_2$.
\end{definition}

\begin{remark}[Linguistic interpretation]
The graded dependent product $\Pi_g(x : A).\, B(x)$ is read:
``for all $x$ of type $A$, $B(x)$ holds at grade~$\geq g$.''  A theory
is only as good as its worst prediction: if $B(x_0)$ fails badly for some
$x_0$, the whole theory is penalized.

The graded dependent sum $\Sigma_g(x : A).\, B(x)$ is read: ``there
exists an $x$ of type $A$ such that $B(x)$ holds at grade~$\geq g$.''
An analysis is a witness---a specific parsing, interpretation, or
rendering---together with evidence of its quality.
\end{remark}

The rules for graded products and sums:

\begin{definition}[Rules for $\Pi_g$ and $\Sigma_g$]
\label{def:pi-sigma-rules}
\leavevmode
\begin{enumerate}[label=(\alph*)]
\item \textbf{$\Pi$-introduction.}
  \[
    \frac{\Gamma, x : A \vdash_g e : B(x)}{
      \Gamma \vdash_g \lambda x.\, e : \Pi_g(x:A).\, B(x)}
  \]
\item \textbf{$\Pi$-elimination.}
  \[
    \frac{\Gamma \vdash_g f : \Pi_g(x:A).\, B(x) \qquad
      \Gamma \vdash_{g'} a : A}{
      \Gamma \vdash_{g \otG g'} f(a) : B(a)}
  \]
\item \textbf{$\Sigma$-introduction.}
  \[
    \frac{\Gamma \vdash_{g_1} a : A \qquad
      \Gamma \vdash_{g_2} b : B(a)}{
      \Gamma \vdash_{g_1 \otG g_2} (a, b) :
      \Sigma_{g_1 \otG g_2}(x:A).\, B(x)}
  \]
\item \textbf{$\Sigma$-elimination.}
  \[
    \frac{\Gamma \vdash_g p : \Sigma_g(x:A).\, B(x)}{
      \Gamma \vdash_{g'} \pi_1(p) : A \qquad
      \Gamma \vdash_{g''} \pi_2(p) : B(\pi_1(p))}
  \]
  where $g' \otG g'' = g$.
\end{enumerate}
\end{definition}

\section{New Type Constructors for Natural Language}
\label{sec:nl-constructors}

Beyond the standard dependent type formers, CTT introduces
domain-specific type constructors for linguistic and poetic phenomena.
Each constructor has formation, introduction, elimination, and
computation rules, and each lives within the four-object framework:
its well-formedness is checked in $\STop$, its proof structure
lives in $\PStr$, its obstructions in $\CSpec$, and its quality
in $\GEval$.

\subsection{Jakobsonian Parallelism}
\label{subsec:jak}

Roman Jakobson identified the ``poetic function'' as the projection
of the principle of equivalence from the axis of selection to the
axis of combination.  We formalize this as a type constructor.

\begin{definition}[Jakobsonian parallelism type]
\label{def:jak}
Given a ``voice'' $v$ (a stratum or combination of strata) and
expressions $e_1, e_2$, the type
\[
  \mathrm{Jak}_v(e_1, e_2) : \Type
\]
is inhabited by evidence of structural equivalence between $e_1$
and $e_2$ in voice $v$.
\begin{enumerate}[label=(\alph*)]
\item \textbf{Formation:}\ $v \in \mathrm{Strata}^+$,
  $e_1 : \mathrm{Expr}$, $e_2 : \mathrm{Expr}$.
\item \textbf{Introduction:}\ a witness
  $w : \mathrm{Jak}_v(e_1, e_2)$ at grade $g$ consists of an
  alignment $\sigma : \mathrm{struct}_v(e_1) \to
  \mathrm{struct}_v(e_2)$ with
  $g = \GId_g(\mathrm{Struct}_v, \mathrm{struct}_v(e_1),
  \mathrm{struct}_v(e_2))$.
\item \textbf{Elimination:}\ from $w : \mathrm{Jak}_v(e_1,e_2)$
  one may extract the alignment $\sigma$ and transfer properties
  from $e_1$ to $e_2$ at the grade of the parallelism.
\end{enumerate}
\end{definition}

\begin{example}
In Shakespeare's Sonnet~73, the three quatrains exhibit syntactic
parallelism (``In me thou see'st \ldots'') at grade $g \approx 0.85$
and semantic parallelism (three metaphors for aging: autumn, twilight,
embers) at grade $g \approx 0.90$.  The Jakobsonian type
$\mathrm{Jak}_{\synt \otG \sem}(Q_1, Q_2)$ captures both dimensions
simultaneously.
\end{example}

\subsection{Iconicity}
\label{subsec:ico}

\begin{definition}[Iconicity type]
\label{def:ico}
The type $\mathrm{Ico}(A, B) : \Type$ is inhabited by evidence that
the phonological structure $A$ \emph{mirrors} the semantic structure
$B$.  Formally, $\mathrm{Ico}(A,B)$ is a graded morphism
\[
  \iota : \mathrm{struct}_{\phon}(A)
  \xrightarrow{g} \mathrm{struct}_{\sem}(B)
\]
that is structure-preserving at grade~$g$.
\end{definition}

\begin{example}
Onomatopoeia (``buzz,'' ``hiss'') achieves $\mathrm{Ico}$ at high
grade.  Sound symbolism (``gl-'' in ``gleam, glitter, glow'')
achieves it at moderate grade.  Arbitrary signs (``dog'') have
$\mathrm{Ico}$ at grade $\approx \gzero$.
\end{example}

\subsection{Poetic Force}
\label{subsec:pf}

\begin{definition}[Poetic force type]
\label{def:pf}
Given a force $\alpha$ (apostrophe, ekphrasis, volta, enjambment,
anaphora, etc.)\ and a content type $A$, the type
\[
  \mathrm{PF}(\alpha, A) : \Type
\]
is inhabited by expressions that convey content of type $A$ with
poetic force $\alpha$ at some grade.
\begin{enumerate}[label=(\alph*)]
\item \textbf{Introduction:}\ given $e : A$ at grade $g_1$ and
  evidence $f$ that $e$ exhibits force $\alpha$ at grade $g_2$,
  the pair $(e, f) : \mathrm{PF}(\alpha, A)$ at grade
  $g_1 \otG g_2$.
\item \textbf{Elimination:}\ from $p : \mathrm{PF}(\alpha, A)$,
  extract the content $\pi_1(p) : A$ and the force evidence
  $\pi_2(p)$.
\end{enumerate}
\end{definition}

\subsection{Rasa Type}
\label{subsec:rasa}

\begin{definition}[Rasa type]
\label{def:rasa}
Given a rasa $r \in \{\text{\'s\.r\.ng\=ara}, \text{h\=asya},
\text{karu\d{n}a}, \text{raudra}, \text{v\=\i ra},
\text{bhay\=anaka}, \text{b\=\i bhatsa},
\text{adbhuta}, \text{\'s\=anta}\}$
and a content type $A$, the type
\[
  \mathrm{Rasa}(r, A) : \Type
\]
is inhabited by evidence that an expression of type $A$ evokes
rasa $r$.  The grade measures the intensity of evocation.
\end{definition}

\begin{remark}
The nine rasas of Bharata's \emph{N\=a\d{t}ya\'s\=astra} correspond
to aesthetic emotions (love, humor, pathos, fury, heroism, terror,
disgust, wonder, peace).  The rasa type formalizes the claim that
aesthetic quality is not a single scalar but a \emph{typed} quantity:
one must specify \emph{which} aesthetic response is evoked and
\emph{how strongly}.
\end{remark}

\subsection{Pustejovskian Qualia}
\label{subsec:qualia}

\begin{definition}[Qualia type]
\label{def:qualia}
For a concept $C$, the \textbf{qualia structure}~\cite{pustejovsky1995} is
\[
  \mathrm{Qualia}(C) \;=\;
  \Sigma_g\bigl(
    f : \mathrm{Formal}(C),\;
    c : \mathrm{Const}(C),\;
    t : \mathrm{Telic}(C),\;
    a : \mathrm{Agent}(C)
  \bigr)
\]
where the four qualia roles are:
\begin{itemize}[nosep]
\item \textbf{Formal:}\ what $C$ is (category, type);
\item \textbf{Constitutive:}\ what $C$ is made of (parts, material);
\item \textbf{Telic:}\ what $C$ is for (purpose, function);
\item \textbf{Agentive:}\ how $C$ comes into being (origin, creation).
\end{itemize}
The grade of the qualia structure is
$g = g_f \otG g_c \otG g_t \otG g_a$.
\end{definition}

\subsection{Discourse Monad}
\label{subsec:dmon}

\begin{definition}[Discourse monad]
\label{def:dmon}
The \textbf{discourse monad} is
\[
  \mathrm{DMon}(A) \;=\; \mathrm{DCtx} \to (A \times \mathrm{DCtx})
\]
where $\mathrm{DCtx}$ is the type of discourse contexts (containing
the discourse referent store, the QUD stack, the common ground, and
the commitment slate).  A term of type $\mathrm{DMon}(A)$ is a
discourse action that, given a context, produces a value of type $A$
and an updated context.

The monad operations have graded types:
\begin{align*}
  \mathsf{return} &: A \to \mathrm{DMon}(A) &&\text{at grade } \gone, \\
  \mathsf{bind} &: \mathrm{DMon}(A) \to (A \to \mathrm{DMon}(B))
  \to \mathrm{DMon}(B) &&\text{at grade } g_1 \otG g_2.
\end{align*}
\end{definition}

\begin{example}
DRT~\cite{kamp1993} is recovered as the special case where
$\mathrm{DCtx}$ contains only a DRS (set of discourse referents +
conditions) and the grade is binary.
\end{example}

\subsection{Force Dynamics}
\label{subsec:fdyn}

\begin{definition}[Force dynamics type]
\label{def:fdyn}
Following Talmy, the \textbf{force dynamics type}
\[
  \mathrm{FDyn}(a, b, r) : \Type
\]
is inhabited by evidence of a force-dynamic interaction between
Agonist $a$ and Antagonist $b$ with result $r \in
\{\text{motion}, \text{rest}, \text{change}, \text{stasis}\}$.
The grade measures the strength of the force-dynamic construal.
\end{definition}

\subsection{Conceptual Blending}
\label{subsec:blend}

\begin{definition}[Blending type]
\label{def:blend}
The \textbf{Fauconnier--Turner blending type}
\[
  \mathrm{Blend}(I_1, I_2, G) : \Type
\]
is inhabited by evidence of a conceptual blend with input spaces
$I_1, I_2$, generic space $G$ (the shared structure), and an emergent
blended space $B$.  Formally:
\begin{enumerate}[label=(\alph*)]
\item projections $\pi_1 : B \to I_1$ and $\pi_2 : B \to I_2$;
\item a generic map $\gamma : G \to B$;
\item emergent structure: $B$ contains elements not in $I_1$ or $I_2$.
\end{enumerate}
The grade measures coherence: how well the blend preserves the
generic structure while generating genuinely emergent meaning.
\end{definition}

\begin{example}
The metaphor ``This surgeon is a butcher'' blends the surgery frame
($I_1$) and the butchery frame ($I_2$) via the shared structure of
cutting ($G$).  The emergent meaning---incompetence---is not present
in either input space.  The blend's grade measures how naturally the
emergent meaning arises.
\end{example}

\section{Instantiation in the Four-Tuple}
\label{sec:instantiation}

Each of the type constructors above inhabits all four objects of the
judgment geometry:

\begin{proposition}[Four-object instantiation]
\label{prop:four-object}
For each type constructor $\mathrm{T}$ in
$\{\GId, \mathrm{Jak}, \mathrm{Ico}, \mathrm{PF}, \mathrm{Rasa},
\mathrm{Qualia}, \mathrm{DMon}, \mathrm{FDyn}, \mathrm{Blend}\}$:
\begin{enumerate}[label=(\roman*)]
\item $\mathrm{T}$ lives in $\PStr$ as a type former with graded
  introduction and elimination rules.
\item The well-formedness of $\mathrm{T}$ is checked in $\STop$:
  the semantic topos provides the context, world knowledge, and
  compositional structure needed to verify that the type is
  well-formed.
\item Obstructions to $\mathrm{T}$ inhabitation are detected in
  $\CSpec$: a non-trivial $H^1$ class in $\CSpec$ indicates that
  the type constructor's constraints cannot be satisfied globally.
\item The quality of a $\mathrm{T}$-inhabitant is computed in
  $\GEval$: the grade $g$ of the typing judgment
  $\Gamma \vdash_g e : \mathrm{T}(\ldots)$ is the output of the
  evaluation function at the relevant stratum.
\end{enumerate}
\end{proposition}

\begin{proof}
By construction.  Each type former's introduction rule produces a
term in $\PStr$.  The premises of the formation rule require
well-formedness conditions checked in $\STop$.  The elimination
rule may fail to apply when local sections do not glue, producing
a class in $\CSpec$.  The grade of the introduction rule is computed
by $\GEval$.
\end{proof}

\begin{remark}[Extensibility]
The list of type constructors is not closed.  Any linguistic or
aesthetic phenomenon that can be cast as a type former with graded
rules can be added to the system.  The open world property
(Theorem~\ref{thm:open-world}) guarantees that the addition is
non-disruptive.  Candidates for future constructors include:
evidentiality types (source of information), honorific types
(social register), and narrative perspective types (focalization).
\end{remark}

\chapter{Stratal Fibration and the 450 Theories}
\label{ch:strata}

Chapters~\ref{ch:theory-of-theories} and~\ref{ch:gdt} developed
theories as typed objects and the graded type theory they inhabit.
This chapter assembles these ingredients into the \emph{stratal
fibration}---the geometric backbone of CTT---and uses it to classify
the full landscape of natural-language theories.

\section{The Stratal Fibration}
\label{sec:stratal-fibration}

\begin{definition}[Stratal fibration]
\label{def:stratal-fibration}
The \textbf{stratal fibration} is a functor
\[
  \pi_s : \mathbf{Lang} \to \mathbf{Strata}
\]
where:
\begin{enumerate}[label=(\alph*)]
\item $\mathbf{Strata}$ is the discrete category on the set
  $\{\phon, \morph, \synt, \sem, \prag, \disc, \info, \pros,
  \poet, \aesth\}$, equipped with the partial order of
  Definition~\ref{def:strata};
\item $\mathbf{Lang}$ is the total category whose objects are
  \emph{local analyses}---typed, graded judgments at a particular
  stratum---and whose morphisms are inter-stratal constraint
  morphisms;
\item the fiber $\pi_s^{-1}(s)$ over stratum $s$ is the category
  of all local sections produced by theories registered at $s$.
\end{enumerate}
\end{definition}

The fibration gives a precise geometric picture: a complete analysis
of an expression is a \emph{section} of the fibration---a choice of
one local analysis at each stratum, subject to the constraint that
neighboring choices are compatible.

\begin{remark}[Relation to the judgment fibration]
The stratal fibration refines the judgment fibration of
Chapter~\ref{ch:fibration} by decomposing the base into ten strata.
The judgment fibration has a single base category (the site); the
stratal fibration partitions the site into stratal sub-sites and
organizes theories by their stratal location.
\end{remark}

\section{GluingData}
\label{sec:gluingdata}

\begin{definition}[GluingData]
\label{def:gluingdata}
A \textbf{GluingData object} for an expression $\varphi$ is a
dependent record
\[
  \mathrm{GluingData}(\varphi) \;=\;
  \bigl\{s \mapsto (A_s, E_s, T_s) : s \in \mathrm{Strata}\bigr\}
\]
where:
\begin{enumerate}[label=(\alph*)]
\item $A_s \in \Type_s$ is the \textbf{type assignment} at
  stratum $s$ (the specification of what the analysis should produce:
  a phonological representation, a parse tree, a logical form, etc.);
\item $E_s : A_s$ is the \textbf{evidence} (the actual analysis:
  the phonological representation, the parse, the logical form);
\item $T_s \in \Grade$ is the \textbf{stratal grade} (how well the
  evidence satisfies the type assignment).
\end{enumerate}
A GluingData object is a \emph{section of the stratal fibration}:
it selects one local analysis per stratum.
\end{definition}

\begin{definition}[Inter-stratal compatibility]
\label{def:inter-stratal}
For strata $s \preceq s'$, the \textbf{inter-stratal compatibility
grade} is
\[
  G_{s,s'}\bigl((A_s, E_s, T_s),\; (A_{s'}, E_{s'}, T_{s'})\bigr)
  \;\in\; \Grade
\]
computed by the inter-stratal constraints $\kappa_{s,s'}$.  It
measures how well the local analyses at $s$ and $s'$ are consistent
with each other.
\end{definition}

\begin{example}[Syntax--semantics interface]
\label{ex:syn-sem}
Consider $s = \synt$, $s' = \sem$.  The type assignment $A_{\synt}$
is a parse tree; $A_{\sem}$ is a logical form.  The inter-stratal
constraint $\kappa_{\synt,\sem}$ checks that the logical form is
compositionally derivable from the parse tree.  If the parse is
$[\text{every}_{\text{NP}} [\text{boy}_{\text{N}}]]
[\text{ran}_{\text{VP}}]$ and the logical form is
$\forall x.\, \mathrm{boy}(x) \to \mathrm{ran}(x)$, the
compatibility grade is high.  If instead the logical form is
$\exists x.\, \mathrm{boy}(x) \wedge \mathrm{ran}(x)$, the
compatibility grade penalizes the quantifier mismatch.
\end{example}

\section{Harmony as Global Section}
\label{sec:harmony-global}

\begin{definition}[Harmony]
\label{def:harmony}
The \textbf{Harmony} of a GluingData object $\sigma$ for expression
$\varphi$ is
\[
  h(\varphi) \;=\;
  \bigotimes_{s \in \mathrm{Strata}} T_s(\varphi)
  \;\otG\;
  \bigotimes_{\substack{s \preceq s' \\ s, s' \in \mathrm{Strata}}}
  G_{s,s'}(\sigma_s, \sigma_{s'}).
\]
An analysis is \textbf{harmonious} if $h(\varphi) \geq \theta$ for
a context-dependent threshold $\theta \in \Grade$.
\end{definition}

Since $\otG$ is addition in log-space, Harmony is the sum of all
stratal grades and inter-stratal compatibility grades.  The key
structural consequence:

\begin{proposition}[Bottleneck principle]
\label{prop:bottleneck}
Let $\sigma$ be a GluingData object with stratal grades
$T_1, \ldots, T_{10}$ and inter-stratal grades
$G_{s,s'}$.  Then
\[
  h(\varphi) \;\leq\; \min\bigl(
  T_1, \ldots, T_{10},\;
  G_{s_1,s_1'}, \ldots, G_{s_k,s_k'}\bigr).
\]
The worst component dominates.
\end{proposition}

\begin{proof}
Since $\Grade \subseteq (-\infty, 0]$ and $\otG = {+}$, the sum
of non-positive terms is at most as large as the smallest summand.
\end{proof}

\begin{remark}[Linguistic consequence]
The bottleneck principle captures a basic intuition:\ a sentence with
perfect syntax, perfect semantics, and a fatal pragmatic violation
(e.g.\ a presupposition failure) is not ``mostly good''---it is
pragmatically infelicitous.  No amount of syntactic well-formedness
can rescue a fatal violation at another stratum.
\end{remark}

\section{Inter-Stratal Obstructions}
\label{sec:inter-stratal-obstructions}

When local analyses are individually well-graded but fail to compose
into a global section, the failure is detected by the first
cohomology of the stratal fibration.

\begin{definition}[Stratal $\Cech$ cohomology]
\label{def:stratal-cech}
Let $\{U_s\}_{s \in \mathrm{Strata}}$ be the canonical covering of
the stratal site (each $U_s$ is the sub-site at stratum $s$).  The
\textbf{stratal \v{C}ech cohomology} is
\[
  H^0\bigl(\{U_s\}, \mathcal{F}\bigr) =
  \ker(\delta^0),
  \qquad
  H^1\bigl(\{U_s\}, \mathcal{F}\bigr) =
  \ker(\delta^1) / \mathrm{im}(\delta^0)
\]
where $\delta^0$ is the coboundary map
\[
  \delta^0(\sigma)_{s,s'} \;=\;
  \sigma_s\big|_{U_s \cap U_{s'}}
  \;-\; \sigma_{s'}\big|_{U_s \cap U_{s'}}
\]
measuring the discrepancy between local analyses on overlaps
(inter-stratal interfaces).
\end{definition}

\begin{theorem}[Obstructions are $H^1$ classes]
\label{thm:obstructions}
A GluingData object $\sigma$ extends to a global section (a fully
harmonious analysis) if and only if the obstruction class
$[\delta^0(\sigma)] \in H^1\bigl(\{U_s\}, \mathcal{F}\bigr)$
vanishes.
\end{theorem}

\begin{proof}
This is the standard sheaf-theoretic gluing criterion applied to the
stratal fibration: sections on the individual opens $U_s$ glue to a
global section if and only if their differences on overlaps are
coboundaries, which is precisely the vanishing of the $H^1$ class.
\end{proof}

\begin{example}[Garden-path effect as obstruction]
\label{ex:garden-path}
Consider ``The horse raced past the barn fell.''  At the $\synt$
stratum, the parser initially analyses ``raced'' as the main verb
(local section $\sigma_{\synt}$).  At the $\sem$ stratum, ``fell''
requires a reduced relative clause reading (local section
$\sigma_{\sem}$).  These local sections disagree on the $\synt$--$\sem$
interface: $\delta^0(\sigma)_{\synt,\sem} \neq 0$, producing a
non-trivial $H^1$ class.  The garden-path effect \emph{is} the
obstruction; reanalysis (reparsing ``raced'' as a passive participle)
is the descent repair.
\end{example}

\begin{construction}[Descent as repair]
\label{con:descent-repair}
Given a non-trivial obstruction $[\omega] \in H^1$:
\begin{enumerate}[label=\textbf{Step \arabic*:}]
\item \textbf{Identify} the worst inter-stratal discrepancy:
  find $(s, s')$ maximizing $|\delta^0(\sigma)_{s,s'}|$.
\item \textbf{Compute the repair frontier}: the set of local
  modifications to $\sigma_s$ or $\sigma_{s'}$ that could reduce
  the discrepancy.
\item \textbf{Apply the best repair}: choose the modification that
  maximizes the resulting Harmony.
\item \textbf{Iterate} until $H^1 = 0$ or a fixpoint is reached.
\end{enumerate}
\end{construction}

\section{The Ten Strata in Detail}
\label{sec:strata-detail}

We now describe each stratum in greater detail, specifying its
sub-site, characteristic types, and representative phenomena.

\paragraph{Phonology ($\phon$).}
The sub-site $(\mathcal{C}_{\phon}, J_{\phon})$ has objects ranging
from individual segments (phonemes with feature bundles such as
$[{+}\text{voice}]$, $[{+}\text{nasal}]$) through syllables
(onset--nucleus--coda structure) and metrical feet (iambs, trochees)
to prosodic words and intonational phrases.  Coverings are given by
the prosodic hierarchy.  Characteristic types include feature bundles,
syllable templates, stress grids, and tone melodies.  The sub-cohomology
$H^1_{\phon}$ detects phonotactic violations (illicit consonant
clusters), stress clashes, and OPC (Obligatory Contour Principle)
violations.

\paragraph{Morphology ($\morph$).}
Objects are morphemes, stems, affixes, and words.  Coverings reflect
the part--whole structure of morphological constituency.  Characteristic
types include inflectional paradigms (person--number--tense--aspect--mood
feature bundles), derivational schemas (nominalization, verbalization),
and compounding templates.  $H^1_{\morph}$ detects paradigm gaps
(defective verbs), blocking effects, and affix-ordering violations.

\paragraph{Syntax ($\synt$).}
Objects range from lexical items and phrases to clauses and sentences.
The sub-site's coverings reflect constituency and dependency structure.
Types include category labels (NP, VP, CP), feature structures
(case, agreement, finiteness), and movement chains (A-movement,
$\bar{A}$-movement, head movement).  $H^1_{\synt}$ detects binding
violations (Principle A/B/C failures), island violations (extraction
from relative clauses or adjuncts), and agreement failures.

\paragraph{Semantics ($\sem$).}
Objects are semantic units: entity terms, predicates, quantifiers,
events, possible worlds, and discourse referents.  Types include
entity types, truth-value types, event types, intensional types
(functions from worlds), and higher-order types for generalized
quantifiers.  The sub-site's coverings reflect compositional scope:
sub-expressions must combine to yield the meaning of the whole.
$H^1_{\sem}$ detects type mismatches (category errors), scope
ambiguities (when multiple scopings have incompatible presuppositions),
and sortal violations.

\paragraph{Pragmatics ($\prag$).}
Objects are speech acts, implicatures, presuppositions, and
conversational moves.  Types include illocutionary-force types
(assertion, question, command), felicity-condition types, and
implicature types (scalar, particularized, generalized).  The
sub-site's coverings reflect the structure of conversational exchange.
$H^1_{\prag}$ detects presupposition failures (``The king of France
is bald'' in a context where France is a republic), Gricean maxim
violations, and infelicitous speech acts.

\paragraph{Discourse ($\disc$).}
Objects are discourse segments, rhetorical relations, and discourse
referents.  Types include RST relation types (elaboration, contrast,
cause, result), SDRT attachment types, and QUD types.  $H^1_{\disc}$
detects incoherent discourse structure (missing rhetorical relations,
dangling anaphora, non-sequiturs) and topic drift.

\paragraph{Information Structure ($\info$).}
Objects are information-structural partitions of clauses: focus,
topic, given, and new material.  Types include focus types
(contrastive, presentational, informational), topic types (aboutness,
contrastive, frame-setting), and givenness types (Gundel's hierarchy:
in focus, activated, familiar, uniquely identifiable, referential,
type identifiable).  $H^1_{\info}$ detects focus--topic mismatches,
violations of the given-before-new principle, and incongruence
between prosodic prominence and information-structural status.

\paragraph{Prosody ($\pros$).}
Objects are prosodic constituents: moras, syllables, feet,
phonological words, phonological phrases, intonational phrases, and
utterances.  Types include stress patterns (primary, secondary,
unstressed), intonation contours (ToBI labels: H*, L*, L+H*, etc.),
and rhythmic templates.  $H^1_{\pros}$ detects stress clashes and
lapses, mismatches between prosodic and syntactic phrasing, and
intonation--meaning incongruence (e.g.\ a falling contour on a
yes/no question).

\paragraph{Poetics ($\poet$).}
Objects are poetic units: feet, lines, couplets, stanzas, cantos,
and entire poems.  Types include meter types (iambic pentameter,
alexandrine, free verse), rhyme types (perfect, slant, eye, internal),
form types (sonnet, villanelle, ghazal, haiku), and figurative types
(metaphor, simile, metonymy, synecdoche).  $H^1_{\poet}$ detects
broken meter, forced rhymes, formal violations (wrong number of lines
in a sonnet), and clich\'{e}d imagery.

\paragraph{Aesthetics ($\aesth$).}
Objects are aesthetic experiences and evaluations.  Types include
rasa types (the nine rasas of Indian aesthetics), sublimity types
(Burkean, Kantian), beauty types (symmetry, proportion, harmony in
the colloquial sense), and affective types (intensity, valence,
arousal).  $H^1_{\aesth}$ detects emotional incoherence (a poem that
begins in the mode of karu\d{n}a and shifts abruptly to h\=asya
without motivating transition), tonal inconsistency, and aesthetic
failures (kitsch, sentimentality).

\section{Classification of 450 Theories}
\label{sec:450-classification}

The following table classifies representative theories by stratum and
depth, indicating for each the theory name, the primary formalism it
contributes, and the inter-stratal interfaces it participates in.

\begin{center}
\small
\renewcommand{\arraystretch}{1.15}
\begin{tabular}{llp{4.5cm}l}
\textbf{Stratum} & \textbf{Representative} & \textbf{Formalism} &
  \textbf{Interfaces} \\
\hline
$\phon$ & SPE, OT, Autosegmental &
  rules / constraints / tiers &
  $\phon$--$\morph$, $\phon$--$\pros$ \\
$\morph$ & DM, WP, Finite-State &
  late insertion / paradigms / FST &
  $\morph$--$\synt$, $\morph$--$\phon$ \\
$\synt$ & Minimalism, HPSG, LFG, CCG, TAG &
  merge / features / c/f-struct / combinators / trees &
  $\synt$--$\sem$, $\synt$--$\info$ \\
$\sem$ & Montague, DRT, Events, Frames &
  IL / DRS / neo-Davidson / FrameNet &
  $\sem$--$\prag$, $\sem$--$\synt$ \\
$\prag$ & Grice, RT, Speech Acts &
  maxims / relevance / felicity &
  $\prag$--$\disc$, $\prag$--$\sem$ \\
$\disc$ & RST, SDRT, QUD &
  rhetorical trees / SDRS / QUD stacks &
  $\disc$--$\info$, $\disc$--$\prag$ \\
$\info$ & Alt.\ Sem., Givenness &
  alternatives / cognitive status &
  $\info$--$\pros$, $\info$--$\synt$ \\
$\pros$ & Metrical, ToBI &
  stress grids / ToBI labels &
  $\pros$--$\phon$, $\pros$--$\poet$ \\
$\poet$ & Gen.\ Metrics, Jakobson, Cog.\ Poetics &
  scansion / parallelism / blending &
  $\poet$--$\aesth$, $\poet$--$\sem$ \\
$\aesth$ & Rasa, Kantian, Sublimity &
  rasas / purposiveness / tension &
  $\aesth$--$\poet$, $\aesth$--$\prag$ \\
\end{tabular}
\end{center}

The 30 theories named explicitly in this monograph represent only the
most prominent.  Within each stratum, the depth parameter
(Definition~\ref{def:theory}) organizes theories by granularity:
depth~$0$ theories handle the stratum's core phenomena; depth~$1$
theories address secondary phenomena; depth~$2$ and beyond handle
specialized sub-domains.  The full registry, maintained
computationally, currently enumerates ${\sim}450$ theories across all
strata and depths.

\begin{remark}[Interaction patterns]
Not all pairs of strata interact equally.  The strongest interfaces
are $\synt$--$\sem$ (the syntax--semantics interface, formalized
since Montague), $\phon$--$\pros$ (the phonology--prosody interface),
and $\prag$--$\disc$ (the pragmatics--discourse interface).  Weaker
but still important interfaces include $\info$--$\pros$ (the
focus--prosody interface: focused constituents receive pitch accents)
and $\poet$--$\aesth$ (the poetics--aesthetics interface: formal
beauty contributes to but does not determine aesthetic value).  The
stratal fibration accommodates all of these by providing inter-stratal
compatibility grades for every pair.
\end{remark}

\begin{proposition}[Fiber completeness]
\label{prop:fiber-complete}
For each stratum $s$, the fiber $\pi_s^{-1}(s)$ is a complete
category: it has all small limits and colimits.  In particular:
\begin{enumerate}[label=(\roman*)]
\item Products in the fiber correspond to \emph{cooperative
  composition} of theories at the same stratum (both must be
  satisfied).
\item Coproducts correspond to \emph{competitive composition}
  (the best theory wins).
\item Equalizers detect agreement between theories.
\item The terminal object is the trivial theory that accepts
  everything at grade $\gone$.
\item The initial object is the vacuous theory that rejects
  everything at grade $\gzero$.
\end{enumerate}
\end{proposition}

\begin{proof}
The fiber $\pi_s^{-1}(s)$ is the category of graded presheaves on
the sub-site $(\mathcal{C}_s, J_s)$.  Presheaf categories are
complete and cocomplete (limits and colimits are computed
pointwise).  The grading structure is preserved because
$(\Grade, \opG, \otG, \gzero, \gone)$ is a complete lattice under
the natural order.
\end{proof}

\chapter{NLU as Term Inference, NLG as Proof Search}
\label{ch:nlu-nlg}

The stratal fibration and graded type theory developed in the
preceding chapters are not merely descriptive: they yield algorithms.
In this chapter we show that natural language understanding (NLU) is
\emph{term inference}---reconstructing the graded typing judgment
for a given utterance---and natural language generation (NLG) is
\emph{proof search}---finding the highest-grade inhabitant of a
given meaning type.

\section{NLU as Term Inference}
\label{sec:nlu}

\begin{definition}[The NLU problem]
\label{def:nlu-problem}
Given an utterance $u$ and a context $\Gamma$, the \textbf{NLU
problem} is to construct a 4-tuple $(\STop_u, \PStr_u, \CSpec_u,
\GEval_u)$ such that
\[
  \Gamma \vdash_g E : A
  \quad\text{with}\quad
  g \;=\; h(u) \;=\; \mathrm{Harm}(\sigma_u)
\]
where $\sigma_u$ is the GluingData section for $u$.
\end{definition}

The four components are constructed as follows:

\paragraph{$\STop_u$: Building the discourse topos.}
The semantic topos for $u$ encodes the context of utterance:
the discourse state (previously introduced referents, the QUD stack,
the common ground), world knowledge (encyclopaedic and situational),
and a model of the speaker (beliefs, intentions, register).
Formally, $\STop_u = \Sh(\mathcal{C}_u, J_u)$ is the category of
sheaves on the contextual site---the Grothendieck topos determined
by the context.

\paragraph{$\PStr_u$: Inferring the type and term.}
The proof structure assigns a type $A$ (the meaning type: what is
being said) and a term $E$ (the evidence: the specific analysis) at
each stratum.  At the $\synt$ stratum, $A$ is the syntactic category
and $E$ is the parse tree.  At the $\sem$ stratum, $A$ is the logical
form type and $E$ is the lambda term.  The full proof structure is
the GluingData section $\sigma_u$.

\paragraph{$\CSpec_u$: Detecting obstructions.}
Obstructions are detected by computing the \v{C}ech cohomology of
the stratal fibration for $\sigma_u$.  A non-trivial $H^1$ class
signals an inter-stratal inconsistency: ambiguity (multiple local
sections at one stratum), presupposition failure (incompatible
pragmatic and semantic sections), or incoherence (discourse structure
that does not match syntactic structure).

\paragraph{$\GEval_u$: Computing quality.}
The graded evaluation computes the Harmony $h(u)$: the $\otG$-product
of all stratal grades and inter-stratal compatibility grades.  This
single number summarizes how well the utterance $u$ conveys its
intended meaning in the given context.

\begin{construction}[Bottom-up type inference across strata]
\label{con:nlu-algorithm}
The NLU algorithm proceeds bottom-up through the stratal hierarchy:
\begin{enumerate}[label=\textbf{Step \arabic*:}]
\item \textbf{Phonological analysis.}  Segment $u$ into phonemes;
  assign feature bundles; compute syllable structure, stress, and
  intonation.  Output:\ $\sigma_{\phon}$ at grade $T_{\phon}$.
\item \textbf{Morphological analysis.}  Identify morphemes; compute
  inflectional and derivational structure.  Output:\ $\sigma_{\morph}$
  at grade $T_{\morph}$.
\item \textbf{Syntactic parsing.}  Build constituent or dependency
  structure.  Multiple candidate parses may be maintained (beam).
  Output:\ $\sigma_{\synt}$ at grade $T_{\synt}$.
\item \textbf{Semantic composition.}  Compositionally construct the
  logical form from the parse.  Output:\ $\sigma_{\sem}$ at grade
  $T_{\sem}$.
\item \textbf{Pragmatic enrichment.}  Compute implicatures, resolve
  presuppositions, determine illocutionary force.  Output:\
  $\sigma_{\prag}$ at grade $T_{\prag}$.
\item \textbf{Discourse integration.}  Attach $u$ to the existing
  discourse structure; update referents and QUD.  Output:\
  $\sigma_{\disc}$ at grade $T_{\disc}$.
\item \textbf{Cross-stratal strata} ($\info$, $\pros$, $\poet$,
  $\aesth$).  Compute information structure, prosodic analysis,
  poetic analysis (if applicable), and aesthetic evaluation.
\item \textbf{Gluing.}  Assemble the GluingData section
  $\sigma_u = (\sigma_s)_{s \in \mathrm{Strata}}$.
  Compute inter-stratal compatibility grades $G_{s,s'}$.
  If $H^1 \neq 0$, invoke the descent repair procedure
  (Construction~\ref{con:descent-repair}).
\item \textbf{Harmony.}  Output $h(u) = \bigotimes_s T_s \otG
  \bigotimes_{s \preceq s'} G_{s,s'}$.
\end{enumerate}
\end{construction}

\begin{definition}[HPLF: Harmonically-balanced Polyphonic Logical Form]
\label{def:hplf}
The output of the NLU algorithm is the $\HPLF$, defined as
\[
  \HPLF(u) \;=\; \bigl(\sigma_u,\; h(u),\; [\omega_u]\bigr)
\]
where $\sigma_u$ is the GluingData section, $h(u)$ is the Harmony,
and $[\omega_u] \in H^1$ is the residual obstruction class (which
may be non-trivial for poetically or aesthetically intentional
``violations'').
\end{definition}

\begin{remark}[Polyphony]
The name ``polyphonic'' reflects the fact that the $\HPLF$ is not a
single-stranded logical form (as in Montague Grammar) but a
\emph{multi-voiced} structure: each stratum contributes one ``voice''
to the analysis, and Harmony measures how well they sing together.
\end{remark}

\section{NLG as Proof Search}
\label{sec:nlg}

\begin{definition}[The NLG problem]
\label{def:nlg-problem}
Given a meaning type $A$, a context $\Gamma$, and a desired Harmony
threshold $\theta$, the \textbf{NLG problem} is to find an utterance
$u$ such that
\[
  \Gamma \vdash_g E_u : A
  \qquad\text{with}\quad
  g \geq \theta.
\]
That is, find a term $E_u$ of type $A$ whose grade is at least
$\theta$.  Under the Curry--Howard correspondence, this is
\emph{proof search}: find the best inhabitant of the type $A$.
\end{definition}

The connection to proof search is precise: if $A$ is a proposition,
then an inhabitant is a proof; if $A$ is a meaning specification,
then an inhabitant is an utterance that conveys that meaning; if $A$
is a poem specification (form + topic + emotion), then an inhabitant
is a poem.

\begin{construction}[$\FMLM$: Fixpoint Masked Language Model]
\label{con:fmlm}
The $\FMLM$ (Fixpoint Masked Language Model) is an iterative
algorithm for NLG that constructs proof terms top-down:
\begin{enumerate}[label=\textbf{Step \arabic*:}]
\item \textbf{Initialize.}  Start with an underspecified template
  $E_0$ in which key positions are masked (denoted $[\mathrm{MASK}]$).
  The template encodes the coarse structure of the target: for a
  sonnet, it might specify 14 lines with metrical slots and a rhyme
  scheme.
  \[
    E_0 \;=\; \lambda\text{-skeleton with masked argument slots}
  \]
\item \textbf{Iterate.}  At each step $t$:
  \begin{enumerate}[label=(\alph*)]
  \item \textbf{Select} the highest-priority masked position
    (determined by the stratal ordering and constraint urgency).
  \item \textbf{Fill} the mask by searching for the
    best-graded term at that position, given the current partial
    term $E_t$ and the context $\Gamma$.
  \item \textbf{Recompute Harmony.}  Update the GluingData section
    and recompute $h(E_{t+1})$.
  \end{enumerate}
  \[
    E_{t+1} \;=\; E_t[\text{best fill} / \mathrm{MASK}_t]
  \]
\item \textbf{Converge.}  The process terminates when no masked
  positions remain, or when $h(E_t) \geq \theta$ and the term is
  stable.
\end{enumerate}
\end{construction}

\begin{theorem}[$\FMLM$ convergence]
\label{thm:fmlm-convergence}
If the Harmony function $h$ is continuous in the grade topology and
each fill operation is monotone (never decreases Harmony by more
than the gain from filling the mask), then the $\FMLM$ converges to
a fixpoint $E_\infty$ with $h(E_\infty) \geq h(E_0)$.
\end{theorem}

\begin{proof}[Proof sketch]
The sequence $h(E_0) \leq h(E_1) \leq \cdots$ is monotone and
bounded above (by $\gone = 0$).  By completeness of $\Grade$ (which
is a closed subset of $\mathbb{R} \cup \{-\infty\}$), the sequence
converges.  Continuity of $h$ ensures the limit is achieved.
\end{proof}

\begin{remark}[Descent as generation]
The $\FMLM$ can be viewed as a descent procedure: the partial term
$E_t$ is a partial section of the stratal fibration, and each fill
operation extends the section to one more stratum or repairs an
inter-stratal inconsistency.  Convergence of the $\FMLM$ is
convergence of the descent.
\end{remark}

\begin{construction}[Beam search over proof terms]
\label{con:beam-search}
In practice, the $\FMLM$ is combined with beam search:
\begin{enumerate}[label=(\alph*)]
\item Maintain $k$ candidate partial terms
  $\{E_t^{(1)}, \ldots, E_t^{(k)}\}$.
\item At each step, expand each candidate by filling one mask, then
  retain the top-$k$ candidates by Harmony.
\item Return the highest-Harmony completed term.
\end{enumerate}
The beam width $k$ trades off quality against computation: larger $k$
explores more of the proof space but costs proportionally more.
\end{construction}

\section{Poetry Generation as Constrained Proof Search}
\label{sec:poetry-generation}

Poetry generation is a special case of NLG in which the meaning type
$A$ encodes not only semantic content but formal constraints (meter,
rhyme, stanza form) and aesthetic requirements (emotional tone,
figurative density).

\begin{definition}[Poem specification type]
\label{def:poem-spec}
A \textbf{poem specification} is a type
\[
  A_{\mathrm{poem}} \;=\;
  \Sigma_g\bigl(
    \text{form} : \mathrm{Form},\;
    \text{topic} : \mathrm{Topic},\;
    \text{emotion} : \mathrm{Rasa}
  \bigr)
\]
where $\mathrm{Form}$ specifies the formal constraints (meter, rhyme
scheme, stanza count, line length), $\mathrm{Topic}$ specifies the
semantic content, and $\mathrm{Rasa}$ specifies the target aesthetic
emotion.
\end{definition}

The poem is the proof term: a well-typed inhabitant of
$A_{\mathrm{poem}}$.  The grade of the proof measures how well the
poem satisfies all constraints simultaneously.

\begin{example}[Sonnet generation]
\label{ex:sonnet}
A Shakespearean sonnet specification:
\begin{align*}
  \mathrm{Form} &= \text{14 lines, iambic pentameter,
    rhyme scheme ABAB CDCD EFEF GG} \\
  \mathrm{Topic} &= \text{passage of time, mortality} \\
  \mathrm{Rasa} &= \text{karu\d{n}a (pathos) with
    a turn to \'s\=anta (peace)}
\end{align*}
Each constraint is a theory slot:
\begin{itemize}[nosep]
\item \emph{Meter} ($\pros$, depth~0):\ each line must scan as
  iambic pentameter.  Grade penalizes trochaic substitutions
  ($g \approx 0.85$) less than spondaic ones ($g \approx 0.7$).
\item \emph{Rhyme} ($\phon$, depth~1):\ line-final words must
  rhyme according to the ABAB scheme.  Perfect rhyme achieves
  $\gone$; slant rhyme achieves $g \approx 0.8$.
\item \emph{Form} ($\poet$, depth~0):\ 14 lines, specific turn
  structure (volta at line~9 or~13).
\item \emph{Imagery} ($\sem$, depth~1):\ metaphorical coherence
  across the three quatrains.
\item \emph{Emotion} ($\aesth$, depth~0):\ tonal arc from pathos
  to peace.
\end{itemize}
\end{example}

\begin{remark}[Obstructions in poetry]
In poetry, some obstructions are \emph{desirable}.  A
``forced rhyme'' is an obstruction at the $\phon$--$\sem$ interface
(the rhyme word does not serve the meaning).  Enjambment is an
obstruction at the $\synt$--$\poet$ interface (the syntactic unit
overflows the line boundary).  In skilled poetry, these obstructions
are \emph{resolved aesthetically}: the forced rhyme creates surprise;
the enjambment creates urgency.  The $\HPLF$ records such
obstructions as non-trivial $H^1$ classes that are \emph{preserved
rather than repaired}---they are features, not bugs.
\end{remark}

\begin{proposition}[Preserved obstructions]
\label{prop:preserved-obstructions}
An obstruction $[\omega] \in H^1$ is \textbf{aesthetically preserved}
if the $\aesth$-stratum grade of preserving $[\omega]$ exceeds the
$\aesth$-stratum grade of repairing it:
\[
  T_{\aesth}(\sigma \text{ with } [\omega]) \;>\;
  T_{\aesth}(\sigma \text{ with } [\omega] = 0).
\]
In this case, the descent procedure
(Construction~\ref{con:descent-repair}) halts: the obstruction is a
local minimum of $H^1$ magnitude but a local \emph{maximum} of
Harmony, because the aesthetic gain outweighs the inter-stratal cost.
\end{proposition}

\begin{proof}
The descent procedure selects repairs that maximize Harmony.  If
repairing $[\omega]$ reduces $T_{\aesth}$ more than it increases the
inter-stratal compatibility grades, the net Harmony decreases, and
the repair is rejected.  The obstruction is therefore preserved at
the fixpoint.
\end{proof}

\begin{remark}[The art of constraint satisfaction]
Poetry generation in CTT is not merely constraint satisfaction: it
is \emph{optimal} constraint satisfaction in a space where some
``violations'' are more valuable than their repairs.  This is the
formal content of the poetic principle that great poetry often
violates expectations in productive ways---what the Russian
Formalists called \emph{defamiliarization} (ostranenie).
\end{remark}

\chapter{The Curry--Howard--Harmony Correspondence}
\label{ch:chh}

The Curry--Howard correspondence---the observation that propositions
are types and proofs are programs---is one of the deepest structural
insights of twentieth-century logic.  In this chapter we extend it
to a three-way correspondence that incorporates the graded, stratal,
and cohomological structure of Judgment Geometry.

\section{The Classical Correspondence}
\label{sec:classical-ch}

We briefly recall the classical Curry--Howard correspondence to fix
notation and set the stage for the extension.

\begin{center}
\renewcommand{\arraystretch}{1.2}
\begin{tabular}{ll}
\textbf{Logic} & \textbf{Type Theory} \\
\hline
Proposition $A$ & Type $A$ \\
Proof of $A$ & Term $e : A$ \\
$A \wedge B$ & $A \times B$ (product type) \\
$A \vee B$ & $A + B$ (sum type) \\
$A \Rightarrow B$ & $A \to B$ (function type) \\
$\forall x.\, P(x)$ & $\Pi(x:A).\, B(x)$ (dependent product) \\
$\exists x.\, P(x)$ & $\Sigma(x:A).\, B(x)$ (dependent sum) \\
$A = B$ & $\mathrm{Id}_A(a,b)$ (identity type) \\
$\top$ (truth) & $\mathbf{1}$ (unit type) \\
$\bot$ (falsity) & $\mathbf{0}$ (empty type) \\
Cut elimination & $\beta$-reduction \\
Consistency & Type safety (progress + preservation) \\
\end{tabular}
\end{center}

The correspondence is a bijection: every derivation in intuitionistic
natural deduction corresponds to a typed lambda term, and every
$\beta$-reduction step corresponds to a cut elimination step.  This
bijection is the foundation of proof assistants, program extraction,
and constructive mathematics.

\section{The Extension to Harmony}
\label{sec:chh-extension}

The classical correspondence is binary: a proposition is either
provable or not; a type is either inhabited or not.  JuGeo adds two
new dimensions---\emph{quality} (grades) and \emph{obstruction
detection} (cohomology)---that have no classical counterpart.  The
Curry--Howard--Harmony (CHH) correspondence extends the table:

\begin{center}
\renewcommand{\arraystretch}{1.2}
\begin{tabular}{lll}
\textbf{Logic} & \textbf{Type Theory} & \textbf{JuGeo / CTT} \\
\hline
Proposition & Type & Specification $A$ \\
Proof & Program & Evidence $E$ (GluingData section) \\
$\wedge$ & $\times$ & $\otG$ (cooperation) \\
$\vee$ & $+$ & $\opG$ (competition) \\
$\Rightarrow$ & $\to$ & Stratal dependency \\
$\forall$ & $\Pi$ & $\Pi_g$ (graded universal) \\
$\exists$ & $\Sigma$ & $\Sigma_g$ (graded existential) \\
$=$ & $\mathrm{Id}$ & $\GId_g$ (graded identity) \\
$\top$ & $\mathbf{1}$ & $\gone$ (perfect grade) \\
$\bot$ & $\mathbf{0}$ & $\gzero$ (failure grade) \\
Cut elimination & $\beta$-reduction & Descent \\
Consistency & Type safety & $H^1 = 0$ \\
--- & --- & Grade $T$ (quality) \\
--- & --- & Harmony $h$ (global quality) \\
\end{tabular}
\end{center}

The last two rows have no classical analogue: the grade and Harmony
are genuinely new data that arise from the graded semiring structure.

\begin{theorem}[Curry--Howard--Harmony correspondence]
\label{thm:chh}
The following identities hold:
\begin{enumerate}[label=(\roman*)]
\item \textbf{Harmony computation is type-checking.}  Computing the
  Harmony $h(\varphi) = \mathrm{Harm}(\sigma_\varphi)$ is
  equivalent to evaluating the grade of the typing judgment
  $\Gamma \vdash_g \sigma_\varphi : \mathrm{Sec}(\pi_s)$, where
  $\mathrm{Sec}(\pi_s)$ is the type of sections of the stratal
  fibration.

\item \textbf{Descent is proof search.}  The descent procedure
  (Construction~\ref{con:descent-repair}) is an algorithm for
  constructing a section of the stratal fibration---i.e., a proof
  term for the specification type.

\item \textbf{$\FMLM$ is constructive existence.}  The $\FMLM$
  fixpoint construction (Construction~\ref{con:fmlm}) is a
  constructive proof of the existence of an inhabitant of the
  target type: the fixpoint $E_\infty$ is the witness.

\item \textbf{Obstructions are preserved $H^1$ classes.}  A
  non-trivial $[\omega] \in H^1$ that is not repaired by descent
  corresponds to a non-trivial inhabitant of the graded identity
  type $\GId_g(A, a, b)$ with $g < \gone$---evidence that $a$ and
  $b$ are ``not quite the same'' in a way that is structurally
  meaningful.

\item \textbf{Grade semiring is decategorification.}  The grade
  $g \in \Grade$ of a typing judgment is the decategorification
  (forgetful image) of the proof term $E$: if we forget the
  structure of $E$ and retain only its quality, we obtain $g$.
\end{enumerate}
\end{theorem}

\begin{proof}
\emph{(i)}
By Definition~\ref{def:harmony}, $h(\varphi) = \bigotimes_s T_s
\otG \bigotimes_{s \preceq s'} G_{s,s'}$.  Each $T_s$ is the grade
of the typing judgment at stratum $s$, and each $G_{s,s'}$ is the
grade of the inter-stratal compatibility check.  The total Harmony
is therefore the grade of the composite judgment that the full
section is well-typed---which is type-checking.

\emph{(ii)}
The descent procedure iteratively modifies local sections to reduce
$H^1$.  Each modification changes the proof term; the goal of making
$H^1 = 0$ is precisely the goal of finding a globally consistent
proof.  By Theorem~\ref{thm:obstructions}, $H^1 = 0$ if and only
if the local sections glue to a global section---a proof.

\emph{(iii)}
The $\FMLM$ starts with a partial term (with masks) and iteratively
fills masks.  Each fill extends the partial proof.  By
Theorem~\ref{thm:fmlm-convergence}, the process converges to a
complete term $E_\infty$.  This term is a constructive witness for
the inhabited-ness of the type $A_{\mathrm{poem}}$: we have
exhibited an element.

\emph{(iv)}
A non-trivial $H^1$ class $[\omega]$ means that local sections agree
at grade $g < \gone$ on their overlap: there exists a graded identity
$p : \GId_g(\ldots)$ witnessing the approximate agreement but no
identity at grade $\gone$.  In poetry, $[\omega]$ is preserved when
it contributes to aesthetic value (Proposition~\ref{prop:preserved-obstructions}).

\emph{(v)}
The map $\Gamma \vdash_g E : A \;\mapsto\; g$ is the
decategorification: it forgets the proof term and retains only the
grade.  This map is a semiring homomorphism from the space of
derivations (with $\otG$ and $\opG$ on derivation trees) to
$(\Grade, \otG, \opG)$.
\end{proof}

\section{Subsumption Theorems}
\label{sec:subsumption}

We now prove that CTT strictly subsumes eight major formalisms.
Each subsumption is an embedding functor from the target formalism
into the graded type theory, and each embedding is faithful: it
preserves all derivations and their computational content.

\begin{theorem}[CTT $\supset$ Montague Grammar]
\label{thm:sub-montague}
Montague Grammar~\cite{montague1973} is a fragment of CTT obtained
by restricting to binary grades ($\Grade = \{\gzero, \gone\}$),
a single stratum ($\sem$), and simply-typed lambda calculus with
generalized quantifiers.
\end{theorem}

\begin{proof}
Define the embedding functor $F_{\mathrm{MG}} : \mathbf{MG} \to
\mathbf{CTT}$ as follows.  For each Montague type $\tau$, set
$F(\tau) = \tau$ (Montague's types are a subset of $\Type_{\sem}$).
For each Montague derivation
$\Gamma \vdash e : \tau$, set
$F(\Gamma \vdash e : \tau) = \Gamma \vdash_{\gone} e : \tau$.
The embedding is faithful because Montague derivations are precisely
the grade-$\gone$ derivations at the $\sem$ stratum.  The embedding
is strict because CTT also admits grades $g \neq \gone$ (e.g.\
metaphorical uses of type-shifting) and analyses at other strata.
\end{proof}

\begin{theorem}[CTT $\supset$ DRT/SDRT]
\label{thm:sub-drt}
Discourse Representation Theory~\cite{kamp1993} and its segmented
extension (SDRT) are fragments of CTT obtained by restricting to
binary grades, the discourse monad $\mathrm{DMon}$, and the $\disc$
stratum.
\end{theorem}

\begin{proof}
A DRS (set of discourse referents + conditions) is a state of the
discourse monad $\mathrm{DMon}$ with binary grades.  DRS
construction rules (introduction of referents, addition of
conditions) correspond to monadic operations $\mathsf{return}$ and
$\mathsf{bind}$.  SDRT's discourse relations correspond to typed
morphisms in the $\disc$-stratum fiber.  The embedding
$F_{\mathrm{DRT}} : \mathbf{DRT} \to \mathbf{CTT}$ maps each DRS
to a discourse-monad term and each accessibility relation to a
morphism in the fiber.  Faithfulness follows because DRT derivations
are in bijection with the binary-grade, single-stratum restriction
of CTT's monadic derivations.
\end{proof}

\begin{theorem}[CTT $\supset$ Priorian tense logic]
\label{thm:sub-prior}
Priorian tense logic (with operators $F$~``it will be the case
that'' and $P$~``it was the case that'') is a fragment of CTT
obtained by restricting to binary grades and binary temporal
operators within the $\sem$ stratum.
\end{theorem}

\begin{proof}
The Priorian operators $F$ and $P$ are modalities on the
$\sem$-stratum types.  Define
$F_{\mathrm{Prior}}(F\varphi) = \Sigma_{\gone}(t :
\mathrm{Time}).\, (t > \mathrm{now}) \times \varphi(t)$
and analogously for $P$.  This embeds tense logic into the graded
dependent sum at grade $\gone$.  Faithfulness follows because the
Priorian axioms ($Fp \to FFp$, etc.)\ are derivable as graded type
isomorphisms in CTT.
\end{proof}

\begin{theorem}[CTT $\supset$ MLTT]
\label{thm:sub-mltt}
Martin-L\"of Type Theory~\cite{mlm1984} is a fragment of CTT obtained
by restricting to binary grades and flat (un-stratified) universes.
\end{theorem}

\begin{proof}
The embedding $F_{\mathrm{MLTT}} : \mathbf{MLTT} \to \mathbf{CTT}$
maps each MLTT judgment $\Gamma \vdash e : A$ to
$\Gamma \vdash_{\gone} e : A$ in the $\sem$ stratum.  The
dependent products, sums, and identity types of MLTT are the
grade-$\gone$ specializations of the graded constructors
(Definitions~\ref{def:graded-pi}, \ref{def:graded-sigma},
\ref{def:gid}).  Computation rules ($\beta$, $\eta$) are preserved
because grade-$\gone$ computation is standard computation.
The embedding is faithful and conservative: it preserves
and reflects derivability.
\end{proof}

\begin{theorem}[CTT $\supset$ OT]
\label{thm:sub-ot}
Optimality Theory~\cite{prince1993} is a fragment of CTT obtained
by restricting to a single stratum and using $\opG$ for candidate
competition and $\otG$ for ranked constraint evaluation.
\end{theorem}

\begin{proof}
An OT grammar is a pair $(\mathrm{Gen}, \mathrm{Con})$ where
$\mathrm{Gen}$ generates candidates and $\mathrm{Con}$ is a ranked
constraint set.  Map each candidate $c$ to a term $E_c$ and each
constraint $C_i$ to a theory $\mathcal{T}_i$ at the relevant
stratum with grade
$\alpha_{\mathcal{T}_i}(c) = -w_i \cdot \mathrm{violations}_i(c)$
where $w_i$ encodes the ranking (higher-ranked constraints have
larger weights).  The OT evaluation
$\mathrm{Eval}(c) = \bigoplus_{c'} \bigotimes_i
\alpha_{\mathcal{T}_i}(c')$
is the Harmony computation restricted to one stratum.
Faithfulness follows because the OT total order on candidates is
recovered by the total order on Harmony values.
\end{proof}

\begin{theorem}[CTT $\supset$ Linear Logic]
\label{thm:sub-linear}
Girard's linear logic~\cite{girard1987} is a fragment of CTT
obtained by restricting to binary grades and interpreting the linear
resource discipline as a graded constraint on usage.
\end{theorem}

\begin{proof}
The embedding uses the grade to track resource usage.
Define $F_{\mathrm{LL}}(A \multimap B) = \Pi_{\gone}(x : A).\, B$
with the constraint that $x$ is used exactly once (enforced by a
theory $\mathcal{T}_{\mathrm{lin}}$ at the $\synt$ stratum that
penalizes non-linear usage with grade $\gzero$).
The multiplicative connectives $\otimes$ and $\oplus$ map to
$\otG$ and $\opG$ respectively.  The exponential $!A$ maps to
unrestricted use (the linear constraint is relaxed).  Faithfulness
follows because the cut-elimination steps of linear logic correspond
to $\beta$-reductions that preserve the linear usage constraint.
\end{proof}

\begin{theorem}[CTT $\supset$ Generative Lexicon]
\label{thm:sub-gl}
Pustejovsky's Generative Lexicon~\cite{pustejovsky1995} is a
fragment of CTT obtained by restricting to binary grades and the
qualia type constructor (Definition~\ref{def:qualia}).
\end{theorem}

\begin{proof}
Each GL lexical entry is a qualia structure
$\mathrm{Qualia}(C) = (f, c, t, a)$ with formal, constitutive,
telic, and agentive roles.  The GL operations of type coercion
and co-composition correspond to graded type coercions:
$\mathrm{coerce}(e, A, B) : \Sigma_g(x : B).\, \GId_g(A, B)$.
The grade of the coercion measures the naturalness of the
semantic shift.  At binary grades, this reduces to Pustejovsky's
binary well-formedness judgments.
\end{proof}

\begin{theorem}[CTT $\supset$ HoTT identity types]
\label{thm:sub-hott}
The identity types of Homotopy Type Theory~\cite{hott2013} are the
$g = \gone$ specialization of graded identity types.
\end{theorem}

\begin{proof}
By Definition~\ref{def:gid}, $\GId_{\gone}(A, a, b)$ has the same
formation, introduction ($\mathsf{refl}$), elimination (transport
$\mathsf{tr}$), and computation ($\beta$-rule) as the HoTT identity
type $\mathrm{Id}_A(a,b)$.  The graded groupoid structure
(Proposition~\ref{prop:graded-groupoid}) at $g = \gone$ recovers
the $\infty$-groupoid structure of HoTT.  Univalence
(Definition~\ref{def:graded-univalence}) at $g = \gone$ recovers
the univalence axiom.  The embedding is faithful because
$\GId_{\gone}$ derivations are in bijection with $\mathrm{Id}_A$
derivations.
\end{proof}

\section{Each Subsumption as a Functor}
\label{sec:subsumption-functors}

The eight subsumption theorems share a common structure: each defines
a functor from the subsumed formalism to CTT.

\begin{definition}[Subsumption functor]
\label{def:subsumption-functor}
A \textbf{subsumption functor} $F : \mathbf{T} \to \mathbf{CTT}$ from
a formalism $\mathbf{T}$ to CTT consists of:
\begin{enumerate}[label=(\alph*)]
\item a map on types:\ $F_{\mathrm{type}} : \mathrm{ob}(\mathbf{T})
  \to \mathrm{ob}(\mathbf{CTT})$;
\item a map on derivations:\ for each derivation
  $\mathcal{D} : \Gamma \vdash_{\mathbf{T}} e : A$, a derivation
  $F(\mathcal{D}) : F(\Gamma) \vdash_{g} F(e) : F(A)$ in CTT;
\item preservation of composition:\ $F$ maps cut/substitution in
  $\mathbf{T}$ to the corresponding operation in CTT.
\end{enumerate}
\end{definition}

\begin{proposition}[Faithfulness]
\label{prop:faithfulness}
Each of the eight embedding functors $F_{\mathrm{MG}}$,
$F_{\mathrm{DRT}}$, $F_{\mathrm{Prior}}$, $F_{\mathrm{MLTT}}$,
$F_{\mathrm{OT}}$, $F_{\mathrm{LL}}$, $F_{\mathrm{GL}}$,
$F_{\mathrm{HoTT}}$ is \textbf{faithful}: it is injective on
hom-sets.  That is, distinct derivations in $\mathbf{T}$ map to
distinct derivations in CTT.
\end{proposition}

\begin{proof}
In each case, the restriction of CTT to the relevant fragment
(binary grades, single stratum, specific type constructors) recovers
the original formalism exactly.  This restriction is a retraction:
composing the embedding with the restriction yields the identity on
$\mathbf{T}$.  A functor with a retraction is faithful.
\end{proof}

\begin{remark}[Strictness of subsumption]
Each subsumption is \emph{strict}: the image of $F$ is a proper
subcategory of $\mathbf{CTT}$.  The ``extra'' structure in CTT
includes:
\begin{itemize}[nosep]
\item continuous grades (all eight formalisms use binary logic);
\item multiple strata (Montague, DRT, OT, GL, HoTT are
  single-stratum);
\item inter-stratal obstructions (none of the eight formalisms have
  cohomological obstruction classes);
\item the domain-specific type constructors ($\mathrm{Jak}$,
  $\mathrm{Ico}$, $\mathrm{PF}$, $\mathrm{Rasa}$, $\mathrm{Blend}$,
  etc.).
\end{itemize}
CTT is thus a genuinely new framework that properly extends all
eight of its predecessors.
\end{remark}

\begin{remark}[The functor perspective]
Viewing each subsumption as a functor has practical value: it means
that any result proved in Montague Grammar, DRT, OT, or any other
subsumed formalism is \emph{automatically a theorem of CTT},
transported along the embedding functor.  This ``import'' mechanism
ensures that CTT inherits the full body of results from its
constituent traditions---it does not start from zero.
\end{remark}

\begin{remark}[Open questions]
Several natural questions remain open:
\begin{enumerate}[label=(\roman*)]
\item Is there a \emph{universal} subsumption functor that
  characterizes CTT as the ``free completion'' of the subsumed
  formalisms?
\item Can the eight embeddings be assembled into a \emph{colimit}
  diagram whose colimit is CTT?
\item What is the precise $2$-categorical structure of the category
  of subsumption functors?
\end{enumerate}
These questions point toward a deeper understanding of CTT's
universal-algebraic status and are left for future work.
\end{remark}

\bigskip
\noindent\rule{\textwidth}{0.4pt}

\medskip
This concludes Part~III.  We have developed the Compositional Theory
of Theories from first principles: theories as typed objects
(Chapter~\ref{ch:theory-of-theories}), graded dependent types
(Chapter~\ref{ch:gdt}), the stratal fibration
(Chapter~\ref{ch:strata}), NLU/NLG as type-theoretic operations
(Chapter~\ref{ch:nlu-nlg}), and the Curry--Howard--Harmony
correspondence (this chapter).  In Parts~IV--VI we turn to detailed
proofs, worked examples, and applications.


% ── Part IV: Proofs ───────────────────────────────────────────
\part{Formal Proofs}
% Part IV: Formal Proofs
% This file is \include'd from main.tex

%%%%%%%%%%%%%%%%%%%%%%%%%%%%%%%%%%%%%%%%%%%%%%%%%%%%%%%%%%%%%%%%%%%%%%%%
%%  Chapter 18 — Proof of Equivalence
%%%%%%%%%%%%%%%%%%%%%%%%%%%%%%%%%%%%%%%%%%%%%%%%%%%%%%%%%%%%%%%%%%%%%%%%
\chapter{Equivalence of Structured and Component-Wise Views}
\label{ch:proof-equiv}

This chapter contains the full proof that the category of structured
(4-tuple) judgments $\mathbf{Judg}_4$ is equivalent to the category of
component-wise judgments $\mathbf{Judg}_8$.  We construct the relevant
categories, the flattening and structuring functors between them, and
verify all required natural isomorphisms.

%% ── §18.1  The Category Judg_8 ──────────────────────────────────────
\section{The category $\mathbf{Judg}_8$}
\label{sec:judg8}

\begin{definition}[Objects of $\mathbf{Judg}_8$]
\label{def:judg8-obj}
An object of $\mathbf{Judg}_8$ is a component tuple
\[
  J = (c,\,\varphi,\,A,\,E,\,O,\,B,\,T,\,\Pi)
\]
where:
\begin{enumerate}[label=(\roman*),nosep]
\item $c \in \ob(\mathcal{C})$ is a coordinate in a small category $\mathcal{C}$
  equipped with a Grothendieck topology~$\mathcal{J}$;
\item $\varphi \in \PSh(\mathcal{C})(c)$ is a local section of a
  presheaf $\mathcal{F}$ at~$c$;
\item $A$ is a type in a dependent type theory $\mathcal{T}$;
\item $E$ is a term such that $\Gamma \vdash E : A$ for some context~$\Gamma$;
\item $O = \{U_i \to X\}_{i \in I}$ is a covering sieve (obligation cover);
\item $B \in \check{H}^1(\mathcal{U}, \mathcal{F})$ is a \v{C}ech
  cohomology class (obstruction);
\item $T \in \Grade$ is a grade;
\item $\Pi$ is a provenance trail (a finite record of evidence sources
  and derivation steps).
\end{enumerate}
These satisfy the \emph{coherence conditions}:
\begin{enumerate}[label=(C\arabic*),nosep]
\item The cover~$O$ refines the topology~$\mathcal{J}$:
  $O \in \mathcal{J}(X)$.
\item The class~$B$ is computed from the presheaf~$\mathcal{F}$
  restricted to the cover~$O$.
\item The grade~$T$ is the evaluation of the trail:
  $T = \mu(\Pi)$, where $\mu$ is the grade-evaluation map.
\item The typing judgment $\Gamma \vdash E : A$ is derivable in~$\mathcal{T}$.
\end{enumerate}
\end{definition}

\begin{definition}[Morphisms of $\mathbf{Judg}_8$]
\label{def:judg8-mor}
A morphism $f : J \to J'$ in $\mathbf{Judg}_8$ is a component tuple
\[
  f = (f_c,\, f_\varphi,\, f_A,\, f_E,\, f_O,\, f_B,\, f_T,\, f_\Pi)
\]
where:
\begin{enumerate}[label=(\roman*),nosep]
\item $f_c : c \to c'$ is a morphism in $\mathcal{C}$;
\item $f_\varphi : \mathcal{F}(c') \to \mathcal{F}(c)$ is the
  restriction along $f_c$, satisfying $f_\varphi(\varphi') = \varphi$;
\item $f_A : A \to A'$ is a type morphism (a term $f_A : A \to A'$
  in context~$\Gamma$);
\item $f_E$ witnesses the transport:
  $\Gamma \vdash f_E : f_A(E) =_{A'} E'$;
\item $f_O$ is a refinement of covers: every $U_i' \in O'$ is covered by
  members of~$O$;
\item $f_B : \check{H}^1(\mathcal{U},\mathcal{F}) \to
  \check{H}^1(\mathcal{U}',\mathcal{F}')$ is the induced map on
  cohomology, satisfying $f_B(B) = B'$;
\item $f_T$ witnesses $T' \leq T$ (grade is non-increasing under morphisms);
\item $f_\Pi$ is an extension of trails: $\Pi' = f_\Pi(\Pi)$.
\end{enumerate}
Composition is componentwise. The identity morphism is the tuple of
identity maps.
\end{definition}

\begin{lemma}
\label{lem:judg8-cat}
$\mathbf{Judg}_8$ is a category.
\end{lemma}

\begin{proof}
We verify the category axioms.
\emph{Identity:} For each object~$J$, the tuple
$(\id_c, \id_{\mathcal{F}(c)}, \id_A, \mathrm{refl}_E, \id_O, \id_B, \mathrm{refl}_T, \id_\Pi)$
satisfies all eight morphism conditions, since identity maps preserve
all structure.

\emph{Composition:} Given $f : J \to J'$ and $g : J' \to J''$, define
$g \circ f$ componentwise:
$(g_c \circ f_c,\; f_\varphi \circ g_\varphi,\; g_A \circ f_A,\;
\ldots)$.
The reversal of order for~$f_\varphi$ follows from contravariance of
presheaves.  Each coherence condition is preserved because each
component morphism preserves its respective condition.

\emph{Associativity:} Follows from associativity of composition in each
ambient category ($\mathcal{C}$, $\Set$, the type theory, etc.).

\emph{Unit laws:} The identity tuple is a two-sided unit for
componentwise composition, since identity is a unit in each factor.
\end{proof}

%% ── §18.2  The Category Judg_4 ──────────────────────────────────────
\section{The category $\mathbf{Judg}_4$}
\label{sec:judg4}

\begin{definition}[Objects of $\mathbf{Judg}_4$]
\label{def:judg4-obj}
An object of $\mathbf{Judg}_4$ is a 4-tuple
\[
  \JJ = (\STop,\; \PStr,\; \CSpec,\; \GEval)
\]
where:
\begin{enumerate}[label=(\roman*),nosep]
\item $\STop = (\mathcal{C}, \mathcal{J}, \mathcal{F}, \Omega,
  a, \gamma, \tau)$ is a \emph{Semantic Topos}:
  a Grothendieck topos $\Sh(\mathcal{C}, \mathcal{J})$ with a
  distinguished presheaf~$\mathcal{F}$, subobject classifier~$\Omega$,
  sheafification~$a$, global-sections geometric morphism~$\gamma$,
  and Lawvere--Tierney topology~$\tau$.
\item $\PStr = (\Gamma, \mathcal{T}, E, A, \mathcal{D})$ is a
  \emph{Proof Structure}: a context~$\Gamma$, a dependent type
  theory~$\mathcal{T}$ with universes, a derivation tree~$\mathcal{D}$
  witnessing $\Gamma \vdash E : A$, together with reduction rules
  ($\beta$, $\eta$) and the full apparatus of dependent types
  ($\Pi$, $\Sigma$, $\mathrm{Id}$, $W$-types).
\item $\CSpec = (\mathcal{U}, C^\bullet, H^\bullet, \partial, \mathrm{SS})$
  is a \emph{Cohomological Spectrum}: a cover~$\mathcal{U}$, the full
  \v{C}ech complex $C^\bullet(\mathcal{U}, \mathcal{F})$, the
  cohomology groups $H^n$ for all $n \geq 0$, connecting maps, and
  the \v{C}ech-to-derived spectral sequence~$\mathrm{SS}$.
\item $\GEval = (G, \opG, \otG, \gzero, \gone, \mu, \multimap, \leq)$
  is a \emph{Graded Evaluation}: a commutative unital quantale $G$
  with evaluation functor~$\mu$, residuation~$\multimap$, and the
  provenance-tracking structure.
\end{enumerate}
These satisfy the \emph{compatibility conditions}:
\begin{enumerate}[label=(K\arabic*),nosep]
\item The cover~$\mathcal{U}$ in $\CSpec$ refines the
  topology~$\mathcal{J}$ in $\STop$.
\item The presheaf~$\mathcal{F}$ in $\STop$ is the coefficient sheaf
  of the \v{C}ech complex in $\CSpec$.
\item The type~$A$ in $\PStr$ classifies sections of~$\mathcal{F}$:
  global sections of~$\mathcal{F}$ correspond to closed terms of type~$A$.
\item The evaluation $\mu : \mathcal{D} \to G$ in $\GEval$ takes
  derivation trees from~$\PStr$ to grades in~$G$.
\end{enumerate}
\end{definition}

\begin{definition}[Morphisms of $\mathbf{Judg}_4$]
\label{def:judg4-mor}
A morphism $F : \JJ \to \JJ'$ in $\mathbf{Judg}_4$ is a 4-tuple
\[
  F = (F_{\STop},\; F_{\PStr},\; F_{\CSpec},\; F_{\GEval})
\]
where:
\begin{enumerate}[label=(\roman*),nosep]
\item $F_{\STop} : \STop \to \STop'$ is a geometric morphism of
  topoi, i.e.\ an adjoint pair $(f^* \dashv f_*)$ with $f^*$
  left-exact, mapping $\mathcal{F} \mapsto \mathcal{F}'$;
\item $F_{\PStr} : \PStr \to \PStr'$ is a context morphism: a
  substitution $\sigma : \Gamma \to \Gamma'$ such that
  $E'[\sigma] \equiv E$ and $A'[\sigma] \equiv A$ (up to
  definitional equality);
\item $F_{\CSpec} : \CSpec \to \CSpec'$ is a chain map
  $C^\bullet(\mathcal{U}, \mathcal{F}) \to
  C^\bullet(\mathcal{U}', \mathcal{F}')$ compatible with coboundary maps;
\item $F_{\GEval} : \GEval \to \GEval'$ is a quantale homomorphism
  $\alpha : G \to G'$ preserving $\opG$, $\otG$, $\gzero$, $\gone$,
  and satisfying $\alpha \circ \mu = \mu' \circ F_{\PStr}$.
\end{enumerate}
These satisfy: $F_{\CSpec}$ is the chain map induced by $F_{\STop}$
on \v{C}ech complexes, and $F_{\GEval}$ is compatible with
$F_{\PStr}$ via the evaluation functors.

Composition is componentwise (composing geometric morphisms, composing
substitutions, composing chain maps, composing quantale homomorphisms).
The identity morphism is the tuple of identity maps.
\end{definition}

\begin{lemma}
\label{lem:judg4-cat}
$\mathbf{Judg}_4$ is a category.
\end{lemma}

\begin{proof}
Composition of geometric morphisms is a geometric morphism
\cite{johnstone2002}.  Composition of substitutions is a substitution.
Composition of chain maps is a chain map.  Composition of quantale
homomorphisms is a quantale homomorphism.  Compatibility conditions
(K1)--(K4) are preserved because each component functor preserves the
relevant structure.  Associativity and unit laws follow from those in
each ambient category.
\end{proof}

%% ── §18.3  The Forgetful Functor ────────────────────────────────────
\section{The forgetful functor $U : \mathbf{Judg}_4 \to \mathbf{Judg}_8$}
\label{sec:forgetful}

\begin{construction}[The functor $U$]
\label{con:forgetful}
Define $U$ as follows.

\emph{On objects.} Given $\JJ = (\STop, \PStr, \CSpec, \GEval)$, set
\[
  U(\JJ) = (c,\, \varphi,\, A,\, E,\, O,\, B,\, T,\, \Pi)
\]
where:
\begin{itemize}[nosep]
\item $c$ is any chosen base object of $\mathcal{C}$ (we fix a global
  choice of base point for the site; for pointed sites this is
  canonical);
\item $\varphi = \mathcal{F}(c)$ is the stalk of the presheaf at~$c$;
\item $(A, E)$ are the type and evidence term from $\PStr$;
\item $O = \mathcal{U}$ is the cover from $\CSpec$;
\item $B = [z] \in H^1(\mathcal{U}, \mathcal{F})$ is the first
  cohomology class from $\CSpec$;
\item $T = \mu(\mathcal{D})$ is the grade of the derivation;
\item $\Pi$ is the provenance record extracted from $\GEval$.
\end{itemize}

\emph{On morphisms.} Given $F = (F_{\STop}, F_{\PStr}, F_{\CSpec},
F_{\GEval})$, extract the underlying component maps:
\begin{itemize}[nosep]
\item $f_c$ is the image of the base point under the inverse image
  functor $f^*$;
\item $f_\varphi$ is the map on stalks induced by $F_{\STop}$;
\item $f_A$ is the type map induced by the substitution in $F_{\PStr}$;
\item $f_E$ is the transport witness from $F_{\PStr}$;
\item $f_O$ is the refinement of covers induced by $F_{\CSpec}$;
\item $f_B$ is the map $H^1(F_{\CSpec})$ on first cohomology;
\item $f_T$ is the grade map from $F_{\GEval}$;
\item $f_\Pi$ is the provenance extension.
\end{itemize}
\end{construction}

\begin{proposition}
\label{prop:U-functor}
$U : \mathbf{Judg}_4 \to \mathbf{Judg}_8$ is a functor.
\end{proposition}

\begin{proof}
\emph{Preservation of identities.}
$U(\id_\JJ) = (\id_c, \id_\varphi, \id_A, \mathrm{refl}_E,
\id_O, \id_B, \mathrm{refl}_T, \id_\Pi)$,
which is the identity in $\mathbf{Judg}_8$ by
\cref{def:judg8-mor}.

\emph{Preservation of composition.}
Let $F : \JJ \to \JJ'$ and $G : \JJ' \to \JJ''$.  Then
$U(G \circ F) = (g_c \circ f_c,\; f_\varphi \circ g_\varphi,\;
\ldots) = U(G) \circ U(F)$,
where the reversal in the second component follows from
contravariance of presheaf restriction and the fact that the inverse
image functor of the composite geometric morphism is the composite of
the inverse image functors.  Each remaining component is computed by
the functoriality of the relevant extraction (chain maps compose,
quantale homomorphisms compose, etc.).
\end{proof}

%% ── §18.4  The Structuring Functor ───────────────────────────────────
\section{The structuring functor $R : \mathbf{Judg}_8 \to \mathbf{Judg}_4$}
\label{sec:enrichment}

\begin{construction}[The functor $R$]
\label{con:enrichment}
Given a component tuple $J = (c, \varphi, A, E, O, B, T, \Pi)$ satisfying
the coherence conditions (C1)--(C4), we construct a 4-tuple
$R(J) = (\STop_J, \PStr_J, \CSpec_J, \GEval_J)$ as follows.

\paragraph{Semantic Topos $\STop_J$.}
The site $(\mathcal{C}, \mathcal{J})$ is given as part of the data.
The presheaf $\mathcal{F}$ is defined by $\varphi \in \mathcal{F}(c)$.
We set $\STop_J$ to be the \emph{smallest sub-topos} of
$\PSh(\mathcal{C})$ that:
\begin{enumerate}[label=(\alph*),nosep]
\item contains the presheaf $\mathcal{F}$,
\item has a topology that refines $\mathcal{J}$,
\item is closed under the internal logic operations
  (finite limits, exponentials, subobject classifier).
\end{enumerate}
Concretely, $\STop_J = \Sh(\mathcal{C}, \mathcal{J}_{\mathrm{min}})$
where $\mathcal{J}_{\mathrm{min}}$ is the coarsest topology on
$\mathcal{C}$ refining~$\mathcal{J}$ for which $\mathcal{F}$ is a
sheaf.  Such a topology exists by the intersection property of
Grothendieck topologies \cite{johnstone2002}.
The remaining components ($\Omega, a, \gamma, \tau$) are canonical:
$\Omega$ is the subobject classifier of $\Sh(\mathcal{C},
\mathcal{J}_{\mathrm{min}})$, $a$ is sheafification,
$\gamma$ is the global-sections geometric morphism to $\Set$,
and $\tau$ is the Lawvere--Tierney topology corresponding to
$\mathcal{J}_{\mathrm{min}}$.

\paragraph{Proof Structure $\PStr_J$.}
From $(A, E)$ and the derivability condition (C4), there exists a
context $\Gamma$ and a derivation tree $\mathcal{D}$ witnessing
$\Gamma \vdash E : A$.  We set $\PStr_J$ to be the
\emph{minimal type theory} containing:
\begin{enumerate}[label=(\alph*),nosep]
\item the types appearing in $\Gamma$ and $A$,
\item the term constructors used in~$E$,
\item identity types $\mathrm{Id}$ for all types in the theory,
\item closure under $\Pi$-types, $\Sigma$-types, and the universe
  $\Type_0$ containing~$A$.
\end{enumerate}
This is well-defined: the intersection of all type theories containing
the given data is again a type theory (closed under the type formation
rules).

\paragraph{Cohomological Spectrum $\CSpec_J$.}
The cover $O = \mathcal{U}$ and the presheaf $\mathcal{F}$ determine
the full \v{C}ech complex
\[
  C^0(\mathcal{U}, \mathcal{F}) \xrightarrow{\delta^0}
  C^1(\mathcal{U}, \mathcal{F}) \xrightarrow{\delta^1}
  C^2(\mathcal{U}, \mathcal{F}) \xrightarrow{\delta^2} \cdots
\]
We take $\CSpec_J$ to be this complex together with all its
cohomology groups $H^n = \ker \delta^n / \mathrm{im}\, \delta^{n-1}$.
The obstruction class~$B$ is recovered as $B = [z] \in H^1$.
The \v{C}ech-to-derived spectral sequence is the standard one
\cite{weibel1994}.

\paragraph{Graded Evaluation $\GEval_J$.}
The grade $T$ lives in a grade set $G$.  We equip $G$ with the
\emph{canonical quantale structure}: define $\opG = \max$,
$\otG = \min$ (the ``bottleneck'' semiring on a totally ordered set),
$\gzero = \bot$, $\gone = \top$.  The evaluation map $\mu$ is defined
by $\mu(\Pi) = T$ (extending to all derivation trees by the
recursive grade computation).
The residuation is $a \multimap b = \begin{cases} \gone & \text{if }
a \leq b, \\ b & \text{otherwise.}\end{cases}$
\end{construction}

\begin{proposition}
\label{prop:R-objects}
For every object $J$ of $\mathbf{Judg}_8$, the tuple $R(J)$ is a
well-defined object of $\mathbf{Judg}_4$.
\end{proposition}

\begin{proof}
We verify the compatibility conditions (K1)--(K4) of
\cref{def:judg4-obj}.

(K1): The cover $\mathcal{U}$ in $\CSpec_J$ is exactly the obligation
cover~$O$, which refines $\mathcal{J}$ by coherence condition (C1),
and $\mathcal{J}_{\mathrm{min}}$ refines~$\mathcal{J}$, so $\mathcal{U}$
refines $\mathcal{J}_{\mathrm{min}}$.

(K2): The presheaf $\mathcal{F}$ used in $\CSpec_J$ is the same
presheaf present in $\STop_J$ by construction.

(K3): The type~$A$ in $\PStr_J$ classifies sections of $\mathcal{F}$
because we have included in the minimal type theory exactly the types
needed to represent sections, and the evidence term~$E$ of type~$A$ is
a witness of a section at coordinate~$c$.

(K4): The evaluation $\mu(\mathcal{D}) = T$ by coherence condition (C3)
and the definition of $\mu$ in $\GEval_J$.
\end{proof}

\begin{construction}[$R$ on morphisms]
\label{con:R-morphisms}
Given a morphism $f = (f_c, f_\varphi, f_A, f_E, f_O, f_B, f_T,
f_\Pi) : J \to J'$ in $\mathbf{Judg}_8$, we define
$R(f) = (F_{\STop}, F_{\PStr}, F_{\CSpec}, F_{\GEval})$ as follows:
\begin{enumerate}[label=(\roman*),nosep]
\item $F_{\STop}$ is the geometric morphism induced by the site
  morphism $f_c : \mathcal{C} \to \mathcal{C}'$ (extended to a
  continuous functor between sites); the inverse image
  $f^*$ sends $\mathcal{F}'$ to $\mathcal{F}$ via $f_\varphi$.
\item $F_{\PStr}$ is the substitution $\sigma$ induced by
  $(f_A, f_E)$: define $\sigma(x) = f_A(x)$ for type variables and
  $f_E$ provides the transport witness.
\item $F_{\CSpec}$ is the chain map induced by $f_O$ (refinement of
  covers) and $f_\varphi$ (the presheaf map): on $n$-cochains,
  the map is the composite of restriction along the refinement
  and the presheaf map.
\item $F_{\GEval}$ is the quantale homomorphism
  $\alpha(g) = f_T(g)$ for $g \in G$, which preserves
  $\opG, \otG$ because $f_T$ is monotone and the quantale
  operations are order-theoretic.
\end{enumerate}
\end{construction}

\begin{proposition}
\label{prop:R-functor}
$R : \mathbf{Judg}_8 \to \mathbf{Judg}_4$ is a functor.
\end{proposition}

\begin{proof}
\emph{Preservation of identities.}  $R(\id_J) = (\id_{\STop_J},
\id_{\PStr_J}, \id_{\CSpec_J}, \id_{\GEval_J})$, since the identity
site morphism induces the identity geometric morphism, the identity
substitution is the identity on contexts, the identity refinement
induces the identity chain map, and the identity on grades is a
quantale homomorphism.

\emph{Preservation of composition.}  Let $f : J \to J'$ and
$g : J' \to J''$.  Each component of $R(g \circ f)$ is computed from
the composite $(g_c \circ f_c, \ldots)$.  The induced geometric
morphism of a composite site morphism is the composite of the
induced geometric morphisms \cite{johnstone2002}, so
$(R(g) \circ R(f))_{\STop} = R(g \circ f)_{\STop}$.  Similarly for
substitutions (composition of substitutions), chain maps (composition
of chain maps), and quantale homomorphisms.
\end{proof}

%% ── §18.5  U ∘ R ≅ id ──────────────────────────────────────────────
\section{The isomorphism $U \circ R \cong \id_{\mathbf{Judg}_8}$}
\label{sec:UR}

\begin{theorem}
\label{thm:UR}
There is a natural isomorphism $\eta : \id_{\mathbf{Judg}_8}
\Rightarrow U \circ R$.
\end{theorem}

\begin{proof}
Let $J = (c, \varphi, A, E, O, B, T, \Pi)$ be an object of
$\mathbf{Judg}_8$.  We must exhibit an isomorphism
$\eta_J : J \xrightarrow{\sim} (U \circ R)(J)$ and verify
naturality.

\emph{The component $\eta_J$.}  By construction,
$(U \circ R)(J) = U(\STop_J, \PStr_J, \CSpec_J, \GEval_J)$.
The forgetful functor~$U$ extracts:
\begin{itemize}[nosep]
\item the base point of $\STop_J$: this is $c$, since $\STop_J$
  was built from the site containing~$c$;
\item the stalk $\mathcal{F}(c)$: this contains $\varphi$ by
  construction of $\STop_J$;
\item $(A, E)$: unchanged, since $\PStr_J$ was the minimal theory
  containing $(A, E)$;
\item $O$: unchanged, since $\CSpec_J$ was built from the cover~$O$;
\item $B$: unchanged, since $H^1$ of the \v{C}ech complex built
  from $(O, \mathcal{F})$ recovers~$B$;
\item $T = \mu(\Pi)$: unchanged by (C3) and the definition of
  $\GEval_J$;
\item $\Pi$: unchanged.
\end{itemize}
Thus $(U \circ R)(J) = J$, and $\eta_J = \id_J$.

\emph{Naturality.}  For any morphism $f : J \to J'$ in
$\mathbf{Judg}_8$, the naturality square
\[
\begin{tikzcd}
  J \ar[r, "\eta_J"] \ar[d, "f"'] &
  (U \circ R)(J) \ar[d, "(U \circ R)(f)"] \\
  J' \ar[r, "\eta_{J'}"'] & (U \circ R)(J')
\end{tikzcd}
\]
commutes trivially, since $\eta_J = \id_J$ and
$\eta_{J'} = \id_{J'}$, so both paths equal~$f$ (using
$(U \circ R)(f) = f$, which follows from the fact that $U$ and $R$
act as identity on the underlying component data).
\end{proof}

%% ── §18.6  R ∘ U ≅ id ──────────────────────────────────────────────
\section{The isomorphism $R \circ U \cong \id_{\mathbf{Judg}_4}$}
\label{sec:RU}

This direction is more substantive: we must show that the canonical
enrichment of the extracted component data is isomorphic to the given
4-tuple.  The key insight is that the ``extra'' structure in a 4-tuple
object is \emph{determined up to canonical isomorphism} by its
underlying components.

\begin{theorem}
\label{thm:RU}
There is a natural isomorphism $\varepsilon : R \circ U
\Rightarrow \id_{\mathbf{Judg}_4}$.
\end{theorem}

\begin{proof}
Let $\JJ = (\STop, \PStr, \CSpec, \GEval)$ be an object of
$\mathbf{Judg}_4$.  We must exhibit an isomorphism
$\varepsilon_\JJ : (R \circ U)(\JJ) \xrightarrow{\sim} \JJ$.

\emph{Step 1: The topos component.}
$(R \circ U)(\JJ)$ has topos $\STop_{U(\JJ)}$, which is the smallest
sub-topos of $\PSh(\mathcal{C})$ containing the presheaf~$\mathcal{F}$
extracted from~$\STop$.  Now $\STop$ itself is a sub-topos of
$\PSh(\mathcal{C})$ containing~$\mathcal{F}$, and $\STop_{U(\JJ)}$ is
the \emph{smallest} such.  Therefore there is a canonical geometric
inclusion
\[
  \iota_{\STop} : \STop_{U(\JJ)} \hookrightarrow \STop.
\]
For this to be an isomorphism, we need $\STop$ to be generated by
$\mathcal{F}$ and the site data---that is, $\STop$ carries no
``extra'' sheaves beyond those determined by the site and the
presheaf.

We claim this holds for the objects of $\mathbf{Judg}_4$: by
compatibility condition (K2), the presheaf $\mathcal{F}$ in $\STop$
is the same as the coefficient sheaf in $\CSpec$, and by (K1), the
topology is determined by the covering data.  Since $\STop$ was
defined as the topos generated by the site and $\mathcal{F}$ (with
the Lawvere--Tierney topology $\tau$ being the unique one making
$\mathcal{F}$ a sheaf for $\mathcal{J}$), we have
$\STop_{U(\JJ)} \cong \STop$.  The isomorphism is the canonical
geometric morphism $\varepsilon_{\STop}$.

\emph{Step 2: The type-theory component.}
$(R \circ U)(\JJ)$ has proof structure $\PStr_{U(\JJ)}$, which is the
minimal type theory containing $(A, E, \Gamma)$ extracted
from~$\PStr$.  Now $\PStr$ contains $(A, E, \Gamma)$ and possibly
additional types and terms.  However, the minimality of $\PStr$ as a
proof structure for the judgment $\JJ$ (every component of $\PStr$ is
used in the derivation~$\mathcal{D}$ or required by the closure
conditions) ensures that $\PStr_{U(\JJ)}$ and $\PStr$ have the same
types and terms up to definitional equality.  Thus
$\varepsilon_{\PStr} : \PStr_{U(\JJ)} \xrightarrow{\sim} \PStr$
is an isomorphism of type theories.

\emph{Step 3: The cohomological component.}
$(R \circ U)(\JJ)$ has spectrum $\CSpec_{U(\JJ)}$, which is the
\v{C}ech complex of $(\mathcal{U}, \mathcal{F})$ with all cohomology
groups.  Since $\CSpec$ is \emph{defined} as this \v{C}ech complex
(compatibility conditions (K1)--(K2) ensure the cover and sheaf
agree), we have $\CSpec_{U(\JJ)} = \CSpec$ on the nose.
Set $\varepsilon_{\CSpec} = \id$.

\emph{Step 4: The grade component.}
$(R \circ U)(\JJ)$ has evaluation $\GEval_{U(\JJ)}$, which is the
canonical quantale on $G$ with $\mu(\Pi) = T$.  The given $\GEval$
has the same underlying set $G$ and the same evaluation $\mu$
(by (K4)).  The quantale structure is determined by the total
order on $G$ (since $\opG = \max$ and $\otG = \min$ are the only
idempotent quantale operations on a totally ordered set that have
$\gone = \top$ and $\gzero = \bot$).  Thus
$\varepsilon_{\GEval} : \GEval_{U(\JJ)} \xrightarrow{\sim} \GEval$
is an isomorphism.

\emph{Naturality.}  Let $F : \JJ \to \JJ'$ be a morphism in
$\mathbf{Judg}_4$.  We must show that the diagram
\[
\begin{tikzcd}
  (R \circ U)(\JJ) \ar[r, "\varepsilon_\JJ"]
    \ar[d, "(R \circ U)(F)"'] &
  \JJ \ar[d, "F"] \\
  (R \circ U)(\JJ') \ar[r, "\varepsilon_{\JJ'}"'] &
  \JJ'
\end{tikzcd}
\]
commutes.  Each component of $\varepsilon$ is either an identity or a
canonical inclusion/isomorphism.  Geometric morphisms compose
functorially, so $F_{\STop} \circ \varepsilon_{\STop} =
\varepsilon_{\STop'} \circ (R \circ U)(F)_{\STop}$ by
functoriality of the sub-topos inclusion.  The type-theory and
grade components are analogous: canonical isomorphisms commute with
structure-preserving maps.  The cohomological component is strictly
an identity, so the square commutes trivially.
\end{proof}

%% ── §18.7  Main Theorem and Corollaries ─────────────────────────────
\section{The equivalence theorem}
\label{sec:equiv-thm}

\begin{theorem}[Categorical equivalence of views]
\label{thm:equiv}
The functors $U : \mathbf{Judg}_4 \to \mathbf{Judg}_8$ and
$R : \mathbf{Judg}_8 \to \mathbf{Judg}_4$ exhibit an equivalence
of categories:
\[
  \mathbf{Judg}_4 \simeq \mathbf{Judg}_8.
\]
\end{theorem}

\begin{proof}
By \cref{thm:UR}, $U \circ R \cong \id_{\mathbf{Judg}_8}$.
By \cref{thm:RU}, $R \circ U \cong \id_{\mathbf{Judg}_4}$.
Hence $U$ and $R$ are quasi-inverse functors, establishing the
equivalence.
\end{proof}

\begin{corollary}[Transfer principle]
\label{cor:transfer}
Let $P$ be any property expressible in the internal language of
$\mathbf{Judg}_8$ (respectively $\mathbf{Judg}_4$).  Then $P$ holds
for a component tuple $J$ if and only if $P$ holds for the corresponding
4-tuple $R(J)$ (respectively, $P$ holds for a 4-tuple $\JJ$ if and
only if $P$ holds for $U(\JJ)$).
\end{corollary}

\begin{proof}
Equivalences of categories preserve and reflect all categorical
properties: limits, colimits, (co)equalizers, pullbacks, exponentials,
and internal-logic statements.  Since $U$ and $R$ form an equivalence,
any property of $J$ in $\mathbf{Judg}_8$ is transferred to $R(J)$ in
$\mathbf{Judg}_4$ via the natural isomorphism $\eta$, and conversely
via $\varepsilon$.
\end{proof}

\begin{corollary}[Faithfulness of $U$]
\label{cor:faithful}
The forgetful functor $U$ is faithful: distinct morphisms in
$\mathbf{Judg}_4$ have distinct underlying component-wise morphisms.
\end{corollary}

\begin{proof}
Any functor that is part of an equivalence is fully faithful.
Faithfulness follows immediately.
\end{proof}

\begin{remark}[Non-injectivity on objects]
\label{rem:non-inj}
Although $U$ is faithful on morphisms, it is not injective on objects
in general: distinct topoi $\STop \neq \STop'$ can project to the
same pair $(c, \varphi)$ if different Grothendieck topologies yield
the same stalks.  The equivalence resolves this through the enrichment
functor $R$, which selects the \emph{canonical} (minimal) enrichment.
Objects that project identically under $U$ are identified up to the
natural isomorphism~$\varepsilon$.
\end{remark}


%% ── Stub for remaining chapters ─────────────────────────────────────
%%%%%%%%%%%%%%%%%%%%%%%%%%%%%%%%%%%%%%%%%%%%%%%%%%%%%%%%%%%%%%%%%%%%%%%%
%%  Chapter 19 — Subsumption Theorems
%%%%%%%%%%%%%%%%%%%%%%%%%%%%%%%%%%%%%%%%%%%%%%%%%%%%%%%%%%%%%%%%%%%%%%%%
\chapter{Subsumption Theorems}
\label{ch:subsumption}

This chapter proves that the JuGeo 4-tuple framework properly
subsumes eight established formalisms.  For each, we give the precise
definition of the subsumed framework, construct a faithful embedding
into $\mathbf{Judg}_4$, prove that the embedding preserves the
relevant structure, and exhibit a concrete example.

\medskip

Throughout, we write $\mathbf{2} = \{\gzero, \gone\}$ for the
two-element Boolean grade semiring with $\opG = \lor$ and
$\otG = \land$, and we write $\mathbf{1}$ for the terminal
category (single object, identity morphism only).

%% ── §19.1  Montague Grammar ────────────────────────────────────────
\section{JuGeo subsumes Montague Grammar}
\label{sec:sub-montague}

\begin{definition}[Montague IL]
\label{def:montague}
Montague's Intensional Logic (IL) is a simply-typed $\lambda$-calculus
with types generated by base types $e$ (entities) and $t$ (truth
values), function types $\alpha \to \beta$, and an intension operator
$\hat{\phantom{x}}$.  A \emph{Montague model} is a triple
$(W, D, \llbracket\cdot\rrbracket)$ where $W$ is a set of possible
worlds, $D$ is a domain of entities, and
$\llbracket\cdot\rrbracket : \mathrm{Term} \to (W \to D_\alpha)$
is an interpretation function.
\end{definition}

\begin{construction}[Embedding $\iota_{\mathrm{MG}} :
  \mathbf{Mont} \to \mathbf{Judg}_4$]
\label{con:montague-embed}
Given a Montague model $(W, D, \llbracket\cdot\rrbracket)$ and a
typed term $M : \alpha$, define:
\begin{itemize}[nosep]
\item $\STop$: Let $\mathcal{C} = \mathcal{O}(W)$ be the poset of
  open subsets of~$W$ (with the usual topology).  The presheaf
  $\mathcal{F}(U) = \{f : U \to D_\alpha\}$ assigns to each open set
  the set of $\alpha$-typed functions on it.  Sheafification gives the
  possible-worlds topos $\Sh(W)$.

\item $\PStr$: The simply-typed $\lambda$-calculus with types built
  from $e$ and $t$ embeds into dependent type theory by setting
  $\Gamma = (w : W)$ and $A_\alpha = D_\alpha$.  The term $M$ becomes
  an evidence term $E = \lambda w.\, \llbracket M \rrbracket(w)$.

\item $\CSpec$: Since $\Sh(W)$ is a topos of sheaves on a
  topological space, and the presheaf $\mathcal{F}$ is already a
  sheaf (function sheaves are sheaves), we have $H^n = 0$ for all
  $n \geq 1$ on any open cover.  Set $\CSpec$ to be the trivial
  \v{C}ech complex.

\item $\GEval$: Set $G = \mathbf{2}$ (binary grades).  A sentence
  $\varphi$ has grade $\gone$ if $\llbracket\varphi\rrbracket(w) =
  \mathrm{true}$ and $\gzero$ otherwise.  The evaluation functor is
  $\mu(\mathcal{D}) = \gone$ if the derivation is valid, $\gzero$
  otherwise.
\end{itemize}
\end{construction}

\begin{theorem}
\label{thm:montague}
The embedding $\iota_{\mathrm{MG}}$ is faithful and preserves
truth conditions.
\end{theorem}

\begin{proof}
\emph{Faithfulness.}  Let $f, g : M \to M'$ be distinct morphisms in
$\mathbf{Mont}$ (i.e., distinct type-preserving maps between
Montague models).  The embedding maps these to geometric morphisms
$\iota(f), \iota(g) : \STop \to \STop'$ of the possible-worlds
topoi.  Since distinct model maps induce distinct inverse image
functors on the sheaf categories (they disagree on at least one
stalk), $\iota(f) \neq \iota(g)$.

\emph{Preservation of truth.}  In Montague semantics,
$\llbracket\varphi\rrbracket(w) = \mathrm{true}$ iff $\varphi$ is
true at world~$w$.  In the JuGeo embedding, this corresponds to the
global section $s(w) = \llbracket\varphi\rrbracket(w)$ being
non-empty, i.e., $\Gamma(\STop, \mathcal{F}_\varphi) \neq \emptyset$.
Since $H^1 = 0$ (trivial cohomology), the fundamental theorem gives
$P \models \varphi$ iff $H^0 \neq 0$, which is exactly the Montague
truth condition.
\end{proof}

\begin{example}
The sentence ``Every man walks'' has Montague translation
$\forall x.\, \mathrm{man}(x) \to \mathrm{walk}(x)$ of type~$t$.
In the embedding: $\STop = \Sh(W)$, $A = \Pi(x:e).\,
\mathrm{man}(x) \to \mathrm{walk}(x)$, $T = \gone$ if the
sentence is true at all worlds, and $H^1 = 0$ (no cohomological
obstruction).
\end{example}

%% ── §19.2  DRT/SDRT ────────────────────────────────────────────────
\section{JuGeo subsumes DRT/SDRT}
\label{sec:sub-drt}

\begin{definition}[DRT]
\label{def:drt}
A \emph{Discourse Representation Structure} (DRS) is a pair
$K = (\mathcal{R}, \mathcal{C})$ where $\mathcal{R}$ is a set of
discourse referents and $\mathcal{C}$ is a set of conditions.
A DRS is \emph{satisfied} by an embedding $f : \mathcal{R} \to D$
into a model domain $D$ if all conditions in $\mathcal{C}$ are true
under~$f$.  SDRT extends DRT with rhetorical relations between
discourse segments.
\end{definition}

\begin{construction}[Embedding $\iota_{\mathrm{DRT}} :
  \mathbf{DRT} \to \mathbf{Judg}_4$]
\label{con:drt-embed}
Given a DRS $K = (\mathcal{R}, \mathcal{C})$:
\begin{itemize}[nosep]
\item $\STop$: The site $\mathcal{C}_{\mathrm{disc}}$ has objects =
  discourse segments $d_1, \ldots, d_n$ and morphisms = anaphoric
  accessibility relations.  The topology is generated by covers
  $\{d_i \to d_{\mathrm{all}}\}$.  The presheaf
  $\mathcal{F}(d_k) = \{f : \mathcal{R}_k \to D\}$ assigns
  possible embeddings.

\item $\PStr$: Define a \emph{discourse monad} type constructor
  $\mathrm{DMon}(A)$ that encapsulates a set of discourse referents
  and a value of type $A$ depending on them.  A DRS becomes a
  term $E_K : \mathrm{DMon}(\mathrm{Prop})$ where
  $\mathrm{Prop} = \Sigma(c : \mathcal{C}).\, \mathrm{Sat}(c, f)$.
  Context: $\Gamma = (D : \Type,\; f : \mathcal{R} \to D)$.

\item $\CSpec$: The \v{C}ech complex is computed on the discourse
  site.  $H^1$ vanishes iff anaphoric references resolve consistently
  across all discourse segments.

\item $\GEval$: $G = \mathbf{2}$.  Grade $\gone$ iff all DRS
  conditions are satisfied.
\end{itemize}
For SDRT: rhetorical relations $R(d_i, d_j)$ become morphisms in the
site category, with the topology encoding which rhetorical structures
constitute complete discourse covers.
\end{construction}

\begin{theorem}
\label{thm:drt}
The embedding $\iota_{\mathrm{DRT}}$ is faithful and preserves
DRS satisfaction.
\end{theorem}

\begin{proof}
\emph{Faithfulness.}  Distinct DRS morphisms (embeddings $f \neq g$)
yield distinct context morphisms in $\PStr$ (the substitutions
$\sigma_f \neq \sigma_g$ disagree on at least one referent),
hence distinct $\mathbf{Judg}_4$ morphisms.

\emph{Preservation of satisfaction.}  A DRS $K$ is satisfied by
$f$ iff all conditions hold under~$f$.  In the embedding, this
corresponds to the existence of a global section of $\mathcal{F}$
over the discourse site---i.e., a compatible family of local
embeddings that agree on shared referents.  The fundamental
theorem identifies this with $H^1 = 0$.  Since DRT
satisfaction is precisely the existence of a consistent global
embedding, the two notions coincide.
\end{proof}

\begin{example}
``A man$_x$ entered.  He$_x$ smiled.''  The discourse site has two
segments $d_1, d_2$ with accessibility $d_1 \to d_2$.  The DRS has
$\mathcal{R} = \{x\}$, $\mathcal{C} = \{\mathrm{man}(x),
\mathrm{enter}(x), \mathrm{smile}(x)\}$.  Local sections assign
embeddings $f_1 : \{x\} \to D$ and $f_2 : \{x\} \to D$.
Compatibility on the overlap $d_1 \cap d_2$ requires $f_1(x) =
f_2(x)$, i.e., the pronoun ``he'' corefers with ``a man.''
\end{example}

%% ── §19.3  MLTT ────────────────────────────────────────────────────
\section{JuGeo subsumes MLTT}
\label{sec:sub-mltt}

\begin{definition}[MLTT]
\label{def:mltt}
Martin-L\"of Type Theory (MLTT) consists of:
\begin{enumerate}[label=(\roman*),nosep]
\item Dependent function types $\Pi(x:A).\,B(x)$;
\item Dependent pair types $\Sigma(x:A).\,B(x)$;
\item Identity types $\mathrm{Id}_A(a,b)$;
\item A universe hierarchy $\Type_0 : \Type_1 : \Type_2 : \cdots$;
\item Inductive types (including $\mathbb{N}$, $W$-types).
\end{enumerate}
Judgments have the form $\Gamma \vdash e : A$.
\end{definition}

\begin{construction}[Embedding $\iota_{\mathrm{MLTT}} :
  \mathbf{MLTT} \to \mathbf{Judg}_4$]
\label{con:mltt-embed}
Given an MLTT judgment $\Gamma \vdash e : A$:
\begin{itemize}[nosep]
\item $\STop$: The site is the terminal category $\mathbf{1}$ with
  the trivial (maximal) topology.  The topos
  $\Sh(\mathbf{1}) \cong \Set$.  The presheaf
  $\mathcal{F}(*) = \llbracket A \rrbracket$ is the set-theoretic
  interpretation of~$A$.

\item $\PStr$: Exactly the MLTT judgment $\Gamma \vdash e : A$ with
  its full derivation tree, type formers, and reduction rules.

\item $\CSpec$: The cover is the trivial cover $\{* \to *\}$.
  The \v{C}ech complex is $C^0 = \mathcal{F}(*)$,
  $C^n = 0$ for $n \geq 1$.  All higher cohomology vanishes:
  $H^n = 0$ for $n \geq 1$.

\item $\GEval$: $G = \mathbf{2}$.  Grade $\gone$ iff
  the judgment is derivable, $\gzero$ otherwise.
\end{itemize}
\end{construction}

\begin{theorem}
\label{thm:mltt}
The embedding $\iota_{\mathrm{MLTT}}$ is faithful, full on the binary
fragment, and preserves derivability.
\end{theorem}

\begin{proof}
\emph{Faithfulness.}  MLTT morphisms are substitutions
$\sigma : \Gamma \to \Gamma'$ with $A'[\sigma] \equiv A$ and
$e'[\sigma] \equiv e$.  These are exactly the context morphisms
in $\PStr$, hence distinct substitutions give distinct
$\mathbf{Judg}_4$ morphisms.

\emph{Fullness on the binary fragment.}  Every $\mathbf{Judg}_4$
morphism in the image of $\iota_{\mathrm{MLTT}}$ has trivial topos
component (since $\STop = \Set$), trivial cohomology, and binary
grade.  The only non-trivial data is the context morphism, which is
an MLTT substitution.  Hence every morphism in the binary-graded,
flat-site fragment of $\mathbf{Judg}_4$ comes from an MLTT morphism.

\emph{Preservation of derivability.}  $\Gamma \vdash e : A$ is
derivable in MLTT iff the embedded judgment has grade $\gone$, since
the grade is $\gone$ exactly when the derivation tree exists.
The cohomology vanishes trivially ($H^1 = 0$ on the terminal site),
so the fundamental theorem reduces to: a global section exists iff
a term of type $A$ exists in context $\Gamma$.
\end{proof}

%% ── §19.4  OT ──────────────────────────────────────────────────────
\section{JuGeo subsumes Optimality Theory}
\label{sec:sub-ot}

\begin{definition}[OT]
\label{def:ot}
An Optimality Theory grammar is a triple
$(\mathrm{Gen}, \mathrm{Con}, \gg)$ where $\mathrm{Gen}$ is a
generator function mapping inputs to candidate sets,
$\mathrm{Con} = \{C_1, \ldots, C_k\}$ is a finite set of ranked
constraints with $C_1 \gg C_2 \gg \cdots \gg C_k$, and the
\emph{optimal candidate} is
$\arg\min_{\mathrm{cand}} (C_1(\mathrm{cand}), \ldots,
C_k(\mathrm{cand}))$ under lexicographic ordering.
\end{definition}

\begin{construction}[Embedding $\iota_{\mathrm{OT}} :
  \mathbf{OT} \to \mathbf{Judg}_4$]
\label{con:ot-embed}
Given an OT grammar $(\mathrm{Gen}, \mathrm{Con}, \gg)$ and an
input~$i$:
\begin{itemize}[nosep]
\item $\STop$: The site has a single stratum (the phonological
  level).  Objects are candidates
  $\mathrm{Gen}(i) = \{o_1, \ldots, o_m\}$.
  Morphisms are identity only.  Topology: trivial.

\item $\PStr$: Context $\Gamma = (i : \mathrm{Input})$.
  For each candidate $o_j$, the type
  $A_j = \Pi(k : [K]).\, C_k(o_j) \leq C_k(o_l)$ (for all
  competitors $o_l$) encodes optimality.  The evidence term
  witnesses the constraint violations.

\item $\CSpec$: Trivial (single stratum, no non-trivial covers).

\item $\GEval$: The grade semiring is $G = \mathbb{N}^k$ with
  lexicographic order.  Define $\otG$ as componentwise addition and
  $\opG$ as componentwise minimum.  The grade of candidate $o_j$ is
  $\mu(o_j) = (C_1(o_j), \ldots, C_k(o_j))$.  The optimal
  candidate has the maximal grade under the reversed lexicographic
  order (fewest violations).

  Equivalently, assign exponential weights $w_k = N^{k-1}$ for
  large $N$ and compute
  $\mu(o_j) = \sum_k w_k \cdot C_k(o_j)$,
  reducing to a single-valued grade with lexicographic
  tie-breaking.
\end{itemize}
\end{construction}

\begin{theorem}
\label{thm:ot}
The embedding $\iota_{\mathrm{OT}}$ is faithful and the optimal
candidate in OT corresponds to the judgment with maximal harmony.
\end{theorem}

\begin{proof}
\emph{Faithfulness.}  Distinct OT grammars have distinct constraint
rankings, which yield distinct grade semirings (different orderings
on $\mathbb{N}^k$), hence distinct $\GEval$ components.

\emph{Optimality $=$ harmony.}  The OT optimal candidate minimizes
violation counts lexicographically.  In the JuGeo embedding, the
harmony of candidate~$o_j$ is
\[
  h(o_j) = \bigotimes_{k=1}^{K} w_k \cdot g_k(o_j),
\]
where $g_k(o_j) = C_k(o_j)$ is the constraint violation count.
The candidate with minimal violation vector has maximal harmony
(since $\otG$ is the bottleneck operation and we use the reversed
ordering).  Thus OT evaluation is a special case of $\GEval$
computation.
\end{proof}

\begin{example}
Input /pat/ with constraints
$\mathrm{Onset} \gg \mathrm{NoCoda} \gg \mathrm{Faith}$.
Candidates: [pat], [pa.ta], [pa].
Grade of [pat]: $(0, 1, 0)$.  Grade of [pa.ta]: $(0, 0, 1)$.
Optimal: [pa.ta] (fewer high-ranked violations).
In JuGeo: harmony $h$ is maximal for [pa.ta].
\end{example}

%% ── §19.5  Linear Logic ────────────────────────────────────────────
\section{JuGeo subsumes Linear Logic}
\label{sec:sub-linear}

\begin{definition}[Linear Logic]
\label{def:linear}
Girard's Linear Logic (LL) is a substructural logic with:
\begin{enumerate}[label=(\roman*),nosep]
\item Multiplicative connectives: $\otimes$ (tensor),
  $\multimap$ (linear implication), $\mathbf{1}$ (unit);
\item Additive connectives: $\oplus$ (plus), $\&$ (with),
  $\mathbf{0}$ (zero), $\top$;
\item Exponentials: $!A$ (of course), $?A$ (why not);
\item No unrestricted weakening or contraction---resources are
  tracked exactly.
\end{enumerate}
\end{definition}

\begin{construction}[Embedding $\iota_{\mathrm{LL}} :
  \mathbf{LL} \to \mathbf{Judg}_4$]
\label{con:linear-embed}
Given an LL sequent $\Gamma \vdash A$:
\begin{itemize}[nosep]
\item $\STop$: Terminal site $\mathbf{1}$ with $\Sh(\mathbf{1})
  \cong \Set$.

\item $\PStr$: The LL proof is embedded into dependent type theory
  via the standard Curry--Howard correspondence for linear logic:
  $A \otimes B$ becomes $\Sigma_{\mathrm{lin}}(x:A).\,B$,
  $A \multimap B$ becomes $\Pi_{\mathrm{lin}}(x:A).\,B$,
  and the exponential $!A$ becomes the unrestricted modality.
  The context~$\Gamma$ tracks linear resources.

\item $\CSpec$: Trivial (flat site).

\item $\GEval$: The grade semiring captures resource usage.
  Set $G = \mathbb{N}$ (counting resource uses).
  $\otG = +$ (combining resources), $\opG = \max$ (alternative
  resource usage), $\gone = 0$ (no resources used),
  $\gzero = \infty$ (impossible).
  The LL constraint ``use each resource exactly once'' becomes
  the grade condition $\mu(\mathcal{D}) = 1$ for each linear
  hypothesis.

  Alternatively, for the full multiplicative-additive structure:
  $G = (\mathbb{N}^n, +, \max, \vec{0}, \vec{\infty})$ where $n$
  counts the number of distinct resources, and the grade tracks
  the exact usage vector.
\end{itemize}
\end{construction}

\begin{theorem}
\label{thm:linear}
The embedding $\iota_{\mathrm{LL}}$ is faithful and preserves
provability.
\end{theorem}

\begin{proof}
\emph{Faithfulness.}  LL proofs are derivation trees with strict
resource discipline.  Distinct proofs have distinct resource usage
patterns, hence distinct grade evaluations.  Since the $\PStr$
component records the full derivation and the $\GEval$ component
records the resource usage, distinct LL morphisms yield distinct
$\mathbf{Judg}_4$ morphisms.

\emph{Preservation of provability.}  An LL sequent
$\Gamma \vdash A$ is provable iff there exists a derivation tree
using each linear hypothesis exactly once.  In the embedding, this
corresponds to the existence of a term $E$ of type $A$ in the
linear context $\Gamma$ with grade $\mu(\mathcal{D}) = \vec{1}$
(each resource used exactly once).  The grade condition
$\mu(\mathcal{D}) = \vec{1}$ is equivalent to the linear logic
resource discipline.
\end{proof}

%% ── §19.6  Generative Lexicon ──────────────────────────────────────
\section{JuGeo subsumes the Generative Lexicon}
\label{sec:sub-gl}

\begin{definition}[Generative Lexicon]
\label{def:gl}
Pustejovsky's Generative Lexicon (GL) assigns to each lexical item
a \emph{qualia structure} $\mathrm{Qualia}(w) = (Q_F, Q_C, Q_T,
Q_A)$ consisting of:
\begin{enumerate}[label=(\roman*),nosep]
\item $Q_F$: formal quale (the type/category);
\item $Q_C$: constitutive quale (parts/material);
\item $Q_T$: telic quale (purpose/function);
\item $Q_A$: agentive quale (origin/creation).
\end{enumerate}
\emph{Type coercion} occurs when a verb selects an argument via
one quale (e.g., ``begin the book'' coerces ``book'' via $Q_T$
to ``reading the book'').
\end{definition}

\begin{construction}[Embedding $\iota_{\mathrm{GL}} :
  \mathbf{GL} \to \mathbf{Judg}_4$]
\label{con:gl-embed}
Given a lexical item $w$ with qualia $\mathrm{Qualia}(w)$:
\begin{itemize}[nosep]
\item $\STop$: The site is the lexical category site, with objects
  = lexical items and morphisms = sense relations (hypernymy,
  meronymy, etc.).

\item $\PStr$: Define the \emph{qualia type constructor}
  \[
    \mathrm{Qualia}(C) = \Sigma(q_F : \mathrm{Form}(C)).\,
    \Sigma(q_C : \mathrm{Const}(C)).\,
    \Sigma(q_T : \mathrm{Telic}(C)).\,
    \mathrm{Agent}(C)
  \]
  as a dependent record type.  Coercion is a function
  $\mathrm{coerce}_T : \mathrm{Qualia}(C) \to C_T$ that
  projects the telic quale and interprets it as a type.

\item $\CSpec$: Trivial for single-word analysis.  For
  compositional semantics across a sentence, the cohomology
  detects coercion conflicts (incompatible qualia across
  composition).

\item $\GEval$: $G = \mathbf{2}$.  Grade $\gone$ if the qualia
  structure is well-formed and coercion succeeds.
\end{itemize}
\end{construction}

\begin{theorem}
\label{thm:gl}
The embedding $\iota_{\mathrm{GL}}$ is faithful and preserves
type coercion.
\end{theorem}

\begin{proof}
\emph{Faithfulness.}  Distinct GL lexical entries have distinct
qualia structures, which yield distinct $\Sigma$-types in $\PStr$.
Distinct coercion operations yield distinct projection functions.

\emph{Preservation of coercion.}  GL coercion ``begin the book''
$\rightsquigarrow$ ``begin reading the book'' is modeled by
selecting $Q_T(\text{book}) = \text{read}$ from the qualia structure.
In the JuGeo embedding, this is the projection
$\pi_T : \mathrm{Qualia}(\text{book}) \to \mathrm{Telic}(\text{book})$
applied to the term $E_{\text{book}} : \mathrm{Qualia}(\text{book})$.
The coerced type $\mathrm{Telic}(\text{book}) = \text{read}(\text{book})$
is the telic quale, and the derivation witnessing
$\Gamma \vdash \mathrm{begin}(\pi_T(E_{\text{book}})) :
\mathrm{Event}$ is exactly the GL coercion mechanism.
\end{proof}

%% ── §19.7  HoTT ───────────────────────────────────────────────────
\section{JuGeo subsumes HoTT}
\label{sec:sub-hott}

\begin{definition}[HoTT identity types]
\label{def:hott}
In Homotopy Type Theory, the identity type $\mathrm{Id}_A(a,b)$ for
$a, b : A$ satisfies:
\begin{enumerate}[label=(\roman*),nosep]
\item Reflexivity: $\mathrm{refl}_a : \mathrm{Id}_A(a,a)$;
\item Symmetry: $p : \mathrm{Id}_A(a,b) \vdash p^{-1} :
  \mathrm{Id}_A(b,a)$;
\item Transitivity: $p : \mathrm{Id}_A(a,b),\; q :
  \mathrm{Id}_A(b,c) \vdash p \cdot q : \mathrm{Id}_A(a,c)$;
\item These form an $\infty$-groupoid structure on identity proofs.
\end{enumerate}
\end{definition}

\begin{construction}[Embedding $\iota_{\mathrm{HoTT}} :
  \mathbf{HoTT} \to \mathbf{Judg}_4$]
\label{con:hott-embed}
Given a HoTT judgment $\Gamma \vdash e : A$:
\begin{itemize}[nosep]
\item $\STop$, $\CSpec$: As for MLTT (\cref{con:mltt-embed}):
  terminal site, trivial cohomology.

\item $\PStr$: The full HoTT type theory, including higher inductive
  types, univalence, and the $\infty$-groupoid structure.

\item $\GEval$: The graded identity type $\GId_g(A,a,b)$ at grade
  $g = \gone$ is set to be exactly the HoTT identity type:
  \[
    \GId_{\gone}(A,a,b) \;\coloneqq\; \mathrm{Id}_A(a,b).
  \]
  At sub-maximal grades $g < \gone$, $\GId_g(A,a,b)$ is a
  relaxed ``near-identity'' type (inhabited when $a$ and $b$ are
  $g$-close but not necessarily equal).  The grade semiring is
  $G = [0,1]$ with $\otG = \min$, $\opG = \max$.
\end{itemize}
\end{construction}

\begin{theorem}
\label{thm:hott}
The embedding $\iota_{\mathrm{HoTT}}$ is faithful and
$\GId_{\gone}(A,a,b) = \mathrm{Id}_A(a,b)$.
\end{theorem}

\begin{proof}
\emph{Faithfulness.}  Follows from faithfulness of the MLTT embedding
(\cref{thm:mltt}), since HoTT extends MLTT.

\emph{Identity type recovery.}
By definition, $\GId_{\gone}(A,a,b) = \mathrm{Id}_A(a,b)$.  We
verify that the graded groupoid structure at grade $\gone$
coincides with the HoTT groupoid structure:
\begin{enumerate}[label=(\roman*),nosep]
\item Reflexivity: $\mathrm{refl}_a : \GId_{\gone}(A,a,a)$ is
  exactly $\mathrm{refl}_a : \mathrm{Id}_A(a,a)$.
\item Symmetry: the graded inverse at $g = \gone$ is HoTT path
  inversion.
\item Transitivity: the graded composition at $g = \gone$ is HoTT
  path concatenation.
\item Higher structure: the graded groupoid laws at $g = \gone$
  are the HoTT groupoid laws (associativity, left/right unit,
  inverse laws hold up to higher identity).
\end{enumerate}
Thus the HoTT $\infty$-groupoid structure is the $g = \gone$ slice
of the graded groupoid.
\end{proof}

\begin{remark}
At grades $g < \gone$, $\GId_g$ provides structure not present in
HoTT: approximate or graded identifications.  This is what makes
the JuGeo framework a \emph{proper} extension (\cref{thm:proper}).
\end{remark}

%% ── §19.8  Prior's Tense Logic ─────────────────────────────────────
\section{JuGeo subsumes Prior's Tense Logic}
\label{sec:sub-prior}

\begin{definition}[Tense Logic]
\label{def:prior}
Prior's basic tense logic extends propositional logic with two
modal operators:
\begin{enumerate}[label=(\roman*),nosep]
\item $P\varphi$ (``it was the case that $\varphi$'');
\item $F\varphi$ (``it will be the case that $\varphi$'').
\end{enumerate}
Models are triples $(T, <, V)$ where $(T, <)$ is a strict linear
order (time) and $V : \mathrm{Prop} \to \mathcal{P}(T)$ is a
valuation.
\end{definition}

\begin{construction}[Embedding $\iota_{\mathrm{TL}} :
  \mathbf{TL} \to \mathbf{Judg}_4$]
\label{con:prior-embed}
Given a tense-logical model $(T, <, V)$ and a formula $\varphi$:
\begin{itemize}[nosep]
\item $\STop$: The site $\mathcal{C}_T$ has objects = intervals
  $[a,b] \subseteq T$ and morphisms = inclusions.  The topology is
  generated by covers $\{[a,c], [c,b]\} \to [a,b]$ (interval
  bisections).  The presheaf $\mathcal{F}([a,b]) = \{t \in [a,b]
  \mid V(\varphi, t) = \mathrm{true}\}$.

\item $\PStr$: The temporal operators become type constructors:
  $P(A) = \Sigma(t' : T).\, (t' < t) \times A(t')$ and
  $F(A) = \Sigma(t' : T).\, (t < t') \times A(t')$.  The
  context is $\Gamma = (t : T)$.

\item $\CSpec$: The \v{C}ech complex on the temporal site detects
  whether a temporally local property holds globally over a time
  interval.

\item $\GEval$: $G = \mathbf{2}$.  Grade $\gone$ iff $\varphi$ is
  true at the evaluation time.
\end{itemize}
\end{construction}

\begin{theorem}
\label{thm:prior}
The embedding $\iota_{\mathrm{TL}}$ is faithful and preserves
tense-logical validity.
\end{theorem}

\begin{proof}
\emph{Faithfulness.}  Distinct tense models $(T, <, V) \neq
(T', <', V')$ yield distinct sites (different interval categories or
different valuations), hence distinct $\STop$ components.  Distinct
tense formulas yield distinct types in $\PStr$.

\emph{Preservation of validity.}  A tense formula $\varphi$ is valid
at time $t$ in model $(T, <, V)$ iff $V(\varphi, t) = \mathrm{true}$.
In the embedding, this is the existence of a global section of
$\mathcal{F}$ over the point $\{t\}$ (which is trivially a cover of
itself), i.e., $\mathcal{F}(\{t\}) \neq \emptyset$.

For temporally extended validity (``$\varphi$ holds throughout
$[a,b]$''), the fundamental theorem applies: $\varphi$ holds on all
of $[a,b]$ iff $H^1 = 0$ on the interval cover, iff local truth at
sub-intervals glues to global truth.

The tense operators are preserved by construction:
$P\varphi$ is true at $t$ iff
$\exists t' < t.\, V(\varphi, t')$ iff the $\Sigma$-type
$\Sigma(t':T).\,(t' < t) \times \mathcal{F}(\{t'\})$ is inhabited.
\end{proof}

%% ── §19.9  Proper Extension ────────────────────────────────────────
\section{JuGeo is a proper extension}
\label{sec:proper}

\begin{theorem}[Proper extension]
\label{thm:proper}
$\mathbf{Judg}_4$ is not equivalent to any of the subsumed
frameworks.  That is, the embeddings $\iota_{\mathrm{MG}},
\iota_{\mathrm{DRT}}, \iota_{\mathrm{MLTT}}, \iota_{\mathrm{OT}},
\iota_{\mathrm{LL}}, \iota_{\mathrm{GL}}, \iota_{\mathrm{HoTT}},
\iota_{\mathrm{TL}}$ are all proper (faithful but not essentially
surjective).
\end{theorem}

\begin{proof}
We exhibit a single judgment $\JJ^* \in \mathbf{Judg}_4$ that lies
outside the image of every embedding.  Define $\JJ^*$ as follows:
\begin{itemize}[nosep]
\item $\STop^*$: a non-trivial site with three strata
  (phonological, semantic, pragmatic) and covers linking all three;
\item $\PStr^*$: a graded judgment
  $\Gamma \vdash_{0.7} E : \mathrm{Poem}$ with fractional grade;
\item $\CSpec^*$: non-trivial $H^1 \neq 0$ (the poem has a forced
  rhyme creating a cohomological obstruction);
\item $\GEval^*$: the grade semiring $[0,1]$ with the bottleneck
  product, evaluating to $T = 0.7$ (the naturalness stratum is the
  bottleneck).
\end{itemize}

This judgment cannot lie in the image of:
\begin{enumerate}[label=(\roman*),nosep]
\item $\iota_{\mathrm{MG}}$: Montague has binary grades and no
  multi-stratal structure ($|G| = 2$, but $\JJ^*$ has $T = 0.7$).
\item $\iota_{\mathrm{DRT}}$: DRT has binary grades and no
  cohomological structure for poetry.
\item $\iota_{\mathrm{MLTT}}$: MLTT has trivial site and binary
  grades.
\item $\iota_{\mathrm{OT}}$: OT has a single stratum and integer
  violation counts (not continuous grades).
\item $\iota_{\mathrm{LL}}$: Linear Logic has resource-counting
  grades ($\mathbb{N}$), not continuous $[0,1]$, and trivial site.
\item $\iota_{\mathrm{GL}}$: GL has binary grades and no
  cohomological structure.
\item $\iota_{\mathrm{HoTT}}$: HoTT has trivial site and no
  multi-stratal harmony.
\item $\iota_{\mathrm{TL}}$: Tense Logic has binary grades and
  temporal site only (not multi-stratal).
\end{enumerate}
Since $\JJ^*$ uses continuous grades, non-trivial cohomology,
\emph{and} multi-stratal structure simultaneously, no single
subsumed framework can express it.
\end{proof}

%%%%%%%%%%%%%%%%%%%%%%%%%%%%%%%%%%%%%%%%%%%%%%%%%%%%%%%%%%%%%%%%%%%%%%%%
%%  Chapter 20 — Descent and Harmony Proofs
%%%%%%%%%%%%%%%%%%%%%%%%%%%%%%%%%%%%%%%%%%%%%%%%%%%%%%%%%%%%%%%%%%%%%%%%
\chapter{Descent and Harmony Proofs}
\label{ch:descent-proofs}

This chapter contains the full proofs of the descent theorem, the
harmony existence and uniqueness theorems, the bottleneck principle,
the $\FMLM$ convergence theorem, and the descent-as-optimisation
correspondence.

%% ── §20.1  The Descent Theorem ──────────────────────────────────────
\section{The descent theorem}
\label{sec:descent-thm}

We first establish the fundamental theorem of Judgment Geometry: the
equivalence between global validity and the vanishing of the first
\v{C}ech cohomology.

\begin{definition}[\v{C}ech complex]
\label{def:cech-complex}
Let $(\mathcal{C}, \mathcal{J})$ be a site, $\mathcal{F}$ a presheaf
on $\mathcal{C}$, and $\mathcal{U} = \{U_i \to X\}_{i \in I}$ a
covering family of an object $X$.  The \v{C}ech complex
$\Cech^\bullet(\mathcal{U}, \mathcal{F})$ is:
\[
  \Cech^n(\mathcal{U}, \mathcal{F})
  = \prod_{i_0 < \cdots < i_n}
    \mathcal{F}(U_{i_0} \times_X \cdots \times_X U_{i_n})
\]
with coboundary maps $\delta^n : \Cech^n \to \Cech^{n+1}$ given by
\[
  (\delta^n s)_{i_0 \ldots i_{n+1}}
  = \sum_{k=0}^{n+1} (-1)^k\, s_{i_0 \ldots \hat{i}_k \ldots i_{n+1}}
    \big|_{U_{i_0} \times_X \cdots \times_X U_{i_{n+1}}}.
\]
The cohomology groups are $\check{H}^n(\mathcal{U}, \mathcal{F})
= \ker \delta^n / \mathrm{im}\, \delta^{n-1}$.
\end{definition}

\begin{definition}[Descent datum]
\label{def:descent-datum}
A \emph{descent datum} for $\mathcal{F}$ on the cover $\mathcal{U}$
is a family of local sections $\{s_i \in \mathcal{F}(U_i)\}_{i \in I}$
satisfying the \emph{cocycle condition}: for all $i, j \in I$,
\[
  s_i|_{U_i \times_X U_j} = s_j|_{U_i \times_X U_j}.
\]
Equivalently, a descent datum is a 0-cocycle $s \in \ker \delta^0$.
\end{definition}

\begin{theorem}[Fundamental Theorem of Judgment Geometry]
\label{thm:fundamental}
Let $(\mathcal{C}, \mathcal{J})$ be a site, $\mathcal{F}$ a presheaf,
$\mathcal{U}$ a cover of $X$, and $\varphi$ a property encoded by
$\mathcal{F}$.  Then:
\[
  X \models \varphi
  \quad\Longleftrightarrow\quad
  \check{H}^1(\mathcal{U}, \mathcal{F}) = 0
  \;\;\text{and}\;\;
  \check{H}^0(\mathcal{U}, \mathcal{F}) \neq 0.
\]
In particular, if $\mathcal{F}$ is a sheaf, then $X \models \varphi$
if and only if a compatible family of local sections (each witnessing
$\varphi$ on its region) exists.
\end{theorem}

\begin{proof}
$(\Rightarrow)$\; Suppose $X \models \varphi$, i.e., there exists a
global section $s \in \mathcal{F}(X)$ witnessing~$\varphi$.
Restricting $s$ to each $U_i$ gives $s_i = s|_{U_i}$.  This family
is automatically compatible:
$s_i|_{U_{ij}} = s|_{U_{ij}} = s_j|_{U_{ij}}$, so $(s_i)$ is a
0-cocycle and $\check{H}^0 \neq 0$.

For $\check{H}^1 = 0$: let $z = (z_{ij}) \in \ker \delta^1$ be any
1-cocycle.  We must show $z$ is a coboundary.  The cocycle condition
gives $z_{ij} + z_{jk} = z_{ik}$ on triple overlaps.  Define
$t_i = z_{i0}$ for a fixed index~$0$.  Then
\[
  (\delta^0 t)_{ij} = t_j - t_i = z_{j0} - z_{i0} = z_{ij}
\]
by the cocycle condition $z_{i0} + z_{0j} = z_{ij}$ (using
$z_{0j} = -z_{j0}$ by antisymmetry).  Hence $z = \delta^0 t$, so
$\check{H}^1 = 0$.

$(\Leftarrow)$\;
Suppose $\check{H}^0 \neq 0$ and $\check{H}^1 = 0$.  The condition
$\check{H}^0 \neq 0$ gives a compatible family $(s_i)_{i \in I}$
with $s_i|_{U_{ij}} = s_j|_{U_{ij}}$ for all $i, j$.  If
$\mathcal{F}$ is a sheaf, the sheaf condition states precisely that
such a compatible family glues to a unique global section
$s \in \mathcal{F}(X)$.  Hence $X \models \varphi$.

If $\mathcal{F}$ is only a presheaf, the vanishing $\check{H}^1 = 0$
is an additional hypothesis ensuring that the obstruction to gluing
(which lives in $\check{H}^1$) vanishes.  Concretely: suppose we have
local sections $s_i$ that are \emph{almost} compatible, differing by
a discrepancy $d_{ij} = s_i|_{U_{ij}} - s_j|_{U_{ij}}$.  The
discrepancy $(d_{ij})$ is a 1-cocycle (it satisfies the cocycle
condition on triple overlaps by a telescoping argument).  If
$\check{H}^1 = 0$, then $(d_{ij}) = (\delta^0 t)_{ij}$ for some
$(t_i)$, and the corrected sections $s_i' = s_i - t_i$ are
compatible: $s_i'|_{U_{ij}} - s_j'|_{U_{ij}} = d_{ij} -
(t_i|_{U_{ij}} - t_j|_{U_{ij}}) = d_{ij} - d_{ij} = 0$.
The corrected family glues (by the sheaf condition applied to the
sheafification of $\mathcal{F}$).
\end{proof}

\begin{corollary}[Verification $= H^0$]
\label{cor:verification}
$X \models \varphi$ if and only if
$\Gamma(X, \mathcal{F}_\varphi) \neq \emptyset$
(the global sections functor is non-empty).  For an acyclic cover
(where \v{C}ech and derived cohomology agree), this is identified
with $\check{H}^0(\mathcal{U}, \mathcal{F}) \neq \emptyset$.
\end{corollary}

\begin{proof}
The global sections functor $\Gamma(X, -)$ is
$\Hom_{\PSh(\mathcal{C})}(\mathbf{y}(X), -)$.  A global section of
$\mathcal{F}_\varphi$ is precisely an element of $\mathcal{F}(X)$
witnessing~$\varphi$.  By \cref{thm:fundamental}, this exists iff
$H^0 \neq 0$ and $H^1 = 0$.  For acyclic covers, the Leray spectral
sequence degenerates: $\check{H}^n(\mathcal{U}, \mathcal{F}) \cong
H^n(X, \mathcal{F})$ for all $n$, so
$\Gamma(X, \mathcal{F}) \cong \check{H}^0$.
\end{proof}

\subsection{The \v{C}ech-to-derived spectral sequence}

\begin{proposition}[\v{C}ech-to-derived comparison]
\label{prop:cech-derived}
For any cover $\mathcal{U}$ of $X$ and any sheaf $\mathcal{F}$, there
is a spectral sequence
\[
  E_2^{p,q} = \check{H}^p(\mathcal{U}, \underline{H}^q(\mathcal{F}))
  \;\Longrightarrow\;
  H^{p+q}(X, \mathcal{F}),
\]
where $\underline{H}^q(\mathcal{F})$ is the presheaf
$U \mapsto H^q(U, \mathcal{F}|_U)$.  When the cover is acyclic
($H^q(U_{i_0 \cdots i_p}, \mathcal{F}) = 0$ for $q > 0$ and all
multi-intersections), this degenerates at $E_2$ and gives
$\check{H}^n(\mathcal{U}, \mathcal{F}) \cong H^n(X, \mathcal{F})$.
\end{proposition}

\begin{proof}
The double complex $C^{p,q} = \Cech^p(\mathcal{U}, \mathcal{I}^q)$,
where $\mathcal{F} \to \mathcal{I}^\bullet$ is an injective
resolution, has two filtrations.  Filtering by rows gives
$E_1^{p,q} = \Cech^p(\mathcal{U}, \underline{H}^q)$, whose
cohomology is $E_2^{p,q} = \check{H}^p(\mathcal{U},
\underline{H}^q)$.  Filtering by columns gives $E_1^{p,q} =
H^q(\prod U_{i_0 \cdots i_p}, \mathcal{F})$; when the cover is
acyclic, $E_1^{p,q} = 0$ for $q > 0$, so the column filtration
degenerates at $E_1$ and both filtrations converge to
$H^{p+q}(X, \mathcal{F})$.  The isomorphism
$\check{H}^n \cong H^n$ follows \cite{weibel1994}.
\end{proof}

%% ── §20.2  Harmony Existence ────────────────────────────────────────
\section{Harmony existence and uniqueness}
\label{sec:harmony}

\begin{definition}[Harmony function]
\label{def:harmony}
Let $\mathcal{T}_1, \ldots, \mathcal{T}_N$ be theories (theory slots)
with stratal grades $g_1, \ldots, g_N$ and weights
$w_1, \ldots, w_N \in \Grade$.  The \emph{harmony function}
$h : \mathrm{Sec}(\mathcal{F}) \to \Grade$ is defined by
\[
  h(\sigma) = \bigotimes_{k=1}^{N} w_k \otG g_k(\sigma),
\]
where $g_k(\sigma)$ is the grade of $\sigma$ with respect to
theory~$\mathcal{T}_k$, and $\mathrm{Sec}(\mathcal{F})$ is the set
of (candidate) sections of the judgment presheaf.
\end{definition}

\begin{definition}[Harmony lattice]
\label{def:harmony-lattice}
The \emph{harmony lattice} is the complete lattice
$(\mathrm{Sec}(\mathcal{F}), \sqsubseteq)$ where
$\sigma \sqsubseteq \sigma'$ iff $h(\sigma) \leq h(\sigma')$
(ordered by harmony value).
\end{definition}

\begin{theorem}[Harmony fixpoint existence]
\label{thm:harmony-fix}
Let $h : \mathrm{Sec}(\mathcal{F}) \to \Grade$ be the harmony
function.  Suppose:
\begin{enumerate}[label=(\roman*),nosep]
\item The number of strata is finite: $N < \infty$.
\item The evaluation $g_k$ is continuous (Scott-continuous) for each
  theory $\mathcal{T}_k$.
\item The grade semiring $(\Grade, \leq)$ is a complete lattice.
\end{enumerate}
Then the refinement operator
$R(\sigma) = \arg\max_{\sigma'} h(\sigma')$ subject to
$\sigma' \sqsubseteq \sigma$ has a fixpoint.
\end{theorem}

\begin{proof}
We apply the Knaster--Tarski fixpoint theorem.  The set
$\mathrm{Sec}(\mathcal{F})$ ordered by $\sqsubseteq$ forms a
complete lattice (it is a sub-lattice of the product
$\prod_{U \in \mathcal{U}} \mathcal{F}(U)$, which is complete because
each $\mathcal{F}(U)$ has a discrete complete order and the product
of complete lattices is complete).

Define the operator $\Phi : \mathrm{Sec}(\mathcal{F}) \to
\mathrm{Sec}(\mathcal{F})$ by
\[
  \Phi(\sigma) = \bigsqcup \{\sigma' \in \mathrm{Sec}(\mathcal{F})
  \mid h(\sigma') \geq h(\sigma)\}.
\]
We verify that $\Phi$ is monotone: if $\sigma_1 \sqsubseteq
\sigma_2$ (i.e., $h(\sigma_1) \leq h(\sigma_2)$), then
$\{\sigma' \mid h(\sigma') \geq h(\sigma_2)\} \subseteq
\{\sigma' \mid h(\sigma') \geq h(\sigma_1)\}$, so
$\Phi(\sigma_2) \sqsubseteq \Phi(\sigma_1)$.  (Note the reversal:
$\Phi$ is monotone on the \emph{opposite} lattice, which suffices for
Knaster--Tarski on the dual.)

By Knaster--Tarski, $\Phi$ has a greatest fixpoint
$\sigma^* = \bigsqcup\{\sigma \mid \sigma \sqsubseteq \Phi(\sigma)\}$.
This fixpoint satisfies $h(\sigma^*) = h(\Phi(\sigma^*))$, meaning
no further refinement improves the harmony: $\sigma^*$ is a
harmony equilibrium.
\end{proof}

\begin{proposition}[Uniqueness of harmony fixpoint]
\label{prop:harmony-unique}
If the harmony function $h$ is \emph{strictly monotone} (i.e., for
each theory $\mathcal{T}_k$ and each section $\sigma$, there is at
most one $\sigma'$ maximising $g_k$ among sections compatible
with $\sigma$), then the fixpoint is unique.
\end{proposition}

\begin{proof}
Suppose $\sigma_1^*$ and $\sigma_2^*$ are two fixpoints.
By strict monotonicity, $\sigma_1^* \sqsubseteq \sigma_2^*$ and
$\sigma_2^* \sqsubseteq \sigma_1^*$, giving
$h(\sigma_1^*) = h(\sigma_2^*)$.  Since each $g_k$ has a unique
maximiser, the sections must agree on each stratum, hence
$\sigma_1^* = \sigma_2^*$.
\end{proof}

%% ── §20.3  The Bottleneck Theorem ───────────────────────────────────
\section{The bottleneck theorem}
\label{sec:bottleneck}

\begin{theorem}[Bottleneck principle]
\label{thm:bottleneck}
Let $h(\sigma) = \bigotimes_{k=1}^{N} w_k \otG g_k(\sigma)$ be the
harmony function.  If $(\Grade, \otG)$ is an idempotent commutative
monoid with $\otG = \min$ (the bottleneck semiring), then:
\[
  h(\sigma) = \min_{k=1}^{N}\; w_k \otG g_k(\sigma),
\]
and the minimum is achieved by at least one stratum $k^*$.
\end{theorem}

\begin{proof}
When $\otG = \min$, the tensor product of a finite family is the
minimum of the family (by induction on $N$: for $N = 1$ the result is
trivial; for $N+1$, $\min(a_1, \ldots, a_N, a_{N+1}) =
\min(\min(a_1, \ldots, a_N), a_{N+1})$ by associativity of $\min$).

Since $\Grade$ is totally ordered and $\{w_k \otG g_k(\sigma)\}$ is
a finite set of elements of $\Grade$, the minimum is achieved by
some index $k^* \in \{1, \ldots, N\}$.  This is the
\emph{bottleneck stratum}: the stratum whose grade most constrains
the overall harmony.
\end{proof}

\begin{corollary}[Bottleneck-driven repair]
\label{cor:bottleneck-repair}
To improve $h(\sigma)$, it is necessary and sufficient to improve the
grade at the bottleneck stratum $k^*$.  Improving any other stratum
$k \neq k^*$ does not change $h(\sigma)$ (until the bottleneck
shifts).
\end{corollary}

\begin{proof}
Since $h(\sigma) = w_{k^*} \otG g_{k^*}(\sigma) =
\min_k(w_k \otG g_k(\sigma))$, increasing $g_k(\sigma)$ for
$k \neq k^*$ does not change the minimum (the minimum is still
$w_{k^*} \otG g_{k^*}(\sigma)$).  Increasing $g_{k^*}(\sigma)$ to
$g_{k^*}'$ raises the harmony to
$\min(w_{k^*} \otG g_{k^*}', \min_{k \neq k^*} w_k \otG g_k)$,
which is $> h(\sigma)$ provided $w_{k^*} \otG g_{k^*}' >
w_{k^*} \otG g_{k^*}$.
\end{proof}

%% ── §20.4  FMLM Convergence ────────────────────────────────────────
\section{$\FMLM$ convergence}
\label{sec:fmlm}

\begin{definition}[$\FMLM$ refinement operator]
\label{def:fmlm-operator}
The \emph{Fill-in by Masked Language Model} ($\FMLM$) procedure
operates on a partially specified section $\sigma_0$ (a ``template''
with masked slots).  The refinement operator
$R : \mathrm{Sec}(\mathcal{F}) \to \mathrm{Sec}(\mathcal{F})$ is:
\begin{enumerate}[label=(\arabic*),nosep]
\item \textbf{Identify bottleneck:} Compute $k^* =
  \arg\min_k\, w_k \otG g_k(\sigma_n)$.
\item \textbf{Generate candidates:} Produce a finite set
  $\mathrm{Cand}(\sigma_n, k^*)$ of candidate replacements for the
  masked slots in stratum $k^*$.
\item \textbf{Grade candidates:} For each candidate $\sigma' \in
  \mathrm{Cand}$, compute $h(\sigma')$.
\item \textbf{Select best:} Set $R(\sigma_n) = \arg\max_{\sigma'}
  h(\sigma')$.
\end{enumerate}
The iteration is $\sigma_{n+1} = R(\sigma_n)$.
\end{definition}

\begin{theorem}[$\FMLM$ convergence]
\label{thm:fmlm}
Let $(\Grade, \leq)$ be a complete lattice with the ascending chain
condition (ACC) (equivalently, $\Grade$ is finite or well-ordered).
Suppose the $\FMLM$ refinement operator $R$ satisfies:
\begin{enumerate}[label=(\roman*),nosep]
\item \emph{Non-degeneracy:} If $\sigma_n$ is not a fixpoint, then
  $\mathrm{Cand}(\sigma_n, k^*)$ contains at least one candidate
  $\sigma'$ with $h(\sigma') > h(\sigma_n)$.
\item \emph{Budget bound:} The total number of iterations is bounded
  by $B < \infty$.
\end{enumerate}
Then the sequence $(\sigma_n)_{n \geq 0}$ converges to a fixpoint
$\sigma^*$ with $h(\sigma^*) \geq h(\sigma_0)$, in at most
$\min(B, |\Grade|)$ steps.
\end{theorem}

\begin{proof}
The harmony sequence $(h(\sigma_n))_{n \geq 0}$ is non-decreasing:
$h(\sigma_{n+1}) = h(R(\sigma_n)) \geq h(\sigma_n)$ because $R$
selects $\arg\max h$ from a candidate set that includes $\sigma_n$
itself (the ``no change'' candidate).

By non-degeneracy, if $\sigma_n$ is not a fixpoint, then
$h(\sigma_{n+1}) > h(\sigma_n)$ (strict increase).

Since $(\Grade, \leq)$ satisfies ACC, every strictly ascending chain
is finite.  Thus the sequence $(h(\sigma_n))$ stabilises after
finitely many steps, say at step $n^* \leq |\Grade|$.  At $n^*$,
$h(\sigma_{n^*+1}) = h(\sigma_{n^*})$, so either $\sigma_{n^*}$ is
a fixpoint of $R$ (the non-degeneracy condition is not triggered)
or the iteration has exhausted the budget $B$.

With the budget bound, the iteration stops at
$n^* = \min(B, |\Grade|)$.

The inequality $h(\sigma^*) \geq h(\sigma_0)$ follows from the
non-decreasing property of the harmony sequence.
\end{proof}

\begin{remark}[Convergence rate]
\label{rem:convergence-rate}
If the grade semiring is $\Grade = [0,1]$ (discretised to $M$
levels), then $|\Grade| = M$ and convergence occurs in at most $M$
steps.  In practice, the bottleneck principle
(\cref{thm:bottleneck}) ensures that each step improves the
bottleneck stratum, and the number of strata $N$ provides a tighter
bound: at most $N \cdot M$ refinements suffice (each stratum can be
the bottleneck at most $M$ times before it is no longer minimal).
\end{remark}

%% ── §20.5  Descent as Optimisation ──────────────────────────────────
\section{Descent as optimisation}
\label{sec:descent-opt}

The final result of this chapter establishes that descent (finding a
global section of a fibration) is equivalent to optimisation (finding
the globally optimal section under the harmony metric).

\begin{theorem}[Descent-optimisation correspondence]
\label{thm:descent-opt}
Let $p : \mathcal{E} \to \mathcal{B}$ be a fibration of categories,
$\mathcal{U} = \{U_i \to X\}$ a cover in $\mathcal{B}$, and
$\mathcal{F}$ the associated presheaf of sections.  Let
$h : \mathrm{Sec}(\mathcal{F}) \to \Grade$ be the harmony function.
Then:
\begin{enumerate}[label=(\roman*)]
\item A \emph{descent datum} $(s_i, \phi_{ij})$ (local sections
  with gluing isomorphisms) corresponds to a \emph{candidate section}
  $\sigma \in \mathrm{Sec}(\mathcal{F})$.
\item Two descent data are \emph{coboundary equivalent}
  ($(s_i, \phi_{ij}) \sim (s_i', \phi_{ij}')$ iff there exist
  isomorphisms $\psi_i : s_i \xrightarrow{\sim} s_i'$ compatible with
  the $\phi$'s) if and only if the corresponding candidate sections
  are \emph{gauge equivalent}: $\sigma \sim_g \sigma'$ iff there
  exists $(t_i)$ with $\sigma' = \sigma + \delta^0 t$.
\item The \emph{cocycle condition} ($\phi_{ij} \circ \phi_{jk} =
  \phi_{ik}$ on triple overlaps) is equivalent to \emph{local
  optimality}: $\sigma$ is a critical point of $h$ restricted to
  each overlap region.
\item A descent datum is \emph{effective} (glues to a global section)
  if and only if the corresponding candidate section is \emph{globally
  optimal}: $h(\sigma) = \max_{\sigma'} h(\sigma')$.
\end{enumerate}
\end{theorem}

\begin{proof}
(i)\; A descent datum assigns to each $U_i$ a section
$s_i \in \mathcal{F}(U_i)$ and to each overlap $U_{ij} = U_i
\times_X U_j$ an isomorphism $\phi_{ij} : s_i|_{U_{ij}}
\xrightarrow{\sim} s_j|_{U_{ij}}$.  This is precisely a 0-cochain
$(s_i) \in \Cech^0(\mathcal{U}, \mathcal{F})$ together with a
1-cochain $(\phi_{ij}) \in \Cech^1(\mathcal{U},
\mathrm{Aut}(\mathcal{F}))$.  The candidate section
$\sigma = (s_i)_{i \in I}$ is the 0-cochain.

(ii)\; A coboundary equivalence is a 0-cochain $(\psi_i)$ with
$\phi_{ij}' = \psi_j \circ \phi_{ij} \circ \psi_i^{-1}$, which is
$(\phi') = (\phi) + \delta^0(\psi)$ in the \v{C}ech complex of
$\mathrm{Aut}(\mathcal{F})$.  On sections, $s_i' = \psi_i(s_i)$,
so $\sigma' = \sigma + \delta^0 t$ where $t_i$ is the displacement
$\psi_i(s_i) - s_i$.  This is gauge equivalence.

(iii)\; The cocycle condition $\phi_{ij} \circ \phi_{jk} = \phi_{ik}$
means the discrepancies $d_{ij} = s_i|_{U_{ij}} - s_j|_{U_{ij}}$
satisfy $d_{ij} + d_{jk} = d_{ik}$, i.e., $(d_{ij}) \in
\ker \delta^1$.  In terms of harmony: the grade $g(\sigma)$
restricted to the overlap $U_{ij}$ depends on the discrepancy
$d_{ij}$.  The cocycle condition says all overlap discrepancies are
coherent, which is local optimality: no local rearrangement of
sections on overlaps can improve the harmony (since the discrepancies
are already minimal---they satisfy the tightest possible constraint).

(iv)\; If the descent datum is effective, there exists a global
section $s \in \mathcal{F}(X)$ with $s|_{U_i} = s_i$.  This global
section has harmony
\[
  h(s) = \bigotimes_k w_k \otG g_k(s) = \gone
\]
(the maximal grade) because there is no discrepancy ($d_{ij} = 0$,
$H^1 = 0$).  Conversely, if $h(\sigma) = \max h$, then all overlap
discrepancies vanish (the bottleneck is $\gone$), the cocycle
condition is satisfied trivially, and the descent datum is effective
by the sheaf condition.
\end{proof}

\begin{corollary}[Equivalence as descent on product site]
\label{cor:equiv-descent}
Two artifacts $P$ and $Q$ are equivalent with respect to a property
$\varphi$ if and only if
$\check{H}^1(\mathcal{U}_{\mathrm{eq}},
\mathcal{F}_{\mathrm{eq}}) = 0$
on the product site $\mathcal{C}(P) \times \mathcal{C}(Q)$, where
$\mathcal{F}_{\mathrm{eq}}(U_P, U_Q) =
\{(s_P, s_Q) \mid s_P|_{U_P} \sim_\varphi s_Q|_{U_Q}\}$.
\end{corollary}

\begin{proof}
The equivalence $P \sim_\varphi Q$ means every local section
witnessing $\varphi$ on a region of $P$ corresponds to a local
section on the corresponding region of $Q$, and these correspondences
are compatible on overlaps.  This is precisely a descent datum on the
product site with the relational presheaf $\mathcal{F}_{\mathrm{eq}}$.
By \cref{thm:fundamental}, the descent datum is effective iff
$H^1 = 0$.
\end{proof}

\begin{proposition}[Support-local invalidation]
\label{prop:support-local}
If evidence changes are supported on a sub-region $V \subseteq X$,
then only judgments in
\[
  \mathrm{Star}_{\mathcal{U}}(V) = \{U_i \in \mathcal{U}
  \mid U_i \times_X V \neq \emptyset\}
\]
may need re-verification.
\end{proposition}

\begin{proof}
Let $\sigma$ be a section and $\sigma'$ the modified section with
changes supported on~$V$.  Then $\sigma|_{U_i} = \sigma'|_{U_i}$ for
all $U_i$ with $U_i \times_X V = \emptyset$.  The discrepancy
$d_{ij}' = \sigma_i'|_{U_{ij}} - \sigma_j'|_{U_{ij}}$ equals
$d_{ij}$ whenever both $U_i$ and $U_j$ are disjoint from $V$.  Thus
the cocycle condition is preserved on all overlaps outside
$\mathrm{Star}(V)$, and only overlaps involving at least one member
of $\mathrm{Star}(V)$ need re-checking.

By the Mayer--Vietoris principle, the cohomology of $X$ decomposes:
the restriction map $H^1(X) \to H^1(X \setminus V)$ is an isomorphism
on the complement of $\mathrm{Star}(V)$, so only the contribution
from $\mathrm{Star}(V)$ can change.
\end{proof}

%%%%%%%%%%%%%%%%%%%%%%%%%%%%%%%%%%%%%%%%%%%%%%%%%%%%%%%%%%%%%%%%%%%%%%%%
%%  Chapter 21 — Metatheoretic Properties
%%%%%%%%%%%%%%%%%%%%%%%%%%%%%%%%%%%%%%%%%%%%%%%%%%%%%%%%%%%%%%%%%%%%%%%%
\chapter{Metatheoretic Properties}
\label{ch:metatheory}

This chapter establishes consistency, conservativity, decidability
results, completeness of the subsumption, and cut elimination for
the graded type system.

%% ── §21.1  Consistency ─────────────────────────────────────────────
\section{Consistency}
\label{sec:consistency}

\begin{theorem}[Consistency of the 4-tuple framework]
\label{thm:consistency}
Assuming ZFC with a Grothendieck universe $\mathbb{U}$, the 4-tuple
framework $(\STop, \PStr, \CSpec, \GEval)$ is consistent: there is no
derivation of $\Gamma \vdash_g E : \bot$ for any context $\Gamma$,
grade $g$, and term $E$.
\end{theorem}

\begin{proof}
We construct a model of the 4-tuple framework in the topos
$\Set_{\mathbb{U}}$ of $\mathbb{U}$-small sets.

\emph{Step 1: Model of $\STop$.}
Take $\mathcal{C} = \mathbf{1}$ (the terminal category) with the
maximal topology.  Then $\Sh(\mathbf{1}) \cong \Set_{\mathbb{U}}$,
which is an elementary topos with natural number object.  Set
$\mathcal{F} = \mathrm{const}_{\{*\}}$ (the constant presheaf on a
singleton).  The subobject classifier is $\Omega = \{0, 1\}$, the
sheafification is the identity, and the global-sections functor
$\gamma^*$ is the identity.  This is a valid Semantic Topos.

\emph{Step 2: Model of $\PStr$.}
Take $\Gamma = ()$ (empty context), $A = \mathbb{1}$ (the unit
type), and $E = * : \mathbb{1}$ (the unique inhabitant).  The type
theory is standard MLTT interpreted in $\Set_{\mathbb{U}}$ via the
standard set-theoretic semantics \cite{mlm1984}.  Under this
interpretation:
\begin{itemize}[nosep]
\item $\Pi$-types are interpreted as dependent products (indexed
  families of functions);
\item $\Sigma$-types as dependent sums (indexed disjoint unions);
\item Identity types $\mathrm{Id}_A(a,b)$ as the singleton
  $\{*\}$ if $a = b$ and $\emptyset$ otherwise;
\item The universe $\Type_0$ is interpreted as $\mathbb{U}$ itself.
\end{itemize}
The empty type $\bot$ is interpreted as $\emptyset$.  Since there is
no set-theoretic function $f : \Gamma \to \emptyset$ when $\Gamma$
is non-empty (or when $\Gamma$ is the terminal object), there is no
closed term of type $\bot$ in the model.

\emph{Step 3: Model of $\CSpec$.}
The trivial cover $\{* \to *\}$ on $\mathbf{1}$ gives
$\Cech^0 = \{*\}$ and $\Cech^n = 0$ for $n \geq 1$.  All cohomology
vanishes: $H^n = 0$.  This is a valid Cohomological Spectrum.

\emph{Step 4: Model of $\GEval$.}
Take $G = \mathbf{2} = \{\gzero, \gone\}$ with $\otG = \land$,
$\opG = \lor$.  The evaluation $\mu(\mathcal{D}) = \gone$ for every
valid derivation.  This is a valid quantale.

\emph{Step 5: Compatibility.}
The compatibility conditions (K1)--(K4) are satisfied trivially in
this model: the trivial cover refines the maximal topology (K1),
the presheaf is shared (K2), the unit type classifies sections
of the constant presheaf (K3), and the evaluation is well-defined
(K4).

\emph{Step 6: Consistency.}
If $\bot$ were derivable in the 4-tuple framework, there would exist
$\Gamma, g, E$ with $\Gamma \vdash_g E : \bot$.  Under the
set-theoretic interpretation, this would give an element of $\emptyset$,
which is impossible.  Hence the framework is consistent.
\end{proof}

%% ── §21.2  Conservativity ──────────────────────────────────────────
\section{Conservativity over MLTT}
\label{sec:conservativity}

\begin{theorem}[Conservative extension]
\label{thm:conservative}
Adding the Grade semiring $\Grade$ to Martin-L\"of Type Theory is
conservative over MLTT: if $\Gamma \vdash_{\gone} E : A$ is
derivable in the graded system and $\Gamma, A, E$ contain no grade
annotations, then $\Gamma \vdash E : A$ is derivable in standard
MLTT.
\end{theorem}

\begin{proof}
We define a \emph{grade-erasure} map $|-| : \mathrm{Judg}_g \to
\mathrm{Judg}$ from graded judgments to ordinary MLTT judgments,
and show it preserves derivability for grade-$\gone$ judgments.

\emph{Grade erasure.}  Define $|-|$ on judgments, types, and terms by:
\begin{align*}
  |\Gamma \vdash_g E : A| &= \Gamma \vdash |E| : |A|, \\
  |\GId_g(A, a, b)| &= \mathrm{Id}_{|A|}(|a|, |b|)
    \quad\text{when } g = \gone, \\
  |\Pi_g(x : A).\, B| &= \Pi(x : |A|).\, |B|
    \quad\text{when } g = \gone, \\
  |\Sigma_g(x : A).\, B| &= \Sigma(x : |A|).\, |B|
    \quad\text{when } g = \gone.
\end{align*}
For $g < \gone$, the graded types have no MLTT counterpart, but we
are only considering $g = \gone$ judgments.

\emph{Preservation of derivations.}  We show by induction on the
derivation of $\Gamma \vdash_{\gone} E : A$ that
$\Gamma \vdash |E| : |A|$ is derivable in MLTT.

\textsc{Case Variable:}
$\dfrac{(x : A) \in \Gamma}{\Gamma \vdash_{\gone} x : A}$.
Erasing: $(x : |A|) \in \Gamma$, so $\Gamma \vdash x : |A|$ by the
MLTT variable rule.

\textsc{Case $\Pi$-introduction:}
$\dfrac{\Gamma, x : A \vdash_{\gone} E : B}
  {\Gamma \vdash_{\gone} \lambda x.\, E : \Pi_{\gone}(x:A).\, B}$.
By induction, $\Gamma, x : |A| \vdash |E| : |B|$.  By the MLTT
$\Pi$-introduction rule,
$\Gamma \vdash \lambda x.\, |E| : \Pi(x:|A|).\, |B|$.

\textsc{Case $\Pi$-elimination:}
$\dfrac{\Gamma \vdash_{\gone} f : \Pi_{\gone}(x:A).\, B \quad
  \Gamma \vdash_{\gone} a : A}
  {\Gamma \vdash_{\gone} f\,a : B[a/x]}$.
By induction, $\Gamma \vdash |f| : \Pi(x:|A|).\,|B|$ and
$\Gamma \vdash |a| : |A|$.  By MLTT $\Pi$-elimination,
$\Gamma \vdash |f|\,|a| : |B|[|a|/x]$.

\textsc{Case Grade-weakening:}
$\dfrac{\Gamma \vdash_g E : A \quad g \leq g'}
  {\Gamma \vdash_{g'} E : A}$.
When $g' = \gone$, we need $g = \gone$ (since $\gone$ is maximal
and $g \leq \gone$ always holds; but if $g < \gone$, the premise
has grade $< \gone$ and the erasure may not apply).  When
$g = \gone$, the premise directly gives the result by induction.

The remaining cases ($\Sigma$-types, identity types, universe rules,
$W$-types) are analogous: at grade $\gone$, every graded type former
reduces to its MLTT counterpart, and every graded rule reduces to
the corresponding MLTT rule.

\emph{Conservativity.}  If $\Gamma \vdash_{\gone} E : A$ is
derivable in the graded system with $\Gamma, A, E$ containing no
grade annotations, then the grade-erasure $|\Gamma \vdash_{\gone}
E : A| = \Gamma \vdash E : A$ is derivable in MLTT by the induction
above.  Hence the graded system proves no new $\gone$-level judgments
about pure MLTT types and terms.
\end{proof}

%% ── §21.3  Decidability ────────────────────────────────────────────
\section{Decidability results}
\label{sec:decidability}

\begin{theorem}[Decidability of grade computation]
\label{thm:dec-grade}
For a finite number of strata $N$ and a decidable grade semiring
$\Grade$, the harmony computation
$h(\sigma) = \bigotimes_{k=1}^{N} w_k \otG g_k(\sigma)$
is decidable.
\end{theorem}

\begin{proof}
The computation requires:
\begin{enumerate}[label=(\arabic*),nosep]
\item Evaluating each $g_k(\sigma)$ for $k = 1, \ldots, N$.  Since
  $\Grade$ is decidable and each $g_k$ is a computable function (by
  hypothesis), each evaluation terminates.
\item Computing $N - 1$ applications of $\otG$.  Since $\otG$ is a
  computable operation on a decidable semiring, each application
  terminates.
\item Computing $N$ applications of the weight multiplication
  $w_k \otG (-)$.  Again computable.
\end{enumerate}
The total computation is a finite sequence of decidable operations,
hence decidable.  The complexity is $O(N \cdot C_g)$ where $C_g$ is
the cost of a single grade evaluation.
\end{proof}

\begin{theorem}[Decidability of harmony]
\label{thm:dec-harmony}
If each theory slot $\mathcal{T}_k$ has a decidable satisfaction
relation (i.e., given $\sigma$ and $\mathcal{T}_k$, it is decidable
whether $g_k(\sigma) \geq t$ for any threshold $t$), then the harmony
comparison $h(\sigma_1) \geq h(\sigma_2)$ is decidable.
\end{theorem}

\begin{proof}
By \cref{thm:dec-grade}, $h(\sigma_1)$ and $h(\sigma_2)$ are both
computable.  The comparison $h(\sigma_1) \geq h(\sigma_2)$ reduces to
$\leq$ on $\Grade$, which is decidable by hypothesis.
\end{proof}

\begin{theorem}[Decidability of graded type-checking]
\label{thm:dec-typecheck}
If the underlying type theory has decidable type-checking (given
$\Gamma, E, A$, it is decidable whether $\Gamma \vdash E : A$),
then graded type-checking ($\Gamma \vdash_g E : A$) is decidable.
\end{theorem}

\begin{proof}
Graded type-checking extends ordinary type-checking with grade
tracking.  Each typing rule of the graded system has the form:
\[
  \dfrac{\Gamma \vdash_{g_1} E_1 : A_1 \quad \cdots \quad
    \Gamma \vdash_{g_n} E_n : A_n \quad g = f(g_1, \ldots, g_n)}
    {\Gamma \vdash_g E : A}
\]
where $f$ is a computable function of the input grades (typically
$\otG$ or $\opG$).  Given a purported derivation, checking each rule
application requires:
\begin{enumerate}[label=(\arabic*),nosep]
\item Checking the underlying MLTT judgment (decidable by hypothesis).
\item Computing the grade from the sub-grades (decidable by
  \cref{thm:dec-grade}).
\item Comparing the computed grade with the asserted grade (decidable
  on $\Grade$).
\end{enumerate}
By induction on the derivation tree, the full check is decidable.
\end{proof}

\begin{theorem}[Undecidability of full descent]
\label{thm:undec-descent}
The descent problem---``given a site $(\mathcal{C}, \mathcal{J})$, a
presheaf $\mathcal{F}$, and a cover $\mathcal{U}$, does
$\check{H}^1(\mathcal{U}, \mathcal{F}) = 0$?''---is undecidable in
general.
\end{theorem}

\begin{proof}
We reduce the halting problem to descent.  Given a Turing machine $M$
and input $w$, construct a site $\mathcal{C}_M$ as follows:
\begin{itemize}[nosep]
\item Objects: the configurations $(q, \tau, h)$ of $M$ on input~$w$
  (state, tape, head position), plus a terminal object~$X$.
\item Morphisms: one-step transitions $c \to c'$ and the unique
  morphism $c \to X$ for each configuration.
\item Cover: $\mathcal{U} = \{c \to X \mid c \text{ reachable}\}$.
\end{itemize}
Define the presheaf $\mathcal{F}(c) = \{*\}$ if $c$ is not a
halting configuration, and $\mathcal{F}(c) = \emptyset$ if $c$ is a
halting configuration.

If $M$ halts on $w$, some configuration $c_h$ has
$\mathcal{F}(c_h) = \emptyset$, and the local sections on
$c_h$'s neighbors cannot be extended through $c_h$.  This creates a
non-trivial 1-cocycle: $\check{H}^1 \neq 0$.

If $M$ does not halt, every configuration has
$\mathcal{F}(c) = \{*\}$, all restrictions are identities, and the
constant family $(*)$ is a compatible descent datum:
$\check{H}^1 = 0$.

Thus $\check{H}^1 = 0$ iff $M$ does not halt on $w$, which is
undecidable.
\end{proof}

\begin{theorem}[Decidability for bounded depth]
\label{thm:dec-bounded}
If the site $\mathcal{C}$ has finite depth $d$ (every chain of
morphisms has length $\leq d$) and $\mathcal{F}$ takes values in
finite sets, then descent is decidable.
\end{theorem}

\begin{proof}
Under these finiteness conditions:
\begin{enumerate}[label=(\arabic*),nosep]
\item The \v{C}ech complex $\Cech^\bullet(\mathcal{U}, \mathcal{F})$
  is a finite complex of finite-dimensional vector spaces (over the
  appropriate coefficient ring).
\item Each coboundary map $\delta^n$ is a finite matrix.
\item $\ker \delta^n$ and $\mathrm{im}\, \delta^{n-1}$ are computable
  by Gaussian elimination (or its analogue for abelian groups).
\item $\check{H}^1 = \ker \delta^1 / \mathrm{im}\, \delta^0$ is
  therefore computable.
\item Checking $\check{H}^1 = 0$ reduces to checking whether
  $\ker \delta^1 = \mathrm{im}\, \delta^0$, which is decidable.
\end{enumerate}
The complexity is polynomial in the size of the \v{C}ech complex:
$O(|\mathcal{U}|^{d+1} \cdot |\mathcal{F}|^3)$ for the Gaussian
elimination step.
\end{proof}

%% ── §21.4  Completeness of the Subsumption ──────────────────────────
\section{Completeness of the subsumption}
\label{sec:completeness}

\begin{theorem}[Proper extension --- detailed version]
\label{thm:proper-detail}
The JuGeo 4-tuple framework is a proper extension of every subsumed
framework: it is not equivalent to any of Montague Grammar, DRT/SDRT,
MLTT, OT, Linear Logic, the Generative Lexicon, HoTT, or Prior's
Tense Logic.
\end{theorem}

\begin{proof}
By \cref{thm:proper}, the separation example $\JJ^*$ lies outside
every embedding.  We sharpen this by exhibiting a judgment that
\emph{requires} all four components simultaneously.

\emph{The separation judgment $\JJ^*$.}
Consider the judgment ``The poem $P$ achieves a unity of meter,
imagery, and meaning across its three stanzas.''  Formalise this as:

\begin{itemize}[nosep]
\item $\STop^*$: Site $\mathcal{C}^*$ with objects
  $\{s_1, s_2, s_3, P\}$ (three stanzas and the full poem), covers
  $\{s_1, s_2, s_3\} \to P$ (the stanzas cover the poem).
  Three strata: $\phon$ (meter), $\sem$ (imagery), $\prag$
  (meaning).  The topos $\Sh(\mathcal{C}^*)$ is non-trivial (not
  equivalent to $\Set$).

\item $\PStr^*$: The graded judgment
  $\Gamma \vdash_{0.7} E_P : \mathrm{Unity}(P)$
  where $\mathrm{Unity}(P)$ is a dependent type encoding the three
  strata of unity.  The grade $0.7$ (not binary) records that the
  meter is excellent ($g_{\phon} = 0.95$), the imagery is strong
  ($g_{\sem} = 0.85$), but the meaning has a forced metaphor
  ($g_{\prag} = 0.7$).

\item $\CSpec^*$: The \v{C}ech complex on the stanza cover yields
  $H^1 \neq 0$: the metaphor in the transition from $s_2$ to $s_3$
  creates a cohomological obstruction (the local meanings do not glue
  smoothly).

\item $\GEval^*$: $G = [0,1]$ with $\otG = \min$.
  $h(\JJ^*) = \min(0.95, 0.85, 0.7) = 0.7$ (the meaning stratum
  is the bottleneck).
\end{itemize}

This judgment simultaneously requires:
\begin{enumerate}[label=(\roman*),nosep]
\item A non-trivial site (rules out MLTT, HoTT, Linear Logic,
  Generative Lexicon);
\item Non-binary grades (rules out Montague, DRT, MLTT, HoTT, Tense
  Logic);
\item Non-trivial $H^1$ (rules out all frameworks that have only
  trivial cohomology);
\item Multi-stratal harmony (rules out OT, which has a single stratum
  with ranked constraints but no cross-stratal harmony).
\end{enumerate}

No single subsumed framework can express all four properties.  Hence
JuGeo is a proper extension.
\end{proof}

%% ── §21.5  Cut Elimination ──────────────────────────────────────────
\section{Cut elimination for the graded system}
\label{sec:cut-elim}

\begin{definition}[Graded sequent calculus]
\label{def:graded-sequent}
The graded sequent calculus has judgments of the form
$\Gamma \vdash_g E : A$ where $\Gamma$ is a context, $g \in \Grade$
is a grade, $E$ is a term, and $A$ is a type.  The \emph{cut rule}
is:
\[
  \dfrac{\Gamma \vdash_g E : A \quad \Gamma, x : A \vdash_{g'} E' : B}
    {\Gamma \vdash_{g \otG g'} E'[E/x] : B}
  \;\;\textsc{(Cut)}
\]
The grade of the conclusion is $g \otG g'$: the tensor product of
the grades of the two premises.
\end{definition}

\begin{definition}[Cut-free derivation]
A derivation is \emph{cut-free} if it does not use the \textsc{(Cut)}
rule.
\end{definition}

\begin{theorem}[Cut elimination]
\label{thm:cut-elim}
If $\Gamma \vdash_g E : A$ has a derivation with cuts, then there
exists a cut-free derivation $\Gamma \vdash_{g'} E' : A$ with
$g' \geq g$ (the grade may increase, but does not decrease).
Moreover, $E' =_\beta E$ (the terms are $\beta$-equal).
\end{theorem}

\begin{proof}
We adapt Gentzen's cut elimination procedure, tracking grades through
each reduction step.

\emph{Measure.}  Define the \emph{cut-rank} of a derivation $\mathcal{D}$
as the pair $(\rho, \nu)$ where $\rho$ is the maximum complexity
(number of connectives) of a cut formula in $\mathcal{D}$, and $\nu$
is the number of cuts of that maximum complexity.  We order cut-ranks
lexicographically.

\emph{Key reduction.}  Consider a topmost cut (one whose premises
are cut-free or contain only lower-rank cuts):
\[
  \dfrac{\dfrac{\vdots}{\Gamma \vdash_{g_1} E_1 : A}
    \quad
    \dfrac{\vdots}{\Gamma, x:A \vdash_{g_2} E_2 : B}}
    {\Gamma \vdash_{g_1 \otG g_2} E_2[E_1/x] : B}
  \;\;\textsc{(Cut)}
\]

\textsc{Case 1: Principal cut on $\Pi$-types.}
The left premise ends with $\Pi$-introduction and the right with
$\Pi$-elimination:
\begin{gather*}
  \dfrac{\Gamma, y:C \vdash_{g_1} E_1 : D}
    {\Gamma \vdash_{g_1} \lambda y.\,E_1 : \Pi(y:C).\,D}
  \quad
  \dfrac{\Gamma \vdash_{g_2} f : \Pi(y:C).\,D \quad
    \Gamma \vdash_{g_3} a : C}
    {\Gamma \vdash_{g_2 \otG g_3} f\,a : D[a/y]}
\end{gather*}
Cutting on $\Pi(y:C).\,D$: the cut reduces to substituting $\lambda
y.\,E_1$ for $f$, giving $(\lambda y.\,E_1)\,a \to_\beta E_1[a/y]$.
The grade of the result is $g_1 \otG g_3$: we lose the $g_2$ factor
(the grade of the eliminated function) but gain back $g_1$ (the
grade of the introduction).  Since $g_2 \leq \gone$ (grades are
bounded by $\gone$) and $g_1 \otG g_3 \geq g_1 \otG g_2 \otG g_3$
(because $g_2 \leq \gone$ and $\otG$ is monotone), the resulting
grade satisfies $g' = g_1 \otG g_3 \geq g_1 \otG g_2 \otG g_3 = g$.

The cut on $\Pi(y:C).\,D$ is replaced by cuts on $C$ and $D$, both
of lower complexity.  Cut-rank decreases.

\textsc{Case 2: Principal cut on $\Sigma$-types.}
Similar: $\Sigma$-introduction paired with projection eliminates the
cut, replacing it with lower-complexity cuts.  The grade bound
$g' \geq g$ holds by monotonicity of $\otG$.

\textsc{Case 3: Principal cut on $\mathrm{Id}$-types.}
The left premise introduces a reflexivity proof
$\mathrm{refl}_a : \mathrm{Id}_A(a,a)$ and the right eliminates it
via $J$ (path induction).  The reduction applies the $J$-computation
rule: $J(\mathrm{refl}_a, d) \to_\beta d$.  The grade computation
$g' = g_1 \otG g_d \geq g_1 \otG g_2 \otG g_d$ follows from
$g_2 \leq \gone$.

\textsc{Case 4: Non-principal cut (commutative case).}
If the cut formula is not the principal formula of one of the
premises, commute the cut upward past the last rule of that premise.
This does not change the cut-rank but moves the cut closer to an
axiom, where it can be eliminated trivially.

\emph{Grade tracking.}  At each reduction step:
\begin{itemize}[nosep]
\item Principal cuts: the reduced derivation has grade $g' \geq g$
  because the eliminated intermediate formula had grade $\leq \gone$,
  and removing it can only increase the tensor product (since
  $a \otG b \leq a$ when $b \leq \gone$ for the bottleneck semiring,
  so removing the $b$ factor gives $a \geq a \otG b$).
\item Commutative cuts: grade is unchanged (the cut is merely
  repositioned).
\end{itemize}

\emph{Termination.}  Each principal-cut reduction strictly decreases
the cut-rank $(\rho, \nu)$: either $\rho$ decreases (the maximum
cut complexity drops) or $\rho$ stays the same and $\nu$ decreases
(the number of maximal-complexity cuts drops by one).  Since
cut-ranks are well-ordered (pairs of natural numbers under
lexicographic order), the process terminates.

\emph{Conclusion.}  After finitely many reduction steps, all cuts
are eliminated.  The resulting derivation
$\Gamma \vdash_{g'} E' : A$ is cut-free, with $g' \geq g$ and
$E' =_\beta E$ (each reduction step is a $\beta$-reduction or a
commutation that preserves $\beta$-equality).
\end{proof}

\begin{corollary}[Subformula property]
\label{cor:subformula}
Every cut-free derivation in the graded sequent calculus has the
subformula property: every type appearing in the derivation is a
subformula of a type in the conclusion or the context.
\end{corollary}

\begin{proof}
The cut rule is the only rule that introduces a type ($A$, the cut
formula) not present in the conclusion.  In a cut-free derivation,
every type is introduced by an introduction or elimination rule,
which only involves subformulas of the conclusion type or context
types.
\end{proof}

\begin{corollary}[Consistency from cut elimination]
\label{cor:consistency-cut}
The graded sequent calculus is consistent: $\vdash_g \bot$ is not
derivable for any $g$.
\end{corollary}

\begin{proof}
By \cref{thm:cut-elim}, if $\vdash_g E : \bot$ were derivable, there
would be a cut-free derivation.  By the subformula property, every
type in the derivation is a subtype of $\bot$ or of a context type.
Since the context is empty and $\bot$ has no introduction rule, no
cut-free derivation of $\bot$ exists.
\end{proof}

%% ── §21.6  Structural Rules ────────────────────────────────────────
\section{Structural rules}
\label{sec:structural}

\begin{theorem}[Admissibility of structural rules]
\label{thm:structural}
The graded sequent calculus admits weakening, exchange, and
contraction, with the following grade effects:
\begin{enumerate}[label=(\roman*)]
\item \textbf{Weakening:} If $\Gamma \vdash_g E : B$, then
  $\Gamma, x : A \vdash_g E : B$ (grade unchanged).
\item \textbf{Exchange:} If $\Gamma, x : A, y : B, \Delta \vdash_g
  E : C$, then $\Gamma, y : B, x : A, \Delta \vdash_g E : C$
  (grade unchanged).
\item \textbf{Contraction:} If $\Gamma, x : A, y : A \vdash_g
  E : B$, then $\Gamma, z : A \vdash_{g'} E[z/x, z/y] : B$ with
  $g' \geq g$ (grade may increase).
\end{enumerate}
\end{theorem}

\begin{proof}
(i) Weakening:  The additional variable $x : A$ is unused in
the derivation of $E : B$.  Each typing rule in the derivation
remains valid with the extended context, since all rules are
monotone in the context.  The grade is unchanged because $x$
contributes no evidence.

(ii) Exchange: The typing rules depend only on the \emph{set} of
context bindings, not their order (contexts are multisets of
typed variables).

(iii) Contraction: Replacing two copies $x, y$ of a variable of
type $A$ with a single copy $z$ identifies two evidence paths.
If $x$ contributed grade $g_x$ and $y$ contributed grade $g_y$,
the contraction replaces $g_x \otG g_y$ with $g_z = g_x \opG g_y
\geq g_x \otG g_y$ (since $a \opG b \geq a \otG b$ in any
semiring where $\opG$ dominates $\otG$, which holds for the
idempotent semiring by the absorption law).
\end{proof}

%% ── §21.7  Source-Independent Soundness ─────────────────────────────
\section{Source-independent soundness}
\label{sec:source-sound}

\begin{theorem}[Source-independent soundness]
\label{thm:source-sound}
If descent succeeds on a family of local sections
$\{s_i\}_{i \in I}$, the resulting global section $s$ is valid
regardless of the source of each $s_i$ (whether produced by an SMT
solver, a runtime test, an LLM, or a human).  The grade of $s$ is
\[
  g(s) = \bigotimes_{i \in I} g(s_i),
\]
the tensor product of the grades of the individual sources.
\end{theorem}

\begin{proof}
The descent theorem (\cref{thm:fundamental}) is a purely mathematical
statement about the \v{C}ech complex: it depends only on the
\emph{values} of the local sections $s_i \in \mathcal{F}(U_i)$ and
their compatibility on overlaps, not on the \emph{process} that
produced them.

Concretely: the gluing condition
$s_i|_{U_{ij}} = s_j|_{U_{ij}}$ is a set-theoretic equality,
verifiable without reference to provenance.  If the equality holds,
the sheaf condition produces a unique global section $s$,
regardless of source.

The grade computation $g(s) = \bigotimes_i g(s_i)$ follows from the
multiplicativity of the grade evaluation: the grade of a composite
derivation is the tensor product of the grades of its parts
(by the definition of $\mu$ on derivation trees).  Each $g(s_i)$
reflects the trust level of its source (e.g., $g = \gone$ for
formal proofs, $g = 0.3$ for LLM suggestions), and the tensor
product computes the overall trust via the bottleneck principle.
\end{proof}

\begin{remark}
Source-independent soundness is the key property that enables
hybrid verification: combining formal proofs (high grade) with
LLM-generated suggestions (low grade) and runtime evidence (medium
grade) into a single judgment whose overall trust is the
conservative combination of all sources.
\end{remark}


% ── Part V: Examples ──────────────────────────────────────────
\part{Worked Examples}
% Part V: Worked Examples
% This file is \include'd from main.tex

\chapter{Program Verification}
\label{ch:ex-verification}

We begin Part~V with the domain closest to JuGeo's origins:
program verification.  We walk through the full 4-tuple
$(\STop,\PStr,\CSpec,\GEval)$ for two concrete programs---a
lock-free queue (where obstructions arise) and a sorting function
(where they vanish).

%% ─────────────────────────────────────────────────────────────
\section{Example 1: A Lock-Free Queue}
\label{sec:lfq}

Consider a Michael--Scott lock-free queue with two operations
\texttt{enqueue} and \texttt{dequeue}, each using a
compare-and-swap (CAS) loop.  The correctness goal is
\emph{linearizability}: every concurrent execution is equivalent to
some sequential one.

%% ── 22.1.1  STop ─────────────────────────────────────────────
\subsection{The Semantic Topos $\STop$}

\begin{construction}[Control-flow site]
\label{con:cfg-site}
Let $\mathcal{C}_Q$ be the category whose
\begin{itemize}[nosep]
  \item \textbf{objects} are program points
        $p_0,p_1,\ldots,p_n$ in the control-flow graph (CFG),
  \item \textbf{morphisms} $p_i\to p_j$ are feasible execution paths
        (sequences of basic-block transitions).
\end{itemize}
The Grothendieck topology~$J$ declares a covering family
of a function node $f$ to be the set of its basic blocks:
\[
  \{B_1\hookrightarrow f,\;B_2\hookrightarrow f,\;\ldots\}
  \;\in\; J(f).
\]
\end{construction}

\begin{definition}[Type presheaf]
\label{def:type-presheaf}
The presheaf $\mathcal{F}_{\mathrm{ty}}:\mathcal{C}_Q^{\op}\to\Set$
assigns to each program point~$p$ the set of variable-type bindings
in scope:
\[
  \mathcal{F}_{\mathrm{ty}}(p)
  \;=\;
  \bigl\{(x:\texttt{Node*}),\;(tail:\texttt{AtomicPtr<Node>}),\;\ldots\bigr\}.
\]
Restriction along a morphism $p\to q$ is the usual scope narrowing.
\end{definition}

The \textbf{subobject classifier} $\Omega$ in the topos
$\Sh(\mathcal{C}_Q,J)$ sends each program point~$p$ to the set of
sieves on~$p$.  Concretely, a sieve on~$p$ is a downward-closed set
of paths arriving at~$p$.  The interpretation: a sieve encodes the
set of \emph{reachable} histories at that point.  An assertion
$\varphi$ at~$p$ is therefore a map
$\llbracket\varphi\rrbracket:1\to\Omega$ selecting the sieve of
paths on which $\varphi$ holds.

\begin{example}[Assertions as sieves]
At program point $p_{\mathrm{cas}}$ (the CAS instruction in
\texttt{enqueue}), the assertion ``\texttt{tail} still points to the
last node'' selects the sieve
$S = \{$paths that have not been preempted by another
\texttt{enqueue}$\}$.  This sieve is a \emph{proper} sub-sieve of
the maximal sieve, reflecting the fact that the assertion can be
invalidated by concurrent interference.
\end{example}

%% ── 22.1.2  PStr ─────────────────────────────────────────────
\subsection{The Proof Structure $\PStr$}

\begin{construction}[Specification type]
\label{con:spec-type}
The specification is the type
\[
  \mathrm{Lin}(Q)
  \;:\equiv\;
  \Pi\,(\sigma:\mathrm{ConcExec}(Q)).\;
  \Sigma\,(\tau:\mathrm{SeqExec}(Q)).\;
  \GId_{\gone}\!\bigl(\mathrm{Obs},\;
    \mathrm{obs}(\sigma),\;\mathrm{obs}(\tau)\bigr),
\]
reading: ``for every concurrent execution~$\sigma$, there exists a
sequential execution~$\tau$ with identical observable behaviour.''
\end{construction}

The \textbf{evidence} inhabiting this type is a \emph{simulation
relation} $R\subseteq\mathrm{State}\times\mathrm{AbsState}$ together
with proofs that $R$ is preserved by every atomic step.

\begin{construction}[Hoare-logic proof tree]
\label{con:hoare-tree}
The derivation tree has the shape of a Hoare-logic proof.
For \texttt{enqueue(x)}:
\[
  \dfrac{
    \dfrac{
      \{\texttt{tail}=t\}
      \;\texttt{new\_node := alloc(x)}\;
      \{\texttt{new\_node}\neq\texttt{null}\}
    }{\text{(Alloc)}}
    \quad
    \dfrac{
      \{\texttt{tail}=t \wedge \texttt{next}=\texttt{null}\}
      \;\texttt{CAS(tail.next, null, new\_node)}\;
      \{R(t,t')\}
    }{\text{(CAS-success)}}
  }{
    \{\mathrm{pre}\}\;\texttt{enqueue(x)}\;\{\mathrm{post}\}
  }
\]
Each horizontal line is a typing rule application
$\Gamma\vdash_{\gone} e:A$; the full tree is the term witnessing
$\mathrm{Lin}(Q)$.
\end{construction}

\textbf{$\beta$-reduction.}\quad
In $\PStr$, $\beta$-reduction corresponds to
\emph{proof simplification}.  For instance, composing the CAS-success
lemma with the swing-tail lemma yields a single invariant-preservation
step:
\[
  (\lambda h.\;\mathrm{swing}(h))\;(\mathrm{cas\text{-}ok}(\sigma))
  \;\;\longrightarrow_\beta\;\;
  \mathrm{swing}(\mathrm{cas\text{-}ok}(\sigma)),
\]
reducing two proof steps to one.

%% ── 22.1.3  CSpec ────────────────────────────────────────────
\subsection{The Cohomological Spectrum $\CSpec$}

\begin{construction}[Thread decomposition as a cover]
\label{con:thread-cover}
Partition the CFG into thread-local sub-graphs
$U_{\mathrm{enq}},U_{\mathrm{deq}}$ (one per operation).  This is a
cover in~$J$.  The \v{C}ech complex is:
\[
  C^0 \;=\;
  \mathrm{Proof}(U_{\mathrm{enq}})
  \;\times\;
  \mathrm{Proof}(U_{\mathrm{deq}}),
  \qquad
  C^1 \;=\;
  \mathrm{Proof}(U_{\mathrm{enq}}\cap U_{\mathrm{deq}}).
\]
Here $C^0$ contains each thread's \emph{local} correctness proof,
and $C^1$ contains proofs about shared state accessed by both threads.
\end{construction}

The coboundary $\delta^0:C^0\to C^1$ computes:
\[
  \delta^0(s_{\mathrm{enq}},s_{\mathrm{deq}})
  \;=\;
  s_{\mathrm{enq}}\big|_{U_{\mathrm{enq}}\cap U_{\mathrm{deq}}}
  \;-\;
  s_{\mathrm{deq}}\big|_{U_{\mathrm{enq}}\cap U_{\mathrm{deq}}}.
\]

\begin{proposition}[Obstruction = data race]
\label{prop:race-obstruction}
$H^1(\mathcal{C}_Q,\mathcal{F}_{\mathrm{inv}})\neq 0$ if and only
if the local invariants of \texttt{enqueue} and \texttt{dequeue}
are inconsistent on shared state---i.e., there exists a data race
or an atomicity violation.
\end{proposition}

\begin{example}[A concrete race]
Suppose a buggy implementation omits the ``swing tail'' step after
a successful CAS.  Then \texttt{enqueue} believes
$\texttt{tail}=\texttt{new\_node}$ on exit, while a concurrent
\texttt{dequeue} still sees the old tail.  The difference
\[
  \delta^0(\texttt{tail}=\texttt{new\_node},\;
           \texttt{tail}=\texttt{old\_tail})
  \;\neq\;0
\]
is a non-trivial 1-cocycle: the obstruction.
Mayer--Vietoris gives the exact sequence
\[
  \cdots\to
  H^0(U_{\mathrm{enq}})\oplus H^0(U_{\mathrm{deq}})
  \xrightarrow{\;\delta^0\;}
  H^0(U_{\mathrm{enq}}\cap U_{\mathrm{deq}})
  \to
  H^1(\mathcal{C}_Q,\mathcal{F}_{\mathrm{inv}})
  \to\cdots
\]
The non-vanishing $H^1$ localises the bug to the shared-state
interface.
\end{example}

\textbf{Repair.}\quad
Adding the swing-tail step makes $\delta^0=0$, whence
$H^1=0$ and the proofs glue into a global linearizability argument.

%% ── 22.1.4  GEval ────────────────────────────────────────────
\subsection{The Graded Evaluation $\GEval$}

\begin{construction}[Verification grade]
\label{con:verif-grade}
The overall grade is a product in the grade semiring:
\[
  G
  \;=\;
  g_{\mathrm{test}} \;\otG\; g_{\mathrm{type}} \;\otG\; g_{\mathrm{perf}},
\]
where each component lies in $\Grade=(-\infty,0]\cup\{-\infty\}$
(with $\gone=0$ perfect, $\gzero=-\infty$ vacuous).
\end{construction}

\begin{itemize}[nosep]
  \item \textbf{Test-coverage grade} $g_{\mathrm{test}}$.
    The trail is: unit tests ($g=-2.1$, covering basic enqueue/dequeue)
    $\to$ integration tests ($g=-1.3$, two-thread interleaving)
    $\to$ stress tests ($g=-0.4$, 64 threads, 10\textsuperscript{6}
    operations).  Convergence of random fuzzing pushes $g_{\mathrm{test}}$
    toward~$\gone$.

  \item \textbf{Type-safety grade} $g_{\mathrm{type}}$.
    With a verified simulation relation checked by a proof assistant,
    $g_{\mathrm{type}}=\gone$ (the proof is machine-checked).

  \item \textbf{Performance grade} $g_{\mathrm{perf}}$.
    Benchmarks show lock-free throughput exceeds a mutex-based queue
    by $3\times$ at 32~threads.  We assign $g_{\mathrm{perf}}=-0.2$
    (minor contention penalty at very high thread counts).
\end{itemize}

The overall grade is
$G = (-0.4)\otG 0 \otG (-0.2) = -0.4$
(bottleneck: test coverage not yet exhaustive).

\medskip

%% ─────────────────────────────────────────────────────────────
\section{Example 2: Insertion Sort}
\label{sec:isort}

To contrast, consider a sequential insertion sort on integer arrays.

\subsection{The 4-Tuple in Brief}

\textbf{$\STop$.}\quad
The site $\mathcal{C}_S$ has program points
$\{p_0,\ldots,p_5\}$ with linear control flow (one loop).  The
topology is trivial: the single function is its own cover.
The type presheaf assigns $\texttt{int[]}$ at every point.

\textbf{$\PStr$.}\quad
The specification type is
$\mathrm{Sorted}(a') \;\wedge\; \mathrm{Perm}(a,a')$.
The proof is a standard loop-invariant argument:
\[
  \dfrac{
    \Gamma\vdash_{\gone}
      \mathrm{inv}(i):\mathrm{Sorted}(a[0..i])
    \quad
    \Gamma\vdash_{\gone}
      \mathrm{step}(i):\mathrm{inv}(i)\Rightarrow\mathrm{inv}(i+1)
  }{
    \Gamma\vdash_{\gone}
      \mathrm{loop}:\mathrm{Sorted}(a[0..n])
  }
\]
$\beta$-reduction unfolds the loop to $n$ applications of
$\mathrm{step}$.

\textbf{$\CSpec$.}\quad
Since the program is sequential, the cover has a single element
(the whole function).  The \v{C}ech complex degenerates:
$C^0=\mathrm{Proof}(U)$, $C^1=0$.  Therefore $H^1=0$---no
obstructions.  The proof glues trivially.

\textbf{$\GEval$.}\quad
$g_{\mathrm{test}}=-0.3$ (exhaustive tests up to $n=10\,000$),
$g_{\mathrm{type}}=\gone$ (mechanised proof of the loop invariant),
$g_{\mathrm{perf}}=-1.5$ ($O(n^2)$ is slow for large~$n$).
Overall $G=-1.5$ (bottleneck: performance).

\begin{remark}[Comparing the two examples]
The lock-free queue has non-trivial $H^1$ that must be resolved;
insertion sort has $H^1=0$ by triviality of the cover.  Both
achieve high type-safety grades, but differ in their test-coverage
and performance profiles.  The 4-tuple captures these differences
precisely: the geometric (cohomological) and evaluative aspects
are orthogonal dimensions of program quality.
\end{remark}

\chapter{Natural Language Semantics}
\label{ch:ex-nl}

We now turn to natural language, the domain where the CTT extension
of JuGeo (Part~III) is most directly exercised.  Three examples
illustrate the 4-tuple at progressively greater interpretive depth:
a classic donkey sentence, a garden-path sentence, and a metaphor.

%% ─────────────────────────────────────────────────────────────
\section{Example 1: Donkey Anaphora}
\label{sec:donkey}

\begin{quote}
\emph{``Every farmer who owns a donkey beats it.''}
\end{quote}

This sentence is the touchstone of formal semantics: the pronoun
\emph{it} must be bound by the indefinite \emph{a donkey}, yet
the two sit in different clauses.  We show how each of the four
objects handles this.

%% ── STop ─────────────────────────────────────────────────────
\subsection{The Semantic Topos $\STop$}

\begin{construction}[NL site]
\label{con:nl-site}
The category $\mathcal{C}_{\mathrm{NL}}$ has:
\begin{itemize}[nosep]
  \item \textbf{Objects}: discourse segments at every granularity---morphemes,
    words, phrases, clauses, sentences, and multi-sentence discourses.
  \item \textbf{Morphisms}: inclusion of a sub-segment into a
    larger segment (e.g.\ the relative clause ``who owns a donkey''
    into the full sentence).
\end{itemize}
The topology~$J$ declares a covering family of a sentence to
be its immediate-constituent phrases:
\[
  \bigl\{\text{NP}_{\text{subj}}\hookrightarrow S,\;
         \text{VP}\hookrightarrow S\bigr\}
  \;\in\;J(S).
\]
\end{construction}

\begin{definition}[Semantic-type presheaf]
The presheaf $\mathcal{F}_{\sem}:\mathcal{C}_{\mathrm{NL}}^{\op}\to\Set$
assigns to each node its semantic type in the stratal universe tower:
\[
  \mathcal{F}_{\sem}(\text{``farmer''}) = U_{\sem},\quad
  \mathcal{F}_{\sem}(\text{``beats''}) = U_{\sem}\to U_{\sem}\to U_{\sem},\quad
  \mathcal{F}_{\sem}(\text{``it''}) = U_{\disc}.
\]
The pronoun lives at the \emph{discourse} stratum $U_{\disc}$
because its referent is determined by cross-clausal context.
\end{definition}

The internal logic of $\Sh(\mathcal{C}_{\mathrm{NL}},J)$ is
intuitionistic: truth of a proposition at a discourse segment~$d$
means truth in \emph{every extension} of~$d$.  This captures the
open-world nature of discourse: adding more sentences can never
revoke what is established.

%% ── PStr ─────────────────────────────────────────────────────
\subsection{The Proof Structure $\PStr$}

We build the semantic representation step by step.

\begin{construction}[Lexicon entries as types]
\label{con:lexicon}
\begin{align*}
  \text{every}   &:\; \Pi(P:U_{\sem}\to\Type).\,
                       \Pi(Q:U_{\sem}\to\Type).\,
                       \bigl(\Pi(x:U_{\sem}).\,P(x)\Rightarrow Q(x)\bigr), \\
  \text{farmer}  &:\; U_{\sem}\to\Prop, \\
  \text{owns}    &:\; U_{\sem}\to U_{\sem}\to\Prop, \\
  \text{a}       &:\; \Pi(P:U_{\sem}\to\Type).\,\Sigma(x:U_{\sem}).\,P(x), \\
  \text{donkey}  &:\; U_{\sem}\to\Prop, \\
  \text{beats}   &:\; U_{\sem}\to U_{\sem}\to\Prop, \\
  \text{it}      &:\; \text{(unresolved discourse referent)}.
\end{align*}
\end{construction}

\textbf{The anaphora problem.}\quad
Under na\"ive composition, ``a donkey'' introduces an existential
$\Sigma(y:U_{\sem}).\,\mathrm{donkey}(y)$ inside the relative clause,
but the pronoun ``it'' sits in the matrix clause and must refer
to~$y$.  The variable~$y$ is not in scope.

\begin{construction}[Resolution via the discourse monad]
\label{con:dmon}
Define the discourse monad
$\mathrm{DMon}(A) = \mathrm{DCtx}\to(A\times\mathrm{DCtx})$
where $\mathrm{DCtx}$ is a list of discourse referents.
The indefinite ``a donkey'' performs a monadic action that
\emph{pushes} a referent~$y$ onto $\mathrm{DCtx}$:
\[
  \llbracket\text{a donkey}\rrbracket
  \;=\;
  \lambda\,\gamma.\;
  \mathrm{let}\;y=\mathrm{fresh}(\gamma)\;\mathrm{in}\;
  (y,\;\gamma\!::\!y).
\]
The pronoun ``it'' is a \emph{lookup}:
$\llbracket\text{it}\rrbracket = \lambda\gamma.\,
  (\mathrm{top}(\gamma),\;\gamma)$.
\end{construction}

Composing inside $\mathrm{DMon}$, the full derivation is:
\[
  \dfrac{
    \dfrac{
      \Gamma\vdash_{\gone} \text{farmer}:U_{\sem}\to\Prop
      \qquad
      \Gamma\vdash_{\gone}
        [\text{who owns a donkey}] :
        \Pi(x).\,\mathrm{DMon}(\mathrm{donkey}(y)\wedge\mathrm{owns}(x,y))
    }{
      \Gamma\vdash_{\gone}
        [\text{farmer who owns a donkey}] :
        \Pi(x).\,\mathrm{DMon}(\mathrm{farmer}(x)\wedge
                                \mathrm{donkey}(y)\wedge
                                \mathrm{owns}(x,y))
    }
  }{
    \Gamma\vdash_{\gone}
      S : \Pi(x).\,
          \bigl(\mathrm{farmer}(x)\wedge\mathrm{donkey}(y)
                \wedge\mathrm{owns}(x,y)\bigr)
          \Rightarrow\mathrm{beats}(x,y)
  }
\]
The monadic plumbing threads $y$ from the relative clause to the
matrix clause, binding ``it'' to the donkey.

%% ── CSpec ────────────────────────────────────────────────────
\subsection{The Cohomological Spectrum $\CSpec$}

\begin{construction}[Clause cover]
\label{con:clause-cover}
Take the cover $\{U_{\mathrm{rel}},U_{\mathrm{mat}}\}$ where
$U_{\mathrm{rel}}$ = the relative clause and
$U_{\mathrm{mat}}$ = the matrix clause.
\end{construction}

\textbf{Before resolution.}\quad
The presheaf of discourse referents restricted to $U_{\mathrm{rel}}$
contains~$y$, but restricted to $U_{\mathrm{mat}}$ does not.
The coboundary is:
\[
  \delta^0(s_{\mathrm{rel}},s_{\mathrm{mat}})
  \;=\;
  s_{\mathrm{rel}}\big|_{U_{\mathrm{rel}}\cap U_{\mathrm{mat}}}
  - s_{\mathrm{mat}}\big|_{U_{\mathrm{rel}}\cap U_{\mathrm{mat}}}
  \;\neq\; 0,
\]
since $y\in s_{\mathrm{rel}}$ but $y\notin s_{\mathrm{mat}}$.
Therefore $H^1\neq 0$: the anaphora is an obstruction to local
gluing.

\textbf{After resolution.}\quad
Switching to the discourse-monad presheaf (which threads $y$
through $\mathrm{DCtx}$) makes $y$ available at
$U_{\mathrm{mat}}$.  Now $\delta^0=0$ and $H^1=0$: the
interpretation glues globally.

\begin{remark}
This is the cohomological content of Dynamic Predicate Logic
and DRT: the discourse monad is exactly the change-of-coefficients
that kills $H^1$.
\end{remark}

%% ── GEval ────────────────────────────────────────────────────
\subsection{The Graded Evaluation $\GEval$}

The sentence receives independent grades at each linguistic stratum:
\begin{align*}
  g_{\phon}  &= -0.05  &&\text{(phonologically well-formed, minor coda cluster)},\\
  g_{\synt}  &= \gone   &&\text{(grammatical)},\\
  g_{\sem}   &= -0.1   &&\text{(coherent once anaphora is resolved)},\\
  g_{\prag}  &= -0.3   &&\text{(pragmatically odd: ``beats'' is harsh)}.
\end{align*}
The overall grade is
$T = g_{\phon}\otG g_{\synt}\otG g_{\sem}\otG g_{\prag}
   = -0.3$
(bottleneck: pragmatic felicity).
In the normalised scale where $\gone=0$ is perfect, a grade of
$-0.3$ corresponds to ``well-formed but stylistically marked.''

\medskip

%% ─────────────────────────────────────────────────────────────
\section{Example 2: Garden-Path Sentence}
\label{sec:garden-path}

\begin{quote}
\emph{``The horse raced past the barn fell.''}
\end{quote}

This sentence is grammatical (a reduced relative clause:
``The horse [that was] raced past the barn fell''),
but readers initially parse ``raced'' as the main verb and
reach ``fell'' with no syntactic slot for it.

\textbf{$\STop$.}\quad
The site is the same NL category $\mathcal{C}_{\mathrm{NL}}$.
The initial parse assigns the cover
$\{$NP:``the horse'', VP:``raced past the barn''$\}$
to the substring before ``fell.''

\textbf{$\PStr$ (first attempt).}\quad
The first derivation tree types ``raced'' as a main verb:
$\Gamma\vdash_{\gone}\text{raced}:U_{\sem}\to U_{\sem}$.
When ``fell'' arrives, the context $\Gamma$ has no open
argument position.  The typing derivation fails:
$\Gamma\nvdash\text{fell}:A$ for any~$A$.

\textbf{$\CSpec$ (first attempt).}\quad
The cover $\{U_{\text{pre-fell}}, U_{\text{fell}}\}$ has
$H^1\neq 0$: the local parse of ``raced past the barn'' does not
extend to include ``fell.''  The obstruction is localised at the
word boundary before ``fell.''

\textbf{Reanalysis.}\quad
The parser \emph{recovers} the cover: ``raced past the barn'' is
re-typed as a reduced relative clause (participial modifier), freeing
a VP slot for ``fell.''
Under the new cover,
$\Gamma\vdash_{\gone}[\text{the horse raced past the barn}]:
  \mathrm{NP}$
and
$\Gamma\vdash_{\gone}\text{fell}:\mathrm{NP}\to S$,
giving $H^1=0$.

\textbf{$\GEval$.}\quad
After reanalysis:
$g_{\synt}=-0.6$ (syntactically legal but highly marked),
$g_{\sem}=\gone$ (semantically clear),
$g_{\prag}=-0.8$ (pragmatically confusing).
Overall $T=-0.8$: a difficult sentence.

\medskip

%% ─────────────────────────────────────────────────────────────
\section{Example 3: Metaphor}
\label{sec:metaphor}

\begin{quote}
\emph{``Juliet is the sun.''}  --- Shakespeare, \emph{Romeo and Juliet}
\end{quote}

\textbf{$\STop$.}\quad
The site is $\mathcal{C}_{\mathrm{NL}}$; the relevant presheaf
is $\mathcal{F}_{\sem}$ assigning semantic types.  At the literal
level, $\mathcal{F}_{\sem}(\text{Juliet})=\texttt{Person}$ and
$\mathcal{F}_{\sem}(\text{sun})=\texttt{Star}$.

\textbf{$\PStr$.}\quad
A literal identity $\GId_{\gone}(\texttt{Entity},\text{Juliet},\text{sun})$
is uninhabited: a person is not a star.  But a \emph{graded} identity
\[
  \GId_g(\texttt{Entity},\text{Juliet},\text{sun})
  \quad\text{with}\quad g\approx -0.35
\]
is inhabited when we restrict to the sub-type of entities possessing
\emph{radiance, warmth, centrality}.  The grade $g$ measures the
proportion of shared structure.

The typing derivation is:
\[
  \dfrac{
    \Gamma\vdash_{\gone}\text{Juliet}:\texttt{Person}
    \qquad
    \Gamma\vdash_{\gone}\text{sun}:\texttt{Star}
    \qquad
    \Gamma\vdash_{-0.35}
      \mathrm{radiance}:\texttt{Person}\cap\texttt{Star}
  }{
    \Gamma\vdash_{-0.35}
      \text{``Juliet is the sun''}
      :\GId_{-0.35}(\texttt{Entity},\text{Juliet},\text{sun})
  }
\]

\textbf{$\CSpec$.}\quad
$H^1=0$ at the metaphorical grade: the identification glues
across the source and target domains because we have restricted
to shared features.  At the literal grade ($g=\gone$), $H^1\neq 0$:
the literal identification fails.

\textbf{$\GEval$.}\quad
$g_{\phon}=\gone$ (simple monosyllables),
$g_{\synt}=\gone$ (subject--copula--predicate),
$g_{\sem}=-0.35$ (metaphorical, not literal),
$g_{\poet}=-0.05$ (vivid imagery, canonical).
The \emph{poetic} grade is near-perfect; the semantic grade reflects
the deliberate non-literalness.  Overall
$T=-0.35$ in the semantic dimension, but the evaluative judgment
is that this is \emph{excellent} poetry: the aesthetic grade overrides
the semantic gap.

\begin{remark}[Graded identity as the algebra of figurative language]
The metaphor example illustrates the central thesis of CTT:
figurative language is not a ``deviation'' from literal meaning
but a principled use of the graded identity type $\GId_g$ at a
grade $\gzero < g < \gone$.  Simile, metonymy, synecdoche, and
irony each correspond to different structural restrictions on
the shared sub-type that witnesses the graded identification.
\end{remark}

\chapter{Poetry Generation and Analysis}
\label{ch:ex-poetry}

Poetry is the domain where all four objects are most densely
populated.  We present three examples: generating a Shakespearean
sonnet, analysing a Keats sonnet, and diagnosing a flawed poem.

%% ─────────────────────────────────────────────────────────────
\section{Example 1: Generating a Shakespearean Sonnet}
\label{sec:sonnet-gen}

\textbf{Prompt.}\; Generate a Shakespearean sonnet on
\emph{the passage of time}.

%% ── STop ─────────────────────────────────────────────────────
\subsection{The Semantic Topos $\STop$}

\begin{construction}[Poetic site]
\label{con:poet-site}
The category $\mathcal{C}_{\mathrm{poet}}$ refines
$\mathcal{C}_{\mathrm{NL}}$ with poetic structure:
\begin{itemize}[nosep]
  \item \textbf{Objects}: syllables, metrical feet, lines,
    quatrains (Q1--Q3), couplet (C), and the full sonnet~$S$.
  \item \textbf{Morphisms}: inclusions
    $\text{syll}\hookrightarrow\text{foot}\hookrightarrow
     \text{line}\hookrightarrow\text{quatrain}\hookrightarrow S$.
\end{itemize}
The topology declares the 14-line decomposition as a cover of~$S$:
\[
  \{l_1,l_2,\ldots,l_{14}\hookrightarrow S\}\in J(S),
\]
and at a coarser level,
$\{Q_1,Q_2,Q_3,C\hookrightarrow S\}\in J(S)$.
\end{construction}

The presheaf $\mathcal{F}_{\mathrm{ling}}$ assigns to each level
the available linguistic material: at the syllable level, phonemes
from the CMU dictionary; at the line level, syntactic frames and
semantic content; at the quatrain level, rhetorical functions
(exposition, complication, turn, resolution).

%% ── PStr ─────────────────────────────────────────────────────
\subsection{The Proof Structure $\PStr$}

\begin{construction}[Sonnet type]
\label{con:sonnet-type}
The specification type is:
\[
  \mathrm{Sonnet}(\theta, \rho, \mu)
  \;:\equiv\;
  \Sigma(l_{1..14}:\mathrm{Line}^{14}).\;
  \mathrm{Theme}(l_{1..14},\theta)
  \;\wedge\;
  \mathrm{Rhyme}(l_{1..14},\rho)
  \;\wedge\;
  \mathrm{Meter}(l_{1..14},\mu),
\]
where $\theta=\text{time}$,
$\rho=\text{ABAB\,CDCD\,EFEF\,GG}$,
$\mu=\text{iambic pentameter}$.
\end{construction}

\textbf{Step 1: Choose a conceit.}\quad
The $\FMLM$ proposes three conceits---time as \emph{river},
\emph{season}, \emph{hourglass}---and scores each against the
theme type.  ``River'' wins ($g=-0.05$) for richness of imagery.

\textbf{Step 2: Generate line by line.}\quad
Each line~$l_i$ is typed in context:
\[
  \Gamma,\;l_1,\ldots,l_{i-1}
  \;\vdash_{g_i}\;
  l_i : \mathrm{Line}(\theta_i,\rho_i,\mu),
\]
where $\theta_i$ is the local sub-theme and $\rho_i$ constrains the
end-rhyme.  The $\FMLM$ fills each line via masked-slot iteration:
\begin{enumerate}[nosep]
  \item Start with a template:
    \texttt{[\_\_\_] of time [\_\_\_] like a [\_\_\_]}.
  \item Generate candidates for each slot from WordNet/FrameNet.
  \item Score: $g_{ij}=\mathrm{grade}(\Gamma\vdash l_i[w_j/M_k]:T_k(w_j))$.
  \item Select $\arg\max_j\,g_{ij}$.
  \item Iterate until grades stabilise.
\end{enumerate}

\textbf{Step 3: Meter check.}\quad
Each line must satisfy $\mathrm{Meter}(l_i,\mu)$: exactly 10
syllables with stress pattern $\cup/\cup/\cup/\cup/\cup/$ (iambic
pentameter).  The $\HPLF$ voice bundle at the $\phon$ stratum
verifies this:
\[
  G_i^{\phon}(l_i)
  \;=\;
  \gone \;\text{if stress pattern matches},
  \qquad
  g < \gone\;\text{otherwise (substitutions allowed at a cost)}.
\]
A feminine ending (11th unstressed syllable) costs $g=-0.1$.

\textbf{Step 4: Rhyme check.}\quad
End-words of paired lines must satisfy
$\GId_{g_r}(U_{\phon},\mathrm{coda}(w_i),\mathrm{coda}(w_j))$
with $g_r$ close to $\gone$.  Perfect rhyme: $g_r=\gone$.
Near-rhyme (e.g.\ ``love''/``prove''): $g_r\approx -0.15$.

\medskip
Suppose the $\FMLM$ produces the opening quatrain:
\begin{verse}
  The river of our years flows ever on,\\
  Through meadows green where youthful pleasures play,\\
  Until the golden afternoon is gone\\
  And twilight steals the colours from the day.
\end{verse}
Each line has been typed at grade $g_i\in[-0.15,\,0]$; the quatrain
as a whole satisfies $\mathrm{Rhyme}(Q_1,\text{ABAB})$ and
$\mathrm{Meter}(Q_1,\mu)$.

%% ── CSpec ────────────────────────────────────────────────────
\subsection{The Cohomological Spectrum $\CSpec$}

With the quatrain cover $\{Q_1,Q_2,Q_3,C\}$:

\begin{example}[Obstruction at the Q2--Q3 boundary]
Suppose line~5 (start of Q2) rhymes correctly with line~7 but
introduces the word ``frozen,'' breaking the river metaphor
(rivers don't freeze in this conceit).  The \v{C}ech coboundary
between Q1 and Q2 is:
\[
  \delta^0(s_{Q_1},s_{Q_2})
  \;=\;
  \text{``river (flowing)''}\big|_{Q_1\cap Q_2}
  \;-\;
  \text{``frozen (static)''}\big|_{Q_1\cap Q_2}
  \;\neq\;0.
\]
$H^1\neq 0$: the metaphor does not glue across quatrains.

\textbf{Repair}: replace ``frozen'' with ``narrowing'' (the
river narrows as time shrinks).  Now $\delta^0=0$ and
$H^1=0$.
\end{example}

\begin{example}[Volta obstruction]
The volta at line~9 must shift the argument (from ``time
passes'' to ``but something endures'').  If Q3 merely
continues the lament without a turn, the structural
constraint $\mathrm{Volta}(l_9)$ is unsatisfied---an
obstruction at the $Q_2$--$Q_3$ boundary in the rhetorical
presheaf.  Repair: rewrite line~9 as a ``But'' or ``Yet''
clause.
\end{example}

%% ── GEval ────────────────────────────────────────────────────
\subsection{The Graded Evaluation $\GEval$}

After repair, the final grades (in the $\HPLF$ framework) are:
\begin{align*}
  g_{\mathrm{meter}}   &= -0.08  &&\text{(one feminine ending in line 12)},\\
  g_{\mathrm{rhyme}}   &= -0.05  &&\text{(one near-rhyme: ``on''/``gone'')},\\
  g_{\mathrm{imagery}} &= -0.12  &&\text{(one clich\'e: ``golden afternoon'')},\\
  g_{\mathrm{arc}}     &= -0.10  &&\text{(volta could be sharper)}.
\end{align*}
The overall harmony:
\[
  \mathrm{Harm}^{\HPLF}
  \;=\;
  g_{\mathrm{meter}} \otG g_{\mathrm{rhyme}} \otG
  g_{\mathrm{imagery}} \otG g_{\mathrm{arc}}
  \;=\; -0.12
\]
(bottleneck: imagery).  The sonnet is competent but not
outstanding; the clich\'e drags the grade down.

\medskip

%% ─────────────────────────────────────────────────────────────
\section{Example 2: Analysing a Keats Sonnet}
\label{sec:keats}

\begin{quote}
\emph{``When I have fears that I may cease to be \\
Before my pen has glean'd my teeming brain\ldots''}\\
\hfill --- Keats, \emph{Sonnet VII}
\end{quote}

\textbf{$\STop$.}\quad
The poetic site is the same $\mathcal{C}_{\mathrm{poet}}$ as above,
with the 14-line cover.  The presheaf assigns Keats's actual
linguistic material at each node.

\textbf{$\PStr$.}\quad
The type $\mathrm{Sonnet}(\text{mortality},\text{ABAB}\ldots,\mu)$
is inhabited by the actual poem.  Key typing observations:
\begin{itemize}[nosep]
  \item ``Glean'd my teeming brain'' uses $\GId_{-0.2}$
    (metaphor: mind = field to be harvested).
  \item ``High-pil\`ed books'' carries
    $\mathrm{Ico}(\text{books},\text{grain})$: an iconicity between
    physical stacking and intellectual accumulation.
  \item The volta at line~9 (``And when I feel\ldots'') introduces
    the beloved, shifting from self to other.
\end{itemize}

\textbf{$\CSpec$.}\quad
$H^1=0$ throughout: every metaphor, every rhyme, every metrical
choice glues perfectly.  The \v{C}ech complex is exact---the
hallmark of a masterwork.

\textbf{$\GEval$.}\quad
$g_{\mathrm{meter}}=\gone$ (flawless iambic pentameter),
$g_{\mathrm{rhyme}}=\gone$ (all perfect rhymes),
$g_{\mathrm{imagery}}=-0.02$ (near-perfect; every image is
fresh),
$g_{\mathrm{arc}}=-0.03$ (the volta is slightly abrupt).
$\mathrm{Harm}^{\HPLF}=-0.03$.  This is essentially a
grade-$\gone$ poem.

\medskip

%% ─────────────────────────────────────────────────────────────
\section{Example 3: A Flawed Poem}
\label{sec:bad-poem}

Consider the following attempt:
\begin{verse}
  The clock goes tick and also tock,\\
  It tells the time just like a clock,\\
  The hours pass both fast and slow,\\
  And that is all I need to know.
\end{verse}

\textbf{$\STop$.}\quad
The site has 4 lines; the cover is $\{l_1,l_2,l_3,l_4\}$.

\textbf{$\PStr$.}\quad
The type is $\mathrm{Quatrain}(\text{time},\text{AABB},\mu')$
for some meter~$\mu'$.  The poem inhabits this type at a low
grade:
\begin{itemize}[nosep]
  \item Line~2 is a tautology (``like a clock'') --- the semantic
    content is $\GId_{\gone}(\text{clock},\text{clock})$, which
    is trivial and adds no information.
  \item Line~4 is a clich\'e (``all I need to know'').
\end{itemize}

\textbf{$\CSpec$.}\quad
$H^1\neq 0$ at the $l_1$--$l_2$ boundary: line~2 repeats
the content of line~1 without advancing the argument.  The
coboundary
$\delta^0(s_1,s_2) = \text{clock}-\text{clock}$ would be zero
\emph{semantically}, but in the poetic presheaf (which values
\emph{novelty}), the redundancy registers as a non-trivial
cocycle.  Similarly, $H^1\neq 0$ at $l_3$--$l_4$ (clich\'e is
an obstruction to freshness).

\textbf{$\GEval$.}\quad
$g_{\mathrm{meter}}=-0.3$ (irregular stress),
$g_{\mathrm{rhyme}}=-0.5$ (``clock''/``clock'' is an identical
rhyme, penalised),
$g_{\mathrm{imagery}}=-2.0$ (no fresh imagery),
$g_{\mathrm{arc}}=-1.5$ (no development).
$\mathrm{Harm}^{\HPLF}=-2.0$ (bottleneck: imagery).
The poem is essentially at the $\gzero$ end of the scale.

\begin{remark}[Diagnosis vs.\ generation]
The 4-tuple does not merely \emph{judge} the bad poem; it
\emph{localises} the failures.  $\CSpec$ points to the
$l_1$--$l_2$ boundary and the $l_3$--$l_4$ boundary;
$\GEval$ identifies imagery as the worst dimension.  A
generation system could use this diagnosis to \emph{repair}
the poem: replace line~2 with a non-tautological elaboration,
replace line~4's clich\'e, and re-check $H^1$.
\end{remark}

\chapter{Mathematical Ideation and Proof}
\label{ch:ex-math}

Mathematics is the domain where $\PStr$ is richest and $\GEval$
most contested.  We work through the irrationality of $\sqrt{2}$
in full 4-tuple detail, then sketch mathematical ideation as
proof search.

%% ─────────────────────────────────────────────────────────────
\section{Example 1: $\sqrt{2}$ Is Irrational}
\label{sec:sqrt2}

\subsection{The Semantic Topos $\STop$}

\begin{construction}[Mathematical site]
\label{con:math-site}
The category $\mathcal{C}_{\mathrm{math}}$ has:
\begin{itemize}[nosep]
  \item \textbf{Objects}: mathematical statements---definitions
    ($\mathrm{gcd}$, ``rational'', ``irrational''), lemmas
    (``if $p^2$ is even then $p$ is even''), and theorems
    ($\sqrt{2}\notin\mathbb{Q}$).
  \item \textbf{Morphisms}: logical dependency.  A morphism
    $L\to T$ exists when lemma~$L$ is used in the proof of
    theorem~$T$.
\end{itemize}
The topology declares a cover of a theorem to be a
\emph{complete} set of dependencies:
\[
  \bigl\{\text{Def}(\gcd),\;
         \text{Lem}(\text{even-square}),\;
         \text{Ax}(\mathbb{Z}\text{-domain})
  \bigr\}
  \;\in\;J\bigl(\sqrt{2}\notin\mathbb{Q}\bigr).
\]
\end{construction}

The presheaf $\mathcal{F}_{\mathrm{obj}}$ assigns to each node the
set of mathematical objects in scope: at $\text{Def}(\gcd)$, the
objects $\{p,q,\gcd(p,q)\}$; at the theorem node, the full
collection $\{p,q,k,\gcd,\sqrt{2}\}$.

\subsection{The Proof Structure $\PStr$}

\begin{construction}[Irrationality type]
The specification is:
\[
  \mathrm{Irrational}(\sqrt{2})
  \;:\equiv\;
  \neg\,\Sigma(p:\mathbb{Z}).\,\Sigma(q:\mathbb{Z}_{>0}).\,
  \GId_{\gone}\!\bigl(\mathbb{R},\;\sqrt{2},\;p/q\bigr).
\]
\end{construction}

The proof term inhabiting this type is a \emph{proof by
contradiction}.  We present the full derivation tree.

\begin{construction}[Proof tree]
\label{con:sqrt2-proof}
\[
\dfrac{
  \dfrac{
    \dfrac{
      [h:\sqrt{2}=p/q,\;\gcd(p,q)=1]
    }{
      2q^2=p^2
    }\text{\scriptsize(algebra)}
    \qquad
    \dfrac{
      2q^2=p^2
    }{
      2\mid p^2
    }\text{\scriptsize(div)}
    \qquad
    \dfrac{
      2\mid p^2
    }{
      2\mid p
    }\text{\scriptsize(prime)}
  }{
    \Gamma,h\vdash_{\gone} e_1:2\mid p
  }
  \quad
  \dfrac{
    \dfrac{
      p=2k
    }{
      2q^2=4k^2
    }\text{\scriptsize(subst)}
    \quad
    \dfrac{
      q^2=2k^2
    }{
      2\mid q
    }\text{\scriptsize(same)}
  }{
    \Gamma,h\vdash_{\gone} e_2:2\mid q
  }
}{
  \dfrac{
    \Gamma,h\vdash_{\gone} (e_1,e_2):2\mid p\;\wedge\;2\mid q
    \qquad
    \Gamma\vdash_{\gone} c:\gcd(p,q)=1
  }{
    \Gamma\vdash_{\gone}\;\bot
  }\text{\scriptsize(contra)}
}
\]
The overall term is
$\lambda h.\;\mathrm{absurd}(\mathrm{gcd\text{-}contra}(e_1,e_2,c))
 :\mathrm{Irrational}(\sqrt{2})$.
\end{construction}

\textbf{$\beta$-reduction.}\quad
Inlining $k=p/2$ into the second branch and simplifying
$2q^2=4(p/2)^2\longrightarrow q^2=2(p/2)^2$ is a $\beta$-step
that shortens the derivation.

\subsection{The Cohomological Spectrum $\CSpec$}

\textbf{Correct proof: $H^1=0$.}\quad
Take the cover
$\{U_1:\text{``$p$ is even''},\;U_2:\text{``$q$ is even''}\}$.
The local proofs (even-square lemma applied to $p^2=2q^2$ and
$q^2=2k^2$) glue at the overlap
$U_1\cap U_2=\{\gcd(p,q)=1\}$, since having both $2\mid p$ and
$2\mid q$ contradicts $\gcd(p,q)=1$.  The \v{C}ech coboundary
vanishes: $\delta^0=0$, so $H^1=0$.

\textbf{Proof with a gap: $H^1\neq 0$.}\quad
Suppose we omit the hypothesis $\gcd(p,q)=1$.  Then the local
proofs ``$2\mid p$'' and ``$2\mid q$'' are both valid, but the
contradiction step has no fuel---$\gcd(p,q)$ is unconstrained.

\begin{proposition}[Gap = non-trivial $H^1$]
The proofs over $U_1$ and $U_2$ define local sections of the
proof presheaf, but $\delta^0(s_1,s_2)\neq 0$ at the overlap
$U_1\cap U_2$: the two local proofs do not compose into a
global contradiction because the necessary compatibility datum
($\gcd=1$) is missing.
\end{proposition}

Repair: adding ``without loss of generality, assume
$\gcd(p,q)=1$'' restores the missing section, making
$\delta^0=0$ and $H^1=0$.

\subsection{The Graded Evaluation $\GEval$}

\begin{align*}
  g_{\mathrm{rigor}}       &= \gone   &&\text{(every step is valid)},\\
  g_{\mathrm{elegance}}    &= -0.3  &&\text{(standard; not the most elegant proof)},\\
  g_{\mathrm{novelty}}     &= -3.0  &&\text{(2500 years old)},\\
  g_{\mathrm{significance}}&= -0.1  &&\text{(foundational result)}.
\end{align*}
Overall $G = \gone\otG(-0.3)\otG(-3.0)\otG(-0.1)=-3.0$
(bottleneck: novelty).  This is correct but unsurprising---exactly
the evaluation a journal referee would give.

\medskip

%% ─────────────────────────────────────────────────────────────
\section{Example 2: Mathematical Ideation}
\label{sec:ideation}

Consider a mathematician who notices:
$1+3=4=2^2$,\;
$1+3+5=9=3^2$,\;
$1+3+5+7=16=4^2$.

\textbf{$\STop$.}\quad
The site $\mathcal{C}_{\mathrm{conj}}$ has objects = observed
equalities and a conjectured theorem node
$T:``\sum_{i=1}^n(2i{-}1)=n^2$''.  The cover of $T$ includes
the examples plus the inductive step (which is not yet proved).

\textbf{$\PStr$.}\quad
The type is $\Pi(n:\mathbb{N}).\,\GId_{\gone}(\mathbb{N},
  \sum_{i=1}^n(2i{-}1),\,n^2)$.
A full term requires induction:
\[
  \dfrac{
    \Gamma\vdash_{\gone}\mathrm{base}:\sum_{i=1}^1(2i{-}1)=1^2
    \qquad
    \Gamma\vdash_{\gone}\mathrm{step}:
      \bigl(\textstyle\sum_{i=1}^n(2i{-}1)=n^2\bigr)
      \Rightarrow
      \bigl(\textstyle\sum_{i=1}^{n+1}(2i{-}1)=(n{+}1)^2\bigr)
  }{
    \Gamma\vdash_{\gone}\mathrm{ind}(\mathrm{base},\mathrm{step}):
      \Pi(n).\,\textstyle\sum_{i=1}^n(2i{-}1)=n^2
  }
\]
Before the inductive step is found, the mathematician has only
the examples.  These are \emph{partial evidence} at grades
$g<\gone$: three witnesses do not constitute a proof.

\textbf{$\CSpec$.}\quad
Before the inductive proof, $H^1\neq 0$: the examples cover
$n=2,3,4$ but do not glue into a proof for all~$n$.  After
finding the inductive step, $H^1=0$: the proof glues globally.
This is a precise formalisation of the difference between
\emph{conjecture} and \emph{theorem}.

\textbf{$\GEval$.}\quad
Before proof: $g_{\mathrm{rigor}}=-2.0$ (empirical only).
After proof: $g_{\mathrm{rigor}}=\gone$,
$g_{\mathrm{elegance}}=-0.1$ (clean one-line induction),
$g_{\mathrm{novelty}}=-2.5$ (well-known),
$g_{\mathrm{significance}}=-0.5$ (useful but not deep).

\begin{remark}[Ideation as proof search]
Mathematical ideation---the process of noticing a pattern and
conjecturing a theorem---is precisely the $\FMLM$ applied to
$\PStr$: generate candidate proof terms (``try induction,''
``try contradiction,'' ``try a clever substitution''), score
each against the type, and iterate.  The grade $g$ rises from
$\gzero$ (no evidence) through intermediate values (empirical
support) to $\gone$ (complete proof).  $\CSpec$ tracks which
parts of the proof are missing at each stage.
\end{remark}

\chapter{Cross-Domain Isomorphisms}
\label{ch:isomorphisms}

The preceding four chapters instantiated the 4-tuple
$(\STop,\PStr,\CSpec,\GEval)$ in four domains: programming,
natural language, poetry, and mathematics.  This chapter
demonstrates that the instantiations are not merely analogous
but \emph{functorially related}: structure-preserving maps
connect every pair.

%% ─────────────────────────────────────────────────────────────
\section{The Category of Domains}
\label{sec:dom-cat}

\begin{definition}[Domain 4-tuple]
A \emph{domain instantiation} is a quadruple
$\mathcal{D}=(\STop_{\mathcal{D}},\PStr_{\mathcal{D}},
 \CSpec_{\mathcal{D}},\GEval_{\mathcal{D}})$
where each component is the restriction of the universal
4-tuple to the domain's site, types, cohomology, and grade
components.
\end{definition}

\begin{definition}[Domain morphism]
A \emph{domain morphism} $\Phi:\mathcal{D}_1\to\mathcal{D}_2$
is a quadruple of functors/natural transformations:
\[
  \Phi = (\Phi_{\STop},\;\Phi_{\PStr},\;\Phi_{\CSpec},\;\Phi_{\GEval})
\]
such that (i)~$\Phi_{\STop}$ is a morphism of sites,
(ii)~$\Phi_{\PStr}$ preserves typing derivations,
(iii)~$\Phi_{\CSpec}$ maps cocycles to cocycles and
coboundaries to coboundaries (hence induces a map on $H^1$),
and (iv)~$\Phi_{\GEval}$ is a semiring homomorphism on grades.
\end{definition}

We write $\mathbf{Dom}$ for the category of domain instantiations
and domain morphisms.  The four domains studied form a cycle:
\[
\begin{tikzcd}
  \mathcal{D}_{\mathrm{prog}}
    \arrow[r,"\Phi_{\mathrm{prog}\to\mathrm{NL}}"]
  & \mathcal{D}_{\mathrm{NL}}
    \arrow[d,"\Phi_{\mathrm{NL}\to\mathrm{poet}}"]
  \\
  \mathcal{D}_{\mathrm{math}}
    \arrow[u,"\Phi_{\mathrm{math}\to\mathrm{prog}}"]
  & \mathcal{D}_{\mathrm{poet}}
    \arrow[l,"\Phi_{\mathrm{poet}\to\mathrm{math}}"]
\end{tikzcd}
\]

%% ─────────────────────────────────────────────────────────────
\section{Programs $\leftrightarrow$ Sentences}
\label{sec:prog-nl}

\begin{construction}[The functor $\Phi_{\mathrm{prog}\to\mathrm{NL}}$]
\label{con:prog-nl}
\begin{itemize}[nosep]
  \item $\Phi_{\STop}$: program points $\mapsto$ discourse segments;
    basic blocks $\mapsto$ phrases; the CFG cover $\mapsto$ constituent
    cover.
  \item $\Phi_{\PStr}$: specifications $\mapsto$ propositions
    (a spec ``output is sorted'' becomes the proposition ``the
    array is sorted''); programs $\mapsto$ utterances (the code
    \emph{is} a statement about what it does).
  \item $\Phi_{\CSpec}$: bugs $\mapsto$ ungrammaticality.
    A data race ($H^1\neq 0$ in the thread cover) maps to a
    binding failure ($H^1\neq 0$ in the clause cover).
  \item $\Phi_{\GEval}$: test quality $\mapsto$ linguistic quality;
    type safety $\mapsto$ grammaticality;
    performance $\mapsto$ pragmatic efficiency.
\end{itemize}
\end{construction}

\begin{theorem}[Faithfulness]
\label{thm:prog-nl-faithful}
$\Phi_{\mathrm{prog}\to\mathrm{NL}}$ is a faithful functor:
distinct programs map to distinct ``sentences'' (no two
structurally different programs have identical NL images).
\end{theorem}

\begin{proof}[Proof sketch]
Faithfulness follows from the injectivity of $\Phi_{\PStr}$ on
derivation trees: different proof trees (Hoare-logic derivations)
yield different typing derivations in $\PStr_{\mathrm{NL}}$
because each atomic Hoare rule maps to a distinct NL typing
rule.
\end{proof}

\begin{remark}[Not full]
The functor is not full: natural-language sentences expressing
emotions, social conventions, or fictional narratives have no
corresponding program.  The image of $\Phi_{\mathrm{prog}\to\mathrm{NL}}$
is the sublanguage of \emph{declarative descriptions of
computational processes}.
\end{remark}

%% ─────────────────────────────────────────────────────────────
\section{Sentences $\leftrightarrow$ Poems}
\label{sec:nl-poet}

\begin{construction}[The functor $\Phi_{\mathrm{NL}\to\mathrm{poet}}$]
\label{con:nl-poet}
\begin{itemize}[nosep]
  \item $\Phi_{\STop}$: $\mathcal{C}_{\mathrm{NL}}\hookrightarrow
    \mathcal{C}_{\mathrm{poet}}$ by adding metrical, rhyme, and
    stanzaic objects.  Every NL site embeds into the poetic site as
    the ``prose sub-site.''
  \item $\Phi_{\PStr}$: a sentence type $A$ maps to
    $A\wedge\mathrm{Meter}(\mu)\wedge\mathrm{Rhyme}(\rho)\wedge
    \mathrm{Form}(f)$---the sentence plus additional constraints.
    Every grammatical sentence is a trivial ``poem'' (free verse
    with no metre or rhyme constraints).
  \item $\Phi_{\CSpec}$: NL obstructions embed as poetic
    obstructions; poetry has \emph{additional} obstructions
    (broken metre, failed rhyme, weak volta).
  \item $\Phi_{\GEval}$: the NL grade components
    $g_{\phon},g_{\synt},g_{\sem},g_{\prag}$ embed into the
    larger poetic grade vector, which adds
    $g_{\mathrm{meter}},g_{\mathrm{rhyme}},g_{\mathrm{imagery}},
    g_{\mathrm{arc}}$.
\end{itemize}
\end{construction}

\begin{theorem}[Faithful embedding]
\label{thm:nl-poet-embed}
$\Phi_{\mathrm{NL}\to\mathrm{poet}}$ is faithful: every
grammatical sentence embeds as a (trivially constrained) poem.
The embedding is not full: poems carry structure that arbitrary
prose does not.
\end{theorem}

The key insight: poetry is \emph{NL plus extra strata in $\PStr$
and extra components in $\GEval$}.  The CTT stratal tower
$U_{\phon}:U_{\morph}:\cdots:U_{\poet}:U_{\aesth}$ makes this
precise---poetry activates the upper strata that prose leaves
dormant.

%% ─────────────────────────────────────────────────────────────
\section{Proofs $\leftrightarrow$ Programs (Curry--Howard)}
\label{sec:math-prog}

\begin{construction}[The functor $\Phi_{\mathrm{math}\to\mathrm{prog}}$]
\label{con:math-prog}
This is the classical Curry--Howard correspondence, enriched with
$\CSpec$ and $\GEval$:
\begin{itemize}[nosep]
  \item $\Phi_{\STop}$: the dependency graph of lemmas $\mapsto$
    the call graph of functions.  A cover (complete set of
    dependencies) $\mapsto$ a cover (complete set of callee
    implementations).
  \item $\Phi_{\PStr}$: propositions $\mapsto$ types;
    proofs $\mapsto$ programs; the derivation tree of a proof
    $\mapsto$ the typing derivation of a program.
  \item $\Phi_{\CSpec}$: a proof gap (missing lemma, $H^1\neq 0$)
    $\mapsto$ a bug (missing case, $H^1\neq 0$).
    Exact correspondence: non-vanishing $H^1$ in the mathematical
    site maps to non-vanishing $H^1$ in the program site.
  \item $\Phi_{\GEval}$: rigor $\mapsto$ type safety;
    elegance $\mapsto$ code cleanliness;
    novelty $\mapsto$ algorithmic innovation;
    significance $\mapsto$ practical impact.
\end{itemize}
\end{construction}

\begin{theorem}[Equivalence on $\CSpec$]
\label{thm:ch-cspec}
Under Curry--Howard,
$H^1(\mathcal{C}_{\mathrm{math}},\mathcal{F}_{\mathrm{proof}})
\cong
H^1(\mathcal{C}_{\mathrm{prog}},\mathcal{F}_{\mathrm{inv}})$.
Proof gaps and program bugs are cohomologically identical.
\end{theorem}

\begin{proof}[Proof sketch]
$\Phi_{\CSpec}$ maps the \v{C}ech complex of the lemma cover to
the \v{C}ech complex of the callee cover.  Since Curry--Howard
is an isomorphism on derivation trees, it induces an isomorphism
on cochain groups, hence on cohomology.
\end{proof}

%% ─────────────────────────────────────────────────────────────
\section{Poems $\leftrightarrow$ Proofs}
\label{sec:poet-math}

This is the most surprising correspondence.

\begin{construction}[The functor $\Phi_{\mathrm{poet}\to\mathrm{math}}$]
\label{con:poet-math}
\begin{itemize}[nosep]
  \item \textbf{Form $\mapsto$ proof structure.}\quad
    The stanzaic form of a sonnet (three quatrains + couplet)
    $\mapsto$ the logical structure of a proof (three lemmas +
    final theorem).
  \item \textbf{Meter $\mapsto$ proof rhythm.}\quad
    The alternation of stressed and unstressed syllables
    $\mapsto$ the alternation of lemma statements and their
    applications in a proof.
  \item \textbf{Rhyme $\mapsto$ structural echo.}\quad
    Rhyming end-words ($\GId_{g_r}$ at the $\phon$ stratum)
    $\mapsto$ reuse of the same lemma at different points in a
    proof (structural parallelism).
  \item \textbf{Volta $\mapsto$ key insight.}\quad
    The turn at line~9 of a sonnet $\mapsto$ the ``aha moment''
    where a proof introduces the decisive idea.
  \item \textbf{Elegance $\mapsto$ elegance.}\quad
    Both domains have an ``elegance'' component in $\GEval$ that
    resists formal definition but is recognised by experts.
\end{itemize}
\end{construction}

\begin{remark}[Why this works]
The correspondence is not metaphorical---it is structural.
Both poems and proofs are terms inhabiting a type in $\PStr$,
both must satisfy global gluing conditions ($H^1=0$ in
$\CSpec$), and both are evaluated by multi-component grades
in $\GEval$.  The functor $\Phi_{\mathrm{poet}\to\mathrm{math}}$
identifies the \emph{formal skeleton} shared by both, while
acknowledging that the \emph{content} differs (linguistic vs.\
logical).
\end{remark}

%% ─────────────────────────────────────────────────────────────
\section{The Cycle and Its Fixed Points}
\label{sec:cycle}

\begin{proposition}[The cycle is idempotent]
\label{prop:cycle-idempotent}
Let $\Psi = \Phi_{\mathrm{math}\to\mathrm{prog}}
            \circ\Phi_{\mathrm{poet}\to\mathrm{math}}
            \circ\Phi_{\mathrm{NL}\to\mathrm{poet}}
            \circ\Phi_{\mathrm{prog}\to\mathrm{NL}}$.
Then $\Psi$ is not the identity on $\mathcal{D}_{\mathrm{prog}}$
(the round trip through four domains loses domain-specific
detail), but $\Psi^2=\Psi$: the image of $\Psi$ is a
\emph{retract}, and applying $\Psi$ again does not lose further
information.
\end{proposition}

\begin{proof}[Proof sketch]
Each functor in the cycle is faithful, hence injective on
morphisms.  The composition $\Psi$ projects onto the
\emph{shared structural skeleton}---the sub-4-tuple consisting
of features present in all four domains.  This sub-4-tuple is
already in the image of $\Psi$, so $\Psi^2$ acts as the
identity on it.
\end{proof}

\begin{definition}[Universal artifact]
An artifact $a$ is \emph{universal} if $\Psi(a)\cong a$---it
is a fixed point of the cycle.  Such an artifact has coherent
realisations in all four domains simultaneously.
\end{definition}

%% ─────────────────────────────────────────────────────────────
\section{Worked Example: Recursion Across Four Domains}
\label{sec:recursion}

We take the concept of \emph{recursion} and exhibit its
4-tuple in each domain, then verify the functorial
correspondences.

\subsection{As a Program: Factorial}

\textbf{$\STop$.}\quad
The CFG contains a self-loop at the recursive call site.

\textbf{$\PStr$.}\quad
Type: $\mathrm{Factorial}:\mathbb{N}\to\mathbb{N}$,
with specification $\mathrm{fact}(n)=n!$.  The term
is $\lambda n.\;\mathrm{if}\;n=0\;\mathrm{then}\;1\;
\mathrm{else}\;n\cdot\mathrm{fact}(n-1)$.
Termination proof: the argument $n$ strictly decreases.

\textbf{$\CSpec$.}\quad
$H^1=0$: the base case and inductive case glue at $n=0$.

\textbf{$\GEval$.}\quad
$g_{\mathrm{rigor}}=\gone$, $g_{\mathrm{perf}}=-0.5$
(stack overflow for large $n$).

\subsection{As a Sentence: Centre-Embedding}

\begin{quote}
\emph{``The cat that the dog that the rat bit chased ran.''}
\end{quote}

\textbf{$\STop$.}\quad
Three nested relative clauses form a recursive phrase structure.

\textbf{$\PStr$.}\quad
The typing derivation nests three NP--VP pairs, each containing
the next.  This is linguistic recursion: the grammar rule
$\text{NP}\to\text{NP}\;\text{RelClause}$ applies to its own
output.

\textbf{$\CSpec$.}\quad
$H^1\neq 0$ for most human parsers: the triple centre-embedding
exceeds working memory, causing a processing breakdown.  The
obstruction is at the innermost clause boundary.

\textbf{$\GEval$.}\quad
$g_{\synt}=\gone$ (grammatical), $g_{\prag}=-2.0$
(incomprehensible to most listeners).

\textbf{Functor check.}\quad
$\Phi_{\mathrm{prog}\to\mathrm{NL}}$ maps the self-loop in
the CFG to the nested clause structure.  The recursive call
corresponds to the recursive application of the NP rule.

\subsection{As a Poem: A Villanelle}

The villanelle repeats two refrain lines (A1 and A2)
throughout, returning to them in a fixed pattern.  The
repetition \emph{is} recursion: the poem calls back to
its own earlier material.

\textbf{$\PStr$.}\quad
Type: $\mathrm{Villanelle}(\theta,\text{ABA}\ldots,\mu)$.
The refrain lines are shared sub-terms, reused at specified
positions---exactly like a recursive call reusing the
function body.

\textbf{$\CSpec$.}\quad
$H^1=0$ when the refrains integrate seamlessly into each
stanza.  $H^1\neq 0$ if a refrain clashes with its local
context (e.g.\ the repeated line contradicts new material).

\textbf{Functor check.}\quad
$\Phi_{\mathrm{NL}\to\mathrm{poet}}$ maps the recursive NP
structure to the refrain structure: both are self-referential
uses of material.

\subsection{As a Proof: Induction}

\textbf{$\PStr$.}\quad
Type: $\Pi(n:\mathbb{N}).\,P(n)$.  The proof term is
$\mathrm{ind}(\mathrm{base},\lambda n.\lambda h.\,
\mathrm{step}(n,h))$ where the step function receives
the inductive hypothesis~$h:P(n)$ and produces $P(n+1)$.
This \emph{is} recursion: the proof of $P(n+1)$ calls
the proof of $P(n)$.

\textbf{$\CSpec$.}\quad
$H^1=0$ when the base case and step glue at $n=0$.

\textbf{Functor check.}\quad
$\Phi_{\mathrm{math}\to\mathrm{prog}}$ (Curry--Howard)
maps the induction term to a recursive program.
$\Phi_{\mathrm{poet}\to\mathrm{math}}$ maps the
villanelle's refrain pattern to the reuse of the inductive
hypothesis.

\medskip

\begin{theorem}[Recursion as a universal artifact]
\label{thm:recursion-universal}
Recursion is a fixed point of the cycle $\Psi$: its 4-tuple
is invariant (up to isomorphism) under the round-trip
$\mathrm{prog}\to\mathrm{NL}\to\mathrm{poet}\to\mathrm{math}
\to\mathrm{prog}$.  The structural skeleton---a self-referential
term with a base case and an inductive/recursive step---is
preserved by every functor in the cycle.
\end{theorem}

\begin{remark}[Other universal artifacts]
Recursion is not the only fixed point.  Other candidates include
\emph{composition} (function composition / sentence conjunction /
stanza sequencing / lemma chaining),
\emph{abstraction} ($\lambda$-abstraction / pronoun binding /
metonymy / universal quantification), and
\emph{symmetry} (invariance under permutation / palindrome /
chiasmus / group-theoretic symmetry).  Each deserves a full
treatment parallel to the recursion example above.
\end{remark}


% ── Part VI: Applications and Horizons ────────────────────────
\part{Applications and Future Directions}
% Part VI: Applications and Future Directions
% This file is \include'd from main.tex

\chapter{LLM-Assisted Software Synthesis}
\label{ch:llm-synthesis}

Large Language Models (LLMs) generate plausible code with remarkable
fluency, yet their outputs carry no structural guarantees: type errors
slip through, modules fail to compose, and subtle semantic bugs hide
behind syntactically correct surfaces.  This chapter shows how
Judgment Geometry supplies the missing \emph{verification and repair
layer} for LLM-generated code.

%% ─── 27.1 ──────────────────────────────────────────────────────
\section{The structural deficit of LLM code}
\label{sec:llm-deficit}

An LLM treats code generation as next-token prediction: given a prompt
$p$, it samples a continuation $c$ from a learned distribution
$P(c \mid p)$.  Nothing in this process enforces that:
\begin{enumerate}[label=(\roman*)]
  \item $c$ type-checks against the project's type system,
  \item $c$ is consistent with specifications in \emph{other} modules,
  \item $c$ satisfies non-local invariants (security policies,
        performance contracts, API versioning).
\end{enumerate}
Testing catches some of these failures, but testing is inherently
\emph{local}: a test suite checks individual behaviors ($H^0$, in our
language) without examining whether local sections \emph{glue} into a
global one ($H^1$).  Formal verification, at the other extreme, is
too rigid---it demands binary pass/fail and cannot exploit partial
correctness.

Judgment Geometry occupies the middle ground: it is \emph{graded}
(partial correctness is meaningful), \emph{cohomological} (it detects
non-local inconsistencies), and \emph{iterative} (the descent algorithm
produces a convergent repair sequence).

%% ─── 27.2 ──────────────────────────────────────────────────────
\section{JuGeo as LLM supervisor}
\label{sec:llm-supervisor}

We instantiate the four objects for LLM-assisted development.

\begin{definition}[LLM Development Site]
\label{def:llm-site}
Let $\mathcal{P}$ be a software project with modules
$M_1,\dots,M_n$.  Define the \emph{LLM development site}
$\STop_{\mathrm{LLM}}$ as follows:
\begin{itemize}
  \item \textbf{Objects}: modules, files, functions, and test suites
        of~$\mathcal{P}$, ordered by containment.
  \item \textbf{Morphisms}: inclusion and dependency maps
        (imports, function calls, type references).
  \item \textbf{Topology}: the coverage $J$ in which a family
        $\{U_i \to U\}$ covers~$U$ when every public symbol of~$U$
        is exercised by at least one~$U_i$.
\end{itemize}
\end{definition}

\begin{definition}[LLM Proof Structure]
\label{def:llm-pstr}
The presheaf $\PStr_{\mathrm{LLM}}$ assigns to each object $U$:
\begin{itemize}
  \item A \emph{specification} $A_U$ (type signature, docstring
        contract, or formal pre/post-condition).
  \item An \emph{evidence bundle} $E_U$: the LLM-generated code
        together with any passing tests.
\end{itemize}
A section $s \in \PStr_{\mathrm{LLM}}(U)$ is a pair $(A_U, E_U)$
asserting ``evidence $E_U$ satisfies specification~$A_U$\!.''
\end{definition}

The \emph{Cohomological Spectrum} $\CSpec_{\mathrm{LLM}}$ is the
\v{C}ech complex of $\PStr_{\mathrm{LLM}}$.  Concretely:
\[
  \Cech^0 = \prod_i \PStr(U_i), \qquad
  \Cech^1 = \prod_{i<j} \PStr(U_i \cap U_j),
\]
where $U_i \cap U_j$ is the shared interface between modules $M_i$
and~$M_j$ (exported types, shared state, API contracts).  A class
$[\alpha] \in \check{H}^1(\STop_{\mathrm{LLM}}, \PStr_{\mathrm{LLM}})$
with $[\alpha]\neq 0$ witnesses an \emph{integration bug}: the
LLM-generated code for $M_i$ and $M_j$ individually type-checks but
disagrees on the overlap.

\begin{example}
Module $M_1$ exports a function \texttt{get\_user} returning a
dictionary with key \texttt{"name"}.  The LLM generates $M_2$
expecting a key \texttt{"username"}.  Both modules pass their local
tests (the LLM generates a mock with the expected key), but the
overlap section records the mismatch.  This is a non-trivial class in
$\check{H}^1$.
\end{example}

The \emph{Graded Evaluation} $\GEval_{\mathrm{LLM}}$ grades each
section along multiple axes:
\begin{equation}
\label{eq:llm-grade}
  g(s) \;=\; g_{\mathrm{type}} \;\otG\;
             g_{\mathrm{test}} \;\otG\;
             g_{\mathrm{style}} \;\otG\;
             g_{\mathrm{perf}},
\end{equation}
where each component lies in $\Grade = (-\infty, 0]$ and
$\otG = +$ aggregates them (recall $\gone = 0$ is perfect).

%% ─── 27.3 ──────────────────────────────────────────────────────
\section{The descent loop for LLM code}
\label{sec:llm-descent}

The central observation is that \emph{iterative LLM repair is exactly
descent}.

\begin{construction}[LLM Descent Algorithm]
\label{con:llm-descent}
Given an LLM $\mathcal{L}$, a project site $\STop$, and a
specification presheaf $\PStr$:
\begin{enumerate}
  \item \textbf{Generate.}\; For each $U_i$, prompt $\mathcal{L}$
        to produce a candidate section
        $s_i^{(0)} \in \PStr(U_i)$.
  \item \textbf{Compute coboundary.}\; Form
        $\delta^0 s^{(0)} \in \Cech^1$; identify all non-zero
        components $\alpha_{ij}$.
  \item \textbf{Repair.}\; For each obstruction $\alpha_{ij}$,
        construct a \emph{repair prompt}
        $p_{ij} = \texttt{``Fix: module } M_i \texttt{ expects \dots\ but }
        M_j \texttt{ provides \dots''}$.
        Feed $p_{ij}$ to~$\mathcal{L}$; obtain $s_i^{(1)}$
        or~$s_j^{(1)}$.
  \item \textbf{Re-evaluate.}\; Compute $\delta^0 s^{(1)}$.
        If $\check{H}^1 = 0$ and
        $\GEval(s^{(1)}) \geq g_{\min}$, accept.
        Otherwise return to step~3.
  \item \textbf{Budget.}\; If the iteration count exceeds a bound
        $N$ or the grade fails to improve for $k$ consecutive
        iterations, halt and report the residual obstructions.
\end{enumerate}
\end{construction}

\begin{proposition}
\label{prop:llm-descent-mono}
If each repair step strictly improves the grade of at least one
obstruction component (i.e., $g(\alpha_{ij}^{(t+1)}) >
g(\alpha_{ij}^{(t)})$ for some~$(i,j)$), and the grade is bounded
above by~$\gone$, then the descent loop terminates in at most
$n(n-1)/2 \cdot \lceil (\gone - g_{\min})/\epsilon \rceil$ steps,
where $\epsilon$ is the minimum per-step improvement.
\end{proposition}

\begin{proof}
There are at most $\binom{n}{2}$ overlap pairs.  Each
can improve at most $(\gone - g_{\min})/\epsilon$ times before
reaching~$\gone$.  Once every overlap grade reaches~$\gone$, the
coboundary vanishes and $\check{H}^1 = 0$.
\end{proof}

\begin{remark}
In practice one does not require $\check{H}^1 = 0$ exactly; a graded
threshold $g_{\min}$ suffices.  This is a key advantage of graded
cohomology over binary verification: the system can declare
``good enough'' and defer minor mismatches.
\end{remark}

%% ─── 27.4 ──────────────────────────────────────────────────────
\section{Prompt engineering as site construction}
\label{sec:prompt-site}

We now observe that the art of prompt engineering admits a precise
geometric interpretation.

\begin{proposition}[Prompts as Covers]
\label{prop:prompt-cover}
Let $\mathcal{T}$ be the target specification (the ``task'' given
to the LLM).  A \emph{prompt decomposition}
$p = \{p_1, \dots, p_k\}$ induces a covering family
$\{U_1, \dots, U_k\}$ of~$\mathcal{T}$ in $\STop$, where each
$U_i$ is the sub-task addressed by prompt~$p_i$.  The LLM's response
to $p_i$ is a local section $s_i \in \PStr(U_i)$.  The full
response glues iff $\check{H}^1(\{U_i\}, \PStr) = 0$.
\end{proposition}

Under this interpretation:
\begin{itemize}
  \item \textbf{Few-shot examples} are \emph{covering sieves}:
        they demonstrate the desired behavior on specific
        sub-regions of the task, guiding the LLM toward sections
        that agree on overlaps.
  \item \textbf{The context window} is the \emph{open cover}
        $\mathcal{U}$: the finite amount of information the LLM can
        attend to at once.
  \item \textbf{Prompt decomposition} is \emph{refinement of the
        cover}: splitting a complex prompt into sub-prompts
        corresponds to refining $\mathcal{U}$ into a finer cover,
        which can only \emph{improve} the cohomological condition
        (by Leray's theorem, refinement does not increase~$H^1$).
  \item \textbf{Chain-of-thought prompting} is \emph{choosing a good
        nerve}: it makes the simplicial structure of the cover
        explicit, so the LLM can track overlaps.
\end{itemize}

%% ─── 27.5 ──────────────────────────────────────────────────────
\section{Multi-agent orchestration via the fibration}
\label{sec:multi-agent}

Modern LLM systems deploy multiple agents, each responsible for one
module or aspect of the codebase.  JuGeo provides the orchestration
framework.

\begin{construction}[Multi-Agent Fibration]
\label{con:multi-agent}
Let $\{A_1, \dots, A_n\}$ be LLM agents, with agent $A_i$
responsible for module~$M_i$.  The \emph{judgment fibration}
$\pi\colon \mathcal{E} \to \STop$ has fiber $\pi^{-1}(U_i)$
consisting of all code artifacts agent~$A_i$ can produce for~$U_i$.
The orchestrator's job is to select a \emph{global section}
$s \in \Gamma(\STop, \mathcal{E})$---i.e., a choice of artifact per
module that glues across all overlaps.  This is precisely the
descent problem.
\end{construction}

The orchestrator proceeds as follows:
\begin{enumerate}
  \item Each agent $A_i$ proposes a local section $s_i$.
  \item The orchestrator computes $\delta^0 s$ and identifies
        obstructions $\alpha_{ij}$.
  \item For each obstruction, the orchestrator sends a
        \emph{compatibility constraint} to agents $A_i$ and~$A_j$.
  \item The agents regenerate; the process iterates.
\end{enumerate}

This is the exact pattern used by systems such as ChatDev and
MetaGPT, but JuGeo gives it a \emph{mathematical foundation}:
convergence, termination, and quality guarantees follow from the
general descent theory of Part~II.

%% ─── 27.6 ──────────────────────────────────────────────────────
\section{Application: full-stack web generation}
\label{sec:webapp-gen}

We illustrate the framework on a concrete application:
generating a complete web application (HTML, CSS, JavaScript,
backend~API).

\begin{example}[Web Application as Stratified Site]
\label{ex:webapp-site}
A web application defines a natural three-stratum site, analogous to
the linguistic strata of Part~III:
\begin{center}
\begin{tabular}{lll}
\textbf{Stratum} & \textbf{Web layer} & \textbf{Linguistic analogue} \\
\hline
$\sigma_1$ & HTML (structure) & $\synt$ (syntax) \\
$\sigma_2$ & CSS (presentation) & $\phon$ (phonology / surface form) \\
$\sigma_3$ & JS (behavior) & $\sem$ (semantics) \\
$\sigma_4$ & API (data) & $\prag$ (pragmatics / context)
\end{tabular}
\end{center}
Cross-stratum descent asks: does the JavaScript match the HTML
structure?  Do the CSS selectors reference elements that actually
exist?  Does the API return data in the format the frontend expects?
\end{example}

The generation loop for a web app is:
\begin{enumerate}
  \item \textbf{Specification.}\;  The user provides a natural-language
        description $\mathcal{T}$.
  \item \textbf{Site construction.}\;  Parse $\mathcal{T}$ into a
        module decomposition $\{U_i\}$ (pages, components, API
        endpoints).
  \item \textbf{LLM generation.}\;  For each $(U_i, \sigma_j)$
        (module $\times$ stratum), the LLM generates a local section.
  \item \textbf{Cross-stratum descent.}\;  Compute
        $\check{H}^1$ across strata: verify that CSS classes match
        HTML elements, that JS event handlers reference existing DOM
        nodes, that API calls match endpoint signatures.
  \item \textbf{Repair.}\;  Feed obstructions back to the LLM as
        repair prompts.
  \item \textbf{Compilation and testing.}\;  Run the project's build
        system and test suite; feed results into $\GEval$.
  \item \textbf{Iterate.}\;  Repeat until $\GEval \geq g_{\min}$.
\end{enumerate}

\begin{remark}
The \texttt{jugeo-webapp} module in the JuGeo codebase implements
a prototype of this loop, using the descent engine from
\texttt{src/jugeo/} to orchestrate HTML/CSS/JS generation and
cross-layer verification.
\end{remark}

%% ─── 27.7 ──────────────────────────────────────────────────────
\section{Comparison with existing approaches}
\label{sec:llm-comparison}

We compare JuGeo-supervised LLM code generation with three
alternatives.

\medskip\noindent
\textbf{Plain LLM generation.}\;
Without JuGeo, an LLM has no mechanism to detect non-local
inconsistencies.  It can only ``hope'' that independently generated
modules cohere.  JuGeo adds the cohomological check ($\CSpec$)
and the descent repair loop, converting hope into a convergent
algorithm.

\medskip\noindent
\textbf{Formal verification.}\;
Systems like Dafny and Lean provide strong guarantees but require
full formal specifications and produce binary verdicts.
JuGeo is \emph{graded}: a section with
$g(s) = -0.3$ is meaningful (``almost correct''),
enabling iterative improvement.  Moreover, JuGeo can incorporate
formal verification as one component of $\GEval$ (set
$g_{\mathrm{formal}} = \gone$ if the prover succeeds,
$g_{\mathrm{formal}} = \gzero$ otherwise), subsuming
formal methods as a special case.

\medskip\noindent
\textbf{Testing alone.}\;
A test suite checks $H^0$: does each module individually behave
correctly?  JuGeo also checks $H^1$: do modules compose correctly?
This is precisely the gap that integration bugs exploit.
By computing the full \v{C}ech complex, JuGeo detects classes
of bugs that no finite test suite can reliably catch---those
arising from the \emph{interaction} of correct components.

\begin{proposition}[Subsumption]
\label{prop:subsumption}
Let $\mathcal{V}$ be any verification methodology that assigns a
binary verdict $v_U \in \{0,1\}$ to each module~$U$.  Then
$\mathcal{V}$ embeds into $\GEval$ via $g_U = v_U \cdot \gone$,
and the $\check{H}^1$ computation of $\CSpec$ subsumes any pairwise
compatibility check that $\mathcal{V}$ performs.
\end{proposition}

\chapter{GOFAI Natural and Creative Language}
\label{ch:gofai-nl}

A central claim of Judgment Geometry is that structured mathematical
theory can \emph{replace} learned statistical models for certain
language tasks.  This chapter develops the claim in full: we show how
the Compositional Theory of Theories (CTT), instantiated with $\sim\!450$
theory slots across ten linguistic strata, drives a chatbot and a
poetry engine that use \emph{no neural model whatsoever}---zero
parameters, zero training data, zero GPU.

%% ─── 28.1 ──────────────────────────────────────────────────────
\section{The GOFAI thesis}
\label{sec:gofai-thesis}

\begin{quote}
\emph{If you have enough theory, you don't need an LLM.}
\end{quote}

The CTT of Part~III provides:
\begin{enumerate}[label=(\arabic*)]
  \item \textbf{450 theory slots} spanning phonology through aesthetics,
        each contributing constraints and generation rules.
  \item \textbf{The $\Grade$ semiring} for combining theory verdicts:
        $\otG$ conjoins independent requirements; $\opG$ selects
        among competing alternatives.
  \item \textbf{The $\FMLM$ decoder}, which fills masked positions
        by iterative fixpoint computation over the full theory
        ensemble.
  \item \textbf{The descent algorithm}, which repairs
        disharmonious output until the global grade crosses a
        threshold.
\end{enumerate}

The result is a system whose every output is \emph{explainable}
(each quality grade traces to a named theory), \emph{controllable}
(change one theory slot, change one aspect of output), and
\emph{compositional} (theories compose via the semiring without
interference).

%% ─── 28.2 ──────────────────────────────────────────────────────
\section{The GOFAI chatbot architecture}
\label{sec:gofai-chatbot}

The chatbot implements a classic pipeline, but every stage is
graded by a JuGeo theory slot.

\subsection{Natural language understanding}
\label{ssec:nlu}

Input processing follows a five-stage NLU pipeline:
\begin{enumerate}
  \item \textbf{Tokenization and POS tagging}\; ($\morph$ stratum).
        Rule-based tokenizer with lexicon lookup; POS tags assigned by
        suffix rules and disambiguation via the highest-grade parse.
  \item \textbf{Syntactic parsing}\; ($\synt$ stratum).
        A slot-based parser identifies clause boundaries, verb frames,
        and modifier attachment.  Each parse receives a grade
        $g_{\synt}$ from theories such as \textsf{ClauseCompleteness}
        and \textsf{InversionLicensing}.
  \item \textbf{Semantic analysis}\; ($\sem$ stratum).
        Frame-semantic analysis using FrameNet frames: the input is
        mapped to a frame with filled roles.  Theories
        \textsf{FrameCoherence} and \textsf{EntityBinding} grade the
        result.
  \item \textbf{Pragmatic interpretation}\; ($\prag$ stratum).
        Speech-act classification (request, assertion, question) and
        Gricean maxim checking.  The theory \textsf{SpeechActFelicity}
        grades the pragmatic parse.
  \item \textbf{Discourse integration}\; ($\disc$ stratum).
        Topic tracking and anaphora resolution update the discourse
        state.  Theories \textsf{TopicCoherence} and
        \textsf{AnaphoraResolution} grade the integration.
\end{enumerate}

The overall NLU grade is the $\otG$-product across strata:
\begin{equation}
\label{eq:nlu-grade}
  g_{\mathrm{NLU}} = g_{\morph} \otG g_{\synt} \otG g_{\sem}
    \otG g_{\prag} \otG g_{\disc}.
\end{equation}

\subsection{Intent classification}
\label{ssec:intent}

The dialogue manager classifies user intent via a pattern-based
system graded by the $\Grade$ semiring.  Each candidate intent
$I \in \{\textsc{PoemRequest}, \textsc{StyleRequest},
\textsc{Feedback}, \textsc{Revision}, \textsc{Question},
\textsc{Greeting}, \textsc{Help}, \textsc{Thanks}\}$ has a set of
indicator patterns; firing patterns contribute grade via~$\opG$
(log-sum-exp), and the highest-grade intent wins:
\begin{equation}
\label{eq:intent-grade}
  I^* = \arg\max_I \; g(I \mid \text{input}).
\end{equation}
No neural classifier is involved: the grade semiring \emph{is} the
classifier.

\subsection{Response generation}
\label{ssec:response-gen}

Given the classified intent $I^*$ and discourse state, the response
pipeline proceeds:
\begin{enumerate}
  \item \textbf{Topic model $\to$ frame selection.}\;
        The current topic (maintained by the $\disc$ stratum) selects
        a FrameNet frame; frame elements become template slots.
  \item \textbf{Template instantiation.}\;
        A response template is chosen from a bank indexed by intent
        and frame; slots are filled from the discourse context and
        knowledge base.
  \item \textbf{Surface realization.}\;
        The filled template is linearized, applying morphological
        rules ($\morph$ stratum) and phonological smoothing ($\phon$
        stratum).
  \item \textbf{$\FMLM$ refinement.}\;
        Any remaining masked positions are filled by the fixpoint
        decoder, which scores candidates against the full theory
        ensemble.
\end{enumerate}

\subsection{Conversation management}
\label{ssec:conversation}

Long-range coherence is maintained by the \emph{discourse monad}
$\mathrm{DMon}$, which tracks:
\begin{itemize}
  \item \textbf{Topic stack}\; ($\disc$ stratum): active topics,
        topic shifts, and topic resumption.
  \item \textbf{Reference resolution}\; ($\sem$ stratum): entity
        mentions and their antecedents.
  \item \textbf{Turn history}\; ($\prag$ stratum): user preferences,
        feedback history, and affective state.
  \item \textbf{User model}: accumulated preferences for poetic form,
        mood, and style, updated after each interaction.
\end{itemize}

The monadic structure ensures that state updates compose cleanly
across turns: each turn is a Kleisli arrow
$\mathrm{Turn} \to \mathrm{DMon}(\mathrm{Response})$.

%% ─── 28.3 ──────────────────────────────────────────────────────
\section{The GOFAI poetry engine}
\label{sec:gofai-poetry}

The poetry engine is the showcase application of the GOFAI approach.
It generates metrically, rhymically, and semantically controlled
verse using no neural model.

\subsection{Form selection}
\label{ssec:form-select}

Each poetic form is a type $F$ in $\PStr$ with specific structural
constraints:
\begin{center}
\begin{tabular}{lll}
\textbf{Form} & \textbf{Constraint type} & \textbf{Theory} \\
\hline
Sonnet & 14 lines, volta at 8 or 12, iambic pentameter
       & \textsf{FormCompliance} \\
Villanelle & 19 lines, two refrains, $\mathrm{ABA}$ tercets
           & \textsf{FormCompliance} \\
Haiku & 5-7-5 syllables, seasonal word (\emph{kigo})
      & \textsf{SyllableCount} \\
Ghazal & couplets with \emph{radif} and \emph{qafia}
       & \textsf{RhymeScheme} \\
Terza rima & $\mathrm{ABA\;BCB\;CDC\;\dots}$ interlocking tercets
           & \textsf{RhymeScheme} \\
Free verse & no fixed meter; scored on rhythm variety
           & \textsf{MeterCompliance}
\end{tabular}
\end{center}
The user may also specify a \emph{poet style} (Keats, Dickinson,
Whitman, Shakespeare, Neruda), which modulates diction, syntax, and
thematic profiles via a style presheaf over the stratal tower.

\subsection{Content generation}
\label{ssec:content-gen}

Content is generated by pure knowledge-base lookup:
\begin{enumerate}
  \item \textbf{Topic $\to$ FrameNet frames.}\;
        The topic word activates one or more FrameNet frames
        (e.g., \emph{love} activates \textsc{Personal\_relationship},
        \textsc{Emotion\_directed}, \textsc{Desiring}).
  \item \textbf{Frame $\to$ lexical units.}\;
        Each frame provides lexical units; these are expanded via
        WordNet synonymy, hypernymy, and semantic similarity.
  \item \textbf{Candidate word pool.}\;
        The union of lexical units and WordNet expansions forms
        the candidate pool, from which the $\FMLM$ decoder draws.
  \item \textbf{Proof strategies.}\;
        Eight proof strategies---\textsf{BaseFact}, \textsf{Causal},
        \textsf{Metaphor}, \textsf{Temporal}, \textsf{Negation},
        \textsf{Relational}, \textsf{Hypernym},
        \textsf{Presupposition}---each produce an $\HPLF$ term and
        a template with bindings.  These provide the propositional
        skeleton of each line.
\end{enumerate}

No neural model is consulted at any point: content derives from
FrameNet, WordNet, CMUDict, and the theory ensemble.

\subsection{Quality control}
\label{ssec:quality}

Each candidate poem is scored by the $\Grade$ semiring along four
primary axes, each backed by specific theory slots:

\medskip\noindent\textbf{Meter.}\;
CMUDict provides stress patterns (sequences of 0, 1, 2 for
unstressed, primary, secondary stress).  The theory
\textsf{MeterCompliance} computes the edit distance between the
observed stress pattern and the target pattern (e.g., iambic
pentameter = $01\,01\,01\,01\,01$), normalized to $[-\infty, 0]$:
\[
  g_{\mathrm{meter}} = -\frac{d_{\mathrm{edit}}(\sigma_{\mathrm{obs}},
    \sigma_{\mathrm{target}})}{\lvert \sigma_{\mathrm{target}} \rvert}.
\]

\medskip\noindent\textbf{Rhyme.}\;
CMUDict phoneme sequences enable perfect rhyme (identical final
stressed vowel and coda), near rhyme (identical vowel, different
coda), and slant rhyme (similar vowel).  The theory
\textsf{RhymeScheme} checks each line-ending pair against the
target scheme:
\[
  g_{\mathrm{rhyme}} = \frac{1}{|\text{pairs}|}\sum_{\text{pairs}}
    g_{\mathrm{match}}(w_i, w_j).
\]

\medskip\noindent\textbf{Imagery.}\;
WordNet Wu--Palmer similarity between successive content words
measures semantic diversity; excessively high similarity (clich\'e)
or excessively low similarity (incoherence) is penalized:
\[
  g_{\mathrm{imagery}} = -\bigl\lvert
    \mathrm{sim}(w_i, w_j) - \tau_{\mathrm{target}} \bigr\rvert,
\]
where $\tau_{\mathrm{target}} \approx 0.4$ balances novelty and
coherence.

\medskip\noindent\textbf{Emotion.}\;
Affect lexicons (valence, arousal, dominance) map each word to an
emotional vector.  The theory \textsf{RasaConsistency} checks that
the poem's emotional trajectory matches the target
mood (\emph{rasa}):
\[
  g_{\mathrm{emotion}} = -\left\lVert
    \vec{e}_{\mathrm{poem}} - \vec{e}_{\mathrm{target}}
  \right\rVert.
\]

The total poem grade is:
\begin{equation}
\label{eq:poem-grade}
  g_{\mathrm{poem}} = g_{\mathrm{meter}} \otG g_{\mathrm{rhyme}}
    \otG g_{\mathrm{imagery}} \otG g_{\mathrm{emotion}}
    \otG g_{\mathrm{style}} \otG g_{\mathrm{form}}.
\end{equation}

\subsection{Beam search as proof search}
\label{ssec:beam-search}

The generation engine maintains a beam of $k$ candidate poems,
scored by their total grade~\eqref{eq:poem-grade}.  At each
step (line or stanza), the $\FMLM$ decoder produces $m$ candidate
continuations for each beam element; the top-$k$ by grade survive.
This is formally identical to \emph{proof search} in the type
theory of Part~III: each candidate poem is a partial proof term,
each continuation extends the proof, and the grade is the
proof's quality.

\begin{proposition}
\label{prop:beam-convergence}
If the per-step grade is bounded below by $g_{\min} > \gzero$ and
the form imposes a finite length~$L$, then beam search terminates
in at most $L$ steps with at least one candidate of grade
$\geq g_{\min}^L$ (where exponentiation is in the $\otG$-monoid,
i.e., $L$-fold $\otG$-product).
\end{proposition}

%% ─── 28.4 ──────────────────────────────────────────────────────
\section{Advantages of GOFAI + JuGeo}
\label{sec:gofai-advantages}

We summarize the structural advantages of the theory-driven approach
over LLM-based generation.

\begin{description}
  \item[Interpretability.]
    Every component of $g_{\mathrm{poem}}$ traces to a named theory
    slot.  A user can ask ``why is the meter grade low?'' and receive
    a specific answer (``line~7 has a trochaic inversion at
    position~3'').  An LLM cannot provide this.
  \item[Controllability.]
    Changing one theory slot changes exactly one aspect of the
    output.  To make the poem more emotionally intense, raise the
    weight of \textsf{RasaConsistency}; to allow more metrical
    freedom, lower the weight of \textsf{MeterCompliance}.
    In an LLM, such fine-grained control is unavailable.
  \item[Compositionality.]
    Theories compose via the $\Grade$ semiring: adding a new
    theory (e.g., \textsf{Alliteration}) does not interfere with
    existing ones.  This is the \emph{open world property} of CTT.
  \item[No training data.]
    The system runs on pure theory: lexicons (WordNet, CMUDict,
    FrameNet), affect dictionaries, and hand-authored constraints.
    No corpus of poems is needed.
  \item[No hallucination.]
    Every factual claim in the poem traces to a knowledge-base
    entry (WordNet gloss, FrameNet frame, or Wikipedia fact).
    The system cannot invent false facts because it has no
    generative model to hallucinate from.
\end{description}

%% ─── 28.5 ──────────────────────────────────────────────────────
\section{Disadvantages and mitigations}
\label{sec:gofai-disadvantages}

\begin{description}
  \item[Fluency.]
    GOFAI-generated text can sound stilted.  Mitigation: large word
    banks ($> 150{,}000$ words from CMUDict), frame-based templates
    from FrameNet, and the $\FMLM$ decoder's multi-axis scoring
    (phonological, semantic, register, diversity) smooth the surface.
  \item[Novelty.]
    Knowledge-base lookup risks producing predictable combinations.
    Mitigation: the \emph{conceptual blending} type constructor
    $\mathrm{Blend}(F_1, F_2)$ of Part~III merges two FrameNet frames
    to produce novel juxtapositions (e.g., blending
    \textsc{Commerce\_sell} with \textsc{Death} for a war poem).
  \item[Scale.]
    Full beam search over 450 theories is expensive.
    Mitigation: theories are lazily evaluated (only active
    theories fire), and the $\FMLM$ decoder caches lexical lookups
    across lines.  In practice, generation of a 14-line sonnet
    takes $<$\,2 seconds on commodity hardware.
\end{description}

%% ─── 28.6 ──────────────────────────────────────────────────────
\section{The hybrid approach}
\label{sec:hybrid}

The GOFAI and LLM approaches are not mutually exclusive; JuGeo
provides the framework for combining them.

\begin{construction}[Hybrid Generation Pipeline]
\label{con:hybrid}
\begin{enumerate}
  \item \textbf{GOFAI generates the structural skeleton.}\;
        The theory ensemble produces a metrically and rhymically
        correct draft with filled semantic roles, graded by
        $\GEval$.  This draft is a \emph{section} in $\PStr$ with
        $\check{H}^1 = 0$ (it glues) but potentially low surface
        fluency.
  \item \textbf{LLM polishes the surface.}\;
        The draft is passed to an LLM with the instruction ``improve
        fluency while preserving meter, rhyme, and meaning.''
        The LLM's output is a new candidate section.
  \item \textbf{JuGeo verifies the result.}\;
        The full theory ensemble re-grades the polished version.
        If $\GEval$ improves and $\check{H}^1$ remains zero, the
        polished version is accepted; otherwise the GOFAI draft is
        preferred.
\end{enumerate}
\end{construction}

This hybrid pipeline gives the best of both worlds: the structural
guarantees of GOFAI and the surface fluency of an LLM, with JuGeo
serving as the arbiter.  The graded evaluation ensures that the
LLM cannot silently break structural invariants---any fluency
improvement that damages meter, rhyme, or semantic coherence will
be caught by the theory ensemble and rejected.

\chapter{Software Generation at Scale}
\label{ch:software-scale}

We now turn from individual-project verification to
\emph{industrial-scale} applications of Judgment Geometry:
continuous integration, multi-repository coherence, specification
mining, refactoring, security analysis, and performance optimization.

%% ─── 29.1 ──────────────────────────────────────────────────────
\section{Continuous verification}
\label{sec:continuous-verif}

Modern software development is organized around
continuous-integration (CI) pipelines.  JuGeo integrates naturally
into this workflow.

\begin{construction}[JuGeo CI Layer]
\label{con:ci-layer}
Augment a standard CI/CD pipeline with a JuGeo verification stage:
\begin{enumerate}
  \item \textbf{Every commit $c$ induces a judgment.}\;
        The diff of $c$ modifies sections on one or more open sets
        $U_i$ in the project site $\STop$.  JuGeo recomputes the
        affected entries of the \v{C}ech complex.
  \item \textbf{Every pull request triggers descent.}\;
        The PR's changeset defines a candidate section~$s'$;
        JuGeo computes $\delta^0 s'$ and reports any new
        obstructions introduced by the change.
  \item \textbf{$\CSpec$ catches integration bugs.}\;
        Standard CI runs unit tests ($H^0$ check).  JuGeo adds the
        $H^1$ check: do the changed modules still glue with their
        neighbors?
  \item \textbf{$\GEval$ tracks quality over time.}\;
        The grade $g(s)$ is computed for every merge to the main
        branch, producing a \emph{quality time series}
        $\{g(s^{(t)})\}_{t \geq 0}$.  Regressions are flagged when
        $g(s^{(t+1)}) < g(s^{(t)}) - \epsilon$ for a configurable
        tolerance~$\epsilon$.
\end{enumerate}
\end{construction}

\begin{remark}
The JuGeo CI layer is \emph{incremental}: only the affected
part of the \v{C}ech complex is recomputed.  If a commit touches
modules $M_i$ and~$M_j$, only the overlap cochains involving $M_i$
and~$M_j$ are re-evaluated.  This makes the computation
proportional to the size of the diff, not the size of the project.
\end{remark}

%% ─── 29.2 ──────────────────────────────────────────────────────
\section{Multi-repository coherence}
\label{sec:multi-repo}

Large organizations operate across monorepos or constellations of
microservice repositories.  JuGeo scales naturally to this setting.

\begin{definition}[Multi-Repository Site]
\label{def:multi-repo-site}
Let $R_1, \dots, R_n$ be repositories.  Define the
\emph{multi-repository site} $\STop_{\mathrm{multi}}$:
\begin{itemize}
  \item Objects: the modules of each $R_i$, plus inter-repository
        API surfaces.
  \item Morphisms: intra-repo dependency + inter-repo API calls.
  \item Topology: a family covers an API surface~$A$ if it
        includes both the provider and all consumers of~$A$.
\end{itemize}
\end{definition}

Under this site:
\begin{itemize}
  \item \textbf{API contracts are $\Cech^1$ conditions.}\;
        The overlap $U_i \cap U_j$ between a service and its
        client is the API contract (request/response schema,
        error codes, latency bounds).  A class
        $[\alpha] \in \check{H}^1$ with $[\alpha]\neq 0$ signals a
        \emph{contract violation}.
  \item \textbf{Breaking changes are $H^1 \neq 0$.}\;
        A breaking API change is precisely a modification that
        introduces a non-trivial class in $\check{H}^1$: the
        provider and consumer disagree on the overlap.
  \item \textbf{Versioning is cover refinement.}\;
        Introducing a new API version $v2$ alongside $v1$ refines
        the cover: consumers of $v1$ and $v2$ live in different
        open sets, and the overlap conditions split accordingly.
        This is sheaf-theoretically natural---it corresponds to
        passing to a finer topology.
\end{itemize}

\begin{proposition}
\label{prop:semver}
Under the multi-repository site, semantic versioning
(major/minor/patch) admits the following cohomological
interpretation:
\begin{enumerate}[label=(\roman*)]
  \item A \emph{patch} change preserves all sections
        ($\check{H}^1$ unchanged).
  \item A \emph{minor} change adds new sections but
        preserves existing ones
        ($\check{H}^1$ can only shrink).
  \item A \emph{major} change may introduce new classes in
        $\check{H}^1$ (existing sections may fail to glue).
\end{enumerate}
\end{proposition}

%% ─── 29.3 ──────────────────────────────────────────────────────
\section{Specification mining}
\label{sec:spec-mining}

A key practical obstacle to verification is that most code lacks
formal specifications.  JuGeo can \emph{infer} specifications from
existing code and tests.

\begin{construction}[Specification Mining via $\PStr$]
\label{con:spec-mining}
Given a module~$U$ with code~$E_U$ and test suite~$T_U$:
\begin{enumerate}
  \item \textbf{Mine candidate types.}\;
        Analyze $E_U$ to extract function signatures, invariants,
        and pre/post-conditions.  Each candidate is a potential
        type $A_U$ in~$\PStr$.
  \item \textbf{Tests as evidence.}\;
        Each passing test $t \in T_U$ is an evidence term
        $e_t : A_U$ in the mined type.  The grade
        $g(e_t : A_U)$ measures how well the test exercises the
        specification.
  \item \textbf{Cross-module consistency.}\;
        Compute $\check{H}^1$ for the mined specifications:
        if module $M_i$'s mined spec says ``\texttt{get\_user}
        returns a string'' but $M_j$'s tests treat the return
        value as a dictionary, $\check{H}^1 \neq 0$ flags the
        inconsistency.
  \item \textbf{Refinement.}\;
        Use the obstructions to refine the mined specifications
        until $\check{H}^1 = 0$.  The result is a consistent,
        machine-inferred specification for the entire project.
\end{enumerate}
\end{construction}

\begin{remark}
Specification mining is dual to code generation: generation starts
from $A$ and produces~$E$; mining starts from $E$ and infers~$A$.
In both cases, $\CSpec$ enforces global consistency.
\end{remark}

%% ─── 29.4 ──────────────────────────────────────────────────────
\section{Refactoring as descent}
\label{sec:refactoring}

Large-scale refactoring---renaming, module extraction, API
migration---is one of the most error-prone activities in software
engineering.  Under JuGeo, refactoring is \emph{exactly} descent.

\begin{proposition}[Refactoring = Descent]
\label{prop:refactoring}
Let $s^{(0)}$ be the current codebase (a section of $\PStr$) and
let $A'$ be the target specification (post-refactoring types and
interfaces).  Refactoring from $s^{(0)}$ to a section $s^{(*)}$
satisfying $A'$ is the descent problem for the presheaf
$\PStr_{A'}$ with initial candidate~$s^{(0)}$.  Specifically:
\begin{enumerate}
  \item The current code $s^{(0)}$ does not satisfy $A'$ on overlaps:
        $\check{H}^1(\STop, \PStr_{A'}) \neq 0$ with $s^{(0)}$ as
        candidate.
  \item Each refactoring step modifies $s^{(t)}$ on some $U_i$,
        producing $s^{(t+1)}$.
  \item Descent converges when $\check{H}^1 = 0$ relative to~$A'$.
\end{enumerate}
\end{proposition}

The $\GEval$ component tracks refactoring progress:
$g(s^{(t)})$ increases monotonically if each step resolves at least
one obstruction.  This provides a \emph{progress metric} for
refactoring projects that is more informative than ``percentage of
files changed.''

%% ─── 29.5 ──────────────────────────────────────────────────────
\section{Security analysis via cohomology}
\label{sec:security}

Security vulnerabilities often arise from \emph{non-local}
interactions: data flows from a trusted source through an untrusted
transformation to a sensitive sink.  This is a cohomological
phenomenon.

\begin{construction}[Security Presheaf]
\label{con:security-presheaf}
Define a presheaf $\mathcal{S}$ on the call-graph site of a program:
\begin{itemize}
  \item To each function $f$, assign the set of
        \emph{security labels} (e.g., public, confidential, secret)
        of its inputs and outputs.
  \item To each call edge $f \to g$, assign the restriction map
        that propagates labels through the call.
\end{itemize}
The sheaf condition requires that labels are consistent along every
path from source to sink.  Failure of the sheaf condition---a
non-trivial class in $\check{H}^1(\mathcal{G}, \mathcal{S})$---is
an \emph{information leak}: data classified ``secret'' reaches a
function labeled ``public'' via an intermediate path that fails to
enforce the label.
\end{construction}

\begin{example}
A web application has a function \texttt{get\_password} labeled
\texttt{secret} and a logging function \texttt{log} labeled
\texttt{public}.  An intermediate function \texttt{format\_user}
calls both.  The overlap cochain records that \texttt{format\_user}
must simultaneously handle \texttt{secret} input and produce
\texttt{public} output---a non-trivial $H^1$ class, flagging a
potential leak path.
\end{example}

Patching the vulnerability is descent: modify the code (repair the
section) until the security labels glue consistently and
$\check{H}^1 = 0$.

%% ─── 29.6 ──────────────────────────────────────────────────────
\section{Performance optimization}
\label{sec:perf-opt}

Performance can be treated as an additional stratum in the JuGeo
framework.

\begin{construction}[Performance Stratum]
\label{con:perf-stratum}
Augment the project site with a \emph{performance stratum}
$\sigma_{\mathrm{perf}}$:
\begin{itemize}
  \item To each function~$f$, assign a performance profile
        (time complexity, memory usage, cache behavior).
  \item The grade $g_{\mathrm{perf}}(f)$ measures how well $f$'s
        observed performance matches its target.
  \item \emph{Hotspots} are obstructions: functions where
        $g_{\mathrm{perf}} \ll \gone$.
  \item \emph{Optimization} is descent: iteratively improve the
        lowest-grade functions until
        $g_{\mathrm{perf}} \geq g_{\min}$.
\end{itemize}
\end{construction}

Cross-stratum descent is particularly powerful here: a performance
optimization that changes a function's interface (e.g., switching
from synchronous to asynchronous I/O) creates an obstruction in
the functional stratum.  JuGeo's cross-stratum descent detects this
and ensures that all callers are updated accordingly.

\begin{remark}
The performance stratum illustrates the \emph{extensibility} of
JuGeo: adding a new stratum requires only defining the appropriate
theory slots and grade functions.  The descent algorithm, the
cohomological machinery, and the $\Grade$ semiring are reused
unchanged.  This is the open-world property of CTT applied to
software engineering.
\end{remark}

\chapter{Future Directions}
\label{ch:future}

We close the monograph with a survey of open problems, speculative
applications, and directions for future research.  Some of these are
natural extensions of the existing framework; others require
fundamentally new mathematics.

%% ─── 30.1 ──────────────────────────────────────────────────────
\section{Higher-categorical Judgment Geometry}
\label{sec:higher-cat}

Throughout this monograph, $\STop$ has been a 1-topos: a category
of sheaves on a Grothendieck site.  The natural generalization
replaces 1-categories with $(\infty,1)$-categories.

\begin{itemize}
  \item \textbf{$\STop$ becomes an $\infty$-topos} in the sense of
        Lurie~\cite{lurie2009}.  The judgment presheaf $\PStr$
        becomes an $\infty$-presheaf valued in $\infty$-groupoids
        rather than sets; this captures not just whether two
        local sections agree on an overlap, but all \emph{higher
        coherences} (homotopies between homotopies, \emph{ad
        infinitum}).
  \item \textbf{$\PStr$ becomes homotopy type theory.}\;
        Types are no longer mere propositions or sets but
        \emph{spaces}; the identity type $\GId_g(a,b)$ of Part~III
        generalizes to a graded path space, and higher identity
        types capture graded homotopies.
  \item \textbf{$\CSpec$ becomes derived $\infty$-cohomology.}\;
        The \v{C}ech complex is replaced by the full derived
        functor of global sections in the $\infty$-topos.  The
        obstruction theory gains access to the full spectrum of
        cohomological invariants, including $H^n$ for all $n \geq 0$.
  \item \textbf{$\GEval$ becomes an $\infty$-graded structure.}\;
        The $\Grade$ semiring is replaced by a
        \emph{graded $\infty$-ring spectrum}, in which grades
        themselves form a homotopy-coherent algebraic structure.
\end{itemize}

\begin{remark}
The motivation for higher-categorical JuGeo is not merely
generality.  In software verification, higher coherences arise
naturally: a refactoring $r_1$ and a refactoring $r_2$ may each
preserve correctness, but their composition $r_2 \circ r_1$ may
not---and tracking \emph{why} requires a homotopy between the
two paths.  In linguistics, higher coherences appear in poetic
ambiguity: a word may have two readings, each reading two
sub-readings, and so forth.  The $\infty$-categorical framework
captures this tower of ambiguity natively.
\end{remark}

The key open problem is:

\begin{quote}
\emph{Does the descent algorithm of Part~II generalize to the
$\infty$-categorical setting with convergence guarantees?}
\end{quote}

We conjecture that it does, provided the $\infty$-graded structure
satisfies an appropriate finiteness condition (e.g., that the
homotopy groups of the grade spectrum stabilize).

%% ─── 30.2 ──────────────────────────────────────────────────────
\section{Quantum Judgment Geometry}
\label{sec:quantum}

Quantum computing introduces programs whose correctness cannot be
verified by classical testing (measuring a quantum state collapses
it).  JuGeo offers an alternative.

\begin{itemize}
  \item \textbf{The site} is the category of quantum circuits,
        with morphisms given by circuit composition and ancilla
        management.
  \item \textbf{Types} are quantum specifications: unitary
        conditions ($U^\dagger U = I$), entanglement invariants,
        and decoherence bounds.
  \item \textbf{Cohomology} captures topological invariants of
        quantum states---analogous to topological quantum
        computation, where the robustness of a quantum gate is
        measured by a topological invariant.
  \item \textbf{Grades} are fidelity measures: the grade
        $g(\rho, \sigma) = \log F(\rho, \sigma)$ where $F$ is the
        quantum fidelity between the ideal state $\sigma$ and the
        actual state~$\rho$.
\end{itemize}

\begin{remark}
The connection between JuGeo and topological quantum computing
is tantalizing: both frameworks use cohomological invariants to
measure the ``robustness'' of a structure.  Whether this analogy
can be made precise---whether quantum error correction codes
arise as descent data in a JuGeo site over quantum circuits---is
an open question.
\end{remark}

%% ─── 30.3 ──────────────────────────────────────────────────────
\section{Biological Judgment Geometry}
\label{sec:biological}

Biological macromolecules can be viewed as ``programs'' whose
``execution'' is physical.  JuGeo provides a verification framework.

\begin{itemize}
  \item \textbf{Codons as morphemes.}\;
        The genetic code maps nucleotide triplets to amino acids,
        exactly as morphological rules map morphemes to meanings.
        The $\morph$ stratum of CTT applies directly.
  \item \textbf{Protein folding as descent.}\;
        A protein's amino acid sequence is a ``local'' description;
        the folded 3D structure is the ``global section'' that the
        sequence must descend to.  Misfolding is $H^1 \neq 0$: the
        local energetic preferences of individual residues fail to
        glue into a globally consistent structure.
  \item \textbf{Mutations as obstructions.}\;
        A point mutation modifies the section on one open set
        (one residue's local environment).  Whether the mutant
        protein still folds correctly is whether the modified
        section still glues---a cohomological question.
  \item \textbf{Fitness as grade.}\;
        Darwinian fitness provides a natural $\GEval$: the grade
        of an organism is its reproductive success, and evolution
        is descent toward higher grades.
\end{itemize}

%% ─── 30.4 ──────────────────────────────────────────────────────
\section{Musical Judgment Geometry: DAW Projects as Programs}
\label{sec:musical}

A digital audio workstation (DAW) project file---a Renoise
\texttt{.xrns}, an Ableton Live \texttt{.als}, an FL~Studio
\texttt{.flp}---is a \emph{complete program}.  It contains tracks
(modules), patterns (functions), notes and effect commands
(instructions), signal-flow routing (data flow), automation curves
(continuous parameter functions), a mixer dependency graph (sends,
returns, sidechains), references to external plugins (libraries),
and a master bus (entry point).  The analogy is not metaphorical:
these project files admit exactly the same 4-tuple structure
$\JJ = (\STop, \PStr, \CSpec, \GEval)$ that governs software
verification.  In this section we develop that structure in detail,
using tracker-style DAWs (especially Renoise) as the running
example while noting how each concept maps to timeline-based
environments such as Ableton Live and FL~Studio.

% ────────────────────────────────────────────────────────────────
\subsection{The site \texorpdfstring{$\STop$}{S} for a DAW project}
\label{ssec:musical-site}

The site is the temporal-structural hierarchy of the project.

\begin{definition}[Musical site]
\label{def:musical-site}
Let $\mathcal{C}_{\mathrm{DAW}}$ be the category whose objects are
the structural units of the project at every scale:
\begin{enumerate}
  \item \textbf{Sample level.}\;
        Individual audio samples at the sample rate
        (e.g.\ 44\,100\,Hz or 48\,000\,Hz).
  \item \textbf{Tick level.}\;
        Tracker ticks or MIDI ticks---the finest rhythmic
        resolution.  In Renoise the default is 12~ticks per line;
        at 8~lines per beat this gives 96~ticks per beat.
  \item \textbf{Line / beat level.}\;
        Individual rows in a Renoise pattern (``lines''),
        corresponding to sixteenth-note or other subdivisions,
        or beats in a timeline DAW.
  \item \textbf{Pattern level.}\;
        Renoise's fundamental compositional unit: a rectangular
        grid of note columns and effect columns, of configurable
        length.  In Ableton this corresponds to a clip; in
        FL~Studio, a pattern clip.
  \item \textbf{Track level.}\;
        Instrument tracks, group tracks, send tracks, and the
        master track.
  \item \textbf{Song level.}\;
        The pattern sequence (Renoise's ``order list''), or the
        arrangement view timeline.
  \item \textbf{Project level.}\;
        The complete \texttt{.xrns}/\texttt{.als}/\texttt{.flp}
        file, including all referenced samples, instruments,
        plugin states, and configuration.
\end{enumerate}
Morphisms are \emph{restriction} maps: a pattern restricts to a
line, a line restricts to a tick, a track restricts to its content
within a single pattern, etc.  Composition of restrictions is
associative; identity restrictions exist at every level.
\end{definition}

The Grothendieck topology on $\mathcal{C}_{\mathrm{DAW}}$ is
defined by covering sieves:

\begin{itemize}
  \item A \emph{cover of the song} is its pattern sequence---the
        ordered collection of pattern instances that tile the
        timeline.
  \item A \emph{cover of a pattern} is its track decomposition---
        the collection of per-track columns within that pattern.
  \item The \emph{overlap} between adjacent patterns in the
        sequence captures crossfade regions, automation
        continuity at pattern boundaries, and sustained notes
        that cross a pattern break.
\end{itemize}

In Renoise, the \textbf{pattern matrix}---the grid of pattern
slots vs.\ tracks---\emph{is} the covering sieve, made visible as
a GUI element.  Each cell $(p, t)$ in the matrix is an open set
$U_{p,t}$, and the pattern matrix is the nerve of the cover.

\begin{definition}[Musical presheaf]
\label{def:musical-presheaf}
The \emph{structure presheaf} $\mathcal{F}$ on
$\mathcal{C}_{\mathrm{DAW}}$ assigns to each open set $U$ the
musical data active on $U$:
\begin{itemize}
  \item At a line: note values, velocities, instrument indices,
        panning, volume, and effect-column commands
        (e.g.\ \texttt{0Exx} retrigger, \texttt{0Bxx} pattern
        break, \texttt{0900} sample offset).
  \item At a track: the DSP effect chain, send amounts, group
        routing.
  \item At an instrument: sample mappings, modulation sets
        (envelopes, LFOs), macro assignments, plugin state.
  \item At a pattern: automation envelopes for any device
        parameter.
\end{itemize}
Restriction maps are the natural projections:
$\mathcal{F}(U) \to \mathcal{F}(V)$ for $V \hookrightarrow U$
restricts the data to the sub-region.
\end{definition}

\begin{remark}[Renoise effect commands as morphisms]
Pattern effect commands such as \texttt{0Bxx} (pattern break) and
\texttt{0Fxx} (set speed/LPB) are morphisms \emph{between} pattern
objects: they alter the flow of control in the pattern sequence,
just as jump instructions alter control flow in a program.
Instrument macros---Renoise's 8 assignable knobs that each control
an arbitrary set of device parameters---are natural transformations
between sub-presheaves, coordinating multiple restriction maps
through a single control surface.
\end{remark}

% ────────────────────────────────────────────────────────────────
\subsection{The proof structure \texorpdfstring{$\PStr$}{P} for
  music}
\label{ssec:musical-proof}

The \emph{type} of a DAW project is its musical specification---the
constraints that the finished piece must satisfy.

\begin{definition}[Musical type]
\label{def:musical-type}
A \emph{musical type} $\tau$ is a tuple of constraints drawn from:
\begin{itemize}
  \item \textbf{Global parameters.}\;
        BPM range, time signature, key, scale/mode.
  \item \textbf{Arrangement structure.}\;
        An ordered list of section labels---e.g.\
        $[\mathsf{intro}, \mathsf{verse}, \mathsf{chorus},
         \mathsf{bridge}, \mathsf{chorus}, \mathsf{outro}]$---with
        target durations.
  \item \textbf{Harmonic rules.}\;
        Chord progression grammars (e.g.\ $\mathrm{ii}$--$\mathrm{V}$--$\mathrm{I}$,
        circle-of-fifths motion), modulation constraints,
        voice-leading rules.
  \item \textbf{Rhythmic constraints.}\;
        Groove template, swing percentage, quantization grid.
  \item \textbf{Timbral constraints.}\;
        Frequency band allocation per track group, stereo-field
        assignments (e.g.\ bass and kick centered, hats and
        percussion panned), dynamic range targets.
  \item \textbf{Reference targets.}\;
        A reference track with measured LUFS loudness,
        spectral profile, and stereo width---the ``type'' the
        final master should inhabit.
\end{itemize}
\end{definition}

The \emph{evidence} inhabiting this type is the actual musical
content: the patterns, notes, automation curves, DSP chains, and
mixer settings that constitute the project.  The \emph{typing
derivation} is the compositional structure that connects each piece
of evidence to its constraint:

\begin{itemize}
  \item Each pattern's content derives from the arrangement type
        for that section.
  \item Each note in a pattern derives from the harmonic type
        active at that moment.
  \item Each DSP device in a chain derives from the timbral type
        for that track's frequency/stereo allocation.
  \item Each automation curve derives from the dynamic contour
        specified by the arrangement arc.
\end{itemize}

\begin{definition}[Renoise-specific types]
\label{def:renoise-types}
We define the following type constructors for Renoise projects:
\begin{align*}
  &\mathsf{PatternType}(\ell, n_t, \mathit{lpb})
    &&\text{--- pattern of length $\ell$ lines, $n_t$ tracks,
           $\mathit{lpb}$ lines per beat,} \\
  &\mathsf{InstrumentType}(S, M, K)
    &&\text{--- instrument with sample set $S$, modulation sets $M$,
           macro map $K$,} \\
  &\mathsf{EffectChainType}(D_1, \ldots, D_k)
    &&\text{--- chain of DSP devices $D_i$ in signal-flow order,} \\
  &\mathsf{AutomationType}(p, e)
    &&\text{--- automation envelope $e$ for device parameter $p$.}
\end{align*}
\end{definition}

A well-typed project is one where every piece of musical content
has a derivation from the specification.  An untyped passage---a
note outside the scale, a pattern with no arrangement role, an
unrouted send---is a type error.

% ────────────────────────────────────────────────────────────────
\subsection{The cohomological spectrum \texorpdfstring{$\CSpec$}{H}
  for music}
\label{ssec:musical-cohomology}

Cohomology detects precisely the obstructions that separate a
collection of locally acceptable musical fragments from a
\emph{finished production}.

\paragraph{$H^0$: Global consistency.}
$H^0(\mathcal{C}_{\mathrm{DAW}}, \mathcal{F})$ is the space of
global sections---the data that is consistent across every open set.
Non-trivial $H^0$ elements include:
\begin{itemize}
  \item The master bus output is valid audio (finite amplitude,
        non-silent, within 0\,dBFS).
  \item All sample references in all instruments resolve to
        existing files.
  \item All plugin identifiers (\texttt{VST2}, \texttt{VST3},
        \texttt{AU}) resolve to installed binaries.
  \item BPM and time signature are consistent across the entire
        song (or the tempo automation is continuous).
\end{itemize}

\paragraph{$H^1$: Obstructions to gluing.}
$H^1 \neq 0$ is where the real problems live---where locally
acceptable sections fail to compose into a coherent whole.  Each
class of $H^1$ corresponds to a specific production problem:

\begin{enumerate}
  \item \textbf{Phase cancellation.}\;
        Two tracks---e.g.\ a layered snare top and snare
        bottom---that sound full in solo but partially cancel when
        summed to the mix bus.  This is a textbook $H^1$ class:
        $\mathcal{F}|_{U_{\mathrm{top}}}$ and
        $\mathcal{F}|_{U_{\mathrm{bot}}}$ are each satisfactory,
        but $\mathcal{F}|_{U_{\mathrm{top}} \cap
        U_{\mathrm{bot}}} = \mathcal{F}|_{\mathrm{sum}}$ has a
        non-trivial obstruction.
  \item \textbf{Frequency masking.}\;
        A kick drum and a bass synth both occupy the
        60--120\,Hz band.  Each sounds powerful in solo; together
        they produce mud.  The overlap
        $U_{\mathrm{kick}} \cap U_{\mathrm{bass}}$ in the
        spectral domain is the source of the obstruction.
  \item \textbf{Key clashes at pattern boundaries.}\;
        Pattern~$A$ is in C~major; pattern~$B$ modulates to
        F$\sharp$~major.  Each pattern is harmonically
        well-formed, but the transition---the overlap in the
        covering sieve---carries an $H^1$ class measuring the
        smoothness of the modulation.
  \item \textbf{Automation discontinuities.}\;
        A filter cutoff automation in pattern~$A$ ends at
        value~$x$; the same parameter in pattern~$B$ begins at
        value~$y \neq x$.  At the pattern boundary, the
        presheaf has a discontinuity---an $H^1$ obstruction in
        the automation sub-presheaf.
  \item \textbf{Sustained-note truncation.}\;
        A note sustains across a pattern break via Renoise's
        \texttt{OFF}/\texttt{---} convention, but the next
        pattern's first line contains a conflicting note-on in
        the same column.  The gluing fails.
  \item \textbf{Sidechain conflicts.}\;
        Two instruments sidechained to the same source
        (e.g.\ kick) with conflicting ducking envelopes: one
        demands fast release, the other demands slow release.
        The shared dependency creates an obstruction in the
        signal-flow presheaf.
  \item \textbf{Tempo automation breaks.}\;
        Pattern at 120~BPM jumps to 140~BPM without a
        ritardando or accelerando transition.  The tempo
        presheaf is discontinuous.
\end{enumerate}

\paragraph{$H^2$ and higher: Structural obstructions.}
Higher cohomology captures arrangement-level issues visible only in
full-song context:
\begin{itemize}
  \item The verse--chorus--verse--chorus form lacks a satisfying
        energy arc (no build, no release).
  \item Mix balance is locally fine section-by-section but the
        overall spectral center of gravity shifts
        unpredictably across the arrangement.
  \item Motivic development stalls---themes introduced in the
        intro never return or transform.
\end{itemize}

\begin{proposition}[Mayer--Vietoris for mixing]
\label{prop:mv-mixing}
Let $A$ and $B$ be two track groups (e.g.\ drums and bass).  Their
spectral/stereo overlap is $A \cap B$.  The Mayer--Vietoris
sequence
\[
  \cdots \to H^0(A \cup B) \to H^0(A) \oplus H^0(B)
  \to H^0(A \cap B) \xrightarrow{\;\partial\;}
  H^1(A \cup B) \to \cdots
\]
gives the connecting homomorphism $\partial$ that sends ``local
spectral compatibility'' to ``global mixing obstruction.''
Concretely: $A$ and $B$ each pass their solo mix checks
($H^0(A), H^0(B)$ trivial), but their shared frequency content
$H^0(A \cap B)$ maps via $\partial$ to a non-trivial $H^1(A \cup
B)$ class---frequency masking.  The standard mixing repair is to
apply complementary EQ curves to $A$ and $B$, carving space in
$A \cap B$ until $\partial = 0$.
\end{proposition}

% ────────────────────────────────────────────────────────────────
\subsection{The graded evaluation \texorpdfstring{$\GEval$}{Q} for
  music}
\label{ssec:musical-grading}

The graded evaluation for a DAW project is multi-stratal, with
each stratum measuring an independent axis of production quality.

\begin{definition}[Musical grade algebra]
\label{def:musical-grade}
Define the grade monoid
$G_{\mathrm{DAW}} = (G, \otG, \gone, \opG, \gzero)$ with strata:
\begin{center}
\renewcommand{\arraystretch}{1.2}
\begin{tabular}{@{}lll@{}}
\textbf{Stratum} & \textbf{What it measures} &
  \textbf{Computation} \\
\hline
$\mathsf{mix}$         & Mix quality
  & LUFS loudness, dynamic range, frequency balance, stereo width \\
$\mathsf{master}$      & Mastering quality
  & Peak level, RMS, crest factor, spectral tilt \\
$\mathsf{arrangement}$ & Structural quality
  & Section lengths, contrast ratios, tension/release arc \\
$\mathsf{harmony}$     & Harmonic quality
  & Voice leading, progression logic, modulation smoothness \\
$\mathsf{rhythm}$      & Rhythmic quality
  & Groove consistency, syncopation interest, quantize accuracy \\
$\mathsf{timbre}$      & Timbral quality
  & Spectral diversity, band allocation, stereo separation \\
$\mathsf{melody}$      & Melodic quality
  & Contour balance, interval distribution, motivic development \\
$\mathsf{production}$  & Production polish
  & Noise floor, artifact absence, latency compensation \\
\end{tabular}
\end{center}
The overall grade is the $\otG$-product across all strata:
\[
  g_{\mathrm{total}} = g_{\mathsf{mix}} \otG\, g_{\mathsf{master}}
    \otG\, g_{\mathsf{arrangement}} \otG\, g_{\mathsf{harmony}}
    \otG\, g_{\mathsf{rhythm}} \otG\, g_{\mathsf{timbre}}
    \otG\, g_{\mathsf{melody}} \otG\, g_{\mathsf{production}}.
\]
\end{definition}

The $\otG$-product enforces the \textbf{bottleneck principle}: a
track with brilliant melodies ($g_{\mathsf{melody}} \approx
\gone$) but a clipping, distorted mix ($g_{\mathsf{mix}} \approx
\gzero$) receives a low total grade, because $\gone \otG \gzero
= \gzero$.  This is the formal expression of the producer's maxim
that a great song with a bad mix never gets played.

\begin{example}[Loudness war as grade collapse]
The ``loudness war'' is the phenomenon of over-compressing masters
to maximize perceived loudness at the cost of dynamic range.  In
our grading, aggressive brick-wall limiting pushes
$g_{\mathsf{master}}$ toward $\gone$ (the LUFS target is met) but
simultaneously drives $g_{\mathsf{mix}} \to \gzero$ (dynamic range
collapses, transients are destroyed).  The $\otG$-product correctly
identifies this as a degradation: the master stratum cannot
compensate for a destroyed mix stratum.
\end{example}

% ────────────────────────────────────────────────────────────────
\subsection{Descent as the production workflow}
\label{ssec:musical-descent}

The actual DAW production workflow is \emph{descent}: a sequence of
refinements that reduce $H^1$ obstructions and increase $\GEval$
grades.

\begin{enumerate}
  \item \textbf{Sketch (initial candidate).}\;
        Rough patterns with placeholder sounds---a bare melodic
        idea in a default sine wave, a four-on-the-floor kick from
        the factory library.  This candidate has many $H^1$
        obstructions: no frequency separation, no arrangement
        logic, placeholder timbres.
  \item \textbf{Arrangement (first descent step).}\;
        Organize patterns into a song structure---intro, build,
        drop, break, drop, outro.  Assign section types from the
        musical type $\tau$.  This resolves structural $H^1$
        classes (the piece now has a coherent macro-form) and
        increases $g_{\mathsf{arrangement}}$.
  \item \textbf{Sound design (second step).}\;
        Replace placeholder instruments with designed sounds:
        layer samples in Renoise's instrument editor, sculpt
        wavetables, program modulation sets and macros.  This
        resolves timbral $H^1$ (the sounds now occupy distinct
        spectral regions) and increases $g_{\mathsf{timbre}}$.
  \item \textbf{Mixing (third step).}\;
        Apply EQ, compression, saturation, and spatial
        processing.  Configure send tracks for reverb and delay.
        Set up sidechain routing.  Automate levels across
        sections.  This resolves frequency-masking, phase, and
        stereo $H^1$ classes, and increases $g_{\mathsf{mix}}$.
  \item \textbf{Mastering (final step).}\;
        Master bus processing: multi-band compression, stereo
        imaging, limiting to target LUFS.  Resolves loudness and
        dynamic-range $H^1$, and pushes $g_{\mathsf{master}}$
        toward $\gone$.
\end{enumerate}

Each step is one iteration of the descent loop from
Definition~\ref{def:musical-site}: produce a candidate section,
compute $\CSpec$ (detect obstructions), compute $\GEval$ (grade
quality), repair the lowest-grade stratum, and repeat.

\begin{remark}[Renoise-specific descent moves]
Renoise's pattern effect commands are \emph{local repair
operations}:
\begin{itemize}
  \item \texttt{0Exx} (retrigger) repairs rhythmic sparseness by
        subdividing a single note event.
  \item \texttt{0Cxx} (set volume) repairs dynamic-range
        obstructions within a pattern.
  \item \texttt{0Gxx} (glide-to-note) repairs melodic
        discontinuities by interpolating between pitches.
  \item DSP chain reordering---EQ before vs.\ after
        compression---changes the $H^1$ structure: pre-EQ
        tightens the compressor's response (repairs spectral
        obstructions first), while post-EQ allows broader
        dynamics but may reintroduce masking.
  \item Instrument macro mappings automate multi-parameter descent:
        a single macro knob can simultaneously adjust filter
        cutoff, resonance, and envelope decay, maintaining
        timbral consistency as one control sweeps a family of
        coherent states.
\end{itemize}
\end{remark}

% ────────────────────────────────────────────────────────────────
\subsection{Music generation via \texorpdfstring{$\FMLM$}{FMLM}}
\label{ssec:musical-fmlm}

JuGeo can \emph{generate} music, not just verify it.  The
$\FMLM$ loop (Fill--Measure--Locate--Mend) applies directly.

\begin{example}[Techno track generation]
\label{ex:techno-gen}
Start with a type:
\[
  \tau = \mathsf{Song}\!\left(
    \begin{aligned}
      &\mathit{genre} = \mathsf{techno},\;
       \mathit{bpm} = 128,\;
       \mathit{key} = \mathrm{A}\,\text{minor}, \\
      &\mathit{structure} = [\mathsf{intro}, \mathsf{build},
        \mathsf{drop}, \mathsf{break}, \mathsf{drop},
        \mathsf{outro}]
    \end{aligned}
  \right).
\]
The $\FMLM$ loop proceeds:
\begin{enumerate}
  \item \textbf{Fill (iteration~1): drums.}\;
        For each section, fill the drum track patterns---kick,
        hats, clap, ride---respecting the rhythmic constraints
        of $\tau$.  After this step, the project has rhythmic
        content but massive $H^1$ (no harmonic content, no
        spectral balance).
  \item \textbf{Measure:}\;
        compute
        $g_{\mathsf{rhythm}} \approx 0.7$,
        $g_{\mathsf{harmony}} = \gzero$,
        $g_{\mathsf{timbre}} \approx 0.2$.
  \item \textbf{Fill (iteration~2): bass line.}\;
        Add a bass track constrained by the key
        ($\mathrm{A}\,\text{minor}$) and the groove (locked to
        kick pattern).  Frequency band: 30--200\,Hz.
  \item \textbf{Measure:}\;
        $g_{\mathsf{harmony}} \approx 0.4$ (bass provides root
        motion), $H^1_{\mathrm{mask}}$ appears between kick and
        bass in the 50--80\,Hz band.
  \item \textbf{Locate + Mend:}\;
        apply sidechain compression from kick to bass
        (resolve masking $H^1$), high-pass the bass at 35\,Hz.
  \item \textbf{Fill (iteration~3): lead/pad.}\;
        Add a pad track (200\,Hz--2\,kHz) and a lead
        (1--8\,kHz), constrained by the harmonic type and the
        remaining spectral space.
  \item \textbf{Fill (iteration~4): FX and transitions.}\;
        Add risers, sweeps, impacts at section boundaries.
        Automate filter cutoffs for builds.  Resolve
        arrangement-level $H^1$ (section transitions are now
        smooth).
  \item \textbf{Final Measure:}\;
        all strata near $\gone$; $H^1 \approx 0$; the project
        is a well-typed, cohomologically trivial,
        high-grade \texttt{.xrns} file.
\end{enumerate}
This is \emph{exactly} how an experienced tracker musician
works---JuGeo formalizes the workflow.
\end{example}

% ────────────────────────────────────────────────────────────────
\subsection{The isomorphism: music \texorpdfstring{$\cong$}{≅}
  software}
\label{ssec:musical-iso}

The structural parallel between a DAW project and a software
project is not a loose analogy but a genuine categorical
equivalence mediated by the 4-tuple.

\begin{proposition}[DAW--software equivalence]
\label{prop:daw-software}
There exists a structure-preserving functor
$\Phi \colon \mathcal{C}_{\mathrm{DAW}} \to
\mathcal{C}_{\mathrm{SW}}$ that maps:
\begin{center}
\renewcommand{\arraystretch}{1.2}
\begin{tabular}{@{}ll@{}}
\textbf{Music (DAW project)} & \textbf{Software project} \\
\hline
Track              & Module / thread \\
Pattern            & Function / procedure \\
Note / effect cmd  & Instruction / statement \\
Signal flow        & Data flow \\
Mixer routing      & Dependency graph \\
Plugin (VST/AU)    & Library (npm, pip, crate) \\
Frequency masking  & Resource contention \\
Phase cancellation & Race condition \\
Clipping (0\,dBFS) & Buffer overflow \\
Sidechain routing  & Semaphore / lock \\
Mastering          & Deployment / release \\
Mix balance        & Load balancing \\
Automation curve   & Configuration schedule \\
Pattern matrix     & Module dependency matrix \\
\end{tabular}
\end{center}
Under $\Phi$, the 4-tuples are compatible:
$\Phi_*(\STop_{\mathrm{DAW}}) = \STop_{\mathrm{SW}}$,
$\Phi_*(\PStr_{\mathrm{DAW}}) = \PStr_{\mathrm{SW}}$, and
similarly for $\CSpec$ and $\GEval$ (with appropriate stratum
relabeling).
\end{proposition}

\begin{remark}
The equivalence is productive in both directions.  Software
engineering concepts suggest musical techniques: \emph{unit
testing} corresponds to \emph{solo monitoring} (check each track
in isolation), \emph{integration testing} to \emph{group bus
monitoring} (check track groups together), and \emph{end-to-end
testing} to \emph{full-mix playback}.  Conversely, musical
concepts suggest software techniques: \emph{sidechain
compression}---dynamically ducking one signal in response to
another---is a pattern for adaptive resource throttling, and
\emph{send/return routing}---sharing a single reverb instance
across many tracks---is a service mesh pattern.
\end{remark}

\begin{remark}
Algorithmic composition has a long history, from Hiller and
Isaacson's \emph{Illiac Suite}~(1957) to modern neural audio
synthesis.  What JuGeo adds is the cohomological dimension: not
merely ``does each pattern sound acceptable in isolation?''
($H^0$) but ``do the patterns compose into a coherent
production?'' ($H^1$).  This is precisely the difference between
a collection of loops and a \emph{finished track}---and it is the
same difference that separates a collection of working functions
from a \emph{shipped program}.
\end{remark}

%% ─── 30.5 ──────────────────────────────────────────────────────
\section{Ethical Judgment Geometry}
\label{sec:ethical}

Ethical reasoning involves multiple, potentially conflicting moral
frameworks---a natural setting for graded, multi-theory judgment.

\begin{itemize}
  \item \textbf{Moral frameworks as theory slots.}\;
        Utilitarianism, deontology, virtue ethics, care ethics,
        and other frameworks each become a theory slot in the CTT.
        Each provides a grade for a given action or policy.
  \item \textbf{Ethical dilemmas as $H^1 \neq 0$.}\;
        An ethical dilemma is precisely a situation where different
        moral frameworks assign contradictory verdicts: the local
        sections (framework-specific judgments) fail to glue.
        The resulting $H^1$ class identifies the \emph{source} of
        the disagreement.
  \item \textbf{Resolution via graded evaluation.}\;
        Rather than forcing a binary choice, $\GEval$ assigns a
        grade that reflects the degree of moral consensus or
        conflict.  A policy with $g \approx \gone$ enjoys broad
        agreement; a policy with $g \ll \gone$ is deeply contested.
  \item \textbf{Pluralism via the open world property.}\;
        New ethical frameworks can be added to the theory ensemble
        without modifying existing ones.  This is genuine moral
        pluralism: no framework is privileged, and the graded
        evaluation tracks how well any action satisfies the full
        ensemble.
\end{itemize}

%% ─── 30.6 ──────────────────────────────────────────────────────
\section{JuGeo for education}
\label{sec:education}

Educational assessment is a domain where graded judgment is
the norm, not the exception.

\begin{itemize}
  \item \textbf{Student work as evidence.}\;
        A student's submission $E$ is evidence for a rubric
        specification~$A$.  The judgment
        $\vdash_g E : A$ says ``$E$ satisfies $A$ to grade~$g$\!.''
  \item \textbf{Partial credit.}\;
        The $\Grade$ semiring provides a principled theory of partial
        credit.  Each rubric criterion is a theory slot with its own
        grade; the total grade is $\otG$-aggregated (all criteria
        matter) or $\opG$-aggregated (best-of, for alternative
        approaches).
  \item \textbf{Feedback as obstruction identification.}\;
        When $g < \gone$, the obstructions identify which
        rubric criteria were not fully met.  This is structured
        feedback: not just ``your grade is 7/10'' but ``the
        logical structure criterion is at $-0.2$ because of a gap
        in the proof of Lemma~3.''
  \item \textbf{Automated grading via descent.}\;
        The grading process itself can be automated:
        parse the submission, compute theory-slot grades, aggregate,
        and report obstructions.  This is structurally identical to
        program verification---the student's work is a ``program''
        and the rubric is its ``specification.''
\end{itemize}

%% ─── 30.7 ──────────────────────────────────────────────────────
\section{Self-improving JuGeo}
\label{sec:self-improving}

In the spirit of G\"odelian self-reference, JuGeo can be applied
to its own development.

\begin{construction}[JuGeo Bootstrap]
\label{con:bootstrap}
Let $\STop_{\mathrm{JG}}$ be the site of JuGeo's own codebase
(the modules in \texttt{src/jugeo/}).  Define:
\begin{itemize}
  \item $\PStr_{\mathrm{JG}}$: JuGeo's specification (the
        mathematical definitions of this monograph) together with
        the implementation as evidence.
  \item $\CSpec_{\mathrm{JG}}$: the \v{C}ech cohomology of the
        implementation relative to the specification.
        $H^1 \neq 0$ identifies modules whose implementation
        deviates from the formal theory.
  \item $\GEval_{\mathrm{JG}}$: the graded evaluation of JuGeo
        by JuGeo---using the very descent algorithm being verified
        to check that the descent algorithm is correctly
        implemented.
\end{itemize}
\end{construction}

This bootstrap is not circular: JuGeo version $n$ verifies JuGeo
version $n+1$, and each version's correctness is established
relative to the previous one.  The process is analogous to
compiler bootstrapping, where compiler version $n$ compiles
version~$n+1$.

%% ─── 30.8 ──────────────────────────────────────────────────────
\section{Open mathematical problems}
\label{sec:open-problems}

We close with a list of open problems that the framework raises.

\begin{enumerate}
  \item \textbf{Decidability.}\;
        Is the full 4-tuple verification problem ($\check{H}^1 = 0$
        and $\GEval \geq g_{\min}$) decidable for any interesting
        fragment of the site category?  For finite sites with
        computable grade functions the answer is trivially yes; for
        sites arising from Turing-complete programming languages
        it is trivially no.  The interesting question is: what is
        the largest class of sites for which the problem remains
        decidable, and does this class cover practically relevant
        programs?

  \item \textbf{Homotopy type of the judgment fibration.}\;
        The judgment fibration $\pi\colon \mathcal{E} \to \STop$
        is a Grothendieck fibration with rich structure.  What are
        its homotopy invariants?  In particular, is the classifying
        space $B\!\operatorname{Aut}(\mathcal{E})$ related to known
        classifying spaces in algebraic topology?

  \item \textbf{Spectral sequences for NLP.}\;
        The \v{C}ech-to-derived spectral sequence converging to
        $H^*(\STop, \PStr)$ provides a computational tool that, to
        our knowledge, has never been applied to NLP\@.  Can the
        differentials $d_r\colon E_r^{p,q} \to E_r^{p+r,q-r+1}$
        be given linguistic interpretations?  Could this yield new
        parsing or generation algorithms that operate stratum-by-stratum?

  \item \textbf{Continuous JuGeo.}\;
        Replace the discrete site with a smooth manifold (or a
        diffeological space).  Sections become smooth functions,
        the \v{C}ech complex becomes the de Rham complex, and
        obstructions become differential forms.  Does this
        ``continuous JuGeo'' yield a useful framework for
        cyber-physical systems, control theory, or signal processing?

  \item \textbf{Categorification of the grade.}\;
        The $\Grade$ semiring is a 0-categorical structure.  Can it
        be categorified---replacing grades by grading categories,
        and the semiring operations by functors---to capture
        ``grades of grades''?  This would give a natural notion of
        \emph{confidence in the grade itself}, addressing the
        meta-question ``how reliable is our quality assessment?''

  \item \textbf{Effective descent bounds.}\;
        Proposition~\ref{prop:llm-descent-mono} gives a worst-case
        bound on descent iterations.  What is the \emph{expected}
        number of iterations under realistic distributions of LLM
        errors?  Is there a phase-transition phenomenon: a critical
        error rate above which descent fails to converge?

  \item \textbf{JuGeo and dependent type theory.}\;
        Part~III introduced graded dependent types.  To what extent
        can JuGeo be \emph{internalized} in a dependent type theory
        (e.g., Agda or Lean~4)?  A full internalization would yield
        machine-checked proofs of JuGeo's own metatheory.
\end{enumerate}

\bigskip

These problems span category theory, homotopy theory, computability,
and applied mathematics.  We hope that the framework developed in
this monograph---simple enough to be stated in a single 4-tuple,
rich enough to organize program verification, natural language
generation, and beyond---will inspire progress on all of them.


% ── Back matter ───────────────────────────────────────────────
\backmatter

\begin{thebibliography}{99}

\bibitem{young2025jg}
H.~Young.
\newblock Judgment Geometry: A Sheaf-Theoretic Framework for
  Universal Program Verification and Generation.
\newblock 2025.

\bibitem{maclane1971}
S.~Mac~Lane.
\newblock \emph{Categories for the Working Mathematician}.
\newblock Springer, 1971.

\bibitem{lurie2009}
J.~Lurie.
\newblock \emph{Higher Topos Theory}.
\newblock Princeton University Press, 2009.

\bibitem{johnstone2002}
P.~Johnstone.
\newblock \emph{Sketches of an Elephant: A Topos Theory Compendium}.
\newblock Oxford University Press, 2002.

\bibitem{mlm1984}
P.~Martin-L\"of.
\newblock \emph{Intuitionistic Type Theory}.
\newblock Bibliopolis, 1984.

\bibitem{hott2013}
The Univalent Foundations Program.
\newblock \emph{Homotopy Type Theory: Univalent Foundations of Mathematics}.
\newblock Institute for Advanced Study, 2013.

\bibitem{weibel1994}
C.~Weibel.
\newblock \emph{An Introduction to Homological Algebra}.
\newblock Cambridge University Press, 1994.

\bibitem{grothendieck1957}
A.~Grothendieck.
\newblock Sur quelques points d'alg\`ebre homologique.
\newblock \emph{T\^ohoku Math.\ J.}, 9:119--221, 1957.

\bibitem{prince1993}
A.~Prince and P.~Smolensky.
\newblock \emph{Optimality Theory}.
\newblock Rutgers University, 1993.

\bibitem{montague1973}
R.~Montague.
\newblock The proper treatment of quantification in ordinary English.
\newblock In \emph{Approaches to Natural Language}, 1973.

\bibitem{kamp1993}
H.~Kamp and U.~Reyle.
\newblock \emph{From Discourse to Logic}.
\newblock Kluwer, 1993.

\bibitem{girard1987}
J.-Y.~Girard.
\newblock Linear logic.
\newblock \emph{Theoretical Computer Science}, 50:1--102, 1987.

\bibitem{pustejovsky1995}
J.~Pustejovsky.
\newblock \emph{The Generative Lexicon}.
\newblock MIT Press, 1995.

\end{thebibliography}

\end{document}
